% Generated by tex_ver/convert.py from ../QFT.md. Edit the Markdown source.
% Regenerate from the QFT directory: /usr/bin/python3 tex_ver/convert.py
% Manual changes to this file are overwritten. QFT.md SHA256: 5534009e0e1c84a53eb54d22b53d082ba550858e30f40d399d2854aab8bf2cc8
\synctex=1
\documentclass[11pt,a4paper,oneside]{report}
\usepackage[T1]{fontenc}
\usepackage{lmodern,CJKutf8,microtype}
\usepackage[a4paper,hmargin=21mm,top=23mm,bottom=24mm,headheight=14pt,headsep=8mm]{geometry}
\usepackage{amsmath,amssymb,slashed,booktabs,enumitem,fancyhdr,needspace}
\usepackage{xcolor}
\newcommand{\qftdisplayarg}[1]{{\displaystyle #1}}
\newcommand{\qftalignarg}[1]{#1}
\NewDocumentCommand{\qftcolorarg}{o m o m m}{\mbox{#5}}
\usepackage[colorlinks,linkcolor=blue!60!black,urlcolor=blue!60!black,bookmarksnumbered,bookmarksdepth=2]{hyperref}
\hypersetup{pdftitle={QFT},pdfauthor={Tianrong Cheng}}
\makeatletter
\ifdefined\pdfnobuiltintounicode
  \AddToHook{cmd/CJK@addcmap/after}{\pdfnobuiltintounicode\font@name}
\fi
\renewcommand\@makechapterhead[1]{{\parindent\z@\raggedright
  \normalfont\large\scshape \@chapapp\space\thechapter\par\nobreak
  \vskip 6pt\normalfont\huge\bfseries\boldmath #1\par\nobreak\vskip 19pt}}
\renewcommand\@makeschapterhead[1]{{\parindent\z@\raggedright
  \normalfont\huge\bfseries #1\par\nobreak\vskip 18pt}}
\renewcommand\section{\@startsection{section}{1}{\z@}{-3ex plus -.6ex}{1.3ex}{\normalfont\Large\bfseries\boldmath}}
\renewcommand\subsection{\@startsection{subsection}{2}{\z@}{-2.3ex plus -.4ex}{.85ex}{\normalfont\large\bfseries\boldmath}}
\newcommand\topic{\@startsection{paragraph}{4}{\z@}{-1.7ex plus -.3ex}{.5ex}{\normalfont\normalsize\bfseries\itshape\boldmath}*}
\newcommand\divtitle[1]{\par\addvspace{1.1ex}\noindent
  {\normalfont\normalsize #1}\par\nobreak\vspace{.35ex}\@afterindentfalse\@afterheading}
\newcommand{\opentocfile}{\newwrite\tf@toc\immediate\openout\tf@toc=\jobname.toc\relax}
\makeatother
\newenvironment{qftinset}{\begin{list}{}{\setlength{\leftmargin}{2em}\setlength{\rightmargin}{0pt}\setlength{\listparindent}{\parindent}}\item[]}{\end{list}}
\newcommand{\qftdiagram}[3]{\begin{center}\leavevmode
  \hbox to #1bp{\vrule width 0pt height #2bp depth 0pt\pdfliteral{#3}\hss}\end{center}}
\setcounter{secnumdepth}{1}
\setcounter{tocdepth}{2}
\setlength{\parindent}{1.4em}
\setlength{\parskip}{0pt}
\setlength{\emergencystretch}{2em}
\setlist{topsep=3pt plus 1pt,itemsep=2pt,parsep=0pt,partopsep=0pt,leftmargin=1.55em}
\raggedbottom
\pagestyle{fancy}\fancyhf{}
\renewcommand{\chaptermark}[1]{\markboth{\chaptername\ \thechapter\enspace #1}{}}
\fancyhead[L]{\small\nouppercase{\leftmark}}
\fancyfoot[C]{\thepage}
\renewcommand{\headrulewidth}{.25pt}
% The two title columns share the same 3.8em left edge.
\newcommand{\tocchapter}[3]{\par\Needspace{4\baselineskip}\addvspace{8pt}%
  {\bfseries\noindent\hyperref[#1]{\makebox[1.5em][l]{#2}#3}\hfill\hyperref[#1]{\pageref*{#1}}\par}\nobreak}
\newcommand{\tocsection}[3]{\par\addvspace{1.2pt}\noindent\hspace*{1.4em}%
  \hyperref[#1]{\makebox[2.4em][l]{#2}#3}\nobreak\kern.35em\dotfill\nobreak\kern.35em\hyperref[#1]{\pageref*{#1}}\par}
\newenvironment{tocgroup}{\par\nobreak\begingroup\leftskip=3.8em\relax
 \rightskip=1em plus 1fil\relax
  \small\parindent=0pt\interlinepenalty=10000\leavevmode}{\par\endgroup}
\newcommand{\tocsep}{\nobreak\,\textperiodcentered\space}
\begin{document}
\begin{CJK*}{UTF8}{gbsn}
\thispagestyle{empty}
\vspace*{.23\textheight}
\begin{center}{\Huge\bfseries QFT\par}
\vspace{1.4em}
{\Large Tianrong Cheng\par}
\vspace{.8em}
{\large \today\par}
\end{center}
\clearpage
\chapter*{Contents}\markboth{Contents}{}\opentocfile
\tocchapter{book-chapter-1}{1}{Classical Field Theory}
\tocsection{book-section-1.1}{1.1}{Lagrangian Formalism}
\begin{tocgroup}\hyperref[book-subsection-section-1.1-1]{Lagrangian}%
\tocsep{}\hyperref[book-subsection-section-1.1-2]{Symmetry Transformations and Noether's Theorem}%
\end{tocgroup}
\tocsection{book-section-1.2}{1.2}{Hamiltonian Formalism}
\begin{tocgroup}\hyperref[book-subsection-section-1.2-1]{Dynamical Systems}%
\tocsep{}\hyperref[book-subsection-section-1.2-2]{Hamiltonian Formalism}%
\tocsep{}\hyperref[book-subsection-section-1.2-3]{Symmetries and Noether's Theorem}%
\end{tocgroup}
\tocsection{book-section-1.3}{1.3}{Definition of Classical Fields}
\begin{tocgroup}\hyperref[book-subsection-section-1.3-1]{Definition of Fields}%
\tocsep{}\hyperref[book-subsection-section-1.3-2]{Finite-Dimensional Representations of $SO(1,3)^+$}%
\end{tocgroup}
\tocsection{book-section-1.4}{1.4}{Dynamics of Classical Fields}
\begin{tocgroup}\hyperref[book-subsection-section-1.4-1]{Lagrangian}%
\tocsep{}\hyperref[book-subsection-section-1.4-2]{Symmetry Transformations and Noether's Theorem}%
\end{tocgroup}
\tocsection{book-section-1.5}{1.5}{Dirac Constraint}
\begin{tocgroup}\hyperref[book-subsection-section-1.5-1]{Singular Lagrangians}%
\end{tocgroup}
\tocsection{book-section-1.6}{1.6}{The Scalar Field}
\begin{tocgroup}\hyperref[book-subsection-section-1.6-1]{Lagrangian}%
\tocsep{}\hyperref[book-subsection-section-1.6-2]{Solutions of the Klein–Gordon Equation}%
\end{tocgroup}
\tocsection{book-section-1.7}{1.7}{The Dirac Field $\psi$}
\begin{tocgroup}\hyperref[book-subsection-section-1.7-1]{Lagrangian}%
\tocsep{}\hyperref[book-subsection-section-1.7-2]{Solutions of the Dirac Equation in the Weyl Representation}%
\tocsep{}\hyperref[book-subsection-section-1.7-3]{Dirac Bilinears}%
\end{tocgroup}
\tocsection{book-section-1.8}{1.8}{The Proca Field}
\begin{tocgroup}\hyperref[book-subsection-section-1.8-1]{Lagrangian}%
\tocsep{}\hyperref[book-subsection-section-1.8-2]{Solutions of the Proca Equation}%
\end{tocgroup}
\tocsection{book-section-1.9}{1.9}{The Vector Field}
\begin{tocgroup}\hyperref[book-subsection-section-1.9-1]{Lagrangian}%
\tocsep{}\hyperref[book-subsection-section-1.9-2]{Solutions of the Maxwell Equations}%
\end{tocgroup}
\tocchapter{book-chapter-2}{2}{Operator Quantization}
\tocsection{book-section-2.1}{2.1}{Pictures of Time Evolution}
\begin{tocgroup}\hyperref[book-subsection-section-2.1-1]{Schrödinger Picture}%
\tocsep{}\hyperref[book-subsection-section-2.1-2]{Heisenberg Picture}%
\end{tocgroup}
\tocsection{book-section-2.2}{2.2}{Canonical Quantization}
\begin{tocgroup}\hyperref[book-subsection-section-2.2-1]{Quantization of the Free Real Scalar Field}%
\end{tocgroup}
\tocsection{book-section-2.3}{2.3}{Quantization with Dirac Brackets}
\begin{tocgroup}\hyperref[book-subsection-section-2.3-1]{Quantization of the Complex Scalar Field}%
\tocsep{}\hyperref[book-subsection-section-2.3-2]{Quantization of the Dirac Field}%
\tocsep{}\hyperref[book-subsection-section-2.3-3]{Quantization of the Proca Field}%
\end{tocgroup}
\tocsection{book-section-2.4}{2.4}{Gupta–Bleuler Quantization}
\tocsection{book-section-2.5}{2.5}{Symmetry}
\begin{tocgroup}\hyperref[book-subsection-section-2.5-1]{Ward Identity}%
\end{tocgroup}
\tocchapter{book-chapter-3}{3}{Path Integrals}
\tocsection{book-section-3.1}{3.1}{Path Integrals in Quantum Mechanics}
\tocsection{book-section-3.2}{3.2}{Path Integrals in Quantum Field Theory}
\tocsection{book-section-3.3}{3.3}{$\phi^3$ Theory}
\tocsection{book-section-3.4}{3.4}{Path Integrals for Fermion Fields}
\begin{tocgroup}\hyperref[book-subsection-section-3.4-1]{Grassmann Algebra}%
\tocsep{}\hyperref[book-subsection-section-3.4-2]{Calculus on the Grassmann Algebra}%
\end{tocgroup}
\tocsection{book-section-3.5}{3.5}{Symmetry in Path Integrals}
\tocsection{book-section-3.6}{3.6}{Scattering Theory in Path Integrals}
\tocchapter{book-chapter-4}{4}{Scattering Theory}
\tocsection{book-section-4.1}{4.1}{Interaction Picture}
\tocsection{book-section-4.2}{4.2}{The S-Matrix and Wick's Theorem}
\tocsection{book-section-4.3}{4.3}{Feynman Diagrams and Feynman Rules}
\tocsection{book-section-4.4}{4.4}{Correlation Functions}
\begin{tocgroup}\hyperref[book-subsection-section-4.4-1]{Perturbative Vacuum Correlation Functions}%
\tocsep{}\hyperref[book-subsection-section-4.4-2]{Nonperturbative Vacuum Correlation Functions}%
\end{tocgroup}
\tocsection{book-section-4.5}{4.5}{LSZ Reduction Formula}
\tocsection{book-section-4.6}{4.6}{Cross Sections and Decay Rates}
\tocchapter{book-chapter-5}{5}{The Standard Model}
\tocsection{book-section-5.1}{5.1}{Yang–Mills Theory}
\begin{tocgroup}\hyperref[book-subsection-section-5.1-1]{Lie Group}%
\tocsep{}\hyperref[book-subsection-section-5.1-2]{Local Gauge Symmetry}%
\end{tocgroup}
\tocsection{book-section-5.2}{5.2}{Fermions and Bosons in the Standard Model}
\begin{tocgroup}\hyperref[book-subsection-section-5.2-1]{General Higgs Lagrangian}%
\tocsep{}\hyperref[book-subsection-section-5.2-2]{Fermions and the Higgs Boson in the Standard Model}%
\tocsep{}\hyperref[book-subsection-section-5.2-3]{Yukawa Coupling and Flavor}%
\end{tocgroup}
\tocsection{book-section-5.3}{5.3}{$W^\pm$, $Z^0$, $\gamma$: Electroweak Theory}
\begin{tocgroup}\hyperref[book-subsection-section-5.3-1]{Gauge Symmetry $SU(2)_L\times U(1)_Y$}%
\tocsep{}\hyperref[book-subsection-section-5.3-2]{Interaction}%
\end{tocgroup}
\tocsection{book-section-5.4}{5.4}{Gluon $g$: Quantum Chromodynamics}
\begin{tocgroup}\hyperref[book-subsection-section-5.4-1]{Gauge Symmetry $SU(3)_c$}%
\end{tocgroup}
\tocchapter{book-chapter-6}{6}{New Physics Beyond the Standard Model}
\tocsection{book-section-6.1}{6.1}{Seesaw Mechanism}
\begin{tocgroup}\hyperref[book-subsection-section-6.1-1]{Type-I Seesaw}%
\tocsep{}\hyperref[book-subsection-section-6.1-2]{Type-II Seesaw}%
\tocsep{}\hyperref[book-subsection-section-6.1-3]{Type-III Seesaw}%
\tocsep{}\hyperref[book-subsection-section-6.1-4]{Leptogenesis}%
\end{tocgroup}
\tocsection{book-section-6.2}{6.2}{$U(1)_{L_\mu-L_\tau}$ Gauge Symmetry}
\begin{tocgroup}\hyperref[book-subsection-section-6.2-1]{$U(1)_{L_\mu-L_\tau}$ Gauge Symmetry Transformations}%
\tocsep{}\hyperref[book-subsection-section-6.2-2]{Lagrangian}%
\end{tocgroup}
\tocchapter{book-chapter-7}{7}{Renormalization}
\tocsection{book-section-7.1}{7.1}{Forest Formula}
\begin{tocgroup}\hyperref[book-subsection-section-7.1-1]{Basic Definitions of Graphs}%
\tocsep{}\hyperref[book-subsection-section-7.1-2]{$\mathcal{L}_{CT}$ and R-Operation Recursion Formula}%
\end{tocgroup}
\tocsection{book-section-7.2}{7.2}{QED Renormalization}
\tocchapter{book-chapter-8}{8}{Cosmology}
\tocsection{book-section-8.1}{8.1}{Geometry of the Universe}
\begin{tocgroup}\hyperref[book-subsection-section-8.1-1]{Experimental Facts}%
\tocsep{}\hyperref[book-subsection-section-8.1-2]{FLRW Metric}%
\end{tocgroup}
\tocsection{book-section-8.2}{8.2}{Evolution of the Universe}
\begin{tocgroup}\hyperref[book-subsection-section-8.2-1]{Dynamics}%
\tocsep{}\hyperref[book-subsection-section-8.2-2]{Dark Matter}%
\tocsep{}\hyperref[book-subsection-section-8.2-3]{Dark Energy}%
\tocsep{}\hyperref[book-subsection-section-8.2-4]{Cosmic Evolution}%
\end{tocgroup}
\tocsection{book-section-8.3}{8.3}{Boltzmann Equation}
\tocchapter{book-chapter-A}{A}{Summary of Feynman Rules}
\tocsection{book-section-A.1}{A.1}{Feynman Diagrams of the Standard Model}
\tocchapter{book-chapter-B}{B}{Mathematical Foundations}
\tocsection{book-section-B.1}{B.1}{Fourier Transform}
\tocsection{book-section-B.2}{B.2}{$\gamma$ Matrices: Weyl Representation}
\tocsection{book-section-B.3}{B.3}{IO(1,3) Group}
\tocsection{book-section-B.4}{B.4}{Green's Functions}

\setcounter{chapter}{0}
\chapter{Classical Field Theory}\label{book-chapter-1}\label{qft-ch1}

\setcounter{section}{0}
\section{Lagrangian Formalism}\label{book-section-1.1}\label{lagrangian-formalism}

\subsection{Lagrangian}\label{book-subsection-section-1.1-1}\label{mechanics-lagrangian}

\topic{Definition}

Consider a system with $n$ degrees of freedom.
\begin{itemize}


\item \textbf{\boldmath Configuration space}: an open set $\mathrm{Conf}\subset\mathbb R^n$.
\begin{itemize}


\item A configuration is a point $q=(q_1,\ldots,q_n)\in\mathrm{Conf}$.

\item The $q_a$ are the generalized coordinates.
\end{itemize}


\item \textbf{\boldmath Paths}: smooth curves, for any $t_0<t_f$,
\begin{equation*}
\begin{aligned}
C:[t_0,t_f]&\to\mathrm{Conf},\\
t&\mapsto C(t),
\end{aligned}
\tag{1.1}\label{eq:1.1}
\end{equation*}
with velocity $\dot C(t):=dC/dt$.

\item \textbf{\boldmath Dynamical structure}: the Lagrangian, a smooth function
\begin{equation*}
\begin{aligned}
L:\mathbb R^n\times\mathbb R^n\times\mathbb R&\to\mathbb R,\\
(x,y,t)&\mapsto L(x,y,t).
\end{aligned}
\tag{1.2}\label{eq:1.2}
\end{equation*}

\item \textbf{\boldmath Action}: the functional on paths
\begin{equation*}
{\displaystyle 
\begin{aligned}
S:\{\text{paths}\}&\to\mathbb R,\\
C&\mapsto S[C]:=\int_{t_0}^{t_f}dt\,L\bigl(C(t),\dot C(t),t\bigr).
\end{aligned}}
\tag{1.3}\label{eq:1.3}
\end{equation*}

\item \textbf{\boldmath Solutions}: paths $\bar C$ at which $S$ is stationary under variations with fixed endpoints,
\begin{equation*}
\frac{d}{ds}S[\bar C+s\eta]\Big|_{s=0}=0,
\tag{1.4}\label{eq:1.4}
\end{equation*}
for every smooth $\eta:[t_0,t_f]\to\mathbb R^n$ with $\eta(t_0)=\eta(t_f)=0$.
\end{itemize}


\topic{Euler–Lagrange Equations}

Write $\frac{\partial L}{\partial x_a}\big|_C$ and $\frac{\partial L}{\partial y_a}\big|_C$ for the partial derivatives of $L$ evaluated at $\bigl(C(t),\dot C(t),t\bigr)$.

The correspondence with the usual notation $q_a,\dot q_a$ is
\begin{equation*}
\begin{aligned}
\frac{\partial L}{\partial \dot q_a}&=\frac{\partial L}{\partial y_a}\Big|_C,\\
\frac{\partial L}{\partial q_a}&=\frac{\partial L}{\partial x_a}\Big|_C.
\end{aligned}
\tag{1.5}\label{eq:1.5}
\end{equation*}
Then
\begin{equation*}
\begin{aligned}
\frac{d}{ds}S[C+s\eta]\Big|_{s=0}
&=\int_{t_0}^{t_f}dt\sum_{a=1}^n\Bigl(\frac{\partial L}{\partial x_a}\Big|_C\,\eta_a+\frac{\partial L}{\partial y_a}\Big|_C\,\dot\eta_a\Bigr)\\
&=\int_{t_0}^{t_f}dt\sum_{a=1}^n\biggl[\frac{\partial L}{\partial x_a}\Big|_C-\frac{d}{dt}\Bigl(\frac{\partial L}{\partial y_a}\Big|_C\Bigr)\biggr]\eta_a\\
&\quad+\biggl[\sum_{a=1}^n\frac{\partial L}{\partial y_a}\Big|_C\,\eta_a\biggr]_{t_0}^{t_f}.
\end{aligned}
\tag{1.6}\label{eq:1.6}
\end{equation*}
The boundary term vanishes because $\eta(t_0)=\eta(t_f)=0$.

Since $\eta$ is arbitrary, the solutions are the paths $\bar C$ with
\begin{equation*}
{\displaystyle \frac{d}{dt}\Bigl(\frac{\partial L}{\partial y_a}\Big|_{\bar C}\Bigr)-\frac{\partial L}{\partial x_a}\Big|_{\bar C}=0,\qquad a=1,\ldots,n.}
\tag{1.7}\label{eq:1.7}
\end{equation*}

\subsection{Symmetry Transformations and Noether's Theorem}\label{book-subsection-section-1.1-2}\label{lagrangian-symmetry}

\topic{Transformations}
\begin{itemize}


\item \textbf{\boldmath Internal transformation}: 

a one-parameter family of diffeomorphisms $\phi_\lambda:\mathrm{Conf}\to\mathrm{Conf}$, $\lambda\in\mathbb R$, with $(\lambda,q)\mapsto\phi_\lambda(q)$ smooth and $\phi_0=\mathrm{id}$, acting on paths pointwise,
\begin{equation*}
\begin{aligned}
\Phi:\mathbb R\times\{\text{paths}\}&\to\{\text{paths}\},\\
(\lambda,C)&\mapsto\Phi_\lambda C:=\phi_\lambda\circ C,
\end{aligned}
\tag{1.8}\label{eq:1.8}
\end{equation*}
The tangent vectors, or velocities, are related by
\begin{equation*}
\frac{d}{dt}(\Phi_\lambda C)_a(t)=\sum_{b=1}^n\frac{\partial(\phi_\lambda)_a}{\partial q_b}\bigl(C(t)\bigr)\,\dot C_b(t),
\tag{1.9}\label{eq:1.9}
\end{equation*}
where $(\phi_\lambda)_a(q)$ is the $a$-th coordinate of $\phi_\lambda(q)$.

\textbf{\boldmath Proof.} The main ingredient is the theorem that, under a smooth map, the image of the tangent vector of a curve is the tangent vector of the image curve.

Write $\frac{\partial\phi_\lambda}{\partial q}(q)$ for the Jacobian matrix with entries $\frac{\partial(\phi_\lambda)_a}{\partial q_b}(q)$, so that the velocity of $\Phi_\lambda C$ is
\[
\frac{d}{dt}(\Phi_\lambda C)(t)=\frac{\partial\phi_\lambda}{\partial q}\bigl(C(t)\bigr)\,\dot C(t)
\]

\item \textbf{\boldmath Time translation}: the shift of paths in time,
\begin{equation*}
\begin{aligned}
\Theta:\mathbb R\times\{\text{paths}\}&\to\{\text{paths}\},\\
(\tau,C)&\mapsto\Theta_\tau C:=C(\cdot-\tau),
\end{aligned}
\tag{1.10}\label{eq:1.10}
\end{equation*}
so that $\Theta_\tau C:[t_0+\tau,t_f+\tau]\to\mathrm{Conf}$, $t\mapsto C(t-\tau)$, with velocity
\begin{equation*}
\frac{d}{dt}(\Theta_\tau C)(t)=\dot C(t-\tau).
\tag{1.11}\label{eq:1.11}
\end{equation*}
\end{itemize}


\topic{Symmetry Transformations}
\begin{itemize}


\item \textbf{\boldmath Symmetry transformation}: an internal transformation $\Phi$, respectively a time translation $\Theta$, under which the action of every path is invariant,
\begin{equation*}
{\displaystyle 
\begin{aligned}
S[\Phi_\lambda C]&=S[C]\quad\text{for every path }C\text{ and every }\lambda\in\mathbb R,\\
S[\Theta_\tau C]&=S[C]\quad\text{for every path }C\text{ and every }\tau\in\mathbb R.
\end{aligned}}
\tag{1.12}\label{eq:1.12}
\end{equation*}
\end{itemize}

Substituting the transformed paths into the action, and shifting $t\to t+\tau$ in the second case,
\begin{equation*}
\begin{aligned}
S[\Phi_\lambda C]&=\int_{t_0}^{t_f}dt\,L\Bigl(\phi_\lambda\bigl(C(t)\bigr),\ \frac{\partial\phi_\lambda}{\partial q}\bigl(C(t)\bigr)\,\dot C(t),\ t\Bigr),\\
S[\Theta_\tau C]&=\int_{t_0+\tau}^{t_f+\tau}dt\,L\bigl(C(t-\tau),\dot C(t-\tau),t\bigr)=\int_{t_0}^{t_f}dt\,L\bigl(C(t),\dot C(t),t+\tau\bigr).
\end{aligned}
\tag{1.13}\label{eq:1.13}
\end{equation*}
Hence symmetry transformations have the following equivalent definition:

$\Phi$, respectively $\Theta$, is a symmetry transformation if and only if, for all $(x,y,t)\in\mathrm{Conf}\times\mathbb R^n\times\mathbb R$ and all $\lambda$, respectively all $\tau$,
\begin{equation*}
{\displaystyle 
\begin{aligned}
L\Bigl(\phi_\lambda(x),\ \frac{\partial\phi_\lambda}{\partial q}(x)\,y,\ t\Bigr)&=L(x,y,t)
\end{aligned}}
\tag{1.14}\label{eq:1.14}
\end{equation*}
\begin{equation*}
{\displaystyle 
L(x,y,t+\tau)=L(x,y,t)\quad(\forall\,\tau\in\mathbb R)
\quad\Longleftrightarrow\quad
\frac{\partial L}{\partial t}(x,y,t)=0.
}
\tag{1.15}\label{eq:1.15}
\end{equation*}

\topic{Conserved Quantities}
\begin{itemize}


\item \textbf{\boldmath Conserved quantity}: a smooth function
\begin{equation*}
\begin{aligned}
\mathcal Q:\mathrm{Conf}\times\mathbb R^n\times\mathbb R&\to\mathbb R,\\
(x,y,t)&\mapsto\mathcal Q(x,y,t),
\end{aligned}
\tag{1.16}\label{eq:1.16}
\end{equation*}
such that $\mathcal Q\bigl(\bar C(t),\dot{\bar C}(t),t\bigr)$ is independent of $t$ for every solution $\bar C$.
\end{itemize}


\topic{Noether's Theorem}

\topic{(A) Time Translation}

If $\Theta$ is a symmetry transformation, i.e. $\partial L/\partial t=0$ on $\mathrm{Conf}\times\mathbb R^n\times\mathbb R$, then the energy
\begin{equation*}
{\displaystyle E(x,y,t):=\sum_{a=1}^n\frac{\partial L}{\partial y_a}(x,y,t)\,y_a-L(x,y,t)}
\tag{1.17}\label{eq:1.17}
\end{equation*}
is a conserved quantity.

\topic{(B) Internal Transformation}

If $\Phi$ is a symmetry transformation, then, with the generator
\begin{equation*}
\begin{aligned}
X:\mathrm{Conf}&\to\mathbb R^n,\\
q&\mapsto\frac{\partial\phi_\lambda(q)}{\partial\lambda}\Big|_{\lambda=0},
\end{aligned}
\tag{1.18}\label{eq:1.18}
\end{equation*}
The generator gives the form of the infinitesimal transformation at $\lambda=0$:
\[
\phi_\lambda(q)=q+\lambda\,X(q)+O(\lambda^2),
\]
i.e. the coordinates transform as $q_a\mapsto q_a+\lambda X_a(q)$.

Correspondingly, the velocity transforms as
\[
y\mapsto y+\lambda\,\frac{\partial X}{\partial q}(q)\,y+O(\lambda^2),
\]
i.e. $y_a\mapsto y_a+\lambda\sum_b\frac{\partial X_a}{\partial q_b}(q)\,y_b$.

the function
\begin{equation*}
{\displaystyle \mathcal Q(x,y,t):=\sum_{a=1}^n\frac{\partial L}{\partial y_a}(x,y,t)\,X_a(x)}
\tag{1.19}\label{eq:1.19}
\end{equation*}
is a conserved quantity.

\topic{Proof}

(A) Let $\bar C$ be a solution. Write $L|_{\bar C}:=L\bigl(\bar C(t),\dot{\bar C}(t),t\bigr)$.

By the chain rule and the Euler–Lagrange equations,
\begin{equation*}
\frac{d}{dt}E\bigl(\bar C(t),\dot{\bar C}(t),t\bigr)
=\frac{d}{dt}\Bigl(\sum_{a=1}^n\frac{\partial L}{\partial y_a}\Big|_{\bar C}\,\dot{\bar C}_a(t)-L\big|_{\bar C}\Bigr)
=-\frac{\partial L}{\partial t}\Big|_{\bar C}
\tag{1.20}\label{eq:1.20}
\end{equation*}
along every solution, which vanishes if $\partial L/\partial t=0$.

(B) Differentiating the first symmetry condition with respect to $\lambda$ at $\lambda=0$,  gives for all $(x,y,t)\in\mathrm{Conf}\times\mathbb R^n\times\mathbb R$
\begin{equation*}
\sum_{a=1}^n\Bigl(\frac{\partial L}{\partial x_a}\,X_a(x)
+\frac{\partial L}{\partial y_a}\sum_{b=1}^n\frac{\partial X_a}{\partial q_b}(x)\,y_b\Bigr)=0,
\tag{1.21}\label{eq:1.21}
\end{equation*}
with the partial derivatives of $L$ evaluated at $(x,y,t)$.

For a solution $\bar C$, the chain rule and the Euler–Lagrange equations give
\begin{equation*}
\frac{d}{dt}\mathcal Q\bigl(\bar C(t),\dot{\bar C}(t),t\bigr)
=\sum_{a=1}^n\Bigl(\frac{\partial L}{\partial x_a}\Big|_{\bar C}\,X_a\bigl(\bar C(t)\bigr)
+\frac{\partial L}{\partial y_a}\Big|_{\bar C}\sum_{b=1}^n\frac{\partial X_a}{\partial q_b}\bigl(\bar C(t)\bigr)\,\dot{\bar C}_b(t)\Bigr),
\tag{1.22}\label{eq:1.22}
\end{equation*}
which is the identity above evaluated at $\bigl(\bar C(t),\dot{\bar C}(t),t\bigr)$, hence zero. $\square$

\topic{Momentum}

Fix $b\in\{1,\ldots,n\}$ and translate the $b$-th coordinate,
\begin{equation*}
\bigl(\phi_\lambda(q)\bigr)_a=q_a+\lambda\,\delta_{ab}.
\tag{1.23}\label{eq:1.23}
\end{equation*}
The Jacobian $\partial\phi_\lambda/\partial q$ is the identity, so the velocity is unchanged, and the symmetry condition reads $L(\phi_\lambda(x),y,t)=L(x,y,t)$ for all $\lambda$, i.e. $\partial L/\partial x_b=0$.

Generator: $X_a(q)=\delta_{ab}$.

Conserved quantity:
\begin{equation*}
\mathcal Q(x,y,t)=\sum_{a=1}^n\frac{\partial L}{\partial y_a}(x,y,t)\,\delta_{ab}=\frac{\partial L}{\partial y_b}(x,y,t),
\tag{1.24}\label{eq:1.24}
\end{equation*}
the momentum $p_b$ conjugate to $q_b$.

\setcounter{section}{1}
\section{Hamiltonian Formalism}\label{book-section-1.2}\label{hamiltonian-formalism}

\subsection{Dynamical Systems}\label{book-subsection-section-1.2-1}\label{dynamical-systems}

\topic{Definition}
\begin{itemize}


\item \textbf{\boldmath State space}: an open set $M\subset\mathbb R^N$.
\begin{itemize}


\item A state is a point $z=(z_1,\ldots,z_N)\in M$.

\item The $z_A$ are the state variables.
\end{itemize}


\item \textbf{\boldmath State transition function}: a smooth function
\begin{equation*}
\begin{aligned}
F:M&\to\mathbb R^N,\\
z&\mapsto F(z).
\end{aligned}
\tag{1.25}\label{eq:1.25}
\end{equation*}

\item \textbf{\boldmath Solutions}: smooth curves $C:I\to M$, with $I\subset\mathbb R$ an open interval, such that
\begin{equation*}
{\displaystyle \dot C_A(t)=F_A\bigl(C(t)\bigr),\qquad A=1,\ldots,N.}
\tag{1.26}\label{eq:1.26}
\end{equation*}
\end{itemize}


\topic{(1) Determinism}

Since $F$ is smooth, for every $t_0\in\mathbb R$ and $z^{(0)}\in M$ there is a unique maximal solution $C$ with $C(t_0)=z^{(0)}$. The state at one time therefore determines the state at every later and every earlier time on its interval, and through each point of $M$ passes exactly one solution curve.

Assume every maximal solution is defined on all of $\mathbb R$. The flow of $F$ is
\begin{equation*}
\begin{aligned}
\Phi_t:M&\to M,\\
z^{(0)}&\mapsto C(t),
\end{aligned}
\tag{1.27}\label{eq:1.27}
\end{equation*}
where $C$ is the solution with $C(0)=z^{(0)}$. It satisfies
\begin{itemize}


\item $\partial_t\Phi_t(z)=F\bigl(\Phi_t(z)\bigr)$ and $\Phi_0=\mathrm{id}$;

\item $\Phi_t\circ\Phi_s=\Phi_{t+s}$, since $F$ does not depend on $t$, so each $\Phi_t$ is a diffeomorphism with inverse $\Phi_{-t}$.
\end{itemize}


\topic{(2) Volume in State Space}

For a region $D\subset M$, write $\mathrm{vol}(D):=\int_Dd^Nz$ and $J_t:=\frac{\partial\Phi_t}{\partial z}$ for the Jacobian matrix of $\Phi_t$. Differentiating $\partial_t\Phi_t=F\circ\Phi_t$ with respect to $z$ and using Jacobi's formula $d\det J=\det J\,\operatorname{tr}(J^{-1}dJ)$,
\begin{equation*}
\begin{aligned}
\partial_tJ_t&=\Bigl(\frac{\partial F}{\partial z}\circ\Phi_t\Bigr)J_t,
\qquad
\partial_t\det J_t=\Bigl(\sum_{A=1}^N\frac{\partial F_A}{\partial z_A}\circ\Phi_t\Bigr)\det J_t,\\
\frac{d}{dt}\mathrm{vol}\bigl(\Phi_t(D)\bigr)
&=\frac{d}{dt}\int_Dd^Nz\,\det J_t
=\int_{\Phi_t(D)}d^Nz\,\sum_{A=1}^N\frac{\partial F_A}{\partial z_A}.
\end{aligned}
\tag{1.28}\label{eq:1.28}
\end{equation*}
Hence $\mathrm{vol}(\Phi_t(D))$ is constant for every $D$ if and only if $\sum_A\partial F_A/\partial z_A=0$.

\subsection{Hamiltonian Formalism}\label{book-subsection-section-1.2-2}\label{hamiltonian-mechanics}

\topic{Definition}

Consider a system with $n$ degrees of freedom.
\begin{itemize}


\item \textbf{\boldmath Phase space}: an open set $M\subset\mathbb R^{2n}$.
\begin{itemize}


\item A state is a point $z=(q_1,\ldots,q_n,p_1,\ldots,p_n)\in M$.

\item The $q_a$ are the generalized coordinates and the $p_a$ the momenta.
\end{itemize}


\item \textbf{\boldmath Observables}: smooth functions
\begin{equation*}
\begin{aligned}
f:M&\to\mathbb R,\\
z&\mapsto f(z).
\end{aligned}
\tag{1.29}\label{eq:1.29}
\end{equation*}

\item \textbf{\boldmath Dynamical structure}:
\begin{itemize}


\item the Poisson bracket $\{\ ,\ \}_{\mathrm{PB}}$ of \eqref{eq:1.87};

\item the Hamiltonian, a smooth function
\begin{equation*}
\begin{aligned}
H:M&\to\mathbb R,\\
z&\mapsto H(z),
\end{aligned}
\tag{1.30}\label{eq:1.30}
\end{equation*}
with no explicit time dependence.
\end{itemize}


\item \textbf{\boldmath Solutions}: smooth curves $C:I\to M$, with $I\subset\mathbb R$ an open interval, such that
\begin{equation*}
\begin{gathered}
\frac{d}{dt}q_a\bigl(C(t)\bigr)=\frac{\partial H}{\partial p_a}\bigl(C(t)\bigr),\quad
\frac{d}{dt}p_a\bigl(C(t)\bigr)=-\frac{\partial H}{\partial q_a}\bigl(C(t)\bigr)
\\[3pt]\Longleftrightarrow\quad
\frac{d}{dt}f\bigl(C(t)\bigr)=\{f,H\}_{\mathrm{PB}}\bigl(C(t)\bigr)\quad\text{for all }f.
\end{gathered}
\tag{1.31}\label{eq:1.31}
\end{equation*}
\end{itemize}


\subsection{Symmetries and Noether's Theorem}\label{book-subsection-section-1.2-3}\label{hamiltonian-symmetry}

\topic{Transformations}
\begin{itemize}


\item \textbf{\boldmath Transformation}: a diffeomorphism with no explicit time dependence,
\begin{equation*}
\begin{aligned}
\phi:M&\to M,\\
z&\mapsto\phi(z),
\end{aligned}
\qquad
Q_a:=q_a\circ\phi,\quad P_a:=p_a\circ\phi,\quad \phi_*f:=f\circ\phi^{-1},
\tag{1.32}\label{eq:1.32}
\end{equation*}
where $Q_a(z),P_a(z)$ are the coordinates of $\phi(z)$, and $\phi_*$ is the pushforward of \eqref{eq:1.64}, with $(\phi_*f)(\phi(z))=f(z)$.
\end{itemize}


\topic{Symmetry Transformations}
\begin{itemize}


\item \textbf{\boldmath Symmetry transformation}: a transformation $\phi$ that preserves both the Poisson bracket and the Hamiltonian,
\begin{equation*}
{\displaystyle 
\begin{aligned}
&\text{(i)}\quad\{f\circ\phi,\,g\circ\phi\}_{\mathrm{PB}}=\{f,g\}_{\mathrm{PB}}\circ\phi\quad\text{for all }f,g,\\
&\text{(ii)}\quad H\circ\phi=H.
\end{aligned}}
\tag{1.33}\label{eq:1.33}
\end{equation*}
\end{itemize}

A transformation satisfying (i) is a canonical transformation. Replacing $f,g$ by $f\circ\phi^{-1},g\circ\phi^{-1}$ shows that (i) is equivalent to $\{\phi_*f,\phi_*g\}_{\mathrm{PB}}=\phi_*\{f,g\}_{\mathrm{PB}}$; (ii) is equivalent to $\phi_*H=H$, i.e. $H(\phi(z))=H(z)$ for all $z\in M$.

Write $z_A$, $A=1,\ldots,2n$, for all phase-space coordinates and $Z_A:=z_A\circ\phi$. By the chain rule,
\begin{equation*}
\begin{aligned}
\{f\circ\phi,\,g\circ\phi\}_{\mathrm{PB}}
&=\sum_{A,B}\Bigl(\frac{\partial f}{\partial z_A}\circ\phi\Bigr)\Bigl(\frac{\partial g}{\partial z_B}\circ\phi\Bigr)\{Z_A,Z_B\}_{\mathrm{PB}},\\
\{f,g\}_{\mathrm{PB}}\circ\phi
&=\sum_{A,B}\Bigl(\frac{\partial f}{\partial z_A}\circ\phi\Bigr)\Bigl(\frac{\partial g}{\partial z_B}\circ\phi\Bigr)\{z_A,z_B\}_{\mathrm{PB}}.
\end{aligned}
\tag{1.34}\label{eq:1.34}
\end{equation*}
The brackets $\{z_A,z_B\}_{\mathrm{PB}}$ are constants, and $f=z_A$, $g=z_B$ pick out a single term, so (i) holds if and only if $\{Z_A,Z_B\}_{\mathrm{PB}}=\{z_A,z_B\}_{\mathrm{PB}}$. By \eqref{eq:1.87}, (i) is equivalent to
\begin{equation*}
{\displaystyle 
\{Q_a,Q_b\}_{\mathrm{PB}}=\{P_a,P_b\}_{\mathrm{PB}}=0,\qquad
\{Q_a,P_b\}_{\mathrm{PB}}=\delta_{ab}.}
\tag{1.35}\label{eq:1.35}
\end{equation*}
A symmetry transformation maps solutions to solutions. Let $C$ be a solution and $C':=\phi\circ C$. Applying \eqref{eq:1.31} to $f\circ\phi$, then (ii) and (i),
\begin{equation*}
\begin{aligned}
\frac{d}{dt}f\bigl(C'(t)\bigr)
&=\{f\circ\phi,\,H\}_{\mathrm{PB}}\bigl(C(t)\bigr)\\
&=\{f\circ\phi,\,H\circ\phi\}_{\mathrm{PB}}\bigl(C(t)\bigr)
=\{f,H\}_{\mathrm{PB}}\bigl(C'(t)\bigr)
\end{aligned}
\tag{1.36}\label{eq:1.36}
\end{equation*}
for every $f$, so $C'$ is a solution by \eqref{eq:1.31}.

\topic{Conserved Quantities}
\begin{itemize}


\item \textbf{\boldmath Conserved quantity}: an observable $f$ such that $f\bigl(C(t)\bigr)$ is independent of $t$ for every solution $C$.
\end{itemize}


\eqref{eq:1.31} is \eqref{eq:1.26} with $N=2n$ and $F_A=\{z_A,H\}_{\mathrm{PB}}$. Since a solution passes through each point of $M$, \eqref{eq:1.31} shows that $f$ is a conserved quantity if and only if
\begin{equation*}
{\displaystyle \{f,H\}_{\mathrm{PB}}=0.}
\tag{1.37}\label{eq:1.37}
\end{equation*}

\topic{Noether's Theorem}

Let $\phi_\lambda$, $\lambda\in\mathbb R$, be a one-parameter subgroup of symmetry transformations:
\begin{itemize}


\item $\phi_\lambda\circ\phi_\mu=\phi_{\lambda+\mu}$ for all $\lambda,\mu\in\mathbb R$; in particular $\phi_0=\mathrm{id}$;

\item $(\lambda,z)\mapsto\phi_\lambda(z)$ is smooth.
\end{itemize}


Define the generator
\begin{equation*}
\begin{aligned}
X:M&\to\mathbb R^{2n},\\
z&\mapsto\frac{\partial\phi_\lambda(z)}{\partial\lambda}\Big|_{\lambda=0},
\end{aligned}
\tag{1.38}\label{eq:1.38}
\end{equation*}
so that $\phi_\lambda$ is the flow of $X$. If $M$ is simply connected, there is a smooth function $G:M\to\mathbb R$, unique up to an additive constant, with
\begin{equation*}
X_A=\{z_A,G\}_{\mathrm{PB}},\qquad A=1,\ldots,2n,
\tag{1.39}\label{eq:1.39}
\end{equation*}
and
\begin{equation*}
{\displaystyle \{G,H\}_{\mathrm{PB}}=0,}
\tag{1.40}\label{eq:1.40}
\end{equation*}
so $G$ is a conserved quantity.

\topic{Proof}

Write $\omega_{AB}:=\{z_A,z_B\}_{\mathrm{PB}}$. By \eqref{eq:1.87} it is constant and invertible, and $\{f,g\}_{\mathrm{PB}}=\sum_{A,B}\omega_{AB}\frac{\partial f}{\partial z_A}\frac{\partial g}{\partial z_B}$. Define $Y$ by $X_A=\sum_B\omega_{AB}Y_B$. Differentiate (i) with $f=z_A$, $g=z_B$, and (ii), with respect to $\lambda$ at $\lambda=0$:
\begin{equation*}
\begin{aligned}
0&=\{X_A,z_B\}_{\mathrm{PB}}+\{z_A,X_B\}_{\mathrm{PB}}
=\sum_{C,D}\omega_{AC}\,\omega_{BD}\Bigl(\frac{\partial Y_D}{\partial z_C}-\frac{\partial Y_C}{\partial z_D}\Bigr),\\
0&=\sum_A\frac{\partial H}{\partial z_A}\,X_A.
\end{aligned}
\tag{1.41}\label{eq:1.41}
\end{equation*}
Since $\omega$ is invertible, $\partial Y_D/\partial z_C=\partial Y_C/\partial z_D$. On a simply connected $M$ this gives $Y_B=\partial G/\partial z_B$ for a smooth $G$, unique up to a constant, i.e. \eqref{eq:1.39}. The second line of \eqref{eq:1.41} is then $\sum_A\frac{\partial H}{\partial z_A}\{z_A,G\}_{\mathrm{PB}}=\{H,G\}_{\mathrm{PB}}=0$, which is \eqref{eq:1.40}. $\square$

\topic{(1) Time Evolution and Energy}

Taking $f=H$ in \eqref{eq:1.37}, $\{H,H\}_{\mathrm{PB}}=0$, so $H$ is a conserved quantity. With $G=H$, \eqref{eq:1.39} gives $X_A=\{z_A,H\}_{\mathrm{PB}}=F_A$, so the subgroup generated by $H$ is the time evolution $\Phi_t$ of \eqref{eq:1.27}.

\topic{(2) Translations and Momentum}

Let $M=\mathbb R^{2n}$ and fix $k\in\mathbb R^n$. Take $\phi_\lambda(q,p)=(q+\lambda k,\,p)$. Then $Q_a=q_a+\lambda k_a$ and $P_a=p_a$ satisfy \eqref{eq:1.35}, and $\phi_\lambda$ is a symmetry transformation if and only if $H(q+\lambda k,p)=H(q,p)$ for all $\lambda$, i.e. $\sum_ak_a\,\partial H/\partial q_a=0$. Here $X=(k,0)$, and \eqref{eq:1.39} holds with
\begin{equation*}
{\displaystyle G=\sum_{a=1}^nk_a\,p_a,}
\tag{1.42}\label{eq:1.42}
\end{equation*}
which is therefore conserved.
\begin{itemize}


\item For $k$ the $b$-th unit vector, the condition is $\partial H/\partial q_b=0$, and $p_b$ is conserved. For example, for $H=\sum_ap_a^2/2m$ every $p_a$ is conserved.

\item For $N$ particles, $q=(\mathbf x_1,\ldots,\mathbf x_N)$ and $p=(\mathbf p_1,\ldots,\mathbf p_N)$, take $k=(\mathbf e,\ldots,\mathbf e)$. The conserved quantity is then $\sum_{A=1}^N\mathbf e\cdot\mathbf p_A$, the component of the total momentum along $\mathbf e$.
\end{itemize}


\setcounter{section}{2}
\section{Definition of Classical Fields}\label{book-section-1.3}\label{classical-field}

\subsection{Definition of Fields}\label{book-subsection-section-1.3-1}\label{field-definition}

\topic{Definition}
\begin{equation*}
\text{Field}:\text{Spacetime}\longrightarrow\text{Vector Space},
\tag{1.43}\label{eq:1.43}
\end{equation*}
\begin{equation*}
{\displaystyle \Psi:\mathbb R^4\longrightarrow\mathbb R^n\ \text{or}\ \mathbb C^n
}
\tag{1.44}\label{eq:1.44}
\end{equation*}
Additional regularity conditions, such as $C^\infty$ smoothness, may be imposed.

The space of fields $\{\Psi\}$ carries a representation of $ISO(1,3)^+$. For $(\Lambda,a)\in ISO(1,3)^+$,
\begin{equation*}
\Psi\xrightarrow{\;(\Lambda,a)\;}U(\Lambda,a)\Psi,
\tag{1.45}\label{eq:1.45}
\end{equation*}
where $U$ is an infinite-dimensional representation that can be constructed from the finite-dimensional representation $D(\Lambda)$:
\begin{equation*}
{\displaystyle 
\begin{aligned}
{}[U(\Lambda,a)\Psi]_\alpha(x)
&:=D_{\alpha\beta}(\Lambda)
[(\Lambda,a)_*\Psi]_\beta(x)\\
&=D_{\alpha\beta}(\Lambda)
\Psi_\beta(\Lambda^{-1}x-\Lambda^{-1}a).
\end{aligned}
}
\tag{1.46}\label{eq:1.46}
\end{equation*}

\subsection{Finite-Dimensional Representations of \texorpdfstring{$SO(1,3)^+$}{SO(1,3)+}}\label{book-subsection-section-1.3-2}\label{lorentz-finite-representations}

$U$ is constructed from $D$ by \eqref{eq:1.46}, and $D$ is recovered from $U$ through $[U(\Lambda,0)\Psi](0)=D(\Lambda)\Psi(0)$. Hence
\begin{equation*}
\begin{aligned}
\Psi_\alpha(x),\ \alpha=1,\ldots,n
&\quad\Longleftrightarrow\quad\text{infinite-dimensional representation }U(\Lambda,a)\\
&\quad\Longleftrightarrow\quad n\text{-dimensional representation }D(\Lambda)\text{ of }SO(1,3)^+.
\end{aligned}
\tag{1.47}\label{eq:1.47}
\end{equation*}

\topic{Lie Algebra Representations}

For $X,Y\in\mathfrak o(1,3)$,
\begin{equation*}
\rho(X):=\frac{d}{dt}D(e^{tX})\Big|_{t=0},\qquad
D(e^{X})=e^{\rho(X)},\qquad
\rho([X,Y])=[\rho(X),\rho(Y)].
\tag{1.48}\label{eq:1.48}
\end{equation*}
$SO(1,3)^+$ is connected, so $D$ is determined by $\rho$. Extend $\rho$ complex-linearly, $\rho(X+iY):=\rho(X)+i\rho(Y)$, to the complexified Lie algebra $\mathfrak o(1,3)_{\mathbb C}$, defined by \eqref{eq:B.32} with $M_4(\mathbb R)$ replaced by $M_4(\mathbb C)$. By \eqref{eq:B.36}, $\mathfrak o(1,3)_{\mathbb C}$ contains every $J^{\mu\nu}$, and $D(e^{-\frac i2\omega_{\mu\nu}J^{\mu\nu}})=e^{-\frac i2\omega_{\mu\nu}\rho(J^{\mu\nu})}$. For the Dirac field, $\rho(J^{\mu\nu})=S^{\mu\nu}$ \eqref{eq:1.106}.

\topic{Decomposition of $\mathfrak o(1,3)_{\mathbb C}$}

From \eqref{eq:B.34},
\begin{equation*}
[J^i,J^j]=\epsilon_{ijk}J^k,\qquad
[J^i,K^j]=\epsilon_{ijk}K^k,\qquad
[K^i,K^j]=-\epsilon_{ijk}J^k.
\tag{1.49}\label{eq:1.49}
\end{equation*}
Define
\begin{equation*}
\begin{aligned}
A^i&:=\frac12\Bigl(\frac12\epsilon_{ijk}J^{jk}+iJ^{0i}\Bigr)=\frac i2\bigl(J^i+iK^i\bigr),\\
B^i&:=\frac12\Bigl(\frac12\epsilon_{ijk}J^{jk}-iJ^{0i}\Bigr)=\frac i2\bigl(J^i-iK^i\bigr),
\end{aligned}
\tag{1.50}\label{eq:1.50}
\end{equation*}
i.e. $A^1=\frac12(J^{23}+iJ^{01})$, $A^2=\frac12(J^{31}+iJ^{02})$, $A^3=\frac12(J^{12}+iJ^{03})$, and $B^i$ is the same with $iJ^{0i}\to-iJ^{0i}$. $A^i$ and $B^i$ are not real, so they lie in $\mathfrak o(1,3)_{\mathbb C}$ but not in $\mathfrak o(1,3)$. By \eqref{eq:1.49},
\begin{equation*}
{\displaystyle [A^i,A^j]=i\epsilon_{ijk}A^k,\qquad
[B^i,B^j]=i\epsilon_{ijk}B^k,\qquad
[A^i,B^j]=0.}
\tag{1.51}\label{eq:1.51}
\end{equation*}
Hence
\begin{equation*}
\mathfrak o(1,3)_{\mathbb C}=\mathfrak g_A\oplus\mathfrak g_B,\qquad
\mathfrak g_A:=\operatorname{span}_{\mathbb C}\{A^1,A^2,A^3\},\quad
\mathfrak g_B:=\operatorname{span}_{\mathbb C}\{B^1,B^2,B^3\}.
\tag{1.52}\label{eq:1.52}
\end{equation*}
The maps $A^i\mapsto\sigma^i/2$ and $B^i\mapsto\sigma^i/2$ identify each of $\mathfrak g_A$ and $\mathfrak g_B$ with the complexification of $\mathfrak{su}(2)\cong\mathfrak o(3)$.

\topic{Irreducible Representations}

Take representations $(\rho_A,V_A)$ of $\mathfrak g_A$ and $(\rho_B,V_B)$ of $\mathfrak g_B$, with $n_A:=\dim V_A$ and $n_B:=\dim V_B$. Their tensor (Kronecker) product $(\rho,V_A\otimes V_B)$ is, for $v_A\in\mathfrak g_A$ and $v_B\in\mathfrak g_B$,
\begin{equation*}
\begin{aligned}
\rho(v_A+v_B)&:=\rho_A(v_A)\otimes I_{n_B}+I_{n_A}\otimes\rho_B(v_B),\\
\rho(v_A+v_B)\bigl(|a\rangle\otimes|b\rangle\bigr)&=\rho_A(v_A)|a\rangle\otimes|b\rangle+|a\rangle\otimes\rho_B(v_B)|b\rangle,\\
e^{\rho(v_A+v_B)}&=e^{\rho_A(v_A)}\otimes e^{\rho_B(v_B)}.
\end{aligned}
\tag{1.53}\label{eq:1.53}
\end{equation*}
The last line holds because the two terms in the first line commute.

The irreducible representations of $\mathfrak{su}(2)$ are $V_j$ with basis $|j,m\rangle$, $m=-j,\ldots,j$, on which $\rho_A(A^3)$ (or $\rho_B(B^3)$) acts as $m$. Here $j=0,\frac12,1,\ldots$ and $\dim V_j=2j+1$. Every finite-dimensional irreducible representation of $\mathfrak g_A\oplus\mathfrak g_B$ is $V_{j_1}\otimes V_{j_2}$, and every finite-dimensional representation is a direct sum of these:
\begin{equation*}
{\displaystyle (j_1,j_2),\qquad j_1,j_2\in\{0,\tfrac12,1,\tfrac32,\ldots\},\qquad\dim=(2j_1+1)(2j_2+1).}
\tag{1.54}\label{eq:1.54}
\end{equation*}

\topic{Rotation by $2\pi$}

Since $A^3+B^3=iJ^3$ by \eqref{eq:1.50}, on $(j_1,j_2)$
\begin{equation*}
e^{2\pi J^3}=I_4,\qquad
e^{2\pi\rho(J^3)}=e^{-2\pi i\rho(A^3+B^3)}=(-1)^{2(j_1+j_2)}I_n.
\tag{1.55}\label{eq:1.55}
\end{equation*}
So $D(e^X):=e^{\rho(X)}$ is a representation of $SO(1,3)^+$ only for integer $j_1+j_2$. For half-integer $j_1+j_2$ it is defined only up to sign; for the Dirac field, $\Lambda_{\frac12}=-I_4$ at $\theta^3=2\pi$.

\setcounter{section}{3}
\section{Dynamics of Classical Fields}\label{book-section-1.4}\label{classical-dynamics}

\subsection{Lagrangian}\label{book-subsection-section-1.4-1}\label{field-theory-lagrangian}

Having defined classical fields, we now turn to their dynamics. In the Lagrangian formalism, this amounts to constructing a Lagrangian density for the fields.
In classical mechanics, the dynamical variables are the generalized coordinates $q_a(t)$, which depend only on time. In field theory, the dynamical variables are the field components $\Psi_\alpha(x)$, which depend on the spacetime point $x=(t,\mathbf x)$.

There are two ways to interpret $x$:
\begin{enumerate}


\item At fixed time $t$, the pair $(\alpha,\mathbf x)$ plays the role of the label $a$ in $q_a(t)$, so $\Psi_\alpha(t,\mathbf x)$ corresponds to $q_a(t)$.

\item More physically, the whole spacetime point $x$ plays the role of $t$ in classical mechanics, so $\Psi_\alpha(x)$ corresponds to $q_a(t)$.
\end{enumerate}


\topic{Definition}
\begin{itemize}


\item \textbf{\boldmath Configuration space}: the fields $\Psi$, with $\Psi(x)\in\mathbb R^n$ or $\mathbb C^n$ at each $x$.

\item \textbf{\boldmath Dynamical structure}: the Lagrangian density, a smooth function
\begin{equation*}
\begin{aligned}
\mathcal L:(\mathbb R^n\ \text{or}\ \mathbb C^n)\times(\mathbb R^n\ \text{or}\ \mathbb C^n)^4&\to\mathbb R,\\
(\Psi_\alpha,\partial_\mu\Psi_\alpha)&\mapsto\mathcal L(\Psi_\alpha,\partial_\mu\Psi_\alpha),
\end{aligned}
\tag{1.56}\label{eq:1.56}
\end{equation*}
in which $\partial_\mu\Psi_\alpha$ are independent variables, with no explicit dependence on $x$.

\item \textbf{\boldmath Action}: the functional on fields, for any bounded region $\Omega\subset\mathbb R^4$,
\begin{equation*}
{\displaystyle S[\Psi]:=\int_\Omega d^4x\,\mathcal L\bigl(\Psi_\alpha(x),\partial_\mu\Psi_\alpha(x)\bigr).}
\tag{1.57}\label{eq:1.57}
\end{equation*}

\item \textbf{\boldmath Solutions}: fields $\Psi$ at which $S$ is stationary under variations vanishing on $\partial\Omega$,
\begin{equation*}
\frac{d}{ds}S[\Psi+s\eta]\Big|_{s=0}=0,
\tag{1.58}\label{eq:1.58}
\end{equation*}
for every smooth $\eta$ with $\eta|_{\partial\Omega}=0$ and every $\Omega$.
\end{itemize}


\topic{Conditions on the Lagrangian Density}

For the local relativistic field theories considered here, we require:
\begin{itemize}


\item \textbf{\boldmath Lorentz invariance:} the action constructed from $\mathcal L$ is a Lorentz scalar.

\item \textbf{\boldmath Reality:} the Lagrangian density $\mathcal L$ is real.

\item \textbf{\boldmath Locality:} $\mathcal L$ depends only on the fields and finitely many of their derivatives at the same spacetime point.
\end{itemize}


Further mathematical restrictions may be imposed as conventional choices rather than universal axioms; ultimately, the essential requirement is agreement with physical reality, so that the Lagrangian describes or explains at least part of the real world.

\topic{Euler–Lagrange Equations}

As in \eqref{eq:1.6}, integrating by parts in \eqref{eq:1.58} and using $\eta|_{\partial\Omega}=0$, $\Psi$ is a solution if and only if
\begin{equation*}
{\displaystyle \partial_\mu\frac{\partial\mathcal L}{\partial(\partial_\mu\Psi_\alpha)}
-\frac{\partial\mathcal L}{\partial\Psi_\alpha}=0.}
\tag{1.59}\label{eq:1.59}
\end{equation*}

\topic{Hamiltonian Formulation}

For a regular Legendre transformation,
\begin{equation*}
\begin{aligned}
\pi_\alpha&:=\frac{\partial\mathcal L}{\partial\dot\Psi_\alpha},
&
\mathcal H&:=\pi_\alpha\dot\Psi_\alpha-\mathcal L,
&
H&=\int d^3x\,\mathcal H,\\
\dot\Psi_\alpha(t,\mathbf x)&=\frac{\delta H}{\delta\pi_\alpha(t,\mathbf x)},
&
\dot\pi_\alpha(t,\mathbf x)&=-\frac{\delta H}{\delta\Psi_\alpha(t,\mathbf x)}.
\end{aligned}
\tag{1.60}\label{eq:1.60}
\end{equation*}

\subsection{Symmetry Transformations and Noether's Theorem}\label{book-subsection-section-1.4-2}\label{classical-noether}

\topic{Transformations}
\begin{itemize}


\item \textbf{\boldmath Transformation}: a one-parameter family $\phi_\lambda$ with $\phi_0=\mathrm{id}$, acting on spacetime and on the field,
\begin{equation*}
\begin{aligned}
\phi_\lambda:\mathbb R^4&\to\mathbb R^4,\\
x&\mapsto x'=\phi_\lambda(x),
\end{aligned}
\qquad
x'^\mu=:x^\mu+\delta x^\mu,\qquad
\Psi\mapsto\Psi',
\tag{1.61}\label{eq:1.61}
\end{equation*}
where the form of $\Psi'$ is given in \eqref{eq:1.64} and \eqref{eq:1.72}: for $\phi_\lambda=(\Lambda,a)$ a Poincaré transformation, $\Psi'=U(\Lambda,a)\Psi$ of \eqref{eq:1.46}; for an internal transformation, $x'=x$ and $\Psi'=g(\lambda)\Psi$.
\end{itemize}


\topic{Symmetry Transformations}
\begin{itemize}


\item \textbf{\boldmath Symmetry transformation}: a transformation under which the action of every field is invariant for every region,
\begin{equation*}
{\displaystyle S'=S\quad\text{for every }\Psi\text{ and every }\Omega,}
\tag{1.62}\label{eq:1.62}
\end{equation*}
with $S'$ the action of $\Psi'$ over $\phi_\lambda(\Omega)$, written out in \eqref{eq:1.66}.
\end{itemize}


\topic{Conserved Quantities}
\begin{itemize}


\item \textbf{\boldmath Conserved quantity}: a functional $\mathcal Q[\Psi,t]$ of the field and its derivatives on the time slice $t$, such that $\mathcal Q[\Psi,t]$ is independent of $t$ for every solution $\Psi$ that, together with its first derivatives, falls off fast enough at spatial infinity.
\end{itemize}


\topic{Noether's Theorem}

If $\phi_\lambda$ is a symmetry transformation, then
\begin{equation*}
{\displaystyle 
\mathcal Q[\Psi,t]:=\int d^3x\left(
\frac{\partial\mathcal L}{\partial\dot\Psi_\alpha}
\frac{\partial\,\delta x^\nu}{\partial\lambda}
\partial_\nu\Psi_\alpha
-\frac{\partial\mathcal L}{\partial\dot\Psi_\alpha}
\frac{\partial\,\delta_0\Psi_\alpha}{\partial\lambda}
-\mathcal L\frac{\partial\,\delta x^0}{\partial\lambda}
\right)
}
\tag{1.63}\label{eq:1.63}
\end{equation*}
is a conserved quantity, where $\delta_0\Psi_\alpha$ is defined in \eqref{eq:1.64}, the derivatives with respect to $\lambda$ are taken at $\lambda=0$, and the fall-off condition makes the integral converge and the surface term in \eqref{eq:1.70} vanish.

\topic{Proof}

\topic{(A) Coordinate Transformation}

Let $x'=\phi_\lambda(x)$, with $\lambda$ infinitesimal. Define two field variations:
\begin{equation*}
\begin{aligned}
\delta_0\Psi_\alpha(x)
&:=\Psi'_\alpha(x')-\Psi_\alpha(x)
=D_{\alpha\beta}(\Lambda)\Psi_\beta(x)-\Psi_\alpha(x),\\
\delta\Psi_\alpha(x)
&:=\Psi'_\alpha(x)-\Psi_\alpha(x)
=[U(\Lambda,a)\Psi]_\alpha(x)-\Psi_\alpha(x).
\end{aligned}
\tag{1.64}\label{eq:1.64}
\end{equation*}
Here $\Psi'$ is the transformed field. Using the pushforward $(\phi_\lambda)_*\Psi:=\Psi\circ\phi_\lambda^{-1}$, we write $\Psi'_\alpha=D_{\alpha\beta}(\Lambda)\,[(\phi_\lambda)_*\Psi]_\beta$. Evaluating at $x'$ gives $\Psi'_\alpha(x')=D_{\alpha\beta}(\Lambda)\Psi_\beta(x)$, which is the last equality in the first line, so the first variation, $\delta_0\Psi_\alpha$, describes the transformation of the field itself. For $\phi_\lambda=(\Lambda,a)$, $\Psi'=U(\Lambda,a)\Psi$ by \eqref{eq:1.46}, which is the last equality in the second line. If $\phi_\lambda$ is not a Poincaré transformation, $D(\Lambda)$ is replaced by the identity and $\delta_0\Psi_\alpha=0$.

To first order in $\lambda$, the two variations are related by
\begin{equation*}
\begin{aligned}
\delta\Psi_\alpha(x)
&=\bigl[\Psi'_\alpha(x')-\Psi_\alpha(x)\bigr]
-\bigl[\Psi'_\alpha(x')-\Psi'_\alpha(x)\bigr]\\
&=\delta_0\Psi_\alpha(x)-\delta x^\mu\partial_\mu\Psi'_\alpha(x)\\
&=\delta_0\Psi_\alpha(x)-\delta x^\mu\partial_\mu\Psi_\alpha(x),
\end{aligned}
\tag{1.65}\label{eq:1.65}
\end{equation*}
where the last step uses $\Psi'_\alpha=\Psi_\alpha+O(\lambda)$.

Over a spacetime region $\Omega$, the actions before and after the transformation are
\begin{equation*}
\begin{aligned}
S&=\int_\Omega d^4x\,
\mathcal L\bigl(\Psi_\alpha(x),\partial_\mu\Psi_\alpha(x)\bigr),\\
S'&=\int_{\phi_\lambda(\Omega)}d^4x'\,
\mathcal L\bigl(\Psi'_\alpha(x'),\partial'_\mu\Psi'_\alpha(x')\bigr),
\end{aligned}
\tag{1.66}\label{eq:1.66}
\end{equation*}
where $\partial'_\mu:=\partial/\partial x'^\mu$. The same function $\mathcal L$ is used in $S'$: the transformation acts on the field configuration and on the region $\Omega$, not on the Lagrangian. No property of $\mathcal L$ is assumed here.

Change the integration variable in $S'$ from $x'$ back to $x$, so that $\phi_\lambda(\Omega)$ becomes $\Omega$. To first order in $\lambda$,
\begin{equation*}
\begin{aligned}
d^4x'
&=\det\!\left(\frac{\partial x'^\mu}{\partial x^\nu}\right)d^4x
=\bigl(1+\partial_\mu\delta x^\mu\bigr)d^4x,\\
\partial'_\mu
&=\frac{\partial x^\nu}{\partial x'^\mu}\partial_\nu
=\partial_\mu-(\partial_\mu\delta x^\nu)\partial_\nu,\\
\Psi'_\alpha(x')
&=\Psi_\alpha(x)+\delta_0\Psi_\alpha(x),\\
\partial'_\mu\Psi'_\alpha(x')
&=\partial_\mu\Psi_\alpha(x)
+\partial_\mu\delta_0\Psi_\alpha(x)
-(\partial_\mu\delta x^\nu)\partial_\nu\Psi_\alpha(x),
\end{aligned}
\tag{1.67}\label{eq:1.67}
\end{equation*}
where the third line is \eqref{eq:1.64}. Hence
\begin{equation*}
S'=\int_\Omega d^4x\,
\bigl(1+\partial_\mu\delta x^\mu\bigr)\,
\mathcal L\bigl(
\Psi_\alpha+\delta_0\Psi_\alpha,\
\partial_\mu\Psi_\alpha
+\partial_\mu\delta_0\Psi_\alpha
-(\partial_\mu\delta x^\nu)\partial_\nu\Psi_\alpha
\bigr).
\tag{1.68}\label{eq:1.68}
\end{equation*}
Expanding \eqref{eq:1.68} to first order in $\lambda$ gives the off-shell identity
\begin{equation*}
\begin{aligned}
S'-S
&=\int_\Omega d^4x\,\biggl[
\mathcal L\,\partial_\mu\delta x^\mu
+\frac{\partial\mathcal L}{\partial\Psi_\alpha}\delta_0\Psi_\alpha\\
&\qquad
+\frac{\partial\mathcal L}{\partial(\partial_\mu\Psi_\alpha)}
\bigl(\partial_\mu\delta_0\Psi_\alpha-(\partial_\mu\delta x^\nu)\partial_\nu\Psi_\alpha\bigr)
\biggr]\\
&=\int_\Omega d^4x\,\biggl\{
\biggl[
\frac{\partial\mathcal L}{\partial\Psi_\alpha}
-\partial_\mu\frac{\partial\mathcal L}{\partial(\partial_\mu\Psi_\alpha)}
\biggr]\delta\Psi_\alpha\\
&\qquad
+\partial_\mu\biggl[
\mathcal L\,\delta x^\mu
+\frac{\partial\mathcal L}{\partial(\partial_\mu\Psi_\alpha)}\delta\Psi_\alpha
\biggr]
\biggr\}.
\end{aligned}
\tag{1.69}\label{eq:1.69}
\end{equation*}
The second equality substitutes $\delta_0\Psi_\alpha=\delta\Psi_\alpha+\delta x^\nu\partial_\nu\Psi_\alpha$ from \eqref{eq:1.65} and uses the product rule together with $\partial_\nu\mathcal L=\frac{\partial\mathcal L}{\partial\Psi_\alpha}\partial_\nu\Psi_\alpha+\frac{\partial\mathcal L}{\partial(\partial_\mu\Psi_\alpha)}\partial_\nu\partial_\mu\Psi_\alpha$; the latter holds because $\mathcal L$ has no explicit $x$-dependence.

By hypothesis, the action is invariant: $S'=S$. Take $\Omega=[t_1,t_2]\times V$ and let $\Psi_\alpha$ satisfy the Euler–Lagrange equations, so that the Euler–Lagrange term in \eqref{eq:1.69} vanishes. Gauss's theorem then gives
\begin{equation*}
\begin{aligned}
0
&=\int_\Omega d^4x\,\partial_\mu\biggl[
\mathcal L\,\delta x^\mu
+\frac{\partial\mathcal L}{\partial(\partial_\mu\Psi_\alpha)}\delta\Psi_\alpha
\biggr]\\
&=\biggl[\int_V d^3x\,\biggl(
\mathcal L\,\delta x^0
+\frac{\partial\mathcal L}{\partial\dot\Psi_\alpha}\delta\Psi_\alpha
\biggr)\biggr]_{t_1}^{t_2}\\
&\quad
+\int_{t_1}^{t_2}dt\oint_{\partial V}dS_i\,\biggl(
\mathcal L\,\delta x^i
+\frac{\partial\mathcal L}{\partial(\partial_i\Psi_\alpha)}\delta\Psi_\alpha
\biggr).
\end{aligned}
\tag{1.70}\label{eq:1.70}
\end{equation*}
By the fall-off condition, the surface term vanishes as $V\to\mathbb R^3$. Since $t_1$ and $t_2$ are arbitrary,
\begin{equation*}
\int d^3x\,\biggl(
\mathcal L\,\delta x^0
+\frac{\partial\mathcal L}{\partial\dot\Psi_\alpha}\delta\Psi_\alpha
\biggr)=\text{const}.
\tag{1.71}\label{eq:1.71}
\end{equation*}
Substituting $\delta\Psi_\alpha=\delta_0\Psi_\alpha-\delta x^\nu\partial_\nu\Psi_\alpha$ from \eqref{eq:1.65}, differentiating with respect to $\lambda$, and reversing the overall sign gives \eqref{eq:1.63} for coordinate transformations.

\topic{(B) Internal Transformation}

Now $\phi_\lambda$ acts trivially on spacetime, and on the field by a constant matrix $g(\lambda)$ with $g(0)=I_n$:
\begin{equation*}
x'^\mu=x^\mu,\qquad
\Psi'_\alpha(x)=g_{\alpha\beta}(\lambda)\,\Psi_\beta(x).
\tag{1.72}\label{eq:1.72}
\end{equation*}
Hence $\delta x^\mu=0$, and the two variations in \eqref{eq:1.64} coincide: $\delta\Psi_\alpha(x)=\delta_0\Psi_\alpha(x)=g_{\alpha\beta}(\lambda)\Psi_\beta(x)-\Psi_\alpha(x)$. Repeating \eqref{eq:1.66}–\eqref{eq:1.71} with $\delta x^\mu=0$: in \eqref{eq:1.67}, $d^4x'=d^4x$ and $\partial'_\mu=\partial_\mu$; the terms containing $\delta x^\mu$ drop out of \eqref{eq:1.68}–\eqref{eq:1.70}; and \eqref{eq:1.71} becomes
\begin{equation*}
\int d^3x\,\frac{\partial\mathcal L}{\partial\dot\Psi_\alpha}\,\delta_0\Psi_\alpha=\text{const}.
\tag{1.73}\label{eq:1.73}
\end{equation*}
Differentiating with respect to $\lambda$ and reversing the overall sign gives \eqref{eq:1.63} with $\delta x^\mu=0$. $\square$

\phantomsection\label{spacetime-conserved-quantities}%
\topic{(1) Spacetime Translations and Energy–Momentum}

Take $\phi_\lambda(x)=x+\lambda\epsilon$ for a fixed $\epsilon\in\mathbb R^4$, so that $\delta x^\mu=\lambda\epsilon^\mu$, $\Lambda=I_4$ and $\delta_0\Psi_\alpha=0$. By \eqref{eq:1.66}–\eqref{eq:1.68}, $S'=S$ holds because $\mathcal L$ has no explicit $x$-dependence, and \eqref{eq:1.63} gives
\begin{equation*}
\mathcal Q=\epsilon^\nu\int d^3x\Bigl(\frac{\partial\mathcal L}{\partial\dot\Psi_\alpha}\partial_\nu\Psi_\alpha-\delta^0{}_\nu\mathcal L\Bigr)=\epsilon^\nu P_\nu,
\tag{1.74}\label{eq:1.74}
\end{equation*}
with the canonical energy–momentum tensor and four-momentum
\begin{equation*}
{\displaystyle 
\begin{aligned}
T^\mu{}_\nu&:=\frac{\partial\mathcal L}{\partial(\partial_\mu\Psi_\alpha)}
\,\partial_\nu\Psi_\alpha-\delta^\mu{}_\nu\mathcal L,
&
\partial_\mu T^\mu{}_\nu&=0,\\
P_\nu&:=\int d^3x\,T^0{}_\nu,
&
\frac{dP_\nu}{dt}&=0,
&
P_0&=H.
\end{aligned}}
\tag{1.75}\label{eq:1.75}
\end{equation*}
Here $\partial_\mu T^\mu{}_\nu=0$ holds on shell by \eqref{eq:1.59}, and $P_0=H$ by \eqref{eq:1.60}.

\setcounter{section}{4}
\section{Dirac Constraint}\label{book-section-1.5}\label{dirac-constraint}

\subsection{Singular Lagrangians}\label{book-subsection-section-1.5-1}\label{singular-lagrangian}

\topic{Definition}

Let $L:\mathbb R^n\times\mathbb R^n\to\mathbb R$ be a smooth Lagrangian as in \eqref{eq:1.2}, without explicit time dependence.
\begin{itemize}


\item \textbf{\boldmath Canonical momenta}:
\begin{equation*}
p_a:=\frac{\partial L}{\partial\dot q_a}(q,\dot q),\qquad a=1,\ldots,n.
\tag{1.76}\label{eq:1.76}
\end{equation*}

\item \textbf{\boldmath Hessian}: the symmetric matrix
\begin{equation*}
W_{ab}:=\frac{\partial^2L}{\partial\dot q_a\,\partial\dot q_b}(q,\dot q),\qquad a,b=1,\ldots,n.
\tag{1.77}\label{eq:1.77}
\end{equation*}

\item \textbf{\boldmath Regular and singular Lagrangians}: $L$ is regular if $\det W\neq0$ on $\mathrm{Conf}\times\mathbb R^n$, and singular otherwise. For a singular $L$ we assume that $\operatorname{rank}W=r<n$ is constant.
\end{itemize}


If $L$ is regular, \eqref{eq:1.76} can be solved for the velocities, $\dot q=\dot q(q,p)$, and
\begin{equation*}
H(q,p):=\sum_{a=1}^np_a\,\dot q_a(q,p)-L\bigl(q,\dot q(q,p)\bigr)
\tag{1.78}\label{eq:1.78}
\end{equation*}
is the Hamiltonian \eqref{eq:1.30} of \hyperref[hamiltonian-formalism]{Section 1.2}. In the rest of this section $L$ is singular.

\topic{Intuitive Physical Picture of Singular Lagrangians}

Take as an example the case of a constant matrix $W$, $L=\frac12\sum_{a,b=1}^nW_{ab}\dot q_a\dot q_b-V(q)$. Then $W$ is the mass matrix, and a $W$ that is not of full rank can be understood intuitively as a vanishing "mass" for motion along some directions. A suitable redefinition of the generalized coordinates brings $W$ to block-diagonal form: the upper-left $r\times r$ block is invertible (and may be taken diagonal), and all other entries vanish, so the kinetic term contains no velocities of the last $n-r$ coordinates. There are two cases, according to whether the potential depends on these $n-r$ coordinates.
\begin{enumerate}


\item In the new coordinates the potential contains only the first $r$ coordinates. The last $n-r$ coordinates do not appear in $L$ at all, and the equations of motion place no restriction on them: the original mathematical description is redundant.

\item In the new coordinates the potential still contains the last $n-r$ coordinates. Their Euler–Lagrange equations contain no accelerations and are only relations among the coordinates: the system is subject to constraints.
\end{enumerate}


\topic{Primary Constraints}

Since $W$ is symmetric, its rank equals the order of its largest invertible principal submatrix, so after reordering the coordinates
\begin{equation*}
\det\bigl(W_{ab}\bigr)_{a,b=1}^{r}\neq0.
\tag{1.79}\label{eq:1.79}
\end{equation*}
Indices $\mu,\nu$ run over the remaining values $r+1,\ldots,n$. By the implicit function theorem, the first $r$ equations of \eqref{eq:1.76} can be solved locally for $\dot q_1,\ldots,\dot q_r$:
\begin{equation*}
\dot q_a=\varphi_a(q,p_b,\dot q_\mu),\qquad a=1,\ldots,r,
\tag{1.80}\label{eq:1.80}
\end{equation*}
where $p_b$ and $\dot q_\mu$ stand for $(p_1,\ldots,p_r)$ and $(\dot q_{r+1},\ldots,\dot q_n)$.

Substituting \eqref{eq:1.80} into the remaining equations of \eqref{eq:1.76} gives
\begin{equation*}
p_\mu=\frac{\partial L}{\partial\dot q_\mu}\bigl(q,\varphi(q,p_b,\dot q_\nu),\dot q_\nu\bigr)=:g_\mu(q,p_b,\dot q_\nu),
\tag{1.81}\label{eq:1.81}
\end{equation*}
and one can further show that $g_\mu$ does not depend on $\dot q_\nu$, as follows.

Differentiating the identity $p_b=\frac{\partial L}{\partial\dot q_b}(q,\varphi,\dot q_\nu)$ with respect to $\dot q_\nu$ gives $\sum_{a=1}^rW_{ba}\frac{\partial\varphi_a}{\partial\dot q_\nu}=-W_{b\nu}$, hence
\begin{equation*}
\frac{\partial}{\partial\dot q_\nu}\frac{\partial L}{\partial\dot q_\mu}\bigl(q,\varphi,\dot q\bigr)
=W_{\mu\nu}+\sum_{a=1}^rW_{\mu a}\frac{\partial\varphi_a}{\partial\dot q_\nu}
=W_{\mu\nu}-\sum_{a,b=1}^rW_{\mu a}\bigl(W^{-1}\bigr)_{ab}W_{b\nu}=0,
\tag{1.82}\label{eq:1.82}
\end{equation*}
with $(W^{-1})_{ab}$ the inverse of the block \eqref{eq:1.79}. The last expression is the Schur complement of that block in $W$, which vanishes since $\operatorname{rank}W=r$.
\begin{itemize}


\item \textbf{\boldmath Primary constraints}:
\begin{equation*}
{\displaystyle \chi_\mu(q,p):=p_\mu-g_\mu(q,p_b)=0,\qquad\mu=r+1,\ldots,n.}
\tag{1.83}\label{eq:1.83}
\end{equation*}
The $n$ equations \eqref{eq:1.76} are equivalent to \eqref{eq:1.80} together with \eqref{eq:1.83}. The constraints \eqref{eq:1.83} involve only $(q,p)$ and define the constraint surface in phase space, the image of the map $(q,\dot q)\mapsto(q,p(q,\dot q))$.

\item \textbf{\boldmath Weak equality}: $F\approx G$ means that the phase-space functions $F$ and $G$ agree on the constraint surface.
\end{itemize}


\topic{Canonical Hamiltonian}

Define
\begin{equation*}
H_c:=\sum_{a=1}^np_a\dot q_a-L(q,\dot q),
\tag{1.84}\label{eq:1.84}
\end{equation*}
with \eqref{eq:1.80} substituted for $\dot q_1,\ldots,\dot q_r$, so that $H_c$ is a function of $(q,p,\dot q_\mu)$. By \eqref{eq:1.80}, \eqref{eq:1.81} and \eqref{eq:1.83},
\begin{equation*}
\frac{\partial H_c}{\partial\dot q_\nu}
=\sum_{a=1}^r\Bigl(p_a-\frac{\partial L}{\partial\dot q_a}\Bigr)\frac{\partial\varphi_a}{\partial\dot q_\nu}
+p_\nu-\frac{\partial L}{\partial\dot q_\nu}
=\chi_\nu(q,p),
\tag{1.85}\label{eq:1.85}
\end{equation*}
which is independent of $\dot q_\mu$ and vanishes on the constraint surface. Hence $H_c$ is independent of $\dot q_\mu$ on the constraint surface, where by \eqref{eq:1.83} it is a function of $(q,p_1,\ldots,p_r)$ alone.

\topic{Total Hamiltonian and Consistency}

Let $H_c$ be the canonical Hamiltonian and $u^A$ the Lagrange multipliers enforcing the primary constraints.
\begin{equation*}
{\displaystyle 
\begin{aligned}
H_T&=H_c+\sum_A u^A(t)\chi_A^{(1)},\\
\dot\chi_A&=\frac{\partial\chi_A}{\partial t}
+\{\chi_A,H_T\}_{\mathrm{PB}}\approx0.
\end{aligned}}
\tag{1.86}\label{eq:1.86}
\end{equation*}
Consistency may yield secondary and higher constraints, fix multipliers, or hold identically. Iterate until closure.

\topic{First-Class and Second-Class Constraints}
\begin{equation*}
\{F,G\}_{\mathrm{PB}}
:=\sum_{a=1}^n\left(
\frac{\partial F}{\partial q_a}\frac{\partial G}{\partial p_a}
-\frac{\partial F}{\partial p_a}\frac{\partial G}{\partial q_a}
\right),
\qquad
\{q_a,p_b\}_{\mathrm{PB}}=\delta_{ab}.
\tag{1.87}\label{eq:1.87}
\end{equation*}
A constraint is first class if its Poisson bracket with every constraint vanishes weakly. After separating all first-class combinations, the remaining second-class constraints $\chi_a^{\mathrm{sc}}$ have an invertible bracket matrix:
\begin{equation*}
\{\chi^{\mathrm{fc}},\chi_A\}_{\mathrm{PB}}\approx0
\quad\text{for all }A,
\qquad
C_{ab}:=\{\chi_a^{\mathrm{sc}},\chi_b^{\mathrm{sc}}\}_{\mathrm{PB}},
\qquad
\det C\ne0.
\tag{1.88}\label{eq:1.88}
\end{equation*}

\topic{Dirac Bracket}
\begin{equation*}
{\displaystyle 
\begin{aligned}
\{F,G\}_{\mathrm D}
&:=\{F,G\}_{\mathrm{PB}}
-\sum_{a,b}\{F,\chi_a^{\mathrm{sc}}\}_{\mathrm{PB}}
(C^{-1})_{ab}\{\chi_b^{\mathrm{sc}},G\}_{\mathrm{PB}},\\
\{F,\chi_a^{\mathrm{sc}}\}_{\mathrm D}&=0.
\end{aligned}}
\tag{1.89}\label{eq:1.89}
\end{equation*}
Second-class constraints may then be imposed strongly. For fields, use functional derivatives, spatial integrals, and inverse integral kernels.

\topic{Physical Degrees of Freedom}

For a regular system with independent constraints and gauge orbits of the expected dimension,
\begin{equation*}
{\displaystyle 
N_{\mathrm{physical}}
=n-N_{\mathrm{first}}-\frac12N_{\mathrm{second}}.
}
\tag{1.90}\label{eq:1.90}
\end{equation*}
For fields, this counts local configuration degrees of freedom per spatial point, subject to boundary and zero-mode qualifications.

\setcounter{section}{5}
\section{The Scalar Field}\label{book-section-1.6}\label{scalar-field}

\subsection{Lagrangian}\label{book-subsection-section-1.6-1}\label{scalar-lagrangian}

\topic{Real Scalar Field}
\begin{equation*}
{\displaystyle 
\mathcal L=\frac12\partial_\mu\phi\,\partial^\mu\phi-V(\phi),
\qquad
\partial_\mu\partial^\mu\phi+V'(\phi)=0.
}
\qquad
V(\phi)=\frac12m^2\phi^2
\ \Longrightarrow\
(\partial_\mu\partial^\mu+m^2)\phi=0.
\tag{1.91}\label{eq:1.91}
\end{equation*}
\begin{equation*}
\pi=\dot\phi,
\qquad
{\displaystyle 
\mathcal H=\frac12\pi^2+\frac12|\boldsymbol\nabla\phi|^2+V(\phi),
\qquad
H=\int d^3x\,\mathcal H.
}
\tag{1.92}\label{eq:1.92}
\end{equation*}

\topic{Complex Scalar Field}

Treat $\phi$ and $\phi^*$ as independent variables:
\begin{equation*}
{\displaystyle 
\begin{aligned}
\mathcal L&=\partial_\mu\phi^*\partial^\mu\phi-m^2\phi^*\phi,\\
(\partial_\mu\partial^\mu+m^2)\phi&=0,\qquad
(\partial_\mu\partial^\mu+m^2)\phi^*=0,\\
\pi_\phi&=\dot\phi^*,\qquad \pi_{\phi^*}=\dot\phi,\\
\mathcal H&=\pi_\phi\pi_{\phi^*}
+|\boldsymbol\nabla\phi|^2+m^2|\phi|^2.
\end{aligned}}
\tag{1.93}\label{eq:1.93}
\end{equation*}

\topic{Conserved Quantities: Spacetime Symmetries}

Translations give the canonical energy–momentum tensor in \eqref{eq:1.75}:
\begin{equation*}
{\displaystyle \begin{aligned}
T^\mu{}_\nu&=\partial^\mu\phi\,\partial_\nu\phi-\delta^\mu{}_\nu\mathcal L
&&\text{(real field)},\\
T^\mu{}_\nu&=\partial^\mu\phi^*\,\partial_\nu\phi
+\partial^\mu\phi\,\partial_\nu\phi^*-\delta^\mu{}_\nu\mathcal L
&&\text{(complex field)},\\
\partial_\mu T^\mu{}_\nu&=0,\qquad
P_\nu=\int d^3x\,T^0{}_\nu,\qquad \frac{dP_\nu}{dt}=0.
\end{aligned}}
\tag{1.94}\label{eq:1.94}
\end{equation*}
For either field, $T^{\mu\nu}:=T^\mu{}_\rho\eta^{\rho\nu}$ is symmetric. Lorentz symmetry gives the rotation and boost charges
\begin{equation*}
{\displaystyle 
J^{\mu\nu}=\int d^3x\,
\bigl(x^\mu T^{0\nu}-x^\nu T^{0\mu}\bigr),
\qquad
\frac{dJ^{\mu\nu}}{dt}=0.
}
\tag{1.95}\label{eq:1.95}
\end{equation*}
Here $P_0$ is the energy, $P^i$ are the spatial momenta, $J^{ij}$ generate rotations, and $J^{0i}$ generate boosts. These conservation laws hold on shell, with convergent integrals and vanishing boundary fluxes.

\topic{Conserved Quantities: Internal and Gauge Symmetries}

The real scalar field has no $U(1)$ phase symmetry of the complex-field type. The complex Lagrangian \eqref{eq:1.93} has a global $U(1)$ symmetry; it is not invariant under an arbitrary local phase $\lambda(x)$ without a gauge field. For constant $\lambda$, the conserved current and charge are
\begin{equation*}
\phi'=e^{i\lambda}\phi,
\qquad
{\displaystyle 
j^\mu=i\bigl(\phi^*\partial^\mu\phi-\phi\,\partial^\mu\phi^*\bigr),
\qquad
\partial_\mu j^\mu=0,
\qquad
Q=\int d^3x\,j^0.
}
\tag{1.96}\label{eq:1.96}
\end{equation*}

\subsection{Solutions of the Klein–Gordon Equation}\label{book-subsection-section-1.6-2}\label{scalar-field-solutions}

For the real field take $V(\phi)=\tfrac12m^2\phi^2$; the complex field in \eqref{eq:1.93} is already free. Define $\Box:=\partial_\mu\partial^\mu$. A plane wave obeys
\begin{equation*}
(\Box+m^2)e^{\mp ip\cdot x}=(m^2-p^2)e^{\mp ip\cdot x}=0,
\qquad
p^2=m^2,\qquad
p^0=E_{\mathbf p}:=\sqrt{\mathbf p^2+m^2}>0.
\tag{1.97}\label{eq:1.97}
\end{equation*}
Using the Lorentz-invariant measure written explicitly below, a real wave packet has the expansion
\begin{equation*}
{\displaystyle 
\phi(x)=\int\frac{d^3p}{(2\pi)^3\,2E_{\mathbf p}}
\left[a(\mathbf p)e^{-ip\cdot x}+a^*(\mathbf p)e^{+ip\cdot x}\right].
}
\tag{1.98}\label{eq:1.98}
\end{equation*}
For a complex field the two frequency amplitudes are independent:
\begin{equation*}
{\displaystyle \begin{aligned}
\phi(x)&=\int\frac{d^3p}{(2\pi)^3\,2E_{\mathbf p}}
\left[a(\mathbf p)e^{-ip\cdot x}+b^*(\mathbf p)e^{+ip\cdot x}\right],\\
\phi^*(x)&=\int\frac{d^3p}{(2\pi)^3\,2E_{\mathbf p}}
\left[a^*(\mathbf p)e^{+ip\cdot x}+b(\mathbf p)e^{-ip\cdot x}\right].
\end{aligned}}
\tag{1.99}\label{eq:1.99}
\end{equation*}
The coefficients are classical complex amplitudes. The real-field condition is $b=a$. These free-field expansions do not apply to a general nonlinear potential $V(\phi)$.

\setcounter{section}{6}
\section{The Dirac Field \texorpdfstring{$\psi$}{ψ}}\label{book-section-1.7}\label{dirac-field}

\subsection{Lagrangian}\label{book-subsection-section-1.7-1}\label{dirac-lagrangian}

The Lagrangian of the Dirac Field has the following form:
\begin{equation*}
{\displaystyle 
\mathcal L=\bar\psi\bigl(i\gamma^\mu\partial_\mu-m\bigr)\psi,
}
\tag{1.100}\label{eq:1.100}
\end{equation*}
where $\psi:\mathbb R^4\to\mathbb C^4$ is the Dirac field, $\bar\psi:=\psi^\dagger\gamma^0$ is its Dirac adjoint, the $\gamma^\mu$ are given by \eqref{eq:B.10}, and $m$ is the mass.

Treating $\psi$ and $\bar\psi$ as independent, the Euler–Lagrange equations give the Dirac equation and its adjoint:
\begin{equation*}
{\displaystyle \begin{aligned}
\bigl(i\gamma^\mu\partial_\mu-m\bigr)\psi&=0,\\
i(\partial_\mu\bar\psi)\gamma^\mu+m\bar\psi&=0.
\end{aligned}}
\tag{1.101}\label{eq:1.101}
\end{equation*}

\topic{Hamiltonian}

Treating $\psi$ and $\bar\psi$ as independent, their conjugate momenta are
\begin{equation*}
{\displaystyle \begin{aligned}
\pi_\psi&:=\frac{\partial\mathcal L}{\partial\dot\psi}
=i\bar\psi\gamma^0=i\psi^\dagger,\\
\pi_{\bar\psi}&:=\frac{\partial\mathcal L}{\partial\dot{\bar\psi}}=0.
\end{aligned}}
\tag{1.102}\label{eq:1.102}
\end{equation*}

\topic{Conserved Quantities: Spacetime Symmetries}

Translations give the canonical energy–momentum tensor and four-momentum
\begin{equation*}
{\displaystyle \begin{aligned}
T^\mu{}_\nu&=i\bar\psi\gamma^\mu\partial_\nu\psi-\delta^\mu{}_\nu\mathcal L,
&\partial_\mu T^\mu{}_\nu&=0,\\
P_\nu&=\int d^3x\,T^0{}_\nu,
&\frac{dP_\nu}{dt}&=0.
\end{aligned}}
\tag{1.103}\label{eq:1.103}
\end{equation*}
With $T^{\mu\nu}:=T^\mu{}_\rho\eta^{\rho\nu}$, Lorentz symmetry gives
\begin{equation*}
{\displaystyle \begin{aligned}
J^{\mu\nu}&=\int d^3x\left[
x^\mu T^{0\nu}-x^\nu T^{0\mu}
+\frac{i}{4}\psi^\dagger[\gamma^\mu,\gamma^\nu]\psi
\right],\\
\frac{dJ^{\mu\nu}}{dt}&=0.
\end{aligned}}
\tag{1.104}\label{eq:1.104}
\end{equation*}
The last term is the spin contribution. $J^{ij}$ and $J^{0i}$ are the rotation and boost charges. Conservation holds on shell, with convergent integrals and vanishing boundary fluxes.

\topic{Conserved Quantities: Internal and Gauge Symmetries}

The free Dirac Lagrangian has a global $U(1)$ symmetry. For constant $\lambda$,
\begin{equation*}
{\displaystyle \begin{aligned}
\psi'&=e^{i\lambda}\psi,\qquad
\bar\psi'=e^{-i\lambda}\bar\psi,\\
j^\mu&=\bar\psi\gamma^\mu\psi,\qquad \partial_\mu j^\mu=0,\\
Q&=\int d^3x\,j^0=\int d^3x\,\psi^\dagger\psi,
\qquad \frac{dQ}{dt}=0.
\end{aligned}}
\tag{1.105}\label{eq:1.105}
\end{equation*}
For a local phase $\lambda(x)$, $\mathcal L'=\mathcal L-j^\mu\partial_\mu\lambda$; local gauge invariance requires coupling to a gauge field.

\subsection{Solutions of the Dirac Equation in the Weyl Representation}\label{book-subsection-section-1.7-2}\label{weyl-spinor-solutions}

\topic{Spinor Representation of the Lorentz Group}

For the Dirac field, $D(\Lambda)$ in \eqref{eq:1.46} is
\begin{equation*}
\Lambda_{\frac12}=\exp\Bigl(-\frac i2\omega_{\mu\nu}S^{\mu\nu}\Bigr),
\qquad
S^{\mu\nu}:=\frac i4[\gamma^\mu,\gamma^\nu],
\tag{1.106}\label{eq:1.106}
\end{equation*}
with the same $\omega_{\mu\nu}$ as in \eqref{eq:B.35}. In the Weyl representation every $S^{\mu\nu}$ is block diagonal:
\begin{equation*}
S^{0i}=-\frac i2\begin{pmatrix}\sigma^i&0\\0&-\sigma^i\end{pmatrix},
\qquad
S^{ij}=\frac12\epsilon^{ijk}\begin{pmatrix}\sigma^k&0\\0&\sigma^k\end{pmatrix}.
\tag{1.107}\label{eq:1.107}
\end{equation*}
The noncompact Lorentz group $SO(1,3)^+$ has no nontrivial finite-dimensional unitary representations; in the Weyl representation, $(S^{0i})^\dagger=-S^{0i}$, so the finite boost matrices for nonzero rapidity are not unitary.

Hence $\Lambda_{\frac12}$ is block diagonal: writing $\psi=\begin{pmatrix}\psi_L\\\psi_R\end{pmatrix}$, the left-handed and right-handed Weyl spinors $\psi_L,\psi_R$ transform independently. For an infinitesimal transformation, with $\omega_{ij}=\epsilon_{ijk}\theta^k$ and $\omega_{0i}=b^i$ as in \eqref{eq:B.34},
\begin{equation*}
\begin{aligned}
\Lambda_{\frac12}
&\simeq I_4-\frac i2\omega_{\mu\nu}S^{\mu\nu}\\
&=I_4-\frac i2\begin{pmatrix}\boldsymbol\theta\cdot\boldsymbol\sigma&0\\0&\boldsymbol\theta\cdot\boldsymbol\sigma\end{pmatrix}
-i\Bigl(-\frac i2\Bigr)\begin{pmatrix}\mathbf b\cdot\boldsymbol\sigma&0\\0&-\mathbf b\cdot\boldsymbol\sigma\end{pmatrix}\\
&=\begin{pmatrix}I_2-i\boldsymbol\theta\cdot\frac{\boldsymbol\sigma}2-\mathbf b\cdot\frac{\boldsymbol\sigma}2&0\\0&I_2-i\boldsymbol\theta\cdot\frac{\boldsymbol\sigma}2+\mathbf b\cdot\frac{\boldsymbol\sigma}2\end{pmatrix},
\end{aligned}
\tag{1.108}\label{eq:1.108}
\end{equation*}
so that
\begin{equation*}
\begin{aligned}
\psi_L&\xrightarrow{\ \Lambda\ }\Bigl(I_2-i\boldsymbol\theta\cdot\frac{\boldsymbol\sigma}2-\mathbf b\cdot\frac{\boldsymbol\sigma}2\Bigr)\psi_L,\\
\psi_R&\xrightarrow{\ \Lambda\ }\Bigl(I_2-i\boldsymbol\theta\cdot\frac{\boldsymbol\sigma}2+\mathbf b\cdot\frac{\boldsymbol\sigma}2\Bigr)\psi_R.
\end{aligned}
\tag{1.109}\label{eq:1.109}
\end{equation*}
In terms of $\psi_L,\psi_R$, the Dirac equation \eqref{eq:1.101} reads
\begin{equation*}
\begin{pmatrix}-m&i\sigma^\mu\partial_\mu\\i\bar\sigma^\mu\partial_\mu&-m\end{pmatrix}
\begin{pmatrix}\psi_L\\\psi_R\end{pmatrix}=0
\quad\Longleftrightarrow\quad
\begin{cases}
i\sigma^\mu\partial_\mu\psi_R-m\psi_L=0,\\
i\bar\sigma^\mu\partial_\mu\psi_L-m\psi_R=0.
\end{cases}
\tag{1.110}\label{eq:1.110}
\end{equation*}
For $m=0$ it reduces to the Weyl equations
\begin{equation*}
\begin{cases}
i\bigl(\partial_t-\boldsymbol\sigma\cdot\nabla\bigr)\psi_L=0,\\
i\bigl(\partial_t+\boldsymbol\sigma\cdot\nabla\bigr)\psi_R=0.
\end{cases}
\tag{1.111}\label{eq:1.111}
\end{equation*}

\topic{Plane-Wave Solutions of the Dirac Equation}

Apply $-i\gamma^\mu\partial_\mu-m$ to the left-hand side of \eqref{eq:1.101}:
\begin{equation*}
\begin{aligned}
\bigl(-i\gamma^\mu\partial_\mu-m\bigr)\bigl(i\gamma^\nu\partial_\nu-m\bigr)\psi
&=\bigl(\gamma^\mu\gamma^\nu\partial_\mu\partial_\nu+m^2\bigr)\psi\\
&=\bigl(\eta^{\mu\nu}\partial_\mu\partial_\nu+m^2\bigr)\psi\\
&=\bigl(\Box+m^2\bigr)\psi,
\end{aligned}
\tag{1.112}\label{eq:1.112}
\end{equation*}
where $\Box:=\eta^{\mu\nu}\partial_\mu\partial_\nu=\partial_t^2-\nabla^2$. Since the left-hand side of \eqref{eq:1.112} vanishes by \eqref{eq:1.101}, every component of $\psi$ satisfies the Klein–Gordon equation
\begin{equation*}
\bigl(\Box+m^2\bigr)\psi=0.
\tag{1.113}\label{eq:1.113}
\end{equation*}
A plane wave $e^{-ip\cdot x}$ satisfies \eqref{eq:1.113} only if $p^2=m^2$. In the Weyl representation \eqref{eq:B.10}, the plane-wave solutions of \eqref{eq:1.101} therefore take the form
\begin{equation*}
\psi(x)=\begin{pmatrix}u_L\\u_R\end{pmatrix}e^{-ip\cdot x},
\qquad
p^0=\pm E_{\mathbf p},\quad
E_{\mathbf p}:=\sqrt{\mathbf p^2+m^2},
\tag{1.114}\label{eq:1.114}
\end{equation*}
where $u_L,u_R\in\mathbb C^2$ are constant. Substituting \eqref{eq:1.114} into \eqref{eq:1.101} gives
\begin{equation*}
\bigl(\slashed{p}-m\bigr)\begin{pmatrix}u_L\\u_R\end{pmatrix}=0,
\qquad
\slashed{p}:=\gamma\cdot p=p^\mu\gamma_\mu.
\tag{1.115}\label{eq:1.115}
\end{equation*}

\topic{(1) Rest Frame}

Take $m>0$ and work in the rest frame, $\mathbf p=\mathbf 0$; there we denote the constants in \eqref{eq:1.114} by $u_L(\mathbf 0)$ and $u_R(\mathbf 0)$. First consider $p^0>0$. Then $p^\mu=(m,\mathbf 0)$, $e^{-ip\cdot x}=e^{-imt}$ and $\slashed{p}=m\gamma^0$, so \eqref{eq:1.115} becomes
\begin{equation*}
\begin{aligned}
\bigl(m\gamma^0-m\bigr)\begin{pmatrix}u_L(\mathbf 0)\\u_R(\mathbf 0)\end{pmatrix}
&=\begin{pmatrix}-mI_2&mI_2\\mI_2&-mI_2\end{pmatrix}
\begin{pmatrix}u_L(\mathbf 0)\\u_R(\mathbf 0)\end{pmatrix}=0\\[6pt]
\Longrightarrow\quad u_L(\mathbf 0)&=u_R(\mathbf 0).
\end{aligned}
\tag{1.116}\label{eq:1.116}
\end{equation*}
Writing $u:=\begin{pmatrix}u_L\\u_R\end{pmatrix}$, there are two independent solutions
\begin{equation*}
u^s(\mathbf 0)=\sqrt m\begin{pmatrix}\xi^s\\\xi^s\end{pmatrix},\qquad s=1,2,
\tag{1.117}\label{eq:1.117}
\end{equation*}
where $\{\xi^s\}$ is a basis of $\mathbb C^2$ with $\xi^{r\dagger}\xi^s=\delta^{rs}$.

\topic{(2) Boost}

Boost along $z$ with rapidity $\eta$, i.e. $\omega_{03}=\eta$ in \eqref{eq:1.106}, so that $\mathbf p=(0,0,p^3)$ and $\tanh\eta=v=p^3/E_{\mathbf p}$. Then
\begin{equation*}
\begin{aligned}
\Lambda_{\frac12}=e^{-i\eta S^{03}}
&=\exp\Bigl[-i\eta\Bigl(-\frac i2\Bigr)\begin{pmatrix}\sigma^3&0\\0&-\sigma^3\end{pmatrix}\Bigr]
=\begin{pmatrix}e^{-\frac12\eta\sigma^3}&0\\0&e^{\frac12\eta\sigma^3}\end{pmatrix}\\
&=\operatorname{diag}\bigl(e^{-\eta/2},\,e^{\eta/2},\,e^{\eta/2},\,e^{-\eta/2}\bigr),
\end{aligned}
\tag{1.118}\label{eq:1.118}
\end{equation*}
\begin{equation*}
\tanh\eta=v
\ \Longrightarrow\
e^{\pm\eta}=\sqrt{\frac{1\pm v}{1\mp v}}=\frac{E_{\mathbf p}\pm p^3}{m},
\qquad
e^{\pm\eta/2}=\sqrt{\frac{E_{\mathbf p}\pm p^3}{m}}.
\tag{1.119}\label{eq:1.119}
\end{equation*}
Hence
\begin{equation*}
u^s(p)=\Lambda_{\frac12}u^s(\mathbf 0)
=\operatorname{diag}\Bigl(\sqrt{E_{\mathbf p}-p^3},\,\sqrt{E_{\mathbf p}+p^3},\,\sqrt{E_{\mathbf p}+p^3},\,\sqrt{E_{\mathbf p}-p^3}\Bigr)
\begin{pmatrix}\xi^s\\\xi^s\end{pmatrix}.
\tag{1.120}\label{eq:1.120}
\end{equation*}
With $\xi^1=\begin{pmatrix}1\\0\end{pmatrix}$ and $\xi^2=\begin{pmatrix}0\\1\end{pmatrix}$,
\begin{equation*}
u^1(p)=\begin{pmatrix}\sqrt{E_{\mathbf p}-p^3}\\0\\\sqrt{E_{\mathbf p}+p^3}\\0\end{pmatrix}
\xrightarrow{\ \eta\to\infty\ }\sqrt{2E_{\mathbf p}}\begin{pmatrix}0\\0\\1\\0\end{pmatrix},
\qquad
u^2(p)=\begin{pmatrix}0\\\sqrt{E_{\mathbf p}+p^3}\\0\\\sqrt{E_{\mathbf p}-p^3}\end{pmatrix}
\xrightarrow{\ \eta\to\infty\ }\sqrt{2E_{\mathbf p}}\begin{pmatrix}0\\1\\0\\0\end{pmatrix}.
\tag{1.121}\label{eq:1.121}
\end{equation*}

\topic{(3) Arbitrary Momentum}

Since $\operatorname{diag}\bigl(\sqrt{E_{\mathbf p}-p^3},\sqrt{E_{\mathbf p}+p^3}\bigr)=\sqrt{E_{\mathbf p}-p^3\sigma^3}$ and $\operatorname{diag}\bigl(\sqrt{E_{\mathbf p}+p^3},\sqrt{E_{\mathbf p}-p^3}\bigr)=\sqrt{E_{\mathbf p}+p^3\sigma^3}$, for arbitrary $\mathbf p$
\begin{equation*}
\sqrt m\,\Lambda_{\frac12}=\begin{pmatrix}\sqrt{p\cdot\sigma}&0\\0&\sqrt{p\cdot\bar\sigma}\end{pmatrix},
\qquad
p\cdot\sigma:=p_\mu\sigma^\mu=p^0I_2-\mathbf p\cdot\boldsymbol\sigma,
\quad
p\cdot\bar\sigma:=p_\mu\bar\sigma^\mu=p^0I_2+\mathbf p\cdot\boldsymbol\sigma,
\tag{1.122}\label{eq:1.122}
\end{equation*}
with positive square roots. The positive-energy solutions are
\begin{equation*}
{\displaystyle 
\psi=u^s(p)\,e^{-ip\cdot x},\qquad
u^s(p)=\begin{pmatrix}\sqrt{p\cdot\sigma}\,\xi^s\\\sqrt{p\cdot\bar\sigma}\,\xi^s\end{pmatrix},\qquad
p^0=E_{\mathbf p}.
}
\tag{1.123}\label{eq:1.123}
\end{equation*}

\topic{Verification}
\begin{equation*}
\begin{aligned}
(p\cdot\bar\sigma)(p\cdot\sigma)
&=\bigl(p^0I_2+p^i\sigma^i\bigr)\bigl(p^0I_2-p^j\sigma^j\bigr)
=(p^0)^2-\frac12\{\sigma^i,\sigma^j\}p^ip^j
=(p^0)^2-\mathbf p^2=m^2,\\[6pt]
\bigl(\slashed{p}-m\bigr)u^s(p)
&=\begin{pmatrix}-m&p\cdot\sigma\\p\cdot\bar\sigma&-m\end{pmatrix}
\begin{pmatrix}\sqrt{p\cdot\sigma}\,\xi^s\\\sqrt{p\cdot\bar\sigma}\,\xi^s\end{pmatrix}
=\begin{pmatrix}\sqrt{p\cdot\sigma}\,\bigl(\sqrt{(p\cdot\sigma)(p\cdot\bar\sigma)}-m\bigr)\xi^s\\\sqrt{p\cdot\bar\sigma}\,\bigl(\sqrt{(p\cdot\bar\sigma)(p\cdot\sigma)}-m\bigr)\xi^s\end{pmatrix}=0.
\end{aligned}
\tag{1.124}\label{eq:1.124}
\end{equation*}

\topic{Negative-Energy Solutions}

For $p^0=-E_{\mathbf p}$, \eqref{eq:1.115} becomes
\begin{equation*}
\bigl(-\gamma^0E_{\mathbf p}-\boldsymbol\gamma\cdot\mathbf p-m\bigr)u(p)=0
\quad\Longleftrightarrow\quad
\bigl(\gamma^0E_{\mathbf p}-\boldsymbol\gamma\cdot(-\mathbf p)+m\bigr)u(p)=0,
\tag{1.125}\label{eq:1.125}
\end{equation*}
which is the positive-energy equation $\bigl(\gamma^0E_{\mathbf p}-\boldsymbol\gamma\cdot\mathbf p-m\bigr)u=0$ with $\mathbf p\to-\mathbf p$ and $m\to-m$. Hence
\begin{equation*}
u^s(p)=\begin{pmatrix}\sqrt{E_{\mathbf p}I_2+\mathbf p\cdot\boldsymbol\sigma}\,\xi^s\\-\sqrt{E_{\mathbf p}I_2-\mathbf p\cdot\boldsymbol\sigma}\,\xi^s\end{pmatrix}
=\begin{pmatrix}\sqrt{-p\cdot\sigma}\,\xi^s\\-\sqrt{-p\cdot\bar\sigma}\,\xi^s\end{pmatrix},
\qquad p^0=-E_{\mathbf p}.
\tag{1.126}\label{eq:1.126}
\end{equation*}

\topic{Antiparticle Solutions}

Define
\begin{equation*}
u^c(p):=\begin{pmatrix}i\sigma^2&0\\0&-i\sigma^2\end{pmatrix}\gamma^0\,u(p),
\qquad
i\sigma^2=\begin{pmatrix}0&1\\-1&0\end{pmatrix}.
\tag{1.127}\label{eq:1.127}
\end{equation*}
For $\mathbf p=(0,0,p^3)$,
\begin{equation*}
i\sigma^2\sqrt{p\cdot\bar\sigma}
=\begin{pmatrix}0&1\\-1&0\end{pmatrix}\begin{pmatrix}\sqrt{E_{\mathbf p}+p^3}&0\\0&\sqrt{E_{\mathbf p}-p^3}\end{pmatrix}
=\begin{pmatrix}0&\sqrt{E_{\mathbf p}-p^3}\\-\sqrt{E_{\mathbf p}+p^3}&0\end{pmatrix}
=\sqrt{p\cdot\sigma}\;i\sigma^2,
\tag{1.128}\label{eq:1.128}
\end{equation*}
and likewise $i\sigma^2\sqrt{p\cdot\sigma}=\sqrt{p\cdot\bar\sigma}\;i\sigma^2$. Applying \eqref{eq:1.127} to \eqref{eq:1.123} and writing $\eta^s:=i\sigma^2\xi^s$,
\begin{equation*}
u^c(p)=\begin{pmatrix}i\sigma^2\sqrt{p\cdot\bar\sigma}\,\xi^s\\-i\sigma^2\sqrt{p\cdot\sigma}\,\xi^s\end{pmatrix}
=\begin{pmatrix}\sqrt{p\cdot\sigma}\,\eta^s\\-\sqrt{p\cdot\bar\sigma}\,\eta^s\end{pmatrix},
\qquad p^0=E_{\mathbf p},
\tag{1.129}\label{eq:1.129}
\end{equation*}
which is \eqref{eq:1.126} at $-p$ with $\xi^s\to\eta^s$. For arbitrary $\mathbf p$, the antiparticle solutions are
\begin{equation*}
{\displaystyle 
\psi=v^s(p)\,e^{+ip\cdot x},\qquad
v^s(p):=\begin{pmatrix}\sqrt{p\cdot\sigma}\,\eta^s\\-\sqrt{p\cdot\bar\sigma}\,\eta^s\end{pmatrix},\qquad
\bigl(\slashed{p}+m\bigr)v^s(p)=0,\qquad
p^0=E_{\mathbf p}.
}
\tag{1.130}\label{eq:1.130}
\end{equation*}

\topic{Properties of the Solutions}

\topic{(1) Orthonormality}

From \eqref{eq:1.123}, $u^\dagger u=2E_{\mathbf p}\,\xi^\dagger\xi$ and $\bar uu=2m\,\xi^\dagger\xi$. With $\xi^{r\dagger}\xi^s=\eta^{r\dagger}\eta^s=\delta^{rs}$,
\begin{equation*}
{\displaystyle 
\begin{aligned}
u^{r\dagger}(p)\,u^s(p)&=2E_{\mathbf p}\,\delta^{rs}, &\qquad \bar u^r(p)\,u^s(p)&=2m\,\delta^{rs},\\
v^{r\dagger}(p)\,v^s(p)&=2E_{\mathbf p}\,\delta^{rs}, &\qquad \bar v^r(p)\,v^s(p)&=-2m\,\delta^{rs},\\
u^{r\dagger}(p)\,v^s(p')&=v^{r\dagger}(p')\,u^s(p)=0, &\qquad p'&:=(p^0,-\mathbf p).
\end{aligned}
}
\tag{1.131}\label{eq:1.131}
\end{equation*}

\topic{(2) Completeness (Spin Sums)}

With $\sum_s\xi^s\xi^{s\dagger}=I_2$,
\begin{equation*}
{\displaystyle 
\sum_s u^s(p)\,\bar u^s(p)=\slashed{p}+m,\qquad
\sum_s v^s(p)\,\bar v^s(p)=\slashed{p}-m.
}
\tag{1.132}\label{eq:1.132}
\end{equation*}

\subsection{Dirac Bilinears}\label{book-subsection-section-1.7-3}\label{dirac-bilinears}

Bilinears have the form $\bar\psi\Gamma\psi$ with $\Gamma$ a $4\times4$ matrix. A basis of the $4\times4$ matrices consists of 16 matrices:

\begin{center}
\renewcommand{\arraystretch}{1.2}
\begin{tabular}{ccc}
\toprule
$\Gamma$ & Number & Replaced by \\
\midrule
$I_4$ & 1 &  \\
$\gamma^\mu$ & 4 &  \\
$S^{\mu\nu}$ & 6 &  \\
$\gamma^{[\mu}\gamma^\nu\gamma^{\rho]}$ & 4 & $\gamma^5\gamma^\mu$ \\
$\gamma^{[\mu}\gamma^\nu\gamma^\rho\gamma^{\sigma]}$ & 1 & $\gamma^5$ \\
\bottomrule
\end{tabular}
\end{center}

where
\begin{equation*}
\gamma^5=\frac{i}{4!}\epsilon_{\mu\nu\rho\sigma}\gamma^\mu\gamma^\nu\gamma^\rho\gamma^\sigma,
\qquad \epsilon_{0123}=+1,
\tag{1.133}\label{eq:1.133}
\end{equation*}
which agrees with \eqref{eq:B.14}. Under $\psi\to\Lambda_{\frac12}\psi$, $\bar\psi\to\bar\psi\Lambda_{\frac12}^{-1}$,
\begin{equation*}
\begin{aligned}
\bar\psi\psi&\xrightarrow{\ \Lambda\ }\bar\psi\psi,\\
\bar\psi\gamma^\mu\psi&\xrightarrow{\ \Lambda\ }\bar\psi\Lambda_{\frac12}^{-1}\gamma^\mu\Lambda_{\frac12}\psi=\Lambda^\mu{}_\nu\,\bar\psi\gamma^\nu\psi,\\
\bar\psi S^{\mu\nu}\psi&\xrightarrow{\ \Lambda\ }\Lambda^\mu{}_\rho\Lambda^\nu{}_\sigma\,\bar\psi S^{\rho\sigma}\psi,\\
\bar\psi\gamma^5\gamma^\mu\psi&\xrightarrow{\ \Lambda\ }\det[\Lambda]\,\Lambda^\mu{}_\nu\,\bar\psi\gamma^5\gamma^\nu\psi\qquad\text{(pseudovector)},\\
\bar\psi\gamma^5\psi&\xrightarrow{\ \Lambda\ }\det[\Lambda]\,\bar\psi\gamma^5\psi\qquad\text{(pseudoscalar)}.
\end{aligned}
\tag{1.134}\label{eq:1.134}
\end{equation*}
\textbf{\boldmath Proof} (for $\gamma^5$)
\begin{equation*}
\begin{aligned}
\Lambda_{\frac12}^{-1}\gamma^5\Lambda_{\frac12}
&=\frac{i}{4!}\epsilon_{\mu\nu\rho\sigma}\Lambda^\mu{}_\alpha\Lambda^\nu{}_\beta\Lambda^\rho{}_\kappa\Lambda^\sigma{}_\lambda\,\gamma^\alpha\gamma^\beta\gamma^\kappa\gamma^\lambda\\
&=\frac{i}{4!}\det[\Lambda]\,\epsilon_{\alpha\beta\kappa\lambda}\gamma^\alpha\gamma^\beta\gamma^\kappa\gamma^\lambda
=\det[\Lambda]\,\gamma^5.
\end{aligned}
\tag{1.135}\label{eq:1.135}
\end{equation*}

\topic{Chirality}
\begin{equation*}
\gamma^5\begin{pmatrix}\psi_L\\0\end{pmatrix}=-\begin{pmatrix}\psi_L\\0\end{pmatrix},
\qquad
\gamma^5\begin{pmatrix}0\\\psi_R\end{pmatrix}=\begin{pmatrix}0\\\psi_R\end{pmatrix}.
\tag{1.136}\label{eq:1.136}
\end{equation*}
The left- and right-handed projection operators are
\begin{equation*}
P_L:=\frac{I_4-\gamma^5}2,\qquad P_R:=\frac{I_4+\gamma^5}2
\quad\Longrightarrow\quad
\begin{cases}
P_L+P_R=I_4,\\
P_LP_R=P_RP_L=0,\\
P_L^n=P_L,\quad P_R^n=P_R.
\end{cases}
\tag{1.137}\label{eq:1.137}
\end{equation*}

\topic{Vector and Axial Currents}
\begin{equation*}
j^\mu:=\bar\psi\gamma^\mu\psi,\qquad
j^{\mu5}:=\bar\psi\gamma^\mu\gamma^5\psi.
\tag{1.138}\label{eq:1.138}
\end{equation*}
The Euler–Lagrange equation of \eqref{eq:1.100} for $\psi$ is $-i\partial_\mu\bar\psi\gamma^\mu-m\bar\psi=0$. Together with \eqref{eq:1.101},
\begin{equation*}
\partial_\mu j^\mu=(\partial_\mu\bar\psi)\gamma^\mu\psi+\bar\psi\gamma^\mu\partial_\mu\psi
=im\bar\psi\psi-im\bar\psi\psi=0
\quad\Longrightarrow\quad
\int d^3x\,\bar\psi\gamma^0\psi=\text{const}.
\tag{1.139}\label{eq:1.139}
\end{equation*}

\topic{Helicity}
\begin{equation*}
h:=\frac{\mathbf p}{|\mathbf p|}\cdot\mathbf S,\qquad
\mathbf S:=\frac12\begin{pmatrix}\boldsymbol\sigma&0\\0&\boldsymbol\sigma\end{pmatrix}.
\tag{1.140}\label{eq:1.140}
\end{equation*}
For $\mathbf p=(0,0,p^3)$ with $p^3>0$, the spinors \eqref{eq:1.121} satisfy $h\,u^1=\frac12u^1$ and $h\,u^2=-\frac12u^2$.

\setcounter{section}{7}
\section{The Proca Field}\label{book-section-1.8}\label{proca-field}

\subsection{Lagrangian}\label{book-subsection-section-1.8-1}\label{proca-lagrangian}

For a real vector field with $m>0$,
\begin{equation*}
F_{\mu\nu}:=\partial_\mu A_\nu-\partial_\nu A_\mu,
\qquad
{\displaystyle 
\mathcal L=-\frac14F_{\mu\nu}F^{\mu\nu}
+\frac12m^2A_\mu A^\mu.
}
\tag{1.141}\label{eq:1.141}
\end{equation*}
\begin{equation*}
{\displaystyle 
\partial_\mu F^{\mu\nu}+m^2A^\nu=0,
\qquad
\partial_\mu A^\mu=0,
\qquad
(\partial_\mu\partial^\mu+m^2)A^\nu=0.
}
\tag{1.142}\label{eq:1.142}
\end{equation*}
Here $\partial_\mu A^\mu=0$ follows from the equations of motion.

\topic{Canonical Momenta and Hamiltonian}

Choose $A_\mu$ as canonical coordinates. All brackets below are at equal time:
\begin{equation*}
\Pi^\mu:=\frac{\partial\mathcal L}{\partial\dot A_\mu}=F^{\mu0},
\qquad
\Pi^0=0,
\qquad
\Pi^i=\dot A_i-\partial_iA_0,
\qquad
\{A_\mu(\mathbf x),\Pi^\nu(\mathbf y)\}_{\mathrm{PB}}
=\delta_\mu{}^\nu\delta^{(3)}(\mathbf x-\mathbf y).
\tag{1.143}\label{eq:1.143}
\end{equation*}
With vanishing spatial surface terms,
\begin{equation*}
{\displaystyle 
\begin{aligned}
H_c&=\int d^3x\left[
\frac12\sum_i(\Pi^i)^2+\frac14\sum_{i,j}(F_{ij})^2
+\frac{m^2}{2}\sum_i(A_i)^2
-\frac{m^2}{2}A_0^2-A_0\partial_i\Pi^i\right],\\
H_T&=H_c+\int d^3x\,u\,\Pi^0.
\end{aligned}}
\tag{1.144}\label{eq:1.144}
\end{equation*}

\topic{Second-Class Constraints}
\begin{equation*}
{\displaystyle 
\begin{aligned}
\chi_1&:=\Pi^0\approx0,
&
\chi_2&:=\partial_i\Pi^i+m^2A_0\approx0,\\
\{\chi_1(\mathbf x),\chi_2(\mathbf y)\}_{\mathrm{PB}}
&=-m^2\delta^{(3)}(\mathbf x-\mathbf y),
&
N_{\mathrm{physical}}&=4-\frac22=3.
\end{aligned}}
\tag{1.145}\label{eq:1.145}
\end{equation*}
Preserving $\chi_2$ fixes $u=\sum_i\partial_iA_i$; there are no further constraints.

Here $\chi_1$ is primary and $\chi_2$ is secondary; both are second class.

\topic{Reduced Hamiltonian}

Eliminating $\Pi^0$ and $A_0$ gives
\begin{equation*}
A_0=-\frac{\partial_i\Pi^i}{m^2},
\qquad
{\displaystyle 
H_{\mathrm{red}}=\int d^3x\left[
\frac12\sum_i(\Pi^i)^2+\frac14\sum_{i,j}(F_{ij})^2
+\frac{m^2}{2}\sum_i(A_i)^2
+\frac{(\partial_i\Pi^i)^2}{2m^2}
\right]\ge0.
}
\tag{1.146}\label{eq:1.146}
\end{equation*}

\topic{Conserved Quantities: Spacetime Symmetries}

For the canonical variables $A_\mu$, translations give
\begin{equation*}
{\displaystyle \begin{aligned}
T^\mu{}_\nu&=-F^{\mu\rho}\partial_\nu A_\rho-\delta^\mu{}_\nu\mathcal L,
&\partial_\mu T^\mu{}_\nu&=0,\\
P_\nu&=\int d^3x\,T^0{}_\nu,
&\frac{dP_\nu}{dt}&=0.
\end{aligned}}
\tag{1.147}\label{eq:1.147}
\end{equation*}
With $T^{\mu\nu}:=T^\mu{}_\rho\eta^{\rho\nu}$, the Lorentz charges include the vector spin contribution:
\begin{equation*}
{\displaystyle 
J^{\mu\nu}=\int d^3x\left[
x^\mu T^{0\nu}-x^\nu T^{0\mu}
+\Pi^\mu A^\nu-\Pi^\nu A^\mu
\right],
\qquad
\frac{dJ^{\mu\nu}}{dt}=0.
}
\tag{1.148}\label{eq:1.148}
\end{equation*}
These are conserved on shell, with convergent integrals and vanishing boundary fluxes.

\topic{Conserved Quantities: Internal and Gauge Symmetries}

For $m>0$, the mass term breaks the Maxwell gauge symmetry $A_\mu\mapsto A_\mu+\partial_\mu\lambda(x)$. The real Proca field has no complex $U(1)$ phase symmetry, so this model has no corresponding phase charge. Its constraint $\partial_i\Pi^i+m^2A_0=0$ is second class, as in \eqref{eq:1.145}.

\subsection{Solutions of the Proca Equation}\label{book-subsection-section-1.8-2}\label{proca-field-solutions}

For $A^\mu(x)=\epsilon^\mu(p)e^{-ip\cdot x}$, the field equation becomes
\begin{equation*}
\begin{aligned}
(m^2-p^2)\epsilon^\mu+p^\mu(p_\nu\epsilon^\nu)&=0,\\
m^2p_\mu\epsilon^\mu&=0.
\end{aligned}
\quad\Longrightarrow\quad
{\displaystyle p^2=m^2,\qquad p_\mu\epsilon^\mu=0}
\tag{1.149}\label{eq:1.149}
\end{equation*}
for a nonzero mode with $m>0$.

\topic{Polarizations}

For $p^2=m^2$, $p^0>0$, there are three polarization vectors, $r=1,2,3$:
\begin{equation*}
{\displaystyle 
p_\mu\epsilon_r^\mu(p)=0,
\qquad
\epsilon_{r\mu}^*(p)\epsilon_s^\mu(p)=-\delta_{rs},
\qquad
\sum_{r=1}^3\epsilon_r^\mu(p)\epsilon_r^{\nu*}(p)
=-\eta^{\mu\nu}+\frac{p^\mu p^\nu}{m^2}.
}
\tag{1.150}\label{eq:1.150}
\end{equation*}
At rest, one may take $\epsilon_r^\mu=(0,\mathbf e_r)$ for an orthonormal spatial basis $\mathbf e_r$. For $\mathbf p\ne0$, choose $\mathbf e_1,\mathbf e_2$ perpendicular to $\mathbf p$, with $\mathbf e_a\cdot\mathbf e_b=\delta_{ab}$. With $E_{\mathbf p}=\sqrt{\mathbf p^2+m^2}$, a basis with two transverse and one longitudinal polarization is
\begin{equation*}
{\displaystyle 
\epsilon_a^\mu(p)=(0,\mathbf e_a)\quad(a=1,2),
\qquad
\epsilon_3^\mu(p)=\left(\frac{|\mathbf p|}{m},
\frac{E_{\mathbf p}}{m}\frac{\mathbf p}{|\mathbf p|}\right).
}
\tag{1.151}\label{eq:1.151}
\end{equation*}
The real field is a superposition of these three modes, with $p^0=E_{\mathbf p}$:
\begin{equation*}
{\displaystyle 
A^\mu(x)=\sum_{r=1}^3\int\frac{d^3p}{(2\pi)^3\,2E_{\mathbf p}}
\left[
a_r(\mathbf p)\epsilon_r^\mu(p)e^{-ip\cdot x}
+a_r^*(\mathbf p)\epsilon_r^{\mu*}(p)e^{+ip\cdot x}
\right].
}
\tag{1.152}\label{eq:1.152}
\end{equation*}
Here $a_r(\mathbf p)$ are classical complex amplitudes; the transversality condition is imposed by the field equation, not by a gauge choice.

\setcounter{section}{8}
\section{The Vector Field}\label{book-section-1.9}\label{vector-field}

\subsection{Lagrangian}\label{book-subsection-section-1.9-1}\label{maxwell-lagrangian}
\begin{equation*}
{\displaystyle 
\mathcal L=-\frac14F_{\mu\nu}F^{\mu\nu},
\qquad
\partial_\mu F^{\mu\nu}=0,
\qquad
A_\mu'=A_\mu+\partial_\mu\lambda(x).
}
\tag{1.153}\label{eq:1.153}
\end{equation*}

\topic{Hamiltonian and First-Class Constraints}
\begin{equation*}
\begin{aligned}
\Pi^\mu&=F^{\mu0},\\
H_c&=\int d^3x\left[
\frac12\sum_i(\Pi^i)^2+\frac14\sum_{i,j}(F_{ij})^2
-A_0\partial_i\Pi^i\right],
&
H_T&=H_c+\int d^3x\,u\,\Pi^0,\\
\chi_1&:=\Pi^0\approx0,
&
\chi_2&:=\partial_i\Pi^i\approx0.
\end{aligned}
\tag{1.154}\label{eq:1.154}
\end{equation*}
Both constraints are first class:
\begin{equation*}
\{\chi_a(\mathbf x),\chi_b(\mathbf y)\}_{\mathrm{PB}}=0,
\qquad a,b=1,2,
\qquad
{\displaystyle N_{\mathrm{physical}}=4-2=2.}
\tag{1.155}\label{eq:1.155}
\end{equation*}

\topic{Coulomb Gauge and Physical Hamiltonian}

In the absence of sources, with fields decaying at spatial infinity,
\begin{equation*}
\sum_i\partial_iA_i=0,
\qquad
A_0=0,
\qquad
\partial_i\Pi^i=0,
\qquad
{\displaystyle 
H_{\mathrm{physical}}
=\int d^3x\left[
\frac12\sum_i(\Pi^i)^2+\frac14\sum_{i,j}(F_{ij})^2
\right]\ge0.
}
\tag{1.156}\label{eq:1.156}
\end{equation*}
The remaining canonical variables are transverse.

\topic{Conserved Quantities: Spacetime Symmetries}

The canonical energy–momentum tensor and four-momentum are
\begin{equation*}
{\displaystyle \begin{aligned}
T^\mu{}_\nu&=-F^{\mu\rho}\partial_\nu A_\rho-\delta^\mu{}_\nu\mathcal L,
&\partial_\mu T^\mu{}_\nu&=0,\\
P_\nu&=\int d^3x\,T^0{}_\nu,
&\frac{dP_\nu}{dt}&=0.
\end{aligned}}
\tag{1.157}\label{eq:1.157}
\end{equation*}
With $T^{\mu\nu}:=T^\mu{}_\rho\eta^{\rho\nu}$, the rotation and boost charges are
\begin{equation*}
{\displaystyle 
J^{\mu\nu}=\int d^3x\left[
x^\mu T^{0\nu}-x^\nu T^{0\mu}
+\Pi^\mu A^\nu-\Pi^\nu A^\mu
\right],
\qquad
\frac{dJ^{\mu\nu}}{dt}=0.
}
\tag{1.158}\label{eq:1.158}
\end{equation*}
These charges are conserved on shell when the integrals converge and boundary fluxes vanish. The canonical densities depend on the gauge; the integrated charges are invariant under gauge transformations whose boundary terms vanish.

\topic{Conserved Quantities: Internal and Gauge Symmetries}

The local transformation in \eqref{eq:1.153} gives a Noether identity and the Gauss constraint. For smooth source-free fields on all of $\mathbb R^3$,
\begin{equation*}
{\displaystyle \begin{aligned}
\partial_\nu\partial_\mu F^{\mu\nu}&\equiv0,
&\partial_i\Pi^i&=0,\\
Q_{\mathrm{el}}&:=\lim_{R\to\infty}\oint_{|\mathbf x|=R}\Pi^i\,dS_i=0.
\end{aligned}}
\tag{1.159}\label{eq:1.159}
\end{equation*}
There is no independent matter $U(1)$ number in this free-field model. Nonzero electric flux requires sources, singularities, or different boundary conditions.

\subsection{Solutions of the Maxwell Equations}\label{book-subsection-section-1.9-2}\label{maxwell-field-solutions}

\topic{Lorenz Gauge}
\begin{equation*}
\partial_\mu A^\mu=0
\quad\Longrightarrow\quad
\partial_\mu\partial^\mu A^\nu=0,
\qquad
A_\mu\mapsto A_\mu+\partial_\mu\lambda,
\qquad
\partial_\mu\partial^\mu\lambda=0
\quad\text{(residual gauge freedom).}
\tag{1.160}\label{eq:1.160}
\end{equation*}

\topic{Plane Waves and Physical Polarizations}

For a plane wave $A^\mu(x)=\epsilon^\mu(p)e^{-ip\cdot x}$ in Lorenz gauge,
\begin{equation*}
{\displaystyle 
p^2=0,\qquad p^0=|\mathbf p|>0,\qquad
p_\mu\epsilon^\mu=0,\qquad
\epsilon^\mu\sim\epsilon^\mu+\lambda p^\mu.
}
\tag{1.161}\label{eq:1.161}
\end{equation*}
In the last expression, $\lambda$ is an arbitrary complex coefficient of a residual plane-wave gauge transformation. The constraint and this equivalence leave two physical polarizations. A Coulomb-gauge basis satisfies
\begin{equation*}
{\displaystyle \begin{aligned}
\epsilon_r^0(p)&=0,\qquad
\mathbf p\cdot\boldsymbol\epsilon_r(p)=0,\qquad
\boldsymbol\epsilon_r^*(p)\cdot\boldsymbol\epsilon_s(p)=\delta_{rs},
\qquad r,s=1,2,\\
\sum_{r=1}^2\epsilon_r^i(p)\epsilon_r^{j*}(p)
&=\delta^{ij}-\frac{p^ip^j}{|\mathbf p|^2}.
\end{aligned}}
\tag{1.162}\label{eq:1.162}
\end{equation*}
For source-free wave packets, the real field in this gauge is, with $p^0=|\mathbf p|$,
\begin{equation*}
{\displaystyle 
A^\mu(x)=\sum_{r=1}^2\int\frac{d^3p}{(2\pi)^3\,2|\mathbf p|}
\left[
a_r(\mathbf p)\epsilon_r^\mu(p)e^{-ip\cdot x}
+a_r^*(\mathbf p)\epsilon_r^{\mu*}(p)e^{+ip\cdot x}
\right].
}
\tag{1.163}\label{eq:1.163}
\end{equation*}
The $a_r(\mathbf p)$ are classical complex amplitudes. Other gauge representatives follow by adding $\partial^\mu\lambda(x)$.

\setcounter{chapter}{1}
\chapter{Operator Quantization}\label{book-chapter-2}\label{qft-ch2}

The first viewpoint in \hyperref[classical-dynamics]{Section 1.4} naturally leads to Operator quantization: spatial position labels the degrees of freedom, while time remains the evolution parameter.

\setcounter{section}{0}
\section{Pictures of Time Evolution}\label{book-section-2.1}\label{pictures-of-time-evolution}

Throughout, the pictures are taken to coincide at $t=0$.

In the Hamiltonian formalism the dynamics is given by the Hamiltonian $\hat H$. The time dependence of an observable therefore has two sources:
\begin{enumerate}


\item The first is the dynamics of the system, i.e. the time evolution generated by $\hat H$, which differs from system to system: a different Hamiltonian gives the observable a different time dependence. 

\item The second is a change in the form of the observable itself, which appears as an explicit time dependence and is described mathematically by the partial derivative with respect to $t$.
\end{enumerate}


For example, the position and momentum operators are the most basic observables. Their form is fixed by their own definition, so they carry no explicit time dependence, but they can still evolve in time. The form of an interaction, by contrast, can change with time: for a charged system, varying the strength of an external magnetic field by hand changes the form of the interaction.

\subsection{Schrödinger Picture}\label{book-subsection-section-2.1-1}\label{schrodinger-picture}

\topic{The dynamics is carried entirely by the states; operators do not evolve in time.}
\begin{equation*}
\frac{d \hat O}{dt} - \frac{\partial \hat O}{\partial t} = 0
\tag{2.1}\label{eq:2.1}
\end{equation*}
Hence
\begin{equation*}
\frac{d\hat q_i}{dt}=\frac{\partial\hat q_i}{\partial t}=0,
\qquad
\frac{d\hat p_i}{dt}=\frac{\partial\hat p_i}{\partial t}=0.
\tag{2.2}\label{eq:2.2}
\end{equation*}
The time dependence of the states comes from the dynamics; as one of the axioms of quantum mechanics, it is given by the celebrated Schrödinger equation
\begin{equation*}
i\frac{\partial}{\partial t}|\psi_S(t)\rangle=\hat H_S|\psi_S(t)\rangle.
\tag{2.3}\label{eq:2.3}
\end{equation*}
Note that the form of $\hat H_S$ in the equation may change with time; accordingly, the solution of the Schrödinger equation takes three forms.

(1) $\frac{\partial \hat H_S}{\partial t} = 0$
\begin{equation*}
|\psi_S(t)\rangle=e^{-i\hat H_S(t-t_0)}\,|\psi_S(t_0)\rangle .
\tag{2.4}\label{eq:2.4}
\end{equation*}
where the exponential of an operator is defined by the power series:
\begin{equation*}
e^{\hat O}:=\sum_{n=0}^{\infty}\frac{1}{n!}\,\hat O^n.
\tag{2.5}\label{eq:2.5}
\end{equation*}
(2) $\frac{\partial\hat H_S}{\partial t}\neq0$, but $[\hat H_S(t),\hat H_S(t')]=0$ for all $t,t'$
\begin{equation*}
|\psi_S(t)\rangle=\exp\Bigl(-i\int_{t_0}^{t}d\tau\,\hat H_S(\tau)\Bigr)\,|\psi_S(t_0)\rangle ,
\tag{2.6}\label{eq:2.6}
\end{equation*}
where
\begin{equation*}
\langle \psi_1 | \int_{t_0}^{t}d\tau\,\hat H_S(\tau) |  \psi_2  \rangle  := \int_{t_0}^{t}d\tau\,\langle\psi_1 |\hat H_S(\tau)|\psi_2\rangle
\tag{2.7}\label{eq:2.7}
\end{equation*}
One can show that
\begin{equation*}
\langle \psi_1 | [\int_{t_0}^{t}d\tau\,\hat H_S(\tau)]^n |  \psi_2  \rangle  = \int_{t_0}^{t}...\int_{t_0}^{t}d\tau_1 ... d\tau_n\,\langle\psi_1 |\hat H_S(\tau_1)...\hat H_S(\tau_n)|\psi_2\rangle
\tag{2.8}\label{eq:2.8}
\end{equation*}
\begin{equation*}
\begin{aligned}
\frac{\partial}{\partial t}\Bigl(\int_{t_0}^{t}d\tau\,\hat H_S(\tau)\Bigr)^n
&=\sum_{k=0}^{n-1}\Bigl(\int_{t_0}^{t}d\tau\,\hat H_S(\tau)\Bigr)^k\,\hat H_S(t)\,\Bigl(\int_{t_0}^{t}d\tau\,\hat H_S(\tau)\Bigr)^{n-1-k}\\
&=n\,\hat H_S(t)\Bigl(\int_{t_0}^{t}d\tau\,\hat H_S(\tau)\Bigr)^{n-1}.
\end{aligned}
\tag{2.9}\label{eq:2.9}
\end{equation*}
Hence
\begin{equation*}
\begin{aligned}
\frac{\partial}{\partial t}\exp\Bigl(-i\int_{t_0}^{t}d\tau\,\hat H_S(\tau)\Bigr)
&=\sum_{n=1}^{\infty}\frac{(-i)^n}{n!}\,n\,\hat H_S(t)\Bigl(\int_{t_0}^{t}d\tau\,\hat H_S(\tau)\Bigr)^{n-1}\\
&=-i\hat H_S(t)\exp\Bigl(-i\int_{t_0}^{t}d\tau\,\hat H_S(\tau)\Bigr).
\end{aligned}
\tag{2.10}\label{eq:2.10}
\end{equation*}
With these formulas one verifies that the solution indeed has the form given above.

(3) $[\hat H_S(t),\hat H_S(t')]\neq0$ for some $t,t'$
\begin{equation*}
|\psi_S(t)\rangle=\mathcal T\exp\Bigl(-i\int_{t_0}^{t}d\tau\,\hat H_S(\tau)\Bigr)\,|\psi_S(t_0)\rangle ,
\tag{2.11}\label{eq:2.11}
\end{equation*}
where the time-ordered exponential is defined by the series (the Dyson series):
\begin{equation*}
\mathcal T\exp\Bigl(-i\int_{t_0}^{t}d\tau\,\hat H_S(\tau)\Bigr)
:=\sum_{n=0}^{\infty}\frac{(-i)^n}{n!}\int_{t_0}^{t}d\tau_1\cdots\int_{t_0}^{t}d\tau_n\,
\mathcal T[\hat H_S(\tau_1)\cdots\hat H_S(\tau_n)],
\tag{2.12}\label{eq:2.12}
\end{equation*}
The time-ordered product $\mathcal T$ orders the operators with time decreasing from left to right:
\begin{equation*}
\mathcal T[\hat H_S(\tau_1)\hat H_S(\tau_2)]
:=\theta(\tau_1-\tau_2)\,\hat H_S(\tau_1)\hat H_S(\tau_2)+\theta(\tau_2-\tau_1)\,\hat H_S(\tau_2)\hat H_S(\tau_1),
\tag{2.13}\label{eq:2.13}
\end{equation*}
The case of $n$ operators is analogous.

One can show that
\begin{equation*}
\frac{\partial}{\partial t}\int_{t_0}^{t}d\tau_1\cdots\int_{t_0}^{t}d\tau_n\,
\mathcal T[\hat H_S(\tau_1)\cdots\hat H_S(\tau_n)]
=n\,\hat H_S(t)\int_{t_0}^{t}d\tau_1\cdots\int_{t_0}^{t}d\tau_{n-1}\,
\mathcal T[\hat H_S(\tau_1)\cdots\hat H_S(\tau_{n-1})],
\tag{2.14}\label{eq:2.14}
\end{equation*}
Hence
\begin{equation*}
\begin{aligned}
\frac{\partial}{\partial t}\mathcal T\exp\Bigl(-i\int_{t_0}^{t}d\tau\,\hat H_S(\tau)\Bigr)
&=\sum_{n=1}^{\infty}\frac{(-i)^n}{n!}\,n\,\hat H_S(t)\int_{t_0}^{t}d\tau_1\cdots\int_{t_0}^{t}d\tau_{n-1}\,
\mathcal T[\hat H_S(\tau_1)\cdots\hat H_S(\tau_{n-1})]\\
&=-i\hat H_S(t)\,\mathcal T\exp\Bigl(-i\int_{t_0}^{t}d\tau\,\hat H_S(\tau)\Bigr).
\end{aligned}
\tag{2.15}\label{eq:2.15}
\end{equation*}
If $[\hat H_S(t),\hat H_S(t')]=0$ for all $t,t'$, then $\mathcal T$ can be dropped and the above reduces to case (2); if they do not commute, the second equality of the power identity in case (2) fails and $\mathcal T$ must be kept.

This series can be obtained by integrating the Schrödinger equation and iterating:
\begin{equation*}
\begin{aligned}
|\psi_S(t)\rangle
&=|\psi_S(t_0)\rangle-i\int_{t_0}^{t}d\tau_1\,\hat H_S(\tau_1)\,|\psi_S(\tau_1)\rangle\\
&=\Bigl[1+(-i)\int_{t_0}^{t}d\tau_1\,\hat H_S(\tau_1)
+(-i)^2\int_{t_0}^{t}d\tau_1\int_{t_0}^{\tau_1}d\tau_2\,\hat H_S(\tau_1)\hat H_S(\tau_2)+\cdots\Bigr]|\psi_S(t_0)\rangle ,
\end{aligned}
\tag{2.16}\label{eq:2.16}
\end{equation*}
where the $n$-th term is integrated over $t\ge\tau_1\ge\cdots\ge\tau_n\ge t_0$, with operators at later times to the left.

$\mathcal T[\hat H_S(\tau_1)\cdots\hat H_S(\tau_n)]$ is symmetric under permutations of $\tau_1,\ldots,\tau_n$; the cube $[t_0,t]^n$ splits into $n!$ regions according to the ordering of the $\tau_i$, each contributing equally, and on $\tau_1>\cdots>\tau_n$ the time-ordered product is the ordered product.

Hence
\begin{equation*}
\int_{t_0}^{t}d\tau_1\int_{t_0}^{\tau_1}d\tau_2\cdots\int_{t_0}^{\tau_{n-1}}d\tau_n\,
\hat H_S(\tau_1)\cdots\hat H_S(\tau_n)
=\frac1{n!}\int_{t_0}^{t}d\tau_1\cdots\int_{t_0}^{t}d\tau_n\,
\mathcal T[\hat H_S(\tau_1)\cdots\hat H_S(\tau_n)],
\tag{2.17}\label{eq:2.17}
\end{equation*}
Substituting this gives the defining series of the time-ordered exponential.

\topic{Evolution Operator}

In each case the solution is linear in $|\psi_S(t_0)\rangle$:
\begin{equation*}
|\psi_S(t)\rangle=\hat U_S(t,t_0)\,|\psi_S(t_0)\rangle,
\tag{2.18}\label{eq:2.18}
\end{equation*}
where
\begin{equation*}
i\frac{\partial}{\partial t}\hat U_S(t,t_0)=\hat H_S(t)\,\hat U_S(t,t_0),
\qquad
\hat U_S(t_0,t_0)=1.
\tag{2.19}\label{eq:2.19}
\end{equation*}
The evolution operator $\hat U_S(t,t_0)$ is the operator exponential acting on $|\psi_S(t_0)\rangle$ in the three cases above.

It is easy to show that
\begin{equation*}
\hat U_S(t_2,t_0) = \hat U_S(t_2,t_1) \hat U_S(t_1,t_0)
\tag{2.20}\label{eq:2.20}
\end{equation*}
\begin{equation*}
\hat U_S(t,t_0)^\dagger  \hat U_S(t,t_0) = \hat 1
\tag{2.21}\label{eq:2.21}
\end{equation*}
\begin{equation*}
\hat U_S(t,t_0)^\dagger = \hat U_S(t_0,t)
\tag{2.22}\label{eq:2.22}
\end{equation*}

\subsection{Heisenberg Picture}\label{book-subsection-section-2.1-2}\label{heisenberg-picture}

\topic{The dynamics is carried entirely by the operators; states do not evolve in time.}

The Heisenberg picture is obtained from the Schrödinger picture by the transformation
\begin{equation*}
{\displaystyle 
\begin{gathered}
|\psi_H(t)\rangle:=\hat U_S(t,0)^\dagger\,|\psi_S(t)\rangle=|\psi_S(0)\rangle,\\[4pt]
\hat O_H(t):=\hat U_S(t,0)^\dagger\,\hat O_S(t)\,\hat U_S(t,0).
\end{gathered}
}
\tag{2.23}\label{eq:2.23}
\end{equation*}
Since the pictures are taken to coincide at $t=0$, we abbreviate $\hat U_S(t,0)$ as $\hat U_S$ below.

The two pictures give the same expectation values
\begin{equation*}
\langle\psi_H(t)| \hat O_H(t) |\psi_H(t)\rangle =
\langle\psi_H(t)|\hat U_S^\dagger \hat U_S \hat O_H(t) \hat U_S^\dagger \hat U_S|\psi_H(t)\rangle=\langle\psi_S(t)|\hat O_S(t)|\psi_S(t)\rangle
\tag{2.24}\label{eq:2.24}
\end{equation*}
We now compute the time dependence of states and operators in the Heisenberg picture
\begin{equation*}
\begin{aligned}
\frac{d\hat O_H}{dt}
&=\Bigl(\frac{\partial\hat U_S}{\partial t}\Bigr)^\dagger\hat O_S\,\hat U_S
+\hat U_S^\dagger\,\frac{d\hat O_S}{dt}\,\hat U_S
+\hat U_S^\dagger\hat O_S\,\frac{\partial\hat U_S}{\partial t}\\
&=i\,\hat U_S^\dagger\bigl[\hat H_S(t),\hat O_S\bigr]\hat U_S
+\hat U_S^\dagger\,\frac{\partial\hat O_S}{\partial t}\,\hat U_S\\
&=i\bigl[\hat H_H(t),\hat O_H\bigr]+\Bigl(\frac{\partial\hat O_S}{\partial t}\Bigr)_H,
\end{aligned}
\tag{2.25}\label{eq:2.25}
\end{equation*}
where $\hat H_H(t):=\hat U_S^\dagger\hat H_S(t)\,\hat U_S$ and $(\partial\hat O_S/\partial t)_H:=\hat U_S^\dagger(\partial\hat O_S/\partial t)\,\hat U_S$.

In cases (1) and (2), $\hat U_S(t,0)$ commutes with $\hat H_S(t)$, so $\hat H_H(t)=\hat H_S(t)$.

\setcounter{section}{1}
\section{Canonical Quantization}\label{book-section-2.2}\label{canonical-quantization}\label{canonical-commutation-relations}

Consider a field $\phi:\mathbb R^4\to\mathbb R^n$ without Dirac constraints, with components $\phi_a$.

The Lagrangian density is a smooth function
\begin{equation*}
\begin{aligned}
\mathcal L:\mathbb R^n\times(\mathbb R^n)^4&\to\mathbb R,\\
(x,y^\mu)&\mapsto\mathcal L(x,y^\mu),
\end{aligned}
\tag{2.26}\label{eq:2.26}
\end{equation*}
where $y^\mu\in\mathbb R^n$, $\mu=0,\ldots,3$.

Write $\mathcal L|_\phi$ and $\frac{\partial\mathcal L}{\partial y^\mu_a}\big|_\phi$ for $\mathcal L$ and its partial derivatives evaluated at $\bigl(\phi(x),\partial_\mu\phi(x)\bigr)$; they are functions of $x$.

The Lagrangian is
\begin{equation*}
L(t):=\int d^3x\,\mathcal L|_\phi(t,\mathbf x).
\tag{2.27}\label{eq:2.27}
\end{equation*}
At each fixed time, the field values form an infinite set of generalized coordinates:
\begin{equation*}
i\longleftrightarrow(a,\mathbf x),
\qquad
q_i(t)\longleftrightarrow\phi_a(t,\mathbf x).
\tag{2.28}\label{eq:2.28}
\end{equation*}

\topic{Conjugate momentum}

The momentum density conjugate to $\phi_a(t,\mathbf x)$ is
\begin{equation*}
\pi_a(t,\mathbf x)
:=
\frac{\partial\mathcal L}{\partial y^0_a}\Big|_\phi(t,\mathbf x).
\tag{2.29}\label{eq:2.29}
\end{equation*}

\topic{Poisson brackets and commutation relations}

The classical fields obey the equal-time canonical Poisson brackets
\begin{equation*}
\begin{aligned}
\{\phi_a(t,\mathbf x),\pi_b(t,\mathbf y)\}_{\mathrm{PB}}
&=\delta_{ab}\delta^{(3)}(\mathbf x-\mathbf y),\\
\{\phi_a(t,\mathbf x),\phi_b(t,\mathbf y)\}_{\mathrm{PB}}
&=\{\pi_a(t,\mathbf x),\pi_b(t,\mathbf y)\}_{\mathrm{PB}}=0.
\end{aligned}
\tag{2.30}\label{eq:2.30}
\end{equation*}
Canonical quantization promotes the fields and their conjugate momenta to operators and replaces these fundamental Poisson brackets according to
\begin{equation*}
\{\ ,\ \}_{\mathrm{PB}}
\longrightarrow
\frac{1}{i}[\ ,\ ].
\tag{2.31}\label{eq:2.31}
\end{equation*}
Thus,
\begin{equation*}
{\displaystyle 
\begin{aligned}
{}[\hat\phi_a(t,\mathbf x),\hat\pi_b(t,\mathbf y)]
&=i\,\delta_{ab}\delta^{(3)}(\mathbf x-\mathbf y),\\
{}[\hat\phi_a(t,\mathbf x),\hat\phi_b(t,\mathbf y)]
&=[\hat\pi_a(t,\mathbf x),\hat\pi_b(t,\mathbf y)]=0.
\end{aligned}}
\tag{2.32}\label{eq:2.32}
\end{equation*}
Here $\delta_{ab}\delta^{(3)}(\mathbf x-\mathbf y)$ generalizes the Kronecker delta $\delta_{ij}$ to the field labels $(a,\mathbf x)$.

\subsection{Quantization of the Free Real Scalar Field}\label{book-subsection-section-2.2-1}\label{free-real-scalar-field-quantization}

For a free field of mass $m$, use the Lorentz-invariant measure $d\Pi_{\mathbf p}$ and positive-energy four-momentum $p$:
\begin{equation*}
\begin{aligned}
E_{\mathbf p}&=\sqrt{\mathbf p^2+m^2},\qquad p^0=E_{\mathbf p},\qquad d\Pi_{\mathbf p}:=\frac{d^3p}{(2\pi)^3\,2E_{\mathbf p}},\\
\hat\phi_I(x)&=\int d\Pi_{\mathbf p}\left(\hat a_{\mathbf p}e^{-ip\cdot x}+\hat a_{\mathbf p}^\dagger e^{+ip\cdot x}\right).
\end{aligned}
\tag{2.33}\label{eq:2.33}
\end{equation*}
The vacuum is normalized to one. The creation operators define covariantly normalized one-particle states:
\begin{equation*}
{\displaystyle \begin{aligned}
{}[\hat a_{\mathbf p},\hat a_{\mathbf q}^\dagger]&=2E_{\mathbf p}(2\pi)^3\delta^{(3)}(\mathbf p-\mathbf q),\qquad
[\hat a_{\mathbf p},\hat a_{\mathbf q}]=[\hat a_{\mathbf p}^\dagger,\hat a_{\mathbf q}^\dagger]=0,\\
\hat a_{\mathbf p}|0\rangle&=0,\qquad |\mathbf p\rangle:=\hat a_{\mathbf p}^\dagger|0\rangle,\\
\langle\mathbf p|\mathbf q\rangle&=2E_{\mathbf p}(2\pi)^3\delta^{(3)}(\mathbf p-\mathbf q).
\end{aligned}}
\tag{2.34}\label{eq:2.34}
\end{equation*}
Normal ordering, denoted by colons, places creation operators to the left of annihilation operators. With the vacuum energy subtracted,
\begin{equation*}
\begin{aligned}
:\hat H_0:&=\int d\Pi_{\mathbf p}\,E_{\mathbf p}\,\hat a_{\mathbf p}^\dagger\hat a_{\mathbf p},\qquad
:\hat{\mathbf P}:=\int d\Pi_{\mathbf p}\,\mathbf p\,\hat a_{\mathbf p}^\dagger\hat a_{\mathbf p},\\
\hat N&=\int d\Pi_{\mathbf p}\,\hat a_{\mathbf p}^\dagger\hat a_{\mathbf p}.
\end{aligned}
\tag{2.35}\label{eq:2.35}
\end{equation*}

\setcounter{section}{2}
\section{Quantization with Dirac Brackets}\label{book-section-2.3}\label{quantization-with-dirac-brackets}

\subsection{Quantization of the Complex Scalar Field}\label{book-subsection-section-2.3-1}\label{complex-scalar-field-quantization}

\subsection{Quantization of the Dirac Field}\label{book-subsection-section-2.3-2}\label{dirac-field-quantization}

\topic{Free Dirac Field}

Using the spinors and their normalization in \hyperref[dirac-field]{Section 1.7}, the operators $\hat b_{\mathbf p}^s$ and $\hat c_{\mathbf p}^s$ annihilate particles and antiparticles, respectively:
\begin{equation*}
\hat\psi_I(x)=\sum_{s=1}^{2}\int d\Pi_{\mathbf p}\left(\hat b_{\mathbf p}^s u^s(p)e^{-ip\cdot x}+\hat c_{\mathbf p}^{s\dagger}v^s(p)e^{+ip\cdot x}\right).
\tag{2.36}\label{eq:2.36}
\end{equation*}
Fermionic quantization uses anticommutators; $\alpha,\beta$ below are Dirac component indices. All mode anticommutators not displayed vanish, and both annihilation operators annihilate $|0\rangle$.
\begin{equation*}
{\displaystyle \begin{aligned}
\{\hat b_{\mathbf p}^r,\hat b_{\mathbf q}^{s\dagger}\}&=\{\hat c_{\mathbf p}^r,\hat c_{\mathbf q}^{s\dagger}\}=2E_{\mathbf p}(2\pi)^3\delta^{rs}\delta^{(3)}(\mathbf p-\mathbf q),\\
\{\hat\psi_{I,\alpha}(t,\mathbf x),\hat\psi_{I,\beta}^\dagger(t,\mathbf y)\}&=\delta_{\alpha\beta}\delta^{(3)}(\mathbf x-\mathbf y),\\
\{\hat\psi_{I,\alpha}(t,\mathbf x),\hat\psi_{I,\beta}(t,\mathbf y)\}&=0.
\end{aligned}}
\tag{2.37}\label{eq:2.37}
\end{equation*}
Fermionic normal ordering includes a minus sign for every interchange of fermionic operators. The normal-ordered energy, momentum, and net fermion number are
\begin{equation*}
\begin{aligned}
:\hat H_0:&=\sum_s\int d\Pi_{\mathbf p}\,E_{\mathbf p}\left(\hat b_{\mathbf p}^{s\dagger}\hat b_{\mathbf p}^s+\hat c_{\mathbf p}^{s\dagger}\hat c_{\mathbf p}^s\right),\\
:\hat{\mathbf P}:&=\sum_s\int d\Pi_{\mathbf p}\,\mathbf p\left(\hat b_{\mathbf p}^{s\dagger}\hat b_{\mathbf p}^s+\hat c_{\mathbf p}^{s\dagger}\hat c_{\mathbf p}^s\right),\\
\hat Q&:=\int d^3x:\hat\psi_I^\dagger\hat\psi_I:
=\sum_s\int d\Pi_{\mathbf p}\left(\hat b_{\mathbf p}^{s\dagger}\hat b_{\mathbf p}^s-\hat c_{\mathbf p}^{s\dagger}\hat c_{\mathbf p}^s\right).
\end{aligned}
\tag{2.38}\label{eq:2.38}
\end{equation*}

\subsection{Quantization of the Proca Field}\label{book-subsection-section-2.3-3}\label{proca-field-quantization}

\setcounter{section}{3}
\section{Gupta–Bleuler Quantization}\label{book-section-2.4}\label{gupta-bleuler-quantization}

\setcounter{section}{4}
\section{Symmetry}\label{book-section-2.5}\label{symmetry}

\subsection{Ward Identity}\label{book-subsection-section-2.5-1}\label{ward-identity}

\setcounter{chapter}{2}
\chapter{Path Integrals}\label{book-chapter-3}\label{qft-ch3}

\setcounter{section}{0}
\section{Path Integrals in Quantum Mechanics}\label{book-section-3.1}\label{path-integrals-in-quantum-mechanics}

\topic{Principle}

Consider a system with $n$ degrees of freedom.
\begin{itemize}


\item \textbf{\boldmath Configuration space}: $\mathrm{Conf}=\mathbb R^n$.

\item \textbf{\boldmath Dynamical structure}: the Lagrangian
\begin{equation*}
\begin{aligned}
L:\mathbb R^n\times\mathbb R^n&\to\mathbb R,\\
(x,y)&\mapsto\sum_{a=1}^n\frac12m_ay_a^2-V(x),
\end{aligned}
\tag{3.1}\label{eq:3.1}
\end{equation*}
in which $m_a>0$ and $V:\mathbb R^n\to\mathbb R$ is smooth.

\item \textbf{\boldmath Paths}: for $t_i<t_f$,
\begin{equation*}
\begin{aligned}
C:[t_i,t_f]&\to\mathrm{Conf},\\
t&\mapsto C(t).
\end{aligned}
\tag{3.2}\label{eq:3.2}
\end{equation*}

\item \textbf{\boldmath Action}: the functional on paths
\begin{equation*}
\begin{aligned}
S:\{\text{paths}\}&\to\mathbb R,\\
C&\mapsto S[C]:=\int_{t_i}^{t_f}dt\,L\bigl(C(t),\dot C(t)\bigr).
\end{aligned}
\tag{3.3}\label{eq:3.3}
\end{equation*}
\end{itemize}


\topic{Quantization}
\begin{itemize}


\item \textbf{\boldmath Initial configuration}:
\begin{equation*}
\begin{aligned}
Q_i:\{t_i\}&\to\mathrm{Conf},\\
t_i&\mapsto q_i.
\end{aligned}
\tag{3.4}\label{eq:3.4}
\end{equation*}

\item \textbf{\boldmath Final configuration}:
\begin{equation*}
\begin{aligned}
Q_f:\{t_f\}&\to\mathrm{Conf},\\
t_f&\mapsto q_f.
\end{aligned}
\tag{3.5}\label{eq:3.5}
\end{equation*}

\item \textbf{\boldmath Transition kernel}: the integral over the paths with these endpoints,
\begin{equation*}
{\displaystyle K(Q_f,Q_i):=\int_{C(t_i)=Q_i(t_i)}^{C(t_f)=Q_f(t_f)}\mathcal DC\,e^{iS[C]},}
\tag{3.6}\label{eq:3.6}
\end{equation*}
also written $\int_{Q_i}^{Q_f}\mathcal DC\,e^{iS[C]}$.
\end{itemize}


\topic{Exact Definition of the Transition Kernel}

Divide $[t_i,t_f]$ into $N+1$ steps.
\begin{itemize}


\item \textbf{\boldmath Measure}:
\begin{equation*}
\begin{gathered}
\Delta t:=\frac{t_f-t_i}{N+1},\qquad t_k:=t_i+k\,\Delta t\quad(k=0,\ldots,N+1),\\
\mathcal DC:=\lim_{N\to\infty}\mathcal N_N\prod_{k=1}^Nd^nC(t_k),\qquad
\mathcal N_N:=\prod_{a=1}^n\Bigl(\frac{m_a}{2\pi i\,\Delta t}\Bigr)^{\frac{N+1}{2}},
\end{gathered}
\tag{3.7}\label{eq:3.7}
\end{equation*}
so that $t_0=t_i$ and $t_{N+1}=t_f$, with $d^nC(t_k):=\prod_{a=1}^ndC_a(t_k)$.

\item \textbf{\boldmath Action}:
\begin{equation*}
S[C]=\lim_{N\to\infty}\Delta t\sum_{k=0}^NL\Bigl(C(t_k),\frac{C(t_{k+1})-C(t_k)}{\Delta t}\Bigr).
\tag{3.8}\label{eq:3.8}
\end{equation*}
\end{itemize}

Thus \eqref{eq:3.6} means
\begin{equation*}
{\displaystyle 
K(Q_f,Q_i)=\lim_{N\to\infty}\mathcal N_N\int\prod_{k=1}^Nd^nC(t_k)\,
\exp\biggl[i\,\Delta t\sum_{k=0}^NL\Bigl(C(t_k),\frac{C(t_{k+1})-C(t_k)}{\Delta t}\Bigr)\biggr],}
\tag{3.9}\label{eq:3.9}
\end{equation*}
with $C(t_0):=q_i$, $C(t_{N+1}):=q_f$, and each $C(t_k)$, $1\le k\le N$, integrated over $\mathrm{Conf}$.

\topic{Functional Derivative}

For a source $J:[t_i,t_f]\to\mathbb R^n$, with $J\cdot C:=\sum_{a=1}^nJ_aC_a$,
\begin{equation*}
K^J(Q_f,Q_i):=\int_{Q_i}^{Q_f}\mathcal DC\,\exp\Bigl(i\int_{t_i}^{t_f}dt\,\bigl[L\bigl(C(t),\dot C(t)\bigr)+J(t)\cdot C(t)\bigr]\Bigr).
\tag{3.10}\label{eq:3.10}
\end{equation*}
In the sliced form \eqref{eq:3.9} the source adds $i\,\Delta t\sum_{k=0}^NJ(t_k)\cdot C(t_k)$ to the exponent, so $J$ enters only through $J(t_0),\ldots,J(t_N)$. For $t_i<t<t_f$ define
\begin{equation*}
\frac{\delta}{\delta J_a(t)}:=\lim_{N\to\infty}\frac{1}{\Delta t}\frac{\partial}{\partial J_a(t_k)},\qquad t_k\to t.
\tag{3.11}\label{eq:3.11}
\end{equation*}
Writing $L_j:=L\bigl(C(t_j),[C(t_{j+1})-C(t_j)]/\Delta t\bigr)$,
\begin{equation*}
\begin{aligned}
\frac{\delta}{\delta J_a(t)}K^J(Q_f,Q_i)
&=\lim_{N\to\infty}\mathcal N_N\int\prod_{j=1}^Nd^nC(t_j)\,
\frac{1}{\Delta t}\frac{\partial}{\partial J_a(t_k)}\exp\biggl[i\,\Delta t\sum_{j=0}^N\bigl(L_j+J(t_j)\cdot C(t_j)\bigr)\biggr]\\
&=\lim_{N\to\infty}\mathcal N_N\int\prod_{j=1}^Nd^nC(t_j)\,
iC_a(t_k)\exp\biggl[i\,\Delta t\sum_{j=0}^N\bigl(L_j+J(t_j)\cdot C(t_j)\bigr)\biggr]\\
&=\int_{Q_i}^{Q_f}\mathcal DC\,iC_a(t)\exp\Bigl(i\int_{t_i}^{t_f}dt'\,\bigl[L\bigl(C(t'),\dot C(t')\bigr)+J(t')\cdot C(t')\bigr]\Bigr).
\end{aligned}
\tag{3.12}\label{eq:3.12}
\end{equation*}

\topic{Example: Free Particle in One Dimension}

Take $n=1$, $V=0$ in \eqref{eq:3.1}:
\begin{equation*}
L(x,y)=\frac12my^2,\qquad Q_i(t_i)=q_i,\qquad Q_f(t_f)=q_f.
\tag{3.13}\label{eq:3.13}
\end{equation*}
By \eqref{eq:3.7}–\eqref{eq:3.9}, with $C(t_0)=q_i$, $C(t_{N+1})=q_f$ and $\mathcal N_N=\bigl(\frac{m}{2\pi i\,\Delta t}\bigr)^{\frac{N+1}{2}}$,
\begin{equation*}
\begin{aligned}
K(Q_f,Q_i)&=\int_{Q_i}^{Q_f}\mathcal DC\,e^{i\int_{t_i}^{t_f}dt\,\frac12m\dot C^2}\\
&=\lim_{N\to\infty}\mathcal N_N\int\prod_{k=1}^{N}dC(t_k)\,e^{\frac{i}{\Delta t}\sum_{k=0}^{N}\frac12m\left(C(t_{k+1})-C(t_k)\right)^2}\\
&=\lim_{N\to\infty}\mathcal N_N\int\prod_{k=1}^{N}dC(t_k)\,
e^{\frac{i}{\Delta t}\left[\frac12mq_i^2+\frac12mq_f^2+mC^2(t_1)+\cdots+mC^2(t_N)\right]}\\
&\qquad{}\times e^{\frac{i}{\Delta t}(-m)\left[q_iC(t_1)+C(t_1)C(t_2)+\cdots+C(t_N)q_f\right]}.
\end{aligned}
\tag{3.14}\label{eq:3.14}
\end{equation*}
Integrate $C(t_1)$, using $\int_{-\infty}^{\infty}dx\,e^{iax^2}=\sqrt{\pi i/a}$ for $a>0$ ($a\to a+i\varepsilon$, $\varepsilon\to0^+$):
\begin{equation*}
\begin{aligned}
\int dC(t_1)\,e^{\frac{i}{\Delta t}\left[mC^2(t_1)-mC(t_1)\left(q_i+C(t_2)\right)\right]}
&=\int dC(t_1)\,e^{\frac{im}{\Delta t}\left(C(t_1)-\frac{q_i+C(t_2)}{2}\right)^2}\,e^{-\frac{im}{4\Delta t}\left(q_i+C(t_2)\right)^2}\\
&=\sqrt{\frac{2\pi i\,\Delta t}{m}\cdot\frac12}\;e^{-\frac{im}{4\Delta t}\left(q_i+C(t_2)\right)^2}.
\end{aligned}
\tag{3.15}\label{eq:3.15}
\end{equation*}
After $C(t_1),\ldots,C(t_{j-1})$, $1\le j\le N$:
\begin{equation*}
\begin{aligned}
K(Q_f,Q_i)&=\lim_{N\to\infty}\mathcal N_N\Bigl(\frac{2\pi i\,\Delta t}{m}\Bigr)^{\frac{j-1}{2}}\frac{1}{\sqrt j}
\int\prod_{k=j}^{N}dC(t_k)\\
&\qquad{}\times 
e^{\frac{im}{2\Delta t}\left[\frac{\left(C(t_j)-q_i\right)^2}{j}+C^2(t_j)+2C^2(t_{j+1})+\cdots+2C^2(t_N)+q_f^2\right]}\\
&\qquad{}\times e^{\frac{im}{2\Delta t}(-2)\left[C(t_j)C(t_{j+1})+\cdots+C(t_N)q_f\right]}.
\end{aligned}
\tag{3.16}\label{eq:3.16}
\end{equation*}
Integrate $C(t_j)$:
\begin{equation*}
\begin{aligned}
&\frac{\left(C(t_j)-q_i\right)^2}{j}+C^2(t_j)-2C(t_j)C(t_{j+1})
&\\[-1pt]
&\quad{}=\frac{j+1}{j}\Bigl(C(t_j)-\frac{q_i+jC(t_{j+1})}{j+1}\Bigr)^2+\frac{\left(C(t_{j+1})-q_i\right)^2}{j+1}-C^2(t_{j+1}),\\
&\int dC(t_j)\,e^{\frac{im}{2\Delta t}\frac{j+1}{j}\left(C(t_j)-\frac{q_i+jC(t_{j+1})}{j+1}\right)^2}\,=\sqrt{\frac{2\pi i\,\Delta t}{m}\cdot\frac{j}{j+1}},
\end{aligned}
\tag{3.17}\label{eq:3.17}
\end{equation*}
which gives \eqref{eq:3.16} with $j+1$; for $j=N$, $C(t_{N+1})=q_f$ and the bracket becomes $\frac{(q_f-q_i)^2}{N+1}$. Hence, with $(N+1)\Delta t=t_f-t_i$,
\begin{equation*}
\begin{aligned}
K(Q_f,Q_i)&=\lim_{N\to\infty}\Bigl(\frac{m}{2\pi i\,\Delta t}\Bigr)^{\frac{N+1}{2}}\Bigl(\frac{2\pi i\,\Delta t}{m}\Bigr)^{\frac N2}\frac{1}{\sqrt{N+1}}\,
e^{\frac{im\left(q_f-q_i\right)^2}{2(N+1)\Delta t}}\\
&={\displaystyle \sqrt{\frac{m}{2\pi i\,(t_f-t_i)}}\;e^{\frac{im\left(q_f-q_i\right)^2}{2(t_f-t_i)}}.}
\end{aligned}
\tag{3.18}\label{eq:3.18}
\end{equation*}
This example shows why the prefactor $\mathcal N_N$ is included in the definition \eqref{eq:3.7} of the path integral. The result \eqref{eq:3.18} is finite, while by its definition $\mathcal N_N$ diverges as $N\to\infty$; without the prefactor, the remaining factor $\bigl(\frac{2\pi i\,\Delta t}{m}\bigr)^{\frac N2}\frac{1}{\sqrt{N+1}}$ in \eqref{eq:3.18} tends to zero, and so would $K(Q_f,Q_i)$.

\topic{Relation to Operator Quantization}

Let $\hat H:=\sum_{a=1}^n\hat p_a^2/(2m_a)+V(\hat q)$ with $[\hat q_a,\hat p_b]=i\delta_{ab}$ and $\hat q_a|q\rangle=q_a|q\rangle$, and let $\hat U_S(t_f,t_i)=e^{-i\hat H(t_f-t_i)}$ be its evolution operator \eqref{eq:2.18}. Then, with $q_i=Q_i(t_i)$ and $q_f=Q_f(t_f)$,
\begin{equation*}
{\displaystyle K(Q_f,Q_i)=\langle q_f|\hat U_S(t_f,t_i)|q_i\rangle=\langle q_f,t_f|q_i,t_i\rangle}
\tag{3.19}\label{eq:3.19}
\end{equation*}
in the Schrödinger and Heisenberg pictures respectively, where $|q,t\rangle:=\hat U_S(t,0)^\dagger|q\rangle$ satisfies $\hat q_H(t)|q,t\rangle=q\,|q,t\rangle$ by \eqref{eq:2.23}.

\setcounter{section}{1}
\section{Path Integrals in Quantum Field Theory}\label{book-section-3.2}\label{path-integrals-in-field-theory}

Following the second viewpoint in \hyperref[classical-dynamics]{Section 1.4}, $t$ is replaced by $x^\mu$.

\topic{Principle}
\begin{itemize}


\item \textbf{\boldmath Configuration space}: $\mathrm{Conf}=\mathbb R^n$.

\item \textbf{\boldmath Dynamical structure}: the Lagrangian density
\begin{equation*}
\begin{aligned}
\mathcal L:\mathbb R^n\times(\mathbb R^n)^4&\to\mathbb R,\\
(x,y^\mu)&\mapsto\sum_{a=1}^n\frac12y_a^\mu y_{a\mu}-V(x),
\end{aligned}
\tag{3.20}\label{eq:3.20}
\end{equation*}
in which $y^\mu\in\mathbb R^n$, $\mu=0,\ldots,3$, and $V:\mathbb R^n\to\mathbb R$ is smooth.

\item \textbf{\boldmath Paths}: for $t_i<t_f$,
\begin{equation*}
\begin{aligned}
\phi:[t_i,t_f]\times\mathbb R^3&\to\mathrm{Conf},\\
x&\mapsto\phi(x),
\end{aligned}
\tag{3.21}\label{eq:3.21}
\end{equation*}
with $\phi(t):=\phi(t,\cdot\,)$ the configuration at time $t$.

\item \textbf{\boldmath Action}: the functional on paths
\begin{equation*}
\begin{aligned}
S:\{\text{paths}\}&\to\mathbb R,\\
\phi&\mapsto S[\phi]:=\int_{\mathbb R^3}d^3x\int_{t_i}^{t_f}dt\,\mathcal L\bigl(\phi(x),\partial_\mu\phi(x)\bigr).
\end{aligned}
\tag{3.22}\label{eq:3.22}
\end{equation*}
\end{itemize}

In the path integral the field takes the place of the path: a path \eqref{eq:3.21} is a classical field \eqref{eq:1.44} restricted to $[t_i,t_f]\times\mathbb R^3$. Paths are therefore denoted by $\phi$, and the integrals below are over field configurations.

Below, $\int d^4x:=\int_{t_i}^{t_f}dt\int_{\mathbb R^3}d^3x$.

\topic{Quantization}
\begin{itemize}


\item \textbf{\boldmath Initial configuration}:
\begin{equation*}
\begin{aligned}
\Phi_i:\{t_i\}\times\mathbb R^3&\to\mathrm{Conf},\\
(t_i,\mathbf x)&\mapsto\Phi_i(t_i,\mathbf x).
\end{aligned}
\tag{3.23}\label{eq:3.23}
\end{equation*}

\item \textbf{\boldmath Final configuration}:
\begin{equation*}
\begin{aligned}
\Phi_f:\{t_f\}\times\mathbb R^3&\to\mathrm{Conf},\\
(t_f,\mathbf x)&\mapsto\Phi_f(t_f,\mathbf x).
\end{aligned}
\tag{3.24}\label{eq:3.24}
\end{equation*}

\item \textbf{\boldmath Transition kernel}:
\begin{equation*}
{\displaystyle K(\Phi_f,\Phi_i):=\int_{\phi(t_i,\mathbf x)=\Phi_i(t_i,\mathbf x)}^{\phi(t_f,\mathbf x)=\Phi_f(t_f,\mathbf x)}\mathcal D\phi\,e^{iS[\phi]},}
\tag{3.25}\label{eq:3.25}
\end{equation*}
also written $\int_{\Phi_i}^{\Phi_f}\mathcal D\phi\,e^{iS[\phi]}$; the boundary conditions hold for every $\mathbf x\in\mathbb R^3$.
\end{itemize}


\topic{Correlation Functions}

As in \eqref{eq:3.19},
\begin{equation*}
K(\Phi_f,\Phi_i)=\langle\Phi_f|\hat U_S(t_f,t_i)|\Phi_i\rangle,
\tag{3.26}\label{eq:3.26}
\end{equation*}
In operator quantization the vacuum expands in this basis, which defines the vacuum functional $\Psi[\Phi]$:
\begin{equation*}
|\mathrm{vac}\rangle=\int\mathcal D\Phi\,|\Phi\rangle\langle\Phi|\mathrm{vac}\rangle=\int\mathcal D\Phi\,|\Phi\rangle\,\Psi[\Phi],\qquad
\Psi[\Phi]:=\langle\Phi|\mathrm{vac}\rangle.
\tag{3.27}\label{eq:3.27}
\end{equation*}
Hence
\begin{equation*}
\begin{aligned}
K(\mathrm{vac}_f,\mathrm{vac}_i)&:=\langle\mathrm{vac}_f|\hat U_S(t_f,t_i)|\mathrm{vac}_i\rangle\\
&=\int\mathcal D\Phi_f\,\mathcal D\Phi_i\,\Psi^*[\Phi_f]\,\Psi[\Phi_i]\,K(\Phi_f,\Phi_i).
\end{aligned}
\tag{3.28}\label{eq:3.28}
\end{equation*}
\begin{itemize}


\item \textbf{\boldmath Correlation function}: 
\begin{equation*}
{\displaystyle 
\begin{aligned}
G_{a_1\cdots a_r}(x_1,\ldots,x_r)&:=\langle \phi_{a_1}(x_1)\cdots\phi_{a_r}(x_r)\rangle_{\mathrm{vac}}\\
&=\frac{\int_{\mathrm{vac}}^{\mathrm{vac}}\mathcal D\phi\,\phi_{a_1}(x_1)\cdots\phi_{a_r}(x_r)\,e^{iS[\phi]}}{\int_{\mathrm{vac}}^{\mathrm{vac}}\mathcal D\phi\,e^{iS[\phi]}},
\end{aligned}}
\tag{3.29}\label{eq:3.29}
\end{equation*}
with $t_i<x_k^0<t_f$.

上述定义中的分母源自于对真空归一化的要求，在上述定义下，\textless{}1\textgreater{}=真空的模=1

\item \textbf{\boldmath Generating functional}: for a source $J:[t_i,t_f]\times\mathbb R^3\to\mathbb R^n$, with $J\cdot\phi$ as in \eqref{eq:3.10},
\begin{equation*}
K^J(\Phi_f,\Phi_i):=\int_{\Phi_i}^{\Phi_f}\mathcal D\phi\,\exp\Bigl(iS[\phi]+i\int d^4x\,J(x)\cdot\phi(x)\Bigr),
\tag{3.30}\label{eq:3.30}
\end{equation*}
and, with the vacuum endpoints as in \eqref{eq:3.28},
\begin{equation*}
{\displaystyle 
Z[J]:=K^J(\mathrm{vac}_f,\mathrm{vac}_i)
=\int\mathcal D\Phi_f\,\mathcal D\Phi_i\,\Psi^*[\Phi_f]\,\Psi[\Phi_i]\,K^J(\Phi_f,\Phi_i),}
\tag{3.31}\label{eq:3.31}
\end{equation*}
so that $Z[0]=K(\mathrm{vac}_f,\mathrm{vac}_i)$.
\end{itemize}


As in \eqref{eq:3.12}, each $\delta/\delta J_a(x)$ acting on \eqref{eq:3.30} brings down $i\phi_a(x)$, so by \eqref{eq:3.31} and \eqref{eq:3.29}
\begin{equation*}
{\displaystyle G_{a_1\cdots a_r}(x_1,\ldots,x_r)=\frac{1}{Z[0]}\left.\Bigl(-i\frac{\delta}{\delta J_{a_1}(x_1)}\Bigr)\cdots\Bigl(-i\frac{\delta}{\delta J_{a_r}(x_r)}\Bigr)Z[J]\right|_{J=0}.}
\tag{3.32}\label{eq:3.32}
\end{equation*}

\topic{Connected Functions}

For a single real scalar field,
\begin{equation*}
\begin{aligned}
G_n(x_1,\ldots,x_n)&:=\langle\Omega|\mathcal T[\hat\phi_H(x_1)\cdots\hat\phi_H(x_n)]|\Omega\rangle
=\left.\frac{1}{Z[0]}\frac{1}{i^n}\frac{\delta^n Z[J]}{\delta J(x_1)\cdots\delta J(x_n)}\right|_{J=0},\\
W[J]&:=-i\ln Z[J],\qquad
G_{n,\mathrm c}=\left.\frac{1}{i^{n-1}}\frac{\delta^nW}{\delta J(x_1)\cdots\delta J(x_n)}\right|_{J=0}.
\end{aligned}
\tag{3.33}\label{eq:3.33}
\end{equation*}
Here $|\Omega\rangle$ is the normalized interacting vacuum and $J(x)$ is a commuting external source. Functional integrals require a regulator; the expressions below define their perturbative use.

\topic{Effective Action and Schwinger--Dyson Equation}

The Legendre transform generates one-particle-irreducible (1PI) vertices, where a 1PI graph stays connected after cutting any single internal line:
\begin{equation*}
\phi_c(x):=\frac{\delta W}{\delta J(x)},\qquad
\Gamma[\phi_c]:=W[J]-\int d^4x\,J\phi_c,\qquad
{\displaystyle \frac{\delta\Gamma}{\delta\phi_c(x)}=-J(x).}
\tag{3.34}\label{eq:3.34}
\end{equation*}
For a translation-invariant scalar functional measure, with functional boundary terms absent,
\begin{equation*}
\left[
\left.\frac{\delta S}{\delta\phi(x)}\right|_{\phi\,\to\,(1/i)\delta/\delta J}
+J(x)\right]Z[J]=0.
\tag{3.35}\label{eq:3.35}
\end{equation*}

\topic{Euclidean Continuation and Fermions}

When the time contour can be rotated without crossing singularities, $t=-i\tau$ gives, for $\mathcal L=\tfrac12\partial_\mu\phi\partial^\mu\phi-V(\phi)$,
\begin{equation*}
e^{iS}\longrightarrow e^{-S_E},\qquad
S_E=\int d\tau\,d^3x\left[\frac12(\partial_\tau\phi)^2
+\frac12(\boldsymbol\nabla\phi)^2+V(\phi)\right].
\tag{3.36}\label{eq:3.36}
\end{equation*}
Fermionic fields in the functional integral are independent Grassmann variables. With a fixed measure ordering, a quadratic Dirac action gives
\begin{equation*}
\int\mathcal D\bar\psi\,\mathcal D\psi\,
e^{i\int d^4x\,\bar\psi(i\gamma^\mu\partial_\mu-m)\psi}\ \propto\ \det(i\gamma^\mu\partial_\mu-m),
\tag{3.37}\label{eq:3.37}
\end{equation*}
where the field-independent normalization is absorbed into $Z[0]$ and the inverse of $i\gamma^\mu\partial_\mu-m$ uses the Feynman prescription.

\topic{Gauge Fixing and Faddeev--Popov Ghosts}

For perturbative $SU(3)$ theory near the trivial gauge configuration, Lorenz gauge introduces the gauge parameter $\xi$ and anticommuting scalar ghost fields $c^a,\bar c^a$:
\begin{equation*}
\begin{aligned}
\mathcal L_{\mathrm{gf}}&=-\frac{1}{2\xi}(\partial^\mu A^a_\mu)^2,\\
\mathcal L_{\mathrm{gh}}&=-\bar c^a\partial^\mu(D_\mu^{ab}c^b),\qquad
D_\mu^{ab}:=\delta^{ab}\partial_\mu+g_s f^{acb}A^c_\mu,\\
\Delta_{\mathrm{FP}}[A]&=\det[-\partial^\mu D_\mu^{ab}].
\end{aligned}
\tag{3.38}\label{eq:3.38}
\end{equation*}
For an Abelian gauge field in this linear gauge, the determinant is field independent and the ghosts decouple. This local perturbative construction does not resolve global gauge-copy ambiguities.

\setcounter{section}{2}
\section{\texorpdfstring{$\phi^3$}{φ3} Theory}\label{book-section-3.3}\label{phi-cubed-theory}

\topic{Action}

Take $n=1$, $\mathrm{Conf}=\mathbb R$ and $V(\phi)=\frac12m^2\phi^2+\frac1{3!}\lambda\phi^3$ in \eqref{eq:3.20}. With $\Box:=\partial_\mu\partial^\mu$ as in \eqref{eq:1.97},
\begin{equation*}
\begin{aligned}
S[\phi]&=\int d^4x\,\mathcal L\bigl(\phi(x),\partial_\mu\phi(x)\bigr)\\
&=\int d^4x\,\Bigl[\frac12\partial_\mu\phi\,\partial^\mu\phi-\frac12m^2\phi^2-\frac1{3!}\lambda\phi^3\Bigr]\\
&=\int d^4x\,\Bigl[-\frac12\phi(\Box+m^2)\phi-\frac1{3!}\lambda\phi^3\Bigr]
+\frac12\int d^3x\,\bigl[\phi\,\partial_0\phi\bigr]_{x^0=t_i}^{x^0=t_f}.
\end{aligned}
\tag{3.39}\label{eq:3.39}
\end{equation*}
$S_0[\phi]$ denotes \eqref{eq:3.39} at $\lambda=0$.

\topic{Feynman Propagator}
\begin{equation*}
\begin{aligned}
(\Box+m^2)\Delta_F(x)&=-i\delta^{(4)}(x),\\
\Delta_F(x)&=\int\frac{d^4p}{(2\pi)^4}\,\frac{-i}{-p^2+m^2-i\varepsilon}\,e^{-ip\cdot x}
=\int\frac{d^4p}{(2\pi)^4}\,\frac{i}{p^2-m^2+i\varepsilon}\,e^{-ip\cdot x},
\end{aligned}
\tag{3.40}\label{eq:3.40}
\end{equation*}
with $\varepsilon\to0^+$; this is the function \eqref{eq:4.12}.

\topic{Free Generating Functional}

From here on $\mathrm{vac}$ is the vacuum of $S_0$, and $J$ vanishes near $x^0=t_i$ and $x^0=t_f$. Let
\begin{equation*}
K_0^J(\Phi_f,\Phi_i):=\int_{\Phi_i}^{\Phi_f}\mathcal D\phi\,\exp\Bigl(iS_0[\phi]+i\int d^4x\,J(x)\phi(x)\Bigr),
\qquad K_0:=K_0^J\big|_{J=0}.
\tag{3.41}\label{eq:3.41}
\end{equation*}
In $K_0^J(\mathrm{vac}_f,\mathrm{vac}_i)$ shift the paths, together with $\Phi_f$ and $\Phi_i$, by
\begin{equation*}
\phi(x)\to\phi(x)+i\int d^4y\,\Delta_F(x-y)J(y),
\tag{3.42}\label{eq:3.42}
\end{equation*}
which leaves the measures in \eqref{eq:3.28} unchanged. By \eqref{eq:3.40}, the integrand of $\int d^4x$ in the exponent becomes
\begin{equation*}
\begin{aligned}
&-\frac12\phi(\Box+m^2)\phi+J\phi\\
&\to-\frac12\phi(\Box+m^2)\phi+J\phi
\underbrace{-\frac i2\int d^4y\,\Delta_F(x-y)J(y)\,(\Box+m^2)\phi(x)}_{-\frac12J\phi}\\
&\qquad+J(x)\,i\int d^4y\,\Delta_F(x-y)J(y)
\underbrace{-\frac12\phi(x)\,i\int d^4y\,\bigl(-i\delta^{(4)}(x-y)\bigr)J(y)}_{-\frac12J\phi}\\
&\qquad\underbrace{+\frac12\int d^4y\,\Delta_F(x-y)J(y)\,(\Box+m^2)\int d^4\bar y\,\Delta_F(x-\bar y)J(\bar y)}_{-\frac i2J\Delta_FJ}\\
&\to-\frac12\phi(\Box+m^2)\phi+\frac i2\int d^4y\,J(y)\Delta_F(x-y)J(x),
\end{aligned}
\tag{3.43}\label{eq:3.43}
\end{equation*}
where the first underbrace holds after integrating by parts in $x$, up to terms at $x^0=t_i,t_f$. Write $\chi(x):=i\int d^4y\,\Delta_F(x-y)J(y)$ for the shift \eqref{eq:3.42}, so that $(\Box+m^2)\chi=J$ by \eqref{eq:3.40}. Including the terms at $x^0=t_i,t_f$ from \eqref{eq:3.39} and from this integration by parts, the exponent of \eqref{eq:3.41} after the shift is exactly
\begin{equation*}
\begin{aligned}
iS_0[\phi+\chi]+i\int d^4x\,J(\phi+\chi)
&=iS_0[\phi]-\frac12\int d^4x\,d^4y\,J(x)\Delta_F(x-y)J(y)\\
&\quad+i\int d^3x\,\Bigl[\bigl(\phi+\tfrac12\chi\bigr)\partial_0\chi\Bigr]_{x^0=t_i}^{x^0=t_f}.
\end{aligned}
\tag{3.44}\label{eq:3.44}
\end{equation*}
Closing the $p^0$ contour of \eqref{eq:3.40} in the lower half-plane for $x^0>0$ and in the upper half-plane for $x^0<0$, around the poles $p^0=\pm(E_{\mathbf p}-i\varepsilon')$ with $E_{\mathbf p}$ as in \eqref{eq:2.33} and $\varepsilon'\to0^+$,
\begin{equation*}
\Delta_F(x)=\int\frac{d^3p}{(2\pi)^3\,2E_{\mathbf p}}\,e^{i\mathbf p\cdot\mathbf x-iE_{\mathbf p}|x^0|}.
\tag{3.45}\label{eq:3.45}
\end{equation*}
Let $\mathcal E:=(-\nabla^2+m^2)^{1/2}$, i.e. $\mathcal E\,e^{i\mathbf p\cdot\mathbf x}=E_{\mathbf p}\,e^{i\mathbf p\cdot\mathbf x}$. Since $J$ vanishes near $t_i$ and $t_f$, \eqref{eq:3.45} gives
\begin{equation*}
\partial_0\chi\big|_{x^0=t_f}=-i\mathcal E\chi\big|_{x^0=t_f},\qquad
\partial_0\chi\big|_{x^0=t_i}=+i\mathcal E\chi\big|_{x^0=t_i},
\tag{3.46}\label{eq:3.46}
\end{equation*}
so $\chi$ has only positive frequencies at $t_f$ and only negative frequencies at $t_i$.

For $S_0$ the vacuum is $|0\rangle$ of \eqref{eq:2.34}. At $t=0$ the Schrödinger-picture operators are $\hat\phi(\mathbf x)=\hat\phi_I(0,\mathbf x)$ and $\hat\pi(\mathbf x)=\partial_0\hat\phi_I(0,\mathbf x)$ by \eqref{eq:4.2} and \eqref{eq:1.92}. Inverting \eqref{eq:2.33}, and acting with $\hat\pi(\mathbf x)$ as $-i\,\delta/\delta\Phi(\mathbf x)$ on $\Psi[\Phi]=\langle\Phi|0\rangle$,
\begin{equation*}
\begin{aligned}
\hat a_{\mathbf p}&=\int d^3x\,e^{-i\mathbf p\cdot\mathbf x}\bigl[E_{\mathbf p}\hat\phi(\mathbf x)+i\hat\pi(\mathbf x)\bigr],\\
0=\langle\Phi|\hat a_{\mathbf p}|0\rangle&=\int d^3x\,e^{-i\mathbf p\cdot\mathbf x}\Bigl[(\mathcal E\Phi)(\mathbf x)+\frac{\delta}{\delta\Phi(\mathbf x)}\Bigr]\Psi[\Phi]\quad\text{for all }\mathbf p,\\
\Psi[\Phi]&=\Psi[0]\exp\Bigl(-\frac12\int d^3x\,\Phi\,\mathcal E\,\Phi\Bigr),
\end{aligned}
\tag{3.47}\label{eq:3.47}
\end{equation*}
where the second line uses $\int d^3x\,f\,\mathcal E\,g=\int d^3x\,(\mathcal Ef)\,g$ and states $\delta\Psi/\delta\Phi(\mathbf x)=-(\mathcal E\Phi)(\mathbf x)\Psi[\Phi]$. With $\chi_{f,i}:=\chi(t_{f,i},\cdot\,)$, the shift turns $\Psi^*[\Phi_f]\Psi[\Phi_i]$ into $\Psi^*[\Phi_f+\chi_f]\Psi[\Phi_i+\chi_i]$, while $\phi(t_{f,i},\cdot\,)=\Phi_{f,i}$. By \eqref{eq:3.46} and \eqref{eq:3.47},
\begin{equation*}
\begin{aligned}
i\int d^3x\,\Bigl[\bigl(\phi+\tfrac12\chi\bigr)\partial_0\chi\Bigr]_{x^0=t_i}^{x^0=t_f}
&=\int d^3x\,\Bigl[\bigl(\Phi_f+\tfrac12\chi_f\bigr)\mathcal E\chi_f+\bigl(\Phi_i+\tfrac12\chi_i\bigr)\mathcal E\chi_i\Bigr],\\
\ln\frac{\Psi^*[\Phi_f+\chi_f]\,\Psi[\Phi_i+\chi_i]}{\Psi^*[\Phi_f]\,\Psi[\Phi_i]}
&=-\int d^3x\,\Bigl[\bigl(\Phi_f+\tfrac12\chi_f\bigr)\mathcal E\chi_f+\bigl(\Phi_i+\tfrac12\chi_i\bigr)\mathcal E\chi_i\Bigr],
\end{aligned}
\tag{3.48}\label{eq:3.48}
\end{equation*}
so the terms at $x^0=t_i,t_f$ in \eqref{eq:3.44} cancel the change of $\Psi^*[\Phi_f]\Psi[\Phi_i]$. Hence
\begin{equation*}
K_0^J(\mathrm{vac}_f,\mathrm{vac}_i)=K_0(\mathrm{vac}_f,\mathrm{vac}_i)\,
\exp\Bigl[-\frac12\int d^4x\,d^4y\,J(x)\Delta_F(x-y)J(y)\Bigr],
\tag{3.49}\label{eq:3.49}
\end{equation*}
and
\begin{equation*}
{\displaystyle 
\begin{aligned}
Z_0[J]&:=K_0^J(\mathrm{vac}_f,\mathrm{vac}_i)
=\int\mathcal D\Phi_f\,\mathcal D\Phi_i\,\Psi^*[\Phi_f]\,\Psi[\Phi_i]\,K_0^J(\Phi_f,\Phi_i)\\
&=Z_0[0]\,\exp\Bigl[-\frac12\int d^4x\,d^4y\,J(x)\Delta_F(x-y)J(y)\Bigr],
\end{aligned}}
\tag{3.50}\label{eq:3.50}
\end{equation*}
with $Z_0[0]=K_0(\mathrm{vac}_f,\mathrm{vac}_i)$ by \eqref{eq:3.49}.

\topic{Two-Point Function}

For $S_0$ and $t_i<z_1^0,z_2^0<t_f$, by \eqref{eq:3.32} and \eqref{eq:3.50},
\begin{equation*}
{\displaystyle 
\begin{aligned}
\langle\phi(z_1)\phi(z_2)\rangle_{\mathrm{vac}}
&=\frac{1}{Z_0[0]}\Bigl(-i\frac{\delta}{\delta J(z_1)}\Bigr)\Bigl(-i\frac{\delta}{\delta J(z_2)}\Bigr)Z_0[J]\Big|_{J=0}\\
&=\Bigl(-i\frac{\delta}{\delta J(z_1)}\Bigr)\Bigl(-i\frac{\delta}{\delta J(z_2)}\Bigr)
\exp\Bigl[-\frac12\int d^4x\,d^4y\,J(x)\Delta_F(x-y)J(y)\Bigr]\Big|_{J=0}\\
&=\frac12\Delta_F(z_1-z_2)+\frac12\Delta_F(z_2-z_1)\\
&=\Delta_F(z_1-z_2).
\end{aligned}}
\tag{3.51}\label{eq:3.51}
\end{equation*}

\topic{Perturbation Expansion}

By \eqref{eq:3.12}, each $-i\,\delta/\delta J(z)$ acting on \eqref{eq:3.41} inserts $\phi(z)$, so
\begin{equation*}
\begin{aligned}
K(\Phi_f,\Phi_i)
&=\int_{\Phi_i}^{\Phi_f}\mathcal D\phi\,\exp\Bigl(iS_0[\phi]+i\int d^4z\,\Bigl(-\frac{\lambda}{3!}\Bigr)\phi^3\Bigr)\\
&=\int_{\Phi_i}^{\Phi_f}\mathcal D\phi\,e^{iS_0[\phi]}\sum_{m=0}^\infty\frac1{m!}\Bigl(i\int d^4z\,\Bigl(-\frac{\lambda}{3!}\Bigr)\phi(z)^3\Bigr)^m\\
&=\sum_{m=0}^\infty\frac1{m!}\Bigl(-\frac{i\lambda}{3!}\Bigr)^m\prod_{k=1}^m\int d^4z_k\Bigl(-i\frac{\delta}{\delta J(z_k)}\Bigr)^3K_0^J(\Phi_f,\Phi_i)\Big|_{J=0}\\
&=\sum_{m=0}^\infty\frac1{m!}\Bigl(\frac{\lambda}{3!}\Bigr)^m\prod_{k=1}^m\int d^4z_k\Bigl(\frac{\delta}{\delta J(z_k)}\Bigr)^3K_0^J(\Phi_f,\Phi_i)\Big|_{J=0}.
\end{aligned}
\tag{3.52}\label{eq:3.52}
\end{equation*}
With the vacuum of $S_0$ at both ends, \eqref{eq:3.49} gives
\begin{equation*}
{\displaystyle 
\begin{aligned}
K(\mathrm{vac}_f,\mathrm{vac}_i)
&=\sum_{m=0}^\infty\frac1{m!}\Bigl(\frac{\lambda}{3!}\Bigr)^m\prod_{k=1}^m\int d^4z_k\Bigl(\frac{\delta}{\delta J(z_k)}\Bigr)^3\\
&\qquad\exp\Bigl[-\frac12\int d^4x\,d^4y\,J(x)\Delta_F(x-y)J(y)\Bigr]\Big|_{J=0}\cdot K_0(\mathrm{vac}_f,\mathrm{vac}_i).
\end{aligned}}
\tag{3.53}\label{eq:3.53}
\end{equation*}

\setcounter{section}{3}
\section{Path Integrals for Fermion Fields}\label{book-section-3.4}\label{path-integrals-for-fermion-fields}

\subsection{Grassmann Algebra}\label{book-subsection-section-3.4-1}\label{grassmann-algebra}

\topic{Algebra}

Let $\mathbb K=\mathbb R$ or $\mathbb C$.
\begin{itemize}


\item \textbf{\boldmath Algebra}: a vector space $\mathcal A$ over $\mathbb K$ with a product
\begin{equation*}
\begin{aligned}
\mathcal A\times\mathcal A&\to\mathcal A,\\
(u,v)&\mapsto uv,
\end{aligned}
\tag{3.54}\label{eq:3.54}
\end{equation*}
that is bilinear: for all $\alpha,\beta\in\mathbb K$ and $u,v,w\in\mathcal A$,
\begin{equation*}
\begin{aligned}
u(\alpha v+\beta w)&=\alpha\,uv+\beta\,uw,\\
(\alpha v+\beta w)u&=\alpha\,vu+\beta\,wu.
\end{aligned}
\tag{3.55}\label{eq:3.55}
\end{equation*}

\item \textbf{\boldmath Unital algebra}: an algebra with an element $1\in\mathcal A$ such that $1u=u=u1$ for all $u\in\mathcal A$.
\begin{itemize}


\item For $\alpha\in\mathbb K$, $\alpha1$ is written $\alpha$.
\end{itemize}


\item \textbf{\boldmath Associative algebra}: an algebra in which $(uv)w=u(vw)$ for all $u,v,w\in\mathcal A$.
\end{itemize}


\topic{Generator}
\begin{itemize}


\item \textbf{\boldmath Generated subalgebra}: for an associative unital algebra $\mathcal A$ over $\mathbb K$ and a subset $\Sigma\subset\mathcal A$ with elements $u_a$, the subalgebra generated by $\Sigma$ is
\begin{equation*}
\mathcal A(\Sigma):=\Bigl\{f_0+\sum_{k\ge1}\sum_{a_1,\ldots,a_k}f_{a_1\cdots a_k}\,u_{a_1}\cdots u_{a_k}\;:\;f_0,\,f_{a_1\cdots a_k}\in\mathbb K\Bigr\},
\tag{3.56}\label{eq:3.56}
\end{equation*}
where only finitely many coefficients are nonzero.

\item \textbf{\boldmath Generating set}: $\Sigma$ is a generating set of $\mathcal A$ if $\mathcal A(\Sigma)=\mathcal A$.
\end{itemize}


\topic{Grassmann Algebra}
\begin{itemize}


\item \textbf{\boldmath Grassmann (exterior) algebra}: an associative unital algebra over $\mathbb K$ with a generating set $\{\zeta_1,\ldots,\zeta_n\}$ such that
\begin{equation*}
\zeta_a\zeta_b+\zeta_b\zeta_a=0,\qquad a,b=1,\ldots,n,
\tag{3.57}\label{eq:3.57}
\end{equation*}
in particular $\zeta_a^2=0$. No further relations are imposed: every relation among the $\zeta_a$ follows from \eqref{eq:3.57}. The algebra is denoted $\mathcal G(\zeta_1,\ldots,\zeta_n)$, or $\mathcal G_n$.

\item \textbf{\boldmath Basis}: the ordered monomials $1$ and $\zeta_{a_1}\cdots\zeta_{a_k}$ with $a_1<\cdots<a_k$ are linearly independent over $\mathbb K$, and
\begin{equation*}
\begin{aligned}
\mathcal G_n&=\operatorname{span}_{\mathbb K}\{1,\ \zeta_{a_1}\cdots\zeta_{a_k}\mid 1\le a_1<\cdots<a_k\le n\},\\
\dim_{\mathbb K}\mathcal G_n&=2^n.
\end{aligned}
\tag{3.58}\label{eq:3.58}
\end{equation*}

\item \textbf{\boldmath Expansion}: every $f\in\mathcal G_n$ has a unique expansion
\begin{equation*}
f=\sum_{k=0}^n\sum_{1\le a_1<\cdots<a_k\le n}f_{a_1\cdots a_k}\,\zeta_{a_1}\cdots\zeta_{a_k},\qquad f_{a_1\cdots a_k}\in\mathbb K,
\tag{3.59}\label{eq:3.59}
\end{equation*}
in which the $k=0$ term is $f_0$.
\end{itemize}


\subsection{Calculus on the Grassmann Algebra}\label{book-subsection-section-3.4-2}\label{grassmann-calculus}

\setcounter{section}{4}
\section{Symmetry in Path Integrals}\label{book-section-3.5}\label{symmetry-in-path-integrals}

\setcounter{section}{5}
\section{Scattering Theory in Path Integrals}\label{book-section-3.6}\label{scattering-theory-in-path-integrals}

\setcounter{chapter}{3}
\chapter{Scattering Theory}\label{book-chapter-4}\label{qft-ch4}

\setcounter{section}{0}
\section{Interaction Picture}\label{book-section-4.1}\label{interaction-picture-dyson}\label{perturbation-theory}\label{interaction-picture}

Let the Hamiltonian be
\begin{equation*}
\hat H=\hat H_0+\hat V
\tag{4.1}\label{eq:4.1}
\end{equation*}
The interaction picture evolves operators with $\hat H_0$ of \eqref{eq:4.1} instead of $\hat H$:
\begin{equation*}
{\displaystyle 
\begin{gathered}
\hat O_I(t):=e^{i\hat H_0t}\,\hat O_S\,e^{-i\hat H_0t},\\[4pt]
|\psi_I(t)\rangle:=e^{i\hat H_0t}\,|\psi_S(t)\rangle .
\end{gathered}
}
\tag{4.2}\label{eq:4.2}
\end{equation*}
The three pictures coincide at $t=0$, and $\langle\psi_I(t)|\hat O_I(t)|\psi_I(t)\rangle=\langle\psi_S(t)|\hat O_S|\psi_S(t)\rangle$.

Differentiating the first of \eqref{eq:4.2},
\begin{equation*}
\frac{d\hat O_I}{dt}=i\,[\hat H_0,\hat O_I].
\tag{4.3}\label{eq:4.3}
\end{equation*}
In particular $\hat H_{0,I}=\hat H_0$, while the interaction term becomes time dependent. Differentiating the second of \eqref{eq:4.2} and using $i\partial_t|\psi_S\rangle=\hat H|\psi_S\rangle$,
\begin{equation*}
i\partial_t|\psi_I(t)\rangle
=e^{i\hat H_0t}\bigl(\hat H-\hat H_0\bigr)|\psi_S(t)\rangle
=e^{i\hat H_0t}\hat Ve^{-i\hat H_0t}\,|\psi_I(t)\rangle,
\tag{4.4}\label{eq:4.4}
\end{equation*}
i.e.
\begin{equation*}
{\displaystyle 
i\partial_t|\psi_I(t)\rangle=\hat V_I(t)\,|\psi_I(t)\rangle,
\qquad
\hat V_I(t):=e^{i\hat H_0t}\,\hat V\,e^{-i\hat H_0t}.
}
\tag{4.5}\label{eq:4.5}
\end{equation*}
Substituting $\hat O_S=e^{-i\hat H_0t}\hat O_I(t)e^{i\hat H_0t}$ into \eqref{eq:2.23},
\begin{equation*}
\hat O_H(t)=\hat U_I(t,0)^\dagger\,\hat O_I(t)\,\hat U_I(t,0),
\qquad
\hat U_I(t,0):=e^{i\hat H_0t}\,\hat U_S(t,0).
\tag{4.6}\label{eq:4.6}
\end{equation*}
For a field operator, \eqref{eq:4.6} reads $\hat\phi_H(t,\mathbf x)=\hat U_I(t,0)^\dagger\,\hat\phi_I(t,\mathbf x)\,\hat U_I(t,0)$, and by \eqref{eq:4.3} $\hat\phi_I(t,\mathbf x)$ obeys the Heisenberg equation of $\hat H_0$ alone.

\topic{Dyson Series}

The state equation \eqref{eq:4.5} is \eqref{eq:2.3} with $\hat H_S$ replaced by $\hat V_I(t)$. The interaction-picture evolution operator $\hat U_I(t,t_0)$, defined by $|\psi_I(t)\rangle=\hat U_I(t,t_0)|\psi_I(t_0)\rangle$, is therefore $\hat U_S(t,t_0)$ of \eqref{eq:2.18} with $\hat H_S\to\hat V_I(t)$: it is unitary, $\hat U_I(t_1,t_2)\hat U_I(t_2,t_3)=\hat U_I(t_1,t_3)$, and $\hat U_I(t,t_0)^\dagger=\hat U_I(t_0,t)$. By \eqref{eq:4.2} and \eqref{eq:2.18},
\begin{equation*}
\hat U_I(t,t_0)=e^{i\hat H_0t}\,\hat U_S(t,t_0)\,e^{-i\hat H_0t_0},
\tag{4.7}\label{eq:4.7}
\end{equation*}
which reduces to the $\hat U_I(t,0)$ of \eqref{eq:4.6} at $t_0=0$; when $\hat H$ has no explicit time dependence, $\hat U_S(t,t_0)=e^{-i\hat H(t-t_0)}$ by \eqref{eq:2.4}. Since $[\hat V_I(t_1),\hat V_I(t_2)]\neq0$ in general, case (3) of \hyperref[schrodinger-picture]{Section 2.1} applies, and by \eqref{eq:2.11}
\begin{equation*}
{\displaystyle 
\hat U_I(t,t_0)=\mathcal T\exp\Bigl(-i\int_{t_0}^{t}dt'\,\hat V_I(t')\Bigr).
}
\tag{4.8}\label{eq:4.8}
\end{equation*}
In field theory, $\hat V=\int d^3x\,\hat{\mathcal H}_{\mathrm{int}}\bigl(\hat\phi_S(\mathbf x)\bigr)$. By \eqref{eq:4.2}, $\hat V_I(t)=\int d^3x\,\hat{\mathcal H}_{\mathrm{int}}\bigl(\hat\phi_I(t,\mathbf x)\bigr)$, and \eqref{eq:4.8} becomes
\begin{equation*}
\hat U_I(t,t_0)=\mathcal T\exp\Bigl(-i\int_{t_0}^{t}dt'\int d^3x\,\hat{\mathcal H}_{\mathrm{int}}\bigl(\hat\phi_I(x)\bigr)\Bigr).
\tag{4.9}\label{eq:4.9}
\end{equation*}
For interactions without derivative couplings, $\hat{\mathcal H}_{\mathrm{int}}=-\hat{\mathcal L}_{\mathrm{int}}$.

\setcounter{section}{1}
\section{The S-Matrix and Wick's Theorem}\label{book-section-4.2}\label{s-matrix-and-wick-theorem}

\topic{Asymptotic States and the S-Matrix}

Assume stable asymptotic particles and the existence of free-particle wave-packet limits. In-states approach the corresponding free states as $t\to-\infty$, and out-states as $t\to+\infty$. In perturbation theory, with adiabatic switching and the vacuum phase removed,
\begin{equation*}
\begin{aligned}
S_{fi}&:=\langle f,\mathrm{out}|i,\mathrm{in}\rangle=\langle f|\hat S|i\rangle,\\
\hat S&=\hat U_I(+\infty,-\infty)
=\mathcal T\exp\left(-i\int_{-\infty}^{+\infty}dt\,\hat V_I(t)\right).
\end{aligned}
\tag{4.10}\label{eq:4.10}
\end{equation*}
Let $p_i$ and $p_f$ denote the total initial and final four-momenta. The connected invariant amplitude $\mathcal M_{fi}$ and the transition operator $\hat T$ are defined by
\begin{equation*}
{\displaystyle \begin{aligned}
\hat S&=1+i\hat T,\qquad
\langle f|\hat T|i\rangle_{\mathrm c}=(2\pi)^4\delta^{(4)}(p_f-p_i)\,\mathcal M_{fi},\\
\hat S^\dagger\hat S&=1\quad\Longrightarrow\quad -i(\hat T-\hat T^\dagger)=\hat T^\dagger\hat T.
\end{aligned}}
\tag{4.11}\label{eq:4.11}
\end{equation*}

\topic{Normal Ordering, Contractions, and Wick's Theorem}

Normal ordering refers to the free vacuum $|0\rangle$. For a free real scalar field, its contraction is the Feynman propagator; $\varepsilon\to0^+$ is understood:
\begin{equation*}
\begin{aligned}
\Delta_F(x-y)&:=\langle0|\mathcal T[\hat\phi_I(x)\hat\phi_I(y)]|0\rangle\\
&=\mathcal T[\hat\phi_I(x)\hat\phi_I(y)]-:\hat\phi_I(x)\hat\phi_I(y):\\
&=\int\frac{d^4p}{(2\pi)^4}\,\frac{i\,e^{-ip\cdot(x-y)}}{p^2-m^2+i\varepsilon}.
\end{aligned}
\tag{4.12}\label{eq:4.12}
\end{equation*}
Wick's theorem expresses a time-ordered product of free fields as the sum over all disjoint contractions, including no contractions. Writing $\hat\phi_j:=\hat\phi_I(x_j)$,
\begin{equation*}
\begin{aligned}
\mathcal T[\hat\phi_1\cdots\hat\phi_n]
&=\sum_{\text{disjoint pairings }P}
\left[\prod_{(i,j)\in P}\Delta_F(x_i-x_j)\right]
:\prod_{k\text{ unpaired in }P}\hat\phi_k:,\\
\langle0|\mathcal T[\hat\phi_1\hat\phi_2\hat\phi_3\hat\phi_4]|0\rangle
&=\Delta_F(x_1-x_2)\Delta_F(x_3-x_4)
+\Delta_F(x_1-x_3)\Delta_F(x_2-x_4)\\
&\quad+\Delta_F(x_1-x_4)\Delta_F(x_2-x_3).
\end{aligned}
\tag{4.13}\label{eq:4.13}
\end{equation*}
For fermions, time ordering and Wick contractions carry the sign of the required permutation of fermionic operators. Free-vacuum expectation values retain only fully contracted terms.

\setcounter{section}{2}
\section{Feynman Diagrams and Feynman Rules}\label{book-section-4.3}\label{feynman-diagrams}

\topic{Scalar Example}

For a real scalar with quartic coupling $\lambda$, the momentum-space rules contribute to $i\mathcal M$. At tree level,
\begin{equation*}
\begin{aligned}
\mathcal L&=\frac12\partial_\mu\phi\,\partial^\mu\phi-\frac12m^2\phi^2-\frac{\lambda}{4!}\phi^4,\qquad
\hat{\mathcal H}_{\mathrm{int}}=\frac{\lambda}{4!}\hat\phi_I^4,\\
\text{internal scalar line}:&\quad\widetilde\Delta_F(p)=\frac{i}{p^2-m^2+i\varepsilon},\qquad
\text{four-scalar vertex}:\quad -i\lambda,\\
 i\mathcal M_{\phi\phi\to\phi\phi}^{\mathrm{tree}}&=-i\lambda.
\end{aligned}
\tag{4.14}\label{eq:4.14}
\end{equation*}

\topic{Propagators and External States}

Define $S_F(x-y):=\langle0|\mathcal T[\hat\psi_I(x)\overline{\hat\psi}_I(y)]|0\rangle$ and $D_{\mu\nu}(x-y):=\langle0|\mathcal T[\hat A_{I,\mu}(x)\hat A_{I,\nu}(y)]|0\rangle$. Their Fourier transforms are
\begin{equation*}
\begin{aligned}
\widetilde S_F(p)&=\frac{i(\slashed{p}+m)}{p^2-m^2+i\varepsilon},\\
\widetilde D_{\mu\nu}^{(\xi)}(p)&=\frac{-i}{p^2+i\varepsilon}\left[\eta_{\mu\nu}-(1-\xi)\frac{p_\mu p_\nu}{p^2+i\varepsilon}\right],
\qquad \mathcal L_{\mathrm{gf}}=-\frac1{2\xi}(\partial_\mu A^\mu)^2,\\
\widetilde D_{\mu\nu}^{\mathrm{Proca}}(p)&=\frac{-i}{p^2-m_V^2+i\varepsilon}\left(\eta_{\mu\nu}-\frac{p_\mu p_\nu}{m_V^2}\right),\qquad m_V>0.
\end{aligned}
\tag{4.15}\label{eq:4.15}
\end{equation*}
Here $\xi$ is the covariant gauge parameter ($\xi=1$ is Feynman gauge), and $m_V$ is the Proca mass. For fields with unit one-particle pole residue, the external factors are
\begin{equation*}
\begin{array}{c|cc}
\text{external state}&\text{incoming}&\text{outgoing}\\ \hline
\text{scalar}&1&1\\
\text{fermion}&u^s(p)&\bar u^s(p)\\
\text{antifermion}&\bar v^s(p)&v^s(p)\\
\text{vector}&\epsilon_\mu^r(p)&\epsilon_\mu^{r*}(p)
\end{array}
\tag{4.16}\label{eq:4.16}
\end{equation*}
The $\epsilon_\mu^r$ are physical polarization vectors. External propagators are amputated; a general field normalization requires the pole-residue factors in \hyperref[lsz-reduction]{Section 4.5}.

\topic{Diagram Rules}

Impose four-momentum conservation at every vertex, integrate each independent loop momentum with $d^4\ell/(2\pi)^4$, and sum connected amputated diagrams. Remove the overall $(2\pi)^4\delta^{(4)}(p_f-p_i)$ to obtain $i\mathcal M$. Divide each graph by its symmetry factor with the external legs held fixed. Each closed fermion loop contributes an additional minus sign; exchanging identical external fermions also changes the sign.

For example, the $s$-channel one-loop four-scalar graph has symmetry factor $2$:
\begin{equation*}
i\mathcal M_s^{\mathrm{1\,loop}}
=\frac{(-i\lambda)^2}{2}\int\frac{d^4\ell}{(2\pi)^4}
\frac{i}{\ell^2-m^2+i\varepsilon}
\frac{i}{(\ell+p_1+p_2)^2-m^2+i\varepsilon},
\qquad s:=(p_1+p_2)^2.
\tag{4.17}\label{eq:4.17}
\end{equation*}
Loop integrals require a regulator and, in a renormalized amplitude, the corresponding counterterms. The full one-loop four-point amplitude also contains the crossed channels.

\setcounter{section}{3}
\section{Correlation Functions}\label{book-section-4.4}\label{correlation-functions}

\subsection{Perturbative Vacuum Correlation Functions}\label{book-subsection-section-4.4-1}\label{perturbative-correlation-functions}

\topic{Vacuum Correlation Functions}

Let $|\Omega\rangle$ be the interacting vacuum and $|0\rangle$ the free vacuum.
\begin{equation*}
\begin{aligned}
G_n(x_1,\ldots,x_n)&:=\frac{\langle\Omega|\mathcal T[\hat\phi_H(x_1)\cdots\hat\phi_H(x_n)]|\Omega\rangle}{\langle\Omega|\Omega\rangle}\\
&=\frac{\langle0|\mathcal T\left[\hat\phi_I(x_1)\cdots\hat\phi_I(x_n)\exp\left(-i\int dt\,\hat V_I(t)\right)\right]|0\rangle}
{\langle0|\mathcal T\exp\left(-i\int dt\,\hat V_I(t)\right)|0\rangle}.
\end{aligned}
\tag{4.18}\label{eq:4.18}
\end{equation*}
The denominator cancels disconnected vacuum bubbles. Graphs contributing to $G_n$ may still be disconnected if each component has external field insertions.

\topic{Generating Functionals and Connected Functions}

For a real scalar source $J(x)$,
\begin{equation*}
\begin{aligned}
Z[J]&:=\frac{\langle\Omega|\mathcal T\exp\left(i\int d^4x\,J(x)\hat\phi_H(x)\right)|\Omega\rangle}{\langle\Omega|\Omega\rangle},\qquad Z[0]=1,\qquad W[J]:=-i\ln Z[J],\\
G_n(x_1,\ldots,x_n)&=\left.\frac1{i^n}\frac{\delta^nZ[J]}{\delta J(x_1)\cdots\delta J(x_n)}\right|_{J=0},\\
G_{n,\mathrm c}(x_1,\ldots,x_n)&=\left.i^{1-n}\frac{\delta^nW[J]}{\delta J(x_1)\cdots\delta J(x_n)}\right|_{J=0}.
\end{aligned}
\tag{4.19}\label{eq:4.19}
\end{equation*}
For $G_1=0$, the four-point function splits into connected and disconnected contributions as
\begin{equation*}
\begin{aligned}
G_4(x_1,x_2,x_3,x_4)
&=G_{4,\mathrm c}(x_1,x_2,x_3,x_4)
+G_2(x_1,x_2)G_2(x_3,x_4)\\
&\quad+G_2(x_1,x_3)G_2(x_2,x_4)
+G_2(x_1,x_4)G_2(x_2,x_3).
\end{aligned}
\tag{4.20}\label{eq:4.20}
\end{equation*}

\subsection{Nonperturbative Vacuum Correlation Functions}\label{book-subsection-section-4.4-2}\label{nonperturbative-correlation-functions}

\setcounter{section}{4}
\section{LSZ Reduction Formula}\label{book-section-4.5}\label{lsz-reduction}

\topic{One-Particle Pole}

Consider a real scalar field with zero vacuum expectation value, an isolated stable one-particle pole of physical mass $m_{\mathrm{phys}}>0$, and nonzero overlap with that particle. Choose the particle-state phase so that
\begin{equation*}
\begin{aligned}
\langle\Omega|\hat\phi_H(0)|\mathbf p\rangle&=\sqrt{Z_{\mathrm{pole}}},\qquad Z_{\mathrm{pole}}>0,\qquad \langle\Omega|\Omega\rangle=1,\\
\widetilde G_2(p)&:=\int d^4x\,e^{ip\cdot x}G_2(x,0)
=\frac{iZ_{\mathrm{pole}}}{p^2-m_{\mathrm{phys}}^2+i\varepsilon}
+\text{terms regular at the pole}.
\end{aligned}
\tag{4.21}\label{eq:4.21}
\end{equation*}
Here $Z_{\mathrm{pole}}$ is the pole residue for the chosen field normalization. Assume the corresponding asymptotic scattering states exist.

\topic{Reduction Formula}

For $r$ incoming momenta $p_a$ and $s$ outgoing momenta $q_b$, all with positive physical energies, define $\widetilde G_{r+s,\mathrm c}$ by Fourier transforming every position with $e^{+ik\cdot x}$. For a connected scattering process with $r+s>2$,
\begin{equation*}
{\displaystyle \begin{aligned}
&i(2\pi)^4\delta^{(4)}\left(\sum_{b=1}^{s}q_b-\sum_{a=1}^{r}p_a\right)\mathcal M_{fi}\\
&\quad=\lim_{\mathrm{on\ shell}}
\left[\prod_{b=1}^{s}\frac{q_b^2-m_{\mathrm{phys}}^2}{i\sqrt{Z_{\mathrm{pole}}}}\right]
\left[\prod_{a=1}^{r}\frac{p_a^2-m_{\mathrm{phys}}^2}{i\sqrt{Z_{\mathrm{pole}}}}\right]
\widetilde G_{r+s,\mathrm c}(q_1,\ldots,q_s,-p_1,\ldots,-p_r).
\end{aligned}}
\tag{4.22}\label{eq:4.22}
\end{equation*}
The limit means extracting the simultaneous external-pole residue with the Feynman prescription before taking all external momenta on shell. Each factor removes one external propagator and its field overlap. A field normalized to unit pole residue has $Z_{\mathrm{pole}}=1$.

\setcounter{section}{5}
\section{Cross Sections and Decay Rates}\label{book-section-4.6}\label{cross-sections-and-decay-rates}

\topic{Phase Space and State Sums}

For total final momentum $P$ and $n$ on-shell outgoing momenta $k_j$,
\begin{equation*}
d\Phi_n(P;k_1,\ldots,k_n):=(2\pi)^4\delta^{(4)}\left(P-\sum_{j=1}^{n}k_j\right)
\prod_{j=1}^{n}\frac{d^3k_j}{(2\pi)^3\,2E_{\mathbf k_j}},\qquad
E_{\mathbf k_j}=\sqrt{\mathbf k_j^2+m_j^2}.
\tag{4.23}\label{eq:4.23}
\end{equation*}
Use $\overline{|\mathcal M|^2}$ for the sum over unobserved final spins and internal quantum numbers, averaged over the equally populated initial states. For two unpolarized initial particles with $N_1,N_2$ such states,
\begin{equation*}
\overline{|\mathcal M|^2}:=\frac1{N_1N_2}\sum_{\text{initial, final states}}|\mathcal M|^2,
\qquad C_{\mathrm{id}}:=\prod_a\nu_a!.
\tag{4.24}\label{eq:4.24}
\end{equation*}
Here $\nu_a$ is the number of identical final particles of species $a$. Divide by $C_{\mathrm{id}}$ when integrating the full phase space of labeled final momenta; omit this factor if the integration region already counts each physical final configuration once. Fixed initial polarizations are not averaged.

\topic{Cross Sections and Decay Rates}

For $1+2\to n$ scattering and decay of a particle of mass $M$ at rest,
\begin{equation*}
{\displaystyle \begin{aligned}
d\sigma&=\frac{\overline{|\mathcal M|^2}}{4\sqrt{(p_1\cdot p_2)^2-m_1^2m_2^2}}\,\frac{d\Phi_n(p_1+p_2)}{C_{\mathrm{id}}},\\
d\Gamma&=\frac{\overline{|\mathcal M|^2}}{2M}\,\frac{d\Phi_n(P)}{C_{\mathrm{id}}},\qquad P=(M,\mathbf0).
\end{aligned}}
\tag{4.25}\label{eq:4.25}
\end{equation*}
In the decay formula, average only over the parent's initial states if an unpolarized rate is intended. For two-body final states, use the single-particle three-momentum magnitudes $|\mathbf p_1|$ and $|\mathbf k_1|$ for the initial and final states in the center-of-mass frame:
\begin{equation*}
\begin{aligned}
\frac{d\sigma_{2\to2}}{d\Omega}
&=\frac1{64\pi^2s\,C_{\mathrm{id}}}\frac{|\mathbf k_1|}{|\mathbf p_1|}\,\overline{|\mathcal M|^2},\qquad s=(p_1+p_2)^2,\\
\frac{d\Gamma_{1\to2}}{d\Omega}
&=\frac{|\mathbf k_1|}{32\pi^2M^2\,C_{\mathrm{id}}}\,\overline{|\mathcal M|^2},\\
\Gamma_{1\to2}&=\frac{|\mathbf k_1|}{8\pi M^2\,C_{\mathrm{id}}}\,\overline{|\mathcal M|^2}\qquad\text{if the squared amplitude is angle independent}.
\end{aligned}
\tag{4.26}\label{eq:4.26}
\end{equation*}
For the decay formulas, $|\mathbf k_1|$ is evaluated in the parent's rest frame. The solid-angle integrals use the same identical-particle counting convention as the full phase space.

\setcounter{chapter}{4}
\chapter{The Standard Model}\label{book-chapter-5}\label{qft-ch5}

\setcounter{section}{0}
\section{Yang–Mills Theory}\label{book-section-5.1}\label{yang-mills}

\subsection{Lie Group}\label{book-subsection-section-5.1-1}\label{lie-group}
\begin{equation*}
\begin{aligned}
U(n)&=\{A\in M_n(\mathbb C)\mid A^\dagger A=I\},\\[6pt]
SU(n)&=\{A\in M_n(\mathbb C)\mid \det A=1,\ A^\dagger A=I\}.
\end{aligned}
\tag{5.1}\label{eq:5.1}
\end{equation*}

\topic{$\mathfrak{u}(n)$}

Consider the one-parameter subgroup
\begin{equation*}
\begin{aligned}
g(\theta)&=e^{i\theta T},\qquad \theta\in\mathbb R,\quad T\in M_n(\mathbb C),\\[6pt]
\left.\frac{d}{d\theta}g(\theta)\right|_{\theta=0}&=iT,\\[6pt]
g^\dagger(\theta)g(\theta)=I\quad\text{for all }\theta\in\mathbb R
&\Longrightarrow i(T-T^\dagger)=0
\Longrightarrow T^\dagger=T.
\end{aligned}
\tag{5.2}\label{eq:5.2}
\end{equation*}
The Lie algebra is
\begin{equation*}
\mathfrak{u}(n)=\{iT\in M_n(\mathbb C)\mid T^\dagger=T\}.
\tag{5.3}\label{eq:5.3}
\end{equation*}
Hermiticity $T^\dagger=T$ means $T_{jj}\in\mathbb R$ and $T_{kj}=T_{jk}^*$ for $j<k$, i.e. $n+2\cdot\frac{n(n-1)}{2}=n^2$ real conditions on the $2n^2$ real parameters of $M_n(\mathbb C)$. Hence $\dim\mathfrak{u}(n)=2n^2-n^2=n^2$.

\topic{$\mathfrak{su}(n)$}

The one-parameter subgroup and its generator are
\begin{equation*}
\begin{aligned}
g(\theta)&=e^{i\theta T},\qquad \theta\in\mathbb R,\quad T\in M_n(\mathbb C),\\[6pt]
\left.\frac{d}{d\theta}g(\theta)\right|_{\theta=0}&=iT.
\end{aligned}
\tag{5.4}\label{eq:5.4}
\end{equation*}
\begin{equation*}
0=\left.\frac{d}{d\theta}\bigl[g^\dagger(\theta)g(\theta)\bigr]\right|_{\theta=0}
=i(T-T^\dagger)
\quad\Longrightarrow\quad T^\dagger=T.
\tag{5.5}\label{eq:5.5}
\end{equation*}
We use the theorem $\det(e^{i\theta T})=e^{i\theta\operatorname{tr}T}$, valid for any $T\in M_n(\mathbb C)$ and $\theta\in\mathbb C$.
\begin{equation*}
0=\left.\frac{d}{d\theta}\det g(\theta)\right|_{\theta=0}
=\left.\frac{d}{d\theta}e^{i\theta\operatorname{tr}T}\right|_{\theta=0}
=i\operatorname{tr}T
\quad\Longrightarrow\quad \operatorname{tr}T=0.
\tag{5.6}\label{eq:5.6}
\end{equation*}
Thus $T$ is traceless and Hermitian, and the Lie algebra is
\begin{equation*}
\mathfrak{su}(n)=\{iT\in M_n(\mathbb C)\mid
\operatorname{tr}T=0,\ T^\dagger=T\}.
\tag{5.7}\label{eq:5.7}
\end{equation*}
Since $\operatorname{tr}T=\sum_jT_{jj}\in\mathbb R$ for Hermitian $T$, tracelessness $\operatorname{tr}T=0$ is one further real condition. Hence $\dim\mathfrak{su}(n)=2n^2-n^2-1=n^2-1$.

Scalar multiplication is taken over $\mathbb R$ to preserve anti-Hermiticity; hence $\mathfrak{su}(n)$ and  $\mathfrak{u}(n)$ are both real Lie algebras.

\phantomsection\label{ym-4}%
\begin{qftinset}

\divtitle{$\circ$ Structure Constants $f^{abc}$}

Let $\{iT^a\}$, with $a=1,\ldots,n^2-1$, be a basis of the real Lie algebra $\mathfrak{su}(n)$.
The structure constants $f^{abc}$ are defined by $[T^a,T^b]=if^{abc}T^c$.

\textbf{\boldmath Claim.} $f^{abc}\in\mathbb R$.

Since the matrices $T^a$ are Hermitian,
\begin{equation*}
\begin{aligned}
([T^a,T^b])^\dagger&=-[T^a,T^b],\qquad
\operatorname{tr}[T^a,T^b]=0,\\[6pt]
[T^a,T^b]&\in\mathfrak{su}(n),\\[6pt]
[T^a,T^b]&=f^{abc}(iT^c),\qquad f^{abc}\in\mathbb R.
\end{aligned}
\tag{5.8}\label{eq:5.8}
\end{equation*}
The expansion coefficients are real because $\{iT^c\}$ is a basis over $\mathbb R$.

\end{qftinset}

\begin{qftinset}

\divtitle{$\circ$ Inner Product on $\mathfrak{su}(n)$}

For $iT_1,iT_2\in\mathfrak{su}(n)$, define
\begin{equation*}
\begin{aligned}
(\cdot,\cdot)&:\mathfrak{su}(n)\times\mathfrak{su}(n)\longrightarrow\mathbb R,\\[6pt]
(iT_1,iT_2)&:=-2\operatorname{tr}\bigl[(iT_1)(iT_2)\bigr]\\
&=2\operatorname{tr}(T_1T_2).
\end{aligned}
\tag{5.9}\label{eq:5.9}
\end{equation*}
\textbf{\boldmath Proof} that \eqref{eq:5.9} satisfies the axioms of an inner product

\emph{(i) Real values.}

Using $T_1^\dagger=T_1$, $T_2^\dagger=T_2$ and cyclicity of the trace,
\begin{equation*}
\begin{aligned}
\bigl[\operatorname{tr}(T_1T_2)\bigr]^*
&=\operatorname{tr}\bigl[(T_1T_2)^\dagger\bigr]
=\operatorname{tr}(T_2^\dagger T_1^\dagger)\\
&=\operatorname{tr}(T_2T_1)
=\operatorname{tr}(T_1T_2),
\end{aligned}
\tag{5.10}\label{eq:5.10}
\end{equation*}
so $(iT_1,iT_2)=2\operatorname{tr}(T_1T_2)\in\mathbb R$.

\emph{(ii) Symmetry.}

By cyclicity of the trace,
\begin{equation*}
\begin{aligned}
(iT_1,iT_2)&=2\operatorname{tr}(T_1T_2)\\
&=2\operatorname{tr}(T_2T_1)=(iT_2,iT_1).
\end{aligned}
\tag{5.11}\label{eq:5.11}
\end{equation*}
\emph{(iii) Bilinearity.}

For $a,b\in\mathbb R$, $a\,iT_1+b\,iT_2=i(aT_1+bT_2)$,
and linearity of the trace gives $(a\,iT_1+b\,iT_2,\,iT_3)=2\operatorname{tr}\bigl[(aT_1+bT_2)T_3\bigr]=a\,(iT_1,iT_3)+b\,(iT_2,iT_3)$. By (ii), the form is also linear in its second argument.

\emph{(iv) Positive definiteness.} Using $T=T^\dagger$,
\begin{equation*}
\begin{aligned}
(iT,iT)&=2\operatorname{tr}(T^2)=2\operatorname{tr}(T^\dagger T)\\
&=2\sum_{j,k}|T_{jk}|^2\ge0.
\end{aligned}
\tag{5.12}\label{eq:5.12}
\end{equation*}
Equality holds if and only if $T_{jk}=0$ for all $j,k$, i.e. $iT=0$. $\square$

\end{qftinset}

\phantomsection\label{ym-orthonormal}%
\begin{qftinset}

\divtitle{$\circ$ Orthonormal Basis of $\mathfrak{su}(n)$ and Total Antisymmetry of $f^{abc}$}

Since \eqref{eq:5.9} is an inner product on $\mathfrak{su}(n)$, we may choose an orthonormal basis $\{iT^a\}$ satisfying $(iT^a,iT^b)=\delta^{ab}$. Equivalently,
\begin{equation*}
{\displaystyle \operatorname{tr}(T^aT^b)=\frac12\delta^{ab}.}
\tag{5.13}\label{eq:5.13}
\end{equation*}
\textbf{\boldmath Total antisymmetry.} With this normalization and cyclicity of the trace, we obtain
\begin{equation*}
\begin{aligned}
\operatorname{tr}([T^a,T^b]T^c)&=\frac{i}{2}f^{abc},\\
f^{abc}&=-2i\operatorname{tr}([T^a,T^b]T^c),\\
f^{bac}&=-f^{abc},\\
f^{acb}
&=-2i\operatorname{tr}(T^aT^cT^b-T^cT^aT^b)\\
&=-2i\operatorname{tr}(T^bT^aT^c-T^aT^bT^c)
=-f^{abc}.
\end{aligned}
\tag{5.14}\label{eq:5.14}
\end{equation*}
\begin{equation*}
{\displaystyle 
f^{abc}=f^{[abc]}\in\mathbb R
}
\tag{5.15}\label{eq:5.15}
\end{equation*}

\end{qftinset}

\subsection{Local Gauge Symmetry}\label{book-subsection-section-5.1-2}\label{ym-5}\label{local-gauge-symmetry}

Fiber bundles provide a useful geometric picture of fields in Yang–Mills theory.
In this picture, an internal space is a fiber which is attached to each spacetime point, and the field takes its value in that fiber.
This makes an additional freedom explicit: we may choose a basis in each fiber,
and this choice can vary from point to point.
It is this local freedom of choice that leads to the  gauge symmetry.

Building on this idea, we now develop the framework of gauge theory below.

The field operator is defined as
\begin{equation*}
\Psi(x):=\begin{pmatrix}
\psi_1(x)\\
\vdots\\
\psi_n(x)
\end{pmatrix},
\tag{5.16}\label{eq:5.16}
\end{equation*}
where each $\psi_J$ is a Dirac spinor and $J=1,\ldots,n$ is the index of the internal space, e.g. color, generation.

Define the Dirac adjoint by
\begin{equation*}
\begin{aligned}
\bar\Psi(x)&:=\bigl(\bar\psi_1(x),\ldots,\bar\psi_n(x)\bigr),\\[6pt]
\bar\psi_J(x)&:=\psi_J^\dagger(x)\gamma^0.
\end{aligned}
\tag{5.17}\label{eq:5.17}
\end{equation*}
The Lagrangian density contains the kinetic term
\begin{equation*}
\mathcal L\supset\bar\Psi\,i\gamma^\mu I_n\partial_\mu\Psi=\sum_{J=1}^{n}\bar\psi_J\,i\gamma^\mu\partial_\mu\psi_J.
\tag{5.18}\label{eq:5.18}
\end{equation*}
Let $R$ be the representation of $SU(n)$ on the internal space. Let $g_{\mathrm{YM}}$ be the coupling and $R(\theta)$ the representation of the local group element.

To compare fields in different internal spaces, we need to establish a relation between these internal spaces.
\begin{itemize}


\item \textbf{\boldmath Wilson line}: for a path $C:[0,1]\to\mathbb R^4$ from $C(0)=x$ to $C(1)=y$, the Wilson line along $C$ maps the internal space at $x$ to that at $y$:
\begin{equation*}
W_C(y,x):=\mathcal P\exp\Bigl(ig_{\mathrm{YM}}\int_0^1ds\,\dot C^\mu(s)A_\mu\bigl(C(s)\bigr)\Bigr),
\tag{5.19}\label{eq:5.19}
\end{equation*}
where $A_\mu:=A_\mu^aT^a$, with $A_\mu^a$ real and $T^a$ the generators in the representation $R$, is the connection and $\dot C^\mu:=dC^\mu/ds$. In general, $W_C(y,x)$ depends on the path $C$, not only on its endpoints $x$ and $y$.

The path-ordered exponential is defined as $\mathcal T\exp$ in \eqref{eq:2.12}–\eqref{eq:2.13} with $\tau\to s$ and $-i\hat H_S(\tau)\to ig_{\mathrm{YM}}\dot C^\mu(s)A_\mu(C(s))$, so $\mathcal P$ puts factors with larger $s$ to the left.

Thus $W_C(y,x)\Psi(x)$ lies in the internal space at $y$, so under a gauge transformation we require it to transform like $\Psi(y)$:
\begin{equation*}
W_C(y,x)\Psi(x)\to R\bigl(\theta(y)\bigr)\,W_C(y,x)\Psi(x).
\tag{5.20}\label{eq:5.20}
\end{equation*}
Since $\Psi(x)\to R\bigl(\theta(x)\bigr)\Psi(x)$, the Wilson line must transform as
\begin{equation*}
W_C(y,x)\to R\bigl(\theta(y)\bigr)\,W_C(y,x)\,R\bigl(\theta(x)\bigr)^{-1}.
\tag{5.21}\label{eq:5.21}
\end{equation*}

\item \textbf{\boldmath Covariant derivative on fields in the representation $R$}: $D_\mu:=\partial_\mu-ig_{\mathrm{YM}}A_\mu$.

Under a gauge transformation, the covariant derivative transforms as
\begin{equation*}
D_\mu\Psi\to R(\theta)\,D_\mu\Psi,\qquad\text{i.e.}\qquad D_\mu\to R(\theta)\,D_\mu\,R(\theta)^{-1}.
\tag{5.22}\label{eq:5.22}
\end{equation*}
\textbf{\boldmath Proof.} For an infinitesimal displacement $\delta x^\mu$ and the straight path $C(s):=x+s\,\delta x$, to first order in $\delta x$, \eqref{eq:5.19} gives
\begin{equation*}
\begin{aligned}
W_C(x+\delta x,x)&=I_n+ig_{\mathrm{YM}}\,\delta x^\mu A_\mu(x),\\
\delta x^\mu D_\mu\Psi(x)&=\Psi(x+\delta x)-W_C(x+\delta x,x)\Psi(x).
\end{aligned}
\tag{5.23}\label{eq:5.23}
\end{equation*}
By \eqref{eq:5.20}, both terms on the right of the second line are multiplied by $R(\theta(x+\delta x))=R(\theta(x))+O(\delta x)$ under a gauge transformation, so $\delta x^\mu D_\mu\Psi(x)\to R(\theta(x))\,\delta x^\mu D_\mu\Psi(x)$ to first order.
Since $\delta x^\mu$ is arbitrary, \eqref{eq:5.22} follows. $\square$

Then, by \eqref{eq:5.22}, $A_\mu$ must transform as
\begin{equation*}
A_\mu\to R(\theta)\,A_\mu\,R(\theta)^{-1}-\frac{i}{g_{\mathrm{YM}}}\bigl(\partial_\mu R(\theta)\bigr)R(\theta)^{-1}.
\tag{5.24}\label{eq:5.24}
\end{equation*}

\item \textbf{\boldmath Field strength}: $[D_\mu,D_\nu]=:-ig_{\mathrm{YM}}F_{\mu\nu}$.

With $F_{\mu\nu}=:F_{\mu\nu}^aT^a$, this gives
\begin{equation*}
\begin{aligned}
F_{\mu\nu}&=\partial_\mu A_\nu-\partial_\nu A_\mu-ig_{\mathrm{YM}}[A_\mu,A_\nu],\\
F_{\mu\nu}^a&=\partial_\mu A_\nu^a-\partial_\nu A_\mu^a+g_{\mathrm{YM}}f^{abc}A_\mu^bA_\nu^c.
\end{aligned}
\tag{5.25}\label{eq:5.25}
\end{equation*}
Under a gauge transformation, the field strength transforms as
\begin{equation*}
F_{\mu\nu}\to R(\theta)\,F_{\mu\nu}\,R(\theta)^{-1}.
\tag{5.26}\label{eq:5.26}
\end{equation*}
\textbf{\boldmath Proof.} By \eqref{eq:5.22}, $[D_\mu,D_\nu]\to R(\theta)\,[D_\mu,D_\nu]\,R(\theta)^{-1}$, and $[D_\mu,D_\nu]=-ig_{\mathrm{YM}}F_{\mu\nu}$. $\square$

Including the gauge field, a gauge-invariant extension of the kinetic term \eqref{eq:5.18} is
\begin{equation*}
{\displaystyle 
\mathcal L\supset\bar\Psi\,i\gamma^\mu D_\mu\Psi-\frac14F_{\mu\nu}^aF^{a\mu\nu}.
}
\tag{5.27}\label{eq:5.27}
\end{equation*}

\item \textbf{\boldmath Covariant derivative on adjoint fields}: for a matrix-valued field $X$ with $X\to R(\theta)\,X\,R(\theta)^{-1}$, such as $F_{\mu\nu}$, $D_\mu X:=\partial_\mu X-ig_{\mathrm{YM}}[A_\mu,X]$.

\item \textbf{\boldmath Equations of motion}: using the definitions above and $j^\nu:=(\bar\Psi\gamma^\nu T^a\Psi)T^a$, the field equations of \eqref{eq:5.27} are
\begin{equation*}
{\displaystyle 
\begin{aligned}
i\gamma^\mu D_\mu\Psi&=0,\\
D_\mu F^{\mu\nu}&=-g_{\mathrm{YM}}j^\nu.
\end{aligned}
}
\tag{5.28}\label{eq:5.28}
\end{equation*}
\end{itemize}

In addition, $F_{\mu\nu}$ satisfies the Bianchi identity $D_\mu F_{\nu\rho}+D_\nu F_{\rho\mu}+D_\rho F_{\mu\nu}=0$.

\setcounter{section}{1}
\section{Fermions and Bosons in the Standard Model}\label{book-section-5.2}\label{sm-fermions}

\subsection{General Higgs Lagrangian}\label{book-subsection-section-5.2-1}\label{higgs-mechanism}\label{higgs-lagrangian}

Let $\Phi:=(\phi_1,\ldots,\phi_n)^T$ be a multiplet of complex scalar fields.

For $\mu_H^2\in\mathbb R$ and $\lambda_H>0$, the Higgs Lagrangian is
\begin{equation*}
{\displaystyle 
\begin{aligned}
\mathcal L_{\mathrm{Higgs}}&=(\partial_\mu\Phi)^\dagger\partial^\mu\Phi-V(\Phi),\\
V(\Phi)&=-\mu_H^2\,\Phi^\dagger\Phi+\lambda_H(\Phi^\dagger\Phi)^2 .
\end{aligned}
}
\tag{5.29}\label{eq:5.29}
\end{equation*}
The Lagrangian \eqref{eq:5.29} depends on $\Phi$ only through $\Phi^\dagger\Phi$ and $(\partial_\mu\Phi)^\dagger\partial^\mu\Phi$.

For a global gauge transformation $U\in U(n)$,
\begin{equation*}
\begin{aligned}
\Phi'&=U\,\Phi,\\
\Phi'^\dagger\Phi'&=\Phi^\dagger U^\dagger U\,\Phi=\Phi^\dagger\Phi,\\
(\partial_\mu\Phi')^\dagger\partial^\mu\Phi'&=(\partial_\mu\Phi)^\dagger U^\dagger U\,\partial^\mu\Phi=(\partial_\mu\Phi)^\dagger\partial^\mu\Phi,
\end{aligned}
\tag{5.30}\label{eq:5.30}
\end{equation*}
so $\mathcal L_{\mathrm{Higgs}}$ is invariant under global $U(n)$ transformations.
\begin{itemize}


\item \textbf{\boldmath Spontaneous Symmetry Breaking: Nonzero Vacuum Expectation Value}

The classical vacuum is the field configuration of minimum energy.

The energy density of \eqref{eq:5.29} is
\begin{equation*}
\mathcal H=(\partial_0\Phi)^\dagger\partial_0\Phi+\sum_{i=1}^3(\partial_i\Phi)^\dagger\partial_i\Phi+V(\Phi).
\tag{5.31}\label{eq:5.31}
\end{equation*}
The first two terms are nonnegative, so the classical vacuum is a constant $\Phi=\Phi_0$ minimizing $V$.

Since $\partial V/\partial(\Phi^\dagger\Phi)=-\mu_H^2+2\lambda_H\Phi^\dagger\Phi$ and $\Phi^\dagger\Phi\ge0$,
\begin{equation*}
{\displaystyle 
\begin{aligned}
\mu_H^2\le0:&\quad\Phi_0=0,\\
\mu_H^2>0:&\quad\Phi_0^\dagger\Phi_0=\frac{v_{\mathrm{SM}}^2}{2},\qquad
v_{\mathrm{SM}}:=\sqrt{\frac{\mu_H^2}{\lambda_H}},\qquad
V(\Phi_0)=-\frac{\lambda_Hv_{\mathrm{SM}}^4}{4}.
\end{aligned}
}
\tag{5.32}\label{eq:5.32}
\end{equation*}
For $\mu_H^2\le0$ the vacuum $\Phi_0=0$ is invariant under all of $U(n)$.

From now on $\mu_H^2>0$: the vacua form the sphere $\Phi_0^\dagger\Phi_0=v_{\mathrm{SM}}^2/2$ in $\mathbb C^n$, and $\Phi=0$ is a local maximum of $V$.
These vacua are degenerate. This degeneracy is called a symmetry-enforced vacuum degeneracy, as opposed to an accidental degeneracy, in which two vacua have the same energy but no symmetry protects this equality.

Physically, however, the vacuum must be unique, so in the mathematical description a physical vacuum has to be chosen by hand. We choose
\begin{equation*}
{\displaystyle 
\Phi_0:=\frac1{\sqrt2}\begin{pmatrix}0\\\vdots\\0\\v_{\mathrm{SM}}\end{pmatrix}
}
\tag{5.33}\label{eq:5.33}
\end{equation*}
and from now on $\Phi_0$ denotes this vacuum.

This choice is arbitrary: any point of the sphere could serve as the physical vacuum.

If the degeneracy is symmetry-enforced, choosing the physical vacuum by hand breaks that symmetry: not all symmetries of the Lagrangian are symmetries of the physical vacuum. Here $\mathcal L_{\mathrm{Higgs}}$ is invariant under all of $U(n)$, but $\Phi_0$ is not. This is spontaneous symmetry breaking.

A symmetry transformation of the Lagrangian is unbroken if it leaves the physical vacuum invariant, and broken otherwise.

Broken generators are denoted $G$, and unbroken ones are still denoted $T$:
\begin{equation*}
{\displaystyle 
\begin{aligned}
T\Phi_0&=0\qquad(\text{unbroken}),\\
G\Phi_0&\neq0\qquad(\text{broken}).
\end{aligned}
}
\tag{5.34}\label{eq:5.34}
\end{equation*}
The $iT$ with $T$ unbroken form a representation of $\mathfrak u(n-1)$, so in $\operatorname{tr}$-orthonormal bases,
\begin{equation*}
\begin{aligned}
G^a\Phi_0&\neq0,\qquad a=1,\ldots,2n-1,\\
T^a\Phi_0&=0,\qquad a=2n,\ldots,n^2,\\
\operatorname{tr}(G^aG^b)&=\operatorname{tr}(T^aT^b)=\frac12\delta^{ab},\qquad\operatorname{tr}(G^aT^b)=0 .
\end{aligned}
\tag{5.35}\label{eq:5.35}
\end{equation*}
Parametrize the $2n$ real fields of $\Phi$ by the Higgs field $h$ and $\phi'^a$, $a=1,\ldots,2n-1$, along the broken generators:
\begin{equation*}
{\displaystyle 
\begin{aligned}
\Phi&=\frac1{\sqrt2}\,e^{i\phi'^a(x)G^a}\begin{pmatrix}0\\\vdots\\0\\v_{\mathrm{SM}}+h\end{pmatrix} .
\end{aligned}
}
\tag{5.36}\label{eq:5.36}
\end{equation*}
\end{itemize}


\phantomsection\label{higgs-mechanism-masses}%
\topic{Higgs Mechanism}

Gauge the $U(n)$ symmetry \eqref{eq:5.30} with the generators $T^a$, $a=1,\ldots,n^2$, of \eqref{eq:5.35} (with $T^a:=G^a$ for $a\le2n-1$), gauge fields $A_\mu^a$, and coupling $g_{\mathrm{YM}}$ as in \eqref{eq:5.22}–\eqref{eq:5.24}. Replacing $\partial_\mu$ by $D_\mu$ in \eqref{eq:5.29} and adding the gauge kinetic term of \eqref{eq:5.27},
\begin{equation*}
{\displaystyle 
\begin{aligned}
\mathcal L&=(D_\mu\Phi)^\dagger D^\mu\Phi-V(\Phi)-\frac14F_{\mu\nu}^aF^{a\mu\nu},\\
D_\mu\Phi&=\partial_\mu\Phi-ig_{\mathrm{YM}}A_\mu\Phi,\qquad A_\mu:=A_\mu^aT^a .
\end{aligned}
}
\tag{5.37}\label{eq:5.37}
\end{equation*}
Substitute the parametrization \eqref{eq:5.36} into \eqref{eq:5.37}, written as $\Phi=(1+h/v_{\mathrm{SM}})\,U(\phi')\Phi_0$ with $U(\phi'):=e^{i\phi'^aG^a}$. With $\mu_H^2=\lambda_Hv_{\mathrm{SM}}^2$ from \eqref{eq:5.32} and $U(\phi')^\dagger U(\phi')=I_n$, the potential depends only on $h$:
\begin{equation*}
\begin{aligned}
\Phi^\dagger\Phi&=\frac{(v_{\mathrm{SM}}+h)^2}2,\\
V&=\lambda_H\Bigl(\Phi^\dagger\Phi-\frac{v_{\mathrm{SM}}^2}2\Bigr)^2-\frac{\lambda_Hv_{\mathrm{SM}}^4}4\\
&=\lambda_Hv_{\mathrm{SM}}^2h^2+\lambda_Hv_{\mathrm{SM}}\,h^3+\frac{\lambda_H}4h^4-\frac{\lambda_Hv_{\mathrm{SM}}^4}4 .
\end{aligned}
\tag{5.38}\label{eq:5.38}
\end{equation*}
For the kinetic term, use $\partial_\mu U(\phi')=U(\phi')\,U(\phi')^\dagger\partial_\mu U(\phi')$. Since $(\partial_\mu U(\phi')^\dagger)U(\phi')=-U(\phi')^\dagger\partial_\mu U(\phi')$, $A^{U(\phi')}_\mu$ is Hermitian, so $\Phi_0^\dagger A^{U(\phi')}_\mu\Phi_0$ is real and the cross terms in $(D_\mu\Phi)^\dagger D^\mu\Phi$ cancel. With $\Phi_0^\dagger\Phi_0=v_{\mathrm{SM}}^2/2$,
\begin{equation*}
\begin{aligned}
D_\mu\Phi&=\frac{\partial_\mu h}{v_{\mathrm{SM}}}\,U(\phi')\Phi_0+\Bigl(1+\frac h{v_{\mathrm{SM}}}\Bigr)\bigl(\partial_\mu U(\phi')-ig_{\mathrm{YM}}A_\mu U(\phi')\bigr)\Phi_0\\
&=\frac{\partial_\mu h}{v_{\mathrm{SM}}}\,U(\phi')\Phi_0-ig_{\mathrm{YM}}\Bigl(1+\frac h{v_{\mathrm{SM}}}\Bigr)U(\phi')A^{U(\phi')}_\mu\Phi_0,\\
A^{U(\phi')}_\mu&:=U(\phi')^\dagger A_\mu U(\phi')+\frac i{g_{\mathrm{YM}}}U(\phi')^\dagger\partial_\mu U(\phi'),\\
(D_\mu\Phi)^\dagger D^\mu\Phi&=\frac12\partial_\mu h\,\partial^\mu h\\
&\quad+\Bigl(1+\frac h{v_{\mathrm{SM}}}\Bigr)^2g_{\mathrm{YM}}^2\,\Phi_0^\dagger A^{U(\phi')}_\mu A^{U(\phi')\mu}\Phi_0 .
\end{aligned}
\tag{5.39}\label{eq:5.39}
\end{equation*}
Substituting \eqref{eq:5.38} and \eqref{eq:5.39} into \eqref{eq:5.37},
\begin{equation*}
{\displaystyle 
\begin{aligned}
\mathcal L={}&\frac12\partial_\mu h\,\partial^\mu h-\frac12m_h^2h^2\\
&-\lambda_Hv_{\mathrm{SM}}\,h^3-\frac{\lambda_H}4h^4+\frac{\lambda_Hv_{\mathrm{SM}}^4}4\\
&+\Bigl(1+\frac h{v_{\mathrm{SM}}}\Bigr)^2g_{\mathrm{YM}}^2\,\Phi_0^\dagger A^{U(\phi')}_\mu A^{U(\phi')\mu}\Phi_0\\
&-\frac14F_{\mu\nu}^aF^{a\mu\nu},\\
m_h^2={}&2\lambda_Hv_{\mathrm{SM}}^2=2\mu_H^2 .
\end{aligned}
}
\tag{5.40}\label{eq:5.40}
\end{equation*}
Define the field $A'_\mu$ and its field strength by
\begin{equation*}
\begin{aligned}
A'_\mu&:=A^{U(\phi')}_\mu=U(\phi')^\dagger A_\mu U(\phi')+\frac i{g_{\mathrm{YM}}}U(\phi')^\dagger\partial_\mu U(\phi'),\\
F'_{\mu\nu}&:=\partial_\mu A'_\nu-\partial_\nu A'_\mu-ig_{\mathrm{YM}}[A'_\mu,A'_\nu]=U(\phi')^\dagger F_{\mu\nu}U(\phi') .
\end{aligned}
\tag{5.41}\label{eq:5.41}
\end{equation*}
Changing the field variables $(\Phi,A_\mu)\to(h,\phi'^a,A'_\mu)$ by \eqref{eq:5.36} and \eqref{eq:5.41}, \eqref{eq:5.40} becomes
\begin{equation*}
{\displaystyle 
\begin{aligned}
\mathcal L={}&\frac12\partial_\mu h\,\partial^\mu h-\frac12m_h^2h^2\\
&-\lambda_Hv_{\mathrm{SM}}\,h^3-\frac{\lambda_H}4h^4+\frac{\lambda_Hv_{\mathrm{SM}}^4}4\\
&+\Bigl(1+\frac h{v_{\mathrm{SM}}}\Bigr)^2g_{\mathrm{YM}}^2\,\Phi_0^\dagger A'_\mu A'^\mu\Phi_0\\
&-\frac14F'^a_{\mu\nu}F'^{a\mu\nu} .
\end{aligned}
}
\tag{5.42}\label{eq:5.42}
\end{equation*}
With $A'_\mu=A'^a_\mu T^a$,  and $T^a\Phi_0=0$ for $a\ge2n$,
\begin{equation*}
\begin{aligned}
\mathcal L={}&\frac12\partial_\mu h\,\partial^\mu h-\frac12m_h^2h^2\\
&-\lambda_Hv_{\mathrm{SM}}\,h^3-\frac{\lambda_H}4h^4+\frac{\lambda_Hv_{\mathrm{SM}}^4}4\\
&+\frac12\Bigl(1+\frac h{v_{\mathrm{SM}}}\Bigr)^2\sum_{a,b=1}^{2n-1}(M^2)^{ab}A'^a_\mu A'^{b\mu}\\
&-\frac14F'^a_{\mu\nu}F'^{a\mu\nu},\\
(M^2)^{ab}:={}&g_{\mathrm{YM}}^2\,\Phi_0^\dagger\{G^a,G^b\}\Phi_0.
\end{aligned}
\tag{5.43}\label{eq:5.43}
\end{equation*}

\subsection{Fermions and the Higgs Boson in the Standard Model}\label{book-subsection-section-5.2-2}\label{sm-fermions-and-higgs}

\begin{center}
\renewcommand{\arraystretch}{1.2}
\begin{tabular}{cccc}
\toprule
 & Generation $\alpha$ & $SU(2)_L$ $i$ & Color $c$ \\
\midrule
$(\ell_L)_{\alpha i}$ & $1,2,3$ & $1,2$ & $1$ \\
$(E_R)_\alpha$ & $1,2,3$ & $1$ & $1$ \\
$(Q_L)_{\alpha ic}$ & $1,2,3$ & $1,2$ & $r,b,g$ \\
$(u_R)_{\alpha c}$ & $1,2,3$ & $1$ & $r,b,g$ \\
$(d_R)_{\alpha c}$ & $1,2,3$ & $1$ & $r,b,g$ \\
$H_i$ & $1$ & $1,2$ & $1$ \\
\bottomrule
\end{tabular}
\end{center}

The index of a singlet (entry $1$) is not written explicitly.

The arbitrary choices in the Standard Model are fixed as follows.
\begin{itemize}


\item \textbf{\boldmath Generators of $SU(3)_c$}: $T_c^a=\frac{\lambda^a}2$ on the color index $c$, where the Gell-Mann matrices $\lambda^a$ ($a=1,\ldots,8$) are
\begin{equation*}
\begin{aligned}
\lambda^1&=\begin{pmatrix}0&1&0\\1&0&0\\0&0&0\end{pmatrix},&
\lambda^2&=\begin{pmatrix}0&-i&0\\i&0&0\\0&0&0\end{pmatrix},&
\lambda^3&=\begin{pmatrix}1&0&0\\0&-1&0\\0&0&0\end{pmatrix},\\[6pt]
\lambda^4&=\begin{pmatrix}0&0&1\\0&0&0\\1&0&0\end{pmatrix},&
\lambda^5&=\begin{pmatrix}0&0&-i\\0&0&0\\i&0&0\end{pmatrix},&
\lambda^6&=\begin{pmatrix}0&0&0\\0&0&1\\0&1&0\end{pmatrix},\\[6pt]
\lambda^7&=\begin{pmatrix}0&0&0\\0&0&-i\\0&i&0\end{pmatrix},&
\lambda^8&=\frac1{\sqrt3}\begin{pmatrix}1&0&0\\0&1&0\\0&0&-2\end{pmatrix},
\end{aligned}
\tag{5.44}\label{eq:5.44}
\end{equation*}
with $c=r,b,g$ corresponding to their rows and columns $1,2,3$.

Up to the antisymmetry \eqref{eq:5.15}, the nonzero structure constants in $[T_c^a,T_c^b]=if^{abc}T_c^c$ are
\begin{equation*}
\begin{aligned}
f^{123}&=1,\\[4pt]
f^{147}=f^{246}=f^{257}=f^{345}&=\tfrac12,\qquad
f^{156}=f^{367}=-\tfrac12,\\[4pt]
f^{458}=f^{678}&=\tfrac{\sqrt3}{2}.
\end{aligned}
\tag{5.45}\label{eq:5.45}
\end{equation*}

\item \textbf{\boldmath Generators of $SU(2)_L$}: $T_L^a=\frac{\sigma^a}2$ on the index $i$, where the Pauli matrices $\sigma^a$ ($a=1,2,3$) are
\begin{equation*}
\sigma^1=\begin{pmatrix}0&1\\1&0\end{pmatrix},\qquad
\sigma^2=\begin{pmatrix}0&-i\\i&0\end{pmatrix},\qquad
\sigma^3=\begin{pmatrix}1&0\\0&-1\end{pmatrix},
\tag{5.46}\label{eq:5.46}
\end{equation*}
so that $[T_L^a,T_L^b]=i\epsilon^{abc}T_L^c$.

\item \textbf{\boldmath Generator of $U(1)_Y$}: the hypercharge $Y$ are

\begin{center}
\renewcommand{\arraystretch}{1.2}
\begin{tabular}{ccccccc}
\toprule
 & $\ell_L$ & $E_R$ & $Q_L$ & $u_R$ & $d_R$ & $H$ \\
\midrule
$Y$ & $-\frac12$ & $-1$ & $\frac16$ & $\frac23$ & $-\frac13$ & $\frac12$ \\
\bottomrule
\end{tabular}
\end{center}

\item \textbf{\boldmath Generators on the fields}: on the fields of the table, the generators of $SU(3)_c$, $SU(2)_L$ and $U(1)_Y$ are
\begin{equation*}
\begin{aligned}
\ell_L&:\quad T_c^a=I_i\otimes0\otimes I_\alpha, &
T_L^a&=\frac{\sigma^a}{2}\otimes I_c\otimes I_\alpha, &
Y&=-\frac12\,I_i\otimes I_c\otimes I_\alpha,\\
E_R&:\quad T_c^a=I_i\otimes0\otimes I_\alpha, &
T_L^a&=0\otimes I_c\otimes I_\alpha, &
Y&=-I_i\otimes I_c\otimes I_\alpha,\\
Q_L&:\quad T_c^a=I_i\otimes\frac{\lambda^a}{2}\otimes I_\alpha, &
T_L^a&=\frac{\sigma^a}{2}\otimes I_c\otimes I_\alpha, &
Y&=\frac16\,I_i\otimes I_c\otimes I_\alpha,\\
u_R&:\quad T_c^a=I_i\otimes\frac{\lambda^a}{2}\otimes I_\alpha, &
T_L^a&=0\otimes I_c\otimes I_\alpha, &
Y&=\frac23\,I_i\otimes I_c\otimes I_\alpha,\\
d_R&:\quad T_c^a=I_i\otimes\frac{\lambda^a}{2}\otimes I_\alpha, &
T_L^a&=0\otimes I_c\otimes I_\alpha, &
Y&=-\frac13\,I_i\otimes I_c\otimes I_\alpha,\\
H&:\quad T_c^a=I_i\otimes0\otimes I_\alpha, &
T_L^a&=\frac{\sigma^a}{2}\otimes I_c\otimes I_\alpha, &
Y&=\frac12\,I_i\otimes I_c\otimes I_\alpha,
\end{aligned}
\tag{5.47}\label{eq:5.47}
\end{equation*}
where $I_i$, $I_c$, $I_\alpha$ are the identities on the indices $i$, $c$, $\alpha$ (the number $1$ on a singlet), and $0$ is the generator on a singlet.

\item \textbf{\boldmath Signs of the couplings}: $g_s,g_2,g_1>0$ in the covariant derivative
\begin{equation*}
D_\mu=\partial_\mu-ig_sA_\mu^aT_c^a-ig_2W_\mu^aT_L^a-ig_1YB_\mu .
\tag{5.48}\label{eq:5.48}
\end{equation*}

\item \textbf{\boldmath Vacuum expectation value of $H$}:
\begin{equation*}
\langle H\rangle=\frac1{\sqrt2}\begin{pmatrix}0\\v_{\mathrm{SM}}\end{pmatrix},\qquad v_{\mathrm{SM}}>0,
\tag{5.49}\label{eq:5.49}
\end{equation*}
which is $\Phi_0$ of \eqref{eq:5.33} with $n=2$.
\end{itemize}


the mass matrix of \eqref{eq:5.43} in the basis $(W'^1_\mu,W'^2_\mu,W'^3_\mu,B'_\mu)$ is
\begin{equation*}
M^2=\frac{v_{\mathrm{SM}}^2}4\begin{pmatrix}g_2^2&0&0&0\\0&g_2^2&0&0\\0&0&g_2^2&-g_1g_2\\0&0&-g_1g_2&g_1^2\end{pmatrix}.
\tag{5.50}\label{eq:5.50}
\end{equation*}

\subsection{Yukawa Coupling and Flavor}\label{book-subsection-section-5.2-3}\label{field-redefinitions-and-flavor}

Besides the gauge interactions in the covariant derivative \eqref{eq:5.48}, the Lagrangian density contains the Yukawa terms, whose couplings $Y_l,Y_d,Y_u$ are complex $3\times3$ matrices carrying only the generation indices:
\begin{equation*}
\begin{aligned}
\mathcal L\supset{}&-\overline{(\ell_L)_\alpha}(Y_l)_{\alpha\beta}H(E_R)_\beta-\overline{(E_R)_\beta}(Y_l)^*_{\alpha\beta}H^\dagger(\ell_L)_\alpha\\
&-\overline{(Q_L)_\alpha}(Y_d)_{\alpha\beta}H(d_R)_\beta-\overline{(d_R)_\beta}(Y_d)^*_{\alpha\beta}H^\dagger(Q_L)_\alpha\\
&-\overline{(Q_L)_\alpha}(Y_u)_{\alpha\beta}H^c(u_R)_\beta-\overline{(u_R)_\beta}(Y_u)^*_{\alpha\beta}(H^c)^\dagger(Q_L)_\alpha,\\[6pt]
H^c:={}&i\sigma^2H^*=\begin{pmatrix}0&1\\-1&0\end{pmatrix}\begin{pmatrix}H_1^*\\H_2^*\end{pmatrix}=\begin{pmatrix}H_2^*\\-H_1^*\end{pmatrix},
\end{aligned}
\tag{5.51}\label{eq:5.51}
\end{equation*}
where $H=(H_1,H_2)^T$ in the index $i$ and the repeated $\alpha,\beta$ are summed.

Each Yukawa term in \eqref{eq:5.51} is invariant under the $SU(3)_c\times SU(2)_L\times U(1)_Y$ gauge transformations.

For example, consider the term $-\overline{(Q_L)_\alpha}(Y_u)_{\alpha\beta}H^c(u_R)_\beta$. A gauge transformation acts on each field $\psi$ of \eqref{eq:5.47} as $\psi\to e^{i\theta_c^aT_c^a+i\theta_L^aT_L^a+i\theta_YY}\psi$ with real $\theta_c^a(x)$, $\theta_L^a(x)$, $\theta_Y(x)$.

With $U_c:=e^{i\theta_c^a\lambda^a/2}$ and $U_L:=e^{i\theta_L^a\sigma^a/2}$,
\begin{equation*}
\begin{aligned}
(Q_L)_\alpha&\to e^{\frac i6\theta_Y}(U_L\otimes U_c)(Q_L)_\alpha,\\
(u_R)_\beta&\to e^{\frac{2i}3\theta_Y}U_c(u_R)_\beta,\\
H&\to e^{\frac i2\theta_Y}U_LH,\\
H^c&\to i\sigma^2\bigl(e^{\frac i2\theta_Y}U_LH\bigr)^*=e^{-\frac i2\theta_Y}\bigl[i\sigma^2U_L^*(i\sigma^2)^\dagger\bigr]i\sigma^2H^*=e^{-\frac i2\theta_Y}U_LH^c,\\
\overline{(Q_L)_\alpha}H^c(u_R)_\beta&\to e^{i(-\frac16-\frac12+\frac23)\theta_Y}\,\overline{(Q_L)_\alpha}(U_L\otimes U_c)^\dagger(U_L\otimes U_c)H^c(u_R)_\beta=\overline{(Q_L)_\alpha}H^c(u_R)_\beta,
\end{aligned}
\tag{5.52}\label{eq:5.52}
\end{equation*}
where $i\sigma^2U_L^*(i\sigma^2)^\dagger=U_L$ because $i\sigma^2(\sigma^a)^*(i\sigma^2)^\dagger=-\sigma^a$ by \eqref{eq:5.46}.

With the generation indices written explicitly, the Lagrangian density of the Standard Model contains the fermion kinetic terms, with the covariant derivative \eqref{eq:5.48} and the generators \eqref{eq:5.47}, and the Yukawa terms \eqref{eq:5.51}:
\begin{equation*}
\begin{aligned}
\mathcal L\supset{}&\overline{(\ell_L)_\alpha}\,i\gamma^\mu\Bigl(\partial_\mu-ig_2W_\mu^a\frac{\sigma^a}2+\frac i2g_1B_\mu\Bigr)(\ell_L)_\alpha\\
&+\overline{(E_R)_\alpha}\,i\gamma^\mu\bigl(\partial_\mu+ig_1B_\mu\bigr)(E_R)_\alpha\\
&+\overline{(Q_L)_\alpha}\,i\gamma^\mu\Bigl(\partial_\mu-ig_sA_\mu^a\frac{\lambda^a}2-ig_2W_\mu^a\frac{\sigma^a}2-\frac i6g_1B_\mu\Bigr)(Q_L)_\alpha\\
&+\overline{(u_R)_\alpha}\,i\gamma^\mu\Bigl(\partial_\mu-ig_sA_\mu^a\frac{\lambda^a}2-\frac{2i}3g_1B_\mu\Bigr)(u_R)_\alpha\\
&+\overline{(d_R)_\alpha}\,i\gamma^\mu\Bigl(\partial_\mu-ig_sA_\mu^a\frac{\lambda^a}2+\frac i3g_1B_\mu\Bigr)(d_R)_\alpha\\[6pt]
&-\overline{(\ell_L)_\alpha}(Y_l)_{\alpha\beta}H(E_R)_\beta-\overline{(E_R)_\beta}(Y_l)^*_{\alpha\beta}H^\dagger(\ell_L)_\alpha\\
&-\overline{(Q_L)_\alpha}(Y_d)_{\alpha\beta}H(d_R)_\beta-\overline{(d_R)_\beta}(Y_d)^*_{\alpha\beta}H^\dagger(Q_L)_\alpha\\
&-\overline{(Q_L)_\alpha}(Y_u)_{\alpha\beta}H^c(u_R)_\beta-\overline{(u_R)_\beta}(Y_u)^*_{\alpha\beta}(H^c)^\dagger(Q_L)_\alpha,
\end{aligned}
\tag{5.53}\label{eq:5.53}
\end{equation*}
where $\sigma^a/2$ acts on the index $i$, $\lambda^a/2$ on the index $c$, and the repeated $\alpha,\beta$ are summed.
\begin{itemize}


\item \textbf{\boldmath Generation transformation}: with a constant $S_\psi\in U(3)$ for each field, acting only on the generation index,
\begin{equation*}
(\psi)_\alpha\to(S_\psi)_{\alpha\beta}(\psi)_\beta,\qquad\psi=\ell_L,E_R,Q_L,u_R,d_R .
\tag{5.54}\label{eq:5.54}
\end{equation*}
\end{itemize}

Inserting $S_\psi S_\psi^\dagger=I_\alpha$ into the kinetic and Yukawa terms of \eqref{eq:5.53},
\begin{equation*}
\begin{aligned}
\overline{(\psi)_\alpha}\,i\gamma^\mu D_\mu(\psi)_\alpha
&=\overline{(\psi)_\alpha}(S_\psi S_\psi^\dagger)_{\alpha\beta}\,i\gamma^\mu D_\mu(\psi)_\beta
=\overline{(S_\psi^\dagger\psi)_\alpha}\,i\gamma^\mu D_\mu(S_\psi^\dagger\psi)_\alpha,\\
\overline{(\ell_L)_\alpha}(Y_l)_{\alpha\beta}H(E_R)_\beta
&=\overline{(\ell_L)_\alpha}(S_{\ell_L}S_{\ell_L}^\dagger Y_lS_{E_R}S_{E_R}^\dagger)_{\alpha\beta}H(E_R)_\beta\\
&=\overline{(S_{\ell_L}^\dagger\ell_L)_\alpha}(S_{\ell_L}^\dagger Y_lS_{E_R})_{\alpha\beta}H(S_{E_R}^\dagger E_R)_\beta,
\end{aligned}
\tag{5.55}\label{eq:5.55}
\end{equation*}
and likewise for the other terms of \eqref{eq:5.51}, where $S_\psi^\dagger$ commutes with $i\gamma^\mu D_\mu$ because $S_\psi$ is constant and acts only on the generation index. Replacing the fields by $S_\psi\psi$ as in \eqref{eq:5.54} turns each $S_\psi^\dagger\psi$ into $\psi$, so the kinetic terms are unchanged and \eqref{eq:5.51} keeps its form with the Yukawa matrices replaced by
\begin{equation*}
Y_l\to S_{\ell_L}^\dagger Y_lS_{E_R},\qquad Y_d\to S_{Q_L}^\dagger Y_dS_{d_R},\qquad Y_u\to S_{Q_L}^\dagger Y_uS_{u_R} .
\tag{5.56}\label{eq:5.56}
\end{equation*}
Since \eqref{eq:5.54} only changes the field variables, Yukawa matrices related by \eqref{eq:5.56} describe the same physics. The parametrization of the Yukawa interactions by $Y_l,Y_d,Y_u$ is therefore redundant: the number of physical parameters is smaller than the number of real parameters in $Y_l,Y_d,Y_u$.

\topic{Parametrization of the Yukawa Couplings: Singular Value Decomposition}

\begin{qftinset}

\divtitle{$\circ$ Leptons}

By the singular value decomposition, the complex $3\times3$ matrix $Y_l$ can be written as
\begin{equation*}
Y_l=S_1\begin{pmatrix}y_1&0&0\\0&y_2&0\\0&0&y_3\end{pmatrix}S_2^\dagger,\qquad S_1,S_2\in U(3),\qquad 0\le y_1\le y_2\le y_3,
\tag{5.57}\label{eq:5.57}
\end{equation*}
where $y_1^2,y_2^2,y_3^2$ are the eigenvalues of $Y_l^\dagger Y_l$.
\begin{itemize}


\item \textbf{\boldmath Orbit}: For fixed $Y_{l,\mathrm{diag}}:=\operatorname{diag}(y_1,y_2,y_3)$, the orbit of $Y_{l,\mathrm{diag}}$ is the image of the map
\begin{equation*}
\begin{aligned}
\mathcal F_l: U(3)\times U(3)&\to M_3(\mathbb C),\\
(S_1,S_2)&\mapsto S_1Y_{l,\mathrm{diag}}S_2^\dagger,
\end{aligned}
\tag{5.58}\label{eq:5.58}
\end{equation*}
where $M_3(\mathbb C)$ is the space of complex $3\times3$ matrices, of real dimension $18$. The question is the real dimension of the orbit.

\item \textbf{\boldmath Stabilizer}: The subgroup of pairs $(S_1,S_2)$ with $S_1Y_{l,\mathrm{diag}}S_2^\dagger=Y_{l,\mathrm{diag}}$.
\end{itemize}


For $0<y_1<y_2<y_3$, the stabilizer condition and its Hermitian conjugate give $S_1Y_{l,\mathrm{diag}}^2S_1^\dagger=Y_{l,\mathrm{diag}}^2=S_2Y_{l,\mathrm{diag}}^2S_2^\dagger$. The diagonal entries of $Y_{l,\mathrm{diag}}^2$ are distinct, so $S_1$ and $S_2$ are diagonal, and since $Y_{l,\mathrm{diag}}$ is invertible, the stabilizer is
\begin{equation*}
S_1=S_2=\operatorname{diag}(e^{i\varphi_1},e^{i\varphi_2},e^{i\varphi_3}),\qquad \varphi_k\in\mathbb R .
\tag{5.59}\label{eq:5.59}
\end{equation*}
The orbit is $U(3)\times U(3)$ modulo the stabilizer, so its real dimension is $2\dim U(3)-3=18-3=15$. The spectrum of $Y_l^\dagger Y_l$ fixes the ordered singular values, so by \eqref{eq:5.57} every $Y_l$ belongs to exactly one orbit, and for $0<y_1<y_2<y_3$ the real parameters of $Y_l$ split as $18=15+3$: $15$ along the orbit and the $3$ singular values. When singular values coincide or vanish, the stabilizer is larger and the orbit has dimension less than $15$.

Under \eqref{eq:5.56}, $S_{\ell_L}$ and $S_{E_R}$ move $Y_l$ along its orbit. In particular, the generation transformation \eqref{eq:5.54} with $S_{\ell_L}=S_1$ and $S_{E_R}=S_2$ gives
\begin{equation*}
{\displaystyle Y_l\to S_{\ell_L}^\dagger Y_lS_{E_R}=\begin{pmatrix}y_1&0&0\\0&y_2&0\\0&0&y_3\end{pmatrix}}\,.
\tag{5.60}\label{eq:5.60}
\end{equation*}
The $15$ real parameters along the orbit are therefore unphysical generation degrees of freedom, and $Y_l$ has $3$ physical parameters, $y_1,y_2,y_3$.

Excluding the unphysical generation degrees of freedom, what remains is flavor: for $y_1<y_2<y_3$ the three generations are physically distinct, and in the basis of \eqref{eq:5.60}, called the flavor basis, the generation index $\alpha=1,2,3$ of $\ell_L$ and $E_R$ labels the flavors $e,\mu,\tau$.

\end{qftinset}

\begin{qftinset}

\divtitle{$\circ$ Quarks}

The matrices $Y_d$ and $Y_u$ have the singular value decompositions
\begin{equation*}
\begin{aligned}
Y_d&=S_3\begin{pmatrix}y_d&0&0\\0&y_s&0\\0&0&y_b\end{pmatrix}S_4^\dagger,&&0\le y_d\le y_s\le y_b,\\
Y_u&=S_5\begin{pmatrix}y_u&0&0\\0&y_c&0\\0&0&y_t\end{pmatrix}S_6^\dagger,&&0\le y_u\le y_c\le y_t,
\end{aligned}
\tag{5.61}\label{eq:5.61}
\end{equation*}
with $S_3,S_4,S_5,S_6\in U(3)$. Unlike $Y_l$, they share $S_{Q_L}$ in \eqref{eq:5.56}, which can remove only one of the left matrices $S_3$ and $S_5$. The generation transformation \eqref{eq:5.54} with $S_{Q_L}=S_3$, $S_{d_R}=S_4$ and $S_{u_R}=S_6$ gives
\begin{equation*}
{\displaystyle Y_d\to\begin{pmatrix}y_d&0&0\\0&y_s&0\\0&0&y_b\end{pmatrix},\qquad Y_u\to V_{\mathrm{CKM}}^\dagger\begin{pmatrix}y_u&0&0\\0&y_c&0\\0&0&y_t\end{pmatrix}},\qquad V_{\mathrm{CKM}}:=S_5^\dagger S_3\in U(3).
\tag{5.62}\label{eq:5.62}
\end{equation*}
The generation transformations that keep the form \eqref{eq:5.62} keep the diagonal matrices $Y_dY_d^\dagger$, $Y_d^\dagger Y_d$ and $Y_u^\dagger Y_u$, whose entries are distinct for $0<y_d<y_s<y_b$ and $0<y_u<y_c<y_t$, so $S_{Q_L},S_{d_R},S_{u_R}$ are diagonal; keeping $Y_d$ also requires $S_{d_R}=S_{Q_L}$. These transformations, and their action on $V_{\mathrm{CKM}}$ by \eqref{eq:5.56}, are
\begin{equation*}
\begin{aligned}
&S_{Q_L}=S_{d_R}=\operatorname{diag}(e^{i\varphi_1},e^{i\varphi_2},e^{i\varphi_3}),\qquad S_{u_R}=\operatorname{diag}(e^{i\varphi'_1},e^{i\varphi'_2},e^{i\varphi'_3}),\\
&(V_{\mathrm{CKM}})_{\alpha\beta}\to e^{i(\varphi_\beta-\varphi'_\alpha)}(V_{\mathrm{CKM}})_{\alpha\beta}.
\end{aligned}
\tag{5.63}\label{eq:5.63}
\end{equation*}
Of the $9$ real parameters of $V_{\mathrm{CKM}}$, $3$ angles and $6$ phases, \eqref{eq:5.63} removes $5$ phases, since $V_{\mathrm{CKM}}$ is unchanged when all six $\varphi_\alpha,\varphi'_\alpha$ are equal. $V_{\mathrm{CKM}}$ therefore has $3$ angles and $1$ phase, and $Y_d,Y_u$ have $6+4=10$ physical parameters out of their $36$ real parameters; with the $3$ of $Y_l$, the Yukawa interactions \eqref{eq:5.51} have $13$ physical parameters out of $54$.

In the basis of \eqref{eq:5.62}, the generation indices of $d_R$ and $u_R$ label the flavors $d,s,b$ and $u,c,t$, and that of $Q_L$ labels $d,s,b$ in $Y_d$; in $Y_u$, the flavors $u,c,t$ are labeled by the index of $V_{\mathrm{CKM}}Q_L$. Unless $V_{\mathrm{CKM}}$ is diagonal, the generation index of $Q_L$ cannot label the flavors of $Y_d$ and $Y_u$ at the same time.

\end{qftinset}

\setcounter{section}{2}
\section{\texorpdfstring{$W^\pm$}{W±}, \texorpdfstring{$Z^0$}{Z0}, \texorpdfstring{$\gamma$}{γ}: Electroweak Theory}\label{book-section-5.3}\label{electroweak-gauge-theory}\label{electroweak-theory}

\subsection{Gauge Symmetry \texorpdfstring{$SU(2)_L\times U(1)_Y$}{SU(2)L× U(1)Y}}\label{book-subsection-section-5.3-1}\label{qed}\label{electroweak-gauge-symmetry}
\begin{itemize}


\item \textbf{\boldmath Generators and covariant derivative}: the generators $T^a$ ($a=1,2,3$) of $SU(2)_L$ and $Y$ of $U(1)_Y$ on the fields of \hyperref[sm-fermions]{Section 5.2}, and the covariant derivative, are
\begin{equation*}
\begin{aligned}
\ell_L&:\quad T^a=\frac{\sigma^a}{2}\otimes I_\alpha\otimes I_s, &
Y&=-\frac12\,I_i\otimes I_\alpha\otimes I_s,\\
E_R&:\quad T^a=0, &
Y&=-I_\alpha\otimes I_s,\\
Q_L&:\quad T^a=\frac{\sigma^a}{2}\otimes I_c\otimes I_\alpha\otimes I_s, &
Y&=\frac16\,I_i\otimes I_c\otimes I_\alpha\otimes I_s,\\
u_R&:\quad T^a=0, &
Y&=\frac23\,I_c\otimes I_\alpha\otimes I_s,\\
d_R&:\quad T^a=0, &
Y&=-\frac13\,I_c\otimes I_\alpha\otimes I_s,\\[6pt]
D_\mu&:=\partial_\mu-ig_2W_\mu^aT^a-ig_1YB_\mu,
\end{aligned}
\tag{5.64}\label{eq:5.64}
\end{equation*}
where $\sigma^a/2$ acts on the $SU(2)_L$ index $i$, and $I_i$, $I_c$, $I_\alpha$, $I_s$ are the identities on the indices $i$, $c$, $\alpha$, $s$.

In $D_\mu$, $W_\mu^a$ and $B_\mu$ are the real $SU(2)_L$ and $U(1)_Y$ gauge fields, $g_2,g_1>0$ are the couplings, and $T^a$, $Y$ are those of the field on which $D_\mu$ acts.
\end{itemize}


\subsection{Interaction}\label{book-subsection-section-5.3-2}\label{electroweak-interaction}
\begin{itemize}


\item \textbf{\boldmath Electric charge, Weinberg angle and the fields $W_\mu^\pm$, $Z_\mu$, $A_\mu^{\mathrm{em}}$}: with $T^a$, $Y$ those of a field of \eqref{eq:5.64},
\begin{equation*}
\begin{aligned}
Q&:=T^3+Y,\qquad T^\pm:=T^1\pm iT^2,\\
\tan\theta_W&:=\frac{g_1}{g_2},\qquad s_W:=\sin\theta_W=\frac{g_1}{\sqrt{g_1^2+g_2^2}},\qquad c_W:=\cos\theta_W=\frac{g_2}{\sqrt{g_1^2+g_2^2}},\\
e&:=g_2s_W=g_1c_W=\frac{g_1g_2}{\sqrt{g_1^2+g_2^2}},\\
W_\mu^\pm&:=\frac{W_\mu^1\mp iW_\mu^2}{\sqrt2},\qquad
\begin{pmatrix}Z_\mu\\A_\mu^{\mathrm{em}}\end{pmatrix}:=\begin{pmatrix}c_W&-s_W\\s_W&c_W\end{pmatrix}\begin{pmatrix}W_\mu^3\\B_\mu\end{pmatrix}.
\end{aligned}
\tag{5.65}\label{eq:5.65}
\end{equation*}
Thus $Q=\operatorname{diag}(0,-1)$ on $\ell_L$, $-1$ on $E_R$, $\operatorname{diag}(\frac23,-\frac13)$ on $Q_L$, $\frac23$ on $u_R$ and $-\frac13$ on $d_R$, where $\operatorname{diag}$ acts on the index $i$ and each entry multiplies the identities on the other indices.
\end{itemize}


The Lagrangian density contains the lepton and quark interaction terms
\begin{equation*}
\begin{aligned}
\mathcal L&\supset\overline{\ell_L}\gamma^\mu\Bigl(g_2W_\mu^a\frac{\sigma^a}2-\frac{g_1}2B_\mu\Bigr)\ell_L-g_1B_\mu\,\overline{E_R}\gamma^\mu E_R,\\
\mathcal L&\supset\overline{Q_L}\gamma^\mu\Bigl(g_2W_\mu^a\frac{\sigma^a}2+\frac{g_1}6B_\mu\Bigr)Q_L+\frac23g_1B_\mu\,\overline{u_R}\gamma^\mu u_R-\frac13g_1B_\mu\,\overline{d_R}\gamma^\mu d_R.
\end{aligned}
\tag{5.66}\label{eq:5.66}
\end{equation*}
Inverting \eqref{eq:5.65},
\begin{equation*}
\begin{aligned}
g_2W_\mu^aT^a+g_1YB_\mu
&=\frac{g_2}{\sqrt2}\bigl(W_\mu^+T^++W_\mu^-T^-\bigr)+A_\mu^{\mathrm{em}}\bigl(g_2s_WT^3+g_1c_WY\bigr)+Z_\mu\bigl(g_2c_WT^3-g_1s_WY\bigr)\\
&=\frac{g_2}{\sqrt2}\bigl(W_\mu^+T^++W_\mu^-T^-\bigr)+eQA_\mu^{\mathrm{em}}+\frac{g_2}{c_W}\bigl(T^3-s_W^2Q\bigr)Z_\mu .
\end{aligned}
\tag{5.67}\label{eq:5.67}
\end{equation*}
Write $\ell_L=(\nu_L,l_L)^T$ and $Q_L=(u_L,d_L)^T$ in the index $i$. Then $\overline{\ell_L}\gamma^\mu T^+\ell_L=\overline{\nu_L}\gamma^\mu l_L$ and $\overline{Q_L}\gamma^\mu T^+Q_L=\overline{u_L}\gamma^\mu d_L$, while $T^\pm=T^3=0$ on $E_R$, $u_R$, $d_R$.

Since $\bar\psi_1\gamma^\mu\psi_2=\overline{P_L\psi_1}\gamma^\mu P_L\psi_2+\overline{P_R\psi_1}\gamma^\mu P_R\psi_2$ by \eqref{eq:1.137}, the $A_\mu^{\mathrm{em}}$ terms combine into the Dirac fields $l:=l_L+E_R$, $u:=u_L+u_R$ and $d:=d_L+d_R$.

Substituting \eqref{eq:5.67} into \eqref{eq:5.66}, the lepton terms are
\begin{equation*}
{\displaystyle 
\begin{aligned}
\mathcal L\supset{}&\frac{g_2}{\sqrt2}\bigl(W_\mu^+\,\overline{\nu_L}\gamma^\mu l_L+W_\mu^-\,\overline{l_L}\gamma^\mu\nu_L\bigr)-eA_\mu^{\mathrm{em}}\,\bar l\gamma^\mu l\\
&+\frac{g_2}{c_W}Z_\mu\Bigl[\frac12\,\overline{\nu_L}\gamma^\mu\nu_L+\Bigl(-\frac12+s_W^2\Bigr)\overline{l_L}\gamma^\mu l_L+s_W^2\,\overline{E_R}\gamma^\mu E_R\Bigr],
\end{aligned}
}
\tag{5.68}\label{eq:5.68}
\end{equation*}
and the quark terms are
\begin{equation*}
{\displaystyle 
\begin{aligned}
\mathcal L\supset{}&\frac{g_2}{\sqrt2}\bigl(W_\mu^+\,\overline{u_L}\gamma^\mu d_L+W_\mu^-\,\overline{d_L}\gamma^\mu u_L\bigr)+eA_\mu^{\mathrm{em}}\Bigl(\frac23\,\bar u\gamma^\mu u-\frac13\,\bar d\gamma^\mu d\Bigr)\\
&+\frac{g_2}{c_W}Z_\mu\Bigl[\Bigl(\frac12-\frac23s_W^2\Bigr)\overline{u_L}\gamma^\mu u_L-\frac23s_W^2\,\overline{u_R}\gamma^\mu u_R\\
&\qquad\qquad+\Bigl(-\frac12+\frac13s_W^2\Bigr)\overline{d_L}\gamma^\mu d_L+\frac13s_W^2\,\overline{d_R}\gamma^\mu d_R\Bigr].
\end{aligned}
}
\tag{5.69}\label{eq:5.69}
\end{equation*}

\setcounter{section}{3}
\section{Gluon \texorpdfstring{$g$}{g}: Quantum Chromodynamics}\label{book-section-5.4}\label{qcd}

\subsection{Gauge Symmetry \texorpdfstring{$SU(3)_c$}{SU(3)c}}\label{book-subsection-section-5.4-1}\label{qcd-gauge-symmetry}
\begin{itemize}


\item \textbf{\boldmath Generators and covariant derivative}: the generators $T^a$ ($a=1,\ldots,8$) of $SU(3)_c$ on the fields of \hyperref[sm-fermions]{Section 5.2}, and the covariant derivative, are
\begin{equation*}
\begin{aligned}
\ell_L&:\quad T^a=0,\\
E_R&:\quad T^a=0,\\
Q_L&:\quad T^a=I_i\otimes\frac{\lambda^a}{2}\otimes I_\alpha\otimes I_s,\\
u_R&:\quad T^a=\frac{\lambda^a}{2}\otimes I_\alpha\otimes I_s,\\
d_R&:\quad T^a=\frac{\lambda^a}{2}\otimes I_\alpha\otimes I_s,\\[6pt]
D_\mu&:=\partial_\mu-ig_sA_\mu^aT^a,
\end{aligned}
\tag{5.70}\label{eq:5.70}
\end{equation*}
where $\lambda^a/2$ acts on the color index $c$, and $I_i$, $I_\alpha$, $I_s$ are the identities on the indices $i$, $\alpha$, $s$.

In $D_\mu$, $A_\mu^a$ are the real $SU(3)_c$ gauge fields, $g_s>0$ is the coupling, and $T^a$ are those of the field on which $D_\mu$ acts.
\end{itemize}


The gauge group of QCD is $SU(3)$. By \eqref{eq:5.7} with $n=3$, $\mathfrak{su}(3)$ has a basis $\{iT^a\}$, $a=1,\ldots,8$, which we take orthonormal in the sense of \eqref{eq:5.13}. The standard choice is $T^a:=\tfrac12\lambda^a$, where the Gell-Mann matrices are

Each $\lambda^a$ is Hermitian and traceless, and $\operatorname{tr}(\lambda^a\lambda^b)=2\delta^{ab}$, so $\operatorname{tr}(T^aT^b)=\tfrac12\delta^{ab}$ as in \eqref{eq:5.13}. By \eqref{eq:5.14}, $f^{abc}=-2i\operatorname{tr}([T^a,T^b]T^c)$; up to the antisymmetry \eqref{eq:5.15}, the nonzero structure constants are

The quadratic Casimir $T^aT^a$ commutes with every $T^b$, since $f^{bac}$ is antisymmetric in $a,c$:
\begin{equation*}
[T^b,T^aT^a]=[T^b,T^a]T^a+T^a[T^b,T^a]=if^{bac}\bigl(T^cT^a+T^aT^c\bigr)=0.
\tag{5.71}\label{eq:5.71}
\end{equation*}
The defining representation of $SU(3)$ is irreducible, so by Schur's lemma $T^aT^a=C_FI_3$. Taking the trace with \eqref{eq:5.13},
\begin{equation*}
3C_F=\operatorname{tr}(T^aT^a)=\tfrac12\cdot8
\quad\Longrightarrow\quad
{\displaystyle T^aT^a=C_FI_3,\qquad C_F=\tfrac43.}
\tag{5.72}\label{eq:5.72}
\end{equation*}

\topic{Quark Fields}

For each flavor $f\in\{u,d,s,c,b,t\}$ the quark field is a color triplet, i.e. the multiplet \eqref{eq:5.16} with $n=3$, with Dirac adjoint as in \eqref{eq:5.17}:
\begin{equation*}
q_f(x):=\begin{pmatrix}\psi_{f,1}(x)\\\psi_{f,2}(x)\\\psi_{f,3}(x)\end{pmatrix},
\qquad
\bar q_f(x):=\bigl(\bar\psi_{f,1}(x),\bar\psi_{f,2}(x),\bar\psi_{f,3}(x)\bigr).
\tag{5.73}\label{eq:5.73}
\end{equation*}
Throughout this section $\gamma^\mu$ acts on the spinor index of each $\psi_{f,a}$ and the $3\times3$ matrices $T^a$ act on the color index, so $\gamma^\mu$ commutes with every matrix built from the $T^a$.

\topic{Local $SU(3)$ Transformation}

Let $\theta^a:\mathbb R^4\to\mathbb R$, $a=1,\ldots,8$, and set
\begin{equation*}
g(x):=e^{i\theta^a(x)T^a}\in SU(3).
\tag{5.74}\label{eq:5.74}
\end{equation*}
This is the internal transformation \eqref{eq:1.72} with $g$ now depending on $x$. The quark fields transform as
\begin{equation*}
q_f'(x)=g(x)\,q_f(x),\qquad
\bar q_f'(x)=\bar q_f(x)\,g^\dagger(x).
\tag{5.75}\label{eq:5.75}
\end{equation*}
Since $\partial_\mu(gq_f)=g\,\partial_\mu q_f+(\partial_\mu g)\,q_f$, the derivative term of \eqref{eq:5.18} is not invariant under \eqref{eq:5.75}:
\begin{equation*}
\bar q_f'\,i\gamma^\mu\partial_\mu q_f'
=\bar q_f\,i\gamma^\mu\partial_\mu q_f
+\bar q_f\,i\gamma^\mu g^\dagger(\partial_\mu g)\,q_f.
\tag{5.76}\label{eq:5.76}
\end{equation*}

\topic{Covariant Derivative and Gluon Fields}

Introduce eight real vector fields $A^a_\mu$, the gluon fields, and define
\begin{equation*}
A_\mu:=A^a_\mu T^a,\qquad
D_\mu:=\partial_\mu-ig_sA_\mu,
\tag{5.77}\label{eq:5.77}
\end{equation*}
where $g_s\in\mathbb R$ is the strong coupling constant and $D_\mu$ acts on color triplets. Let $A_\mu'$ denote the transformed gluon field and $D_\mu':=\partial_\mu-ig_sA_\mu'$. We require $D_\mu'q_f'=g\,D_\mu q_f$:
\begin{equation*}
\begin{aligned}
D_\mu'(gq_f)&=g\,\partial_\mu q_f+(\partial_\mu g)\,q_f-ig_sA_\mu'\,g\,q_f,\\[6pt]
g\,D_\mu q_f&=g\,\partial_\mu q_f-ig_s\,gA_\mu\,q_f.
\end{aligned}
\tag{5.78}\label{eq:5.78}
\end{equation*}
Equating the two for all $q_f$ gives $A_\mu'g=gA_\mu-\frac{i}{g_s}\partial_\mu g$, i.e.
\begin{equation*}
{\displaystyle 
A_\mu'=gA_\mu g^\dagger-\frac{i}{g_s}(\partial_\mu g)\,g^\dagger
=gA_\mu g^\dagger+\frac{i}{g_s}\,g\,\partial_\mu g^\dagger.
}
\tag{5.79}\label{eq:5.79}
\end{equation*}
Differentiating $gg^\dagger=I_3$ gives $\bigl((\partial_\mu g)g^\dagger\bigr)^\dagger=-(\partial_\mu g)g^\dagger$, and Jacobi's formula $\partial_\mu\det g=\det g\,\operatorname{tr}(g^\dagger\partial_\mu g)$ with $\det g=1$ gives $\operatorname{tr}\bigl((\partial_\mu g)g^\dagger\bigr)=0$. Hence $A_\mu'$ is again Hermitian and traceless, $A_\mu'=A_\mu'^aT^a$ with real $A_\mu'^a$.

Since $D_\mu'(gq_f)=g\,D_\mu q_f$ holds for all $q_f$, as operators on color triplets
\begin{equation*}
D_\mu'=g\,D_\mu\,g^\dagger.
\tag{5.80}\label{eq:5.80}
\end{equation*}
For infinitesimal $\theta^a$, $g=I_3+i\theta^aT^a+O(\theta^2)$ and \eqref{eq:5.79} gives
\begin{equation*}
\begin{aligned}
A_\mu'&=A_\mu+i\theta^b[T^b,A_\mu]+\frac1{g_s}(\partial_\mu\theta^a)T^a+O(\theta^2)\\
&=A_\mu-f^{bca}\theta^bA^c_\mu T^a+\frac1{g_s}(\partial_\mu\theta^a)T^a+O(\theta^2),\\[6pt]
\delta A^a_\mu:=A_\mu'^a-A_\mu^a
&=\frac1{g_s}\partial_\mu\theta^a+f^{abc}A^b_\mu\theta^c+O(\theta^2).
\end{aligned}
\tag{5.81}\label{eq:5.81}
\end{equation*}

\topic{Field Strength}

Define $F_{\mu\nu}$ by $[D_\mu,D_\nu]=:-ig_sF_{\mu\nu}$. Since $[\partial_\mu,A_\nu]q_f=(\partial_\mu A_\nu)q_f$,
\begin{equation*}
\begin{aligned}
{}[D_\mu,D_\nu]
&=[\partial_\mu-ig_sA_\mu,\ \partial_\nu-ig_sA_\nu]\\
&=-ig_s\bigl(\partial_\mu A_\nu-\partial_\nu A_\mu\bigr)-g_s^2[A_\mu,A_\nu],\\[6pt]
F_{\mu\nu}&=\partial_\mu A_\nu-\partial_\nu A_\mu-ig_s[A_\mu,A_\nu].
\end{aligned}
\tag{5.82}\label{eq:5.82}
\end{equation*}
Expanding in the basis, $F_{\mu\nu}=F^a_{\mu\nu}T^a$ with
\begin{equation*}
{\displaystyle 
F^a_{\mu\nu}=\partial_\mu A^a_\nu-\partial_\nu A^a_\mu+g_sf^{abc}A^b_\mu A^c_\nu.
}
\tag{5.83}\label{eq:5.83}
\end{equation*}
By \eqref{eq:5.80}, $[D_\mu',D_\nu']=g[D_\mu,D_\nu]g^\dagger$, so
\begin{equation*}
F_{\mu\nu}'=gF_{\mu\nu}g^\dagger,
\qquad
\delta F^a_{\mu\nu}=f^{abc}F^b_{\mu\nu}\theta^c+O(\theta^2).
\tag{5.84}\label{eq:5.84}
\end{equation*}
By cyclicity of the trace and \eqref{eq:5.13},
\begin{equation*}
\operatorname{tr}\bigl(F_{\mu\nu}'F'^{\mu\nu}\bigr)
=\operatorname{tr}\bigl(F_{\mu\nu}F^{\mu\nu}\bigr)
=\tfrac12F^a_{\mu\nu}F^{a\mu\nu}.
\tag{5.85}\label{eq:5.85}
\end{equation*}

\topic{The QCD Lagrangian}
\begin{equation*}
{\displaystyle 
\mathcal L_{\mathrm{QCD}}
=\sum_{f}\bar q_f\bigl(i\gamma^\mu D_\mu-m_f\bigr)q_f
-\frac12\operatorname{tr}\bigl(F_{\mu\nu}F^{\mu\nu}\bigr)
=\sum_{f}\bar q_f\bigl(i\gamma^\mu D_\mu-m_f\bigr)q_f
-\frac14F^a_{\mu\nu}F^{a\mu\nu},
}
\tag{5.86}\label{eq:5.86}
\end{equation*}
where the sum runs over the six flavors in \eqref{eq:5.73}, $m_f$ is the mass of flavor $f$, and $D_\mu$ acts on the $q_f$ to its right. Under \eqref{eq:5.75} and \eqref{eq:5.79}, $D_\mu'q_f'=g\,D_\mu q_f$ by construction, so
\begin{equation*}
\bar q_f'\bigl(i\gamma^\mu D_\mu'-m_f\bigr)q_f'
=\bar q_f\,g^\dagger\bigl(i\gamma^\mu g\,D_\mu-m_f\,g\bigr)q_f
=\bar q_f\bigl(i\gamma^\mu D_\mu-m_f\bigr)q_f.
\tag{5.87}\label{eq:5.87}
\end{equation*}
Together with \eqref{eq:5.85}, $\mathcal L_{\mathrm{QCD}}$ is invariant under the local transformations \eqref{eq:5.75}, \eqref{eq:5.79}.

Substituting \eqref{eq:5.77} and \eqref{eq:5.83} into \eqref{eq:5.86} separates the free and interaction terms:
\begin{equation*}
\begin{aligned}
\mathcal L_{\mathrm{QCD}}
={}&\sum_f\bar q_f\bigl(i\gamma^\mu\partial_\mu-m_f\bigr)q_f
-\frac14\bigl(\partial_\mu A^a_\nu-\partial_\nu A^a_\mu\bigr)\bigl(\partial^\mu A^{a\nu}-\partial^\nu A^{a\mu}\bigr)\\[6pt]
&+g_sA^a_\mu\sum_f\bar q_f\gamma^\mu T^aq_f
-\frac{g_s}{2}f^{abc}\bigl(\partial_\mu A^a_\nu-\partial_\nu A^a_\mu\bigr)A^{b\mu}A^{c\nu}
-\frac{g_s^2}{4}f^{abc}f^{ade}A^b_\mu A^c_\nu A^{d\mu}A^{e\nu}.
\end{aligned}
\tag{5.88}\label{eq:5.88}
\end{equation*}

\topic{Equations of Motion}

Treating $q_f$, $\bar q_f$ and $A^a_\mu$ as independent, the Euler–Lagrange equation for $\bar q_f$ is
\begin{equation*}
\bigl(i\gamma^\mu D_\mu-m_f\bigr)q_f=0.
\tag{5.89}\label{eq:5.89}
\end{equation*}
For $A^a_\nu$, define the color current $j^{a\nu}:=\sum_f\bar q_f\gamma^\nu T^aq_f$. From \eqref{eq:5.86} with \eqref{eq:5.83} and \eqref{eq:5.88},
\begin{equation*}
\begin{aligned}
\frac{\partial\mathcal L_{\mathrm{QCD}}}{\partial(\partial_\mu A^a_\nu)}
&=-\frac12F^{b\rho\sigma}\,\delta^{ab}\bigl(\delta^\mu_\rho\delta^\nu_\sigma-\delta^\mu_\sigma\delta^\nu_\rho\bigr)
=-F^{a\mu\nu},\\[6pt]
\frac{\partial\mathcal L_{\mathrm{QCD}}}{\partial A^a_\nu}
&=g_sj^{a\nu}-\frac12F^{b\rho\sigma}\,g_s\bigl(f^{bad}\delta^\nu_\rho A^d_\sigma+f^{bca}A^c_\rho\delta^\nu_\sigma\bigr)\\
&=g_sj^{a\nu}-g_sf^{abc}A^c_\mu F^{b\mu\nu}\\
&=g_sj^{a\nu}+g_sf^{abc}A^b_\mu F^{c\mu\nu}.
\end{aligned}
\tag{5.90}\label{eq:5.90}
\end{equation*}
The Euler–Lagrange equation $\partial_\mu\bigl(\partial\mathcal L_{\mathrm{QCD}}/\partial(\partial_\mu A^a_\nu)\bigr)=\partial\mathcal L_{\mathrm{QCD}}/\partial A^a_\nu$ then reads
\begin{equation*}
{\displaystyle 
\partial_\mu F^{a\mu\nu}+g_sf^{abc}A^b_\mu F^{c\mu\nu}=-g_sj^{a\nu}.
}
\tag{5.91}\label{eq:5.91}
\end{equation*}
The left-hand side is the $a$-component of $\partial_\mu F^{\mu\nu}-ig_s[A_\mu,F^{\mu\nu}]$, the covariant derivative of $F^{\mu\nu}$ in the adjoint representation; with $j^\nu:=j^{a\nu}T^a$, \eqref{eq:5.91} is
\begin{equation*}
D_\mu F^{\mu\nu}:=\partial_\mu F^{\mu\nu}-ig_s[A_\mu,F^{\mu\nu}]=-g_s\,j^\nu.
\tag{5.92}\label{eq:5.92}
\end{equation*}

\setcounter{chapter}{5}
\chapter{New Physics Beyond the Standard Model}\label{book-chapter-6}\label{qft-ch6}

\setcounter{section}{0}
\section{Seesaw Mechanism}\label{book-section-6.1}\label{seesaw-mechanism}

\subsection{Type-I Seesaw}\label{book-subsection-section-6.1-1}\label{type-i-seesaw}

\topic{Fields and Lagrangian}

Extend the Standard Model by $n$ right-handed neutrino fields.
\begin{itemize}


\item \textbf{\boldmath Right-handed neutrinos}: $N_{IR}$, $I=1,\ldots,n$, are four-component chiral spinors with $P_RN_{IR}=N_{IR}$ and are singlets of $SU(3)_c\times SU(2)_L\times U(1)_Y$.

\item \textbf{\boldmath Flavor multiplets}: $N_R:=(N_{1R},\ldots,N_{nR})^T$ and $\nu_L:=(\nu_{eL},\nu_{\mu L},\nu_{\tau L})^T$. The lepton doublets $\ell_{\alpha L}$, charged singlets $E_{\alpha R}$ and Higgs doublet $H$ are those of Section 5.3.
\end{itemize}


The conjugate Higgs doublet and the additional terms in the Lagrangian are
\begin{equation*}
H^c:=i\sigma^2H^*
=\begin{pmatrix}H^{0*}\\-H^-\end{pmatrix}
\sim\Bigl(\mathbf2,-\frac12\Bigr),\qquad H^-:=(H^+)^*,
\tag{6.1}\label{eq:6.1}
\end{equation*}
\begin{equation*}
{\displaystyle 
\begin{aligned}
\mathcal L_{\mathrm I}
={}&\mathcal L_{\mathrm{SM}}
+\overline{N_R}\,i\gamma^\mu\partial_\mu N_R\\
&-\left(\overline{\ell_{\alpha L}}\,(Y_\nu)_{\alpha I}H^cN_{IR}
+\frac12\overline{N_R^c}\,M_RN_R+\mathrm{h.c.}\right),\\
Y_\nu&\in\mathbb C^{3\times n},\qquad M_R\in\mathbb C^{n\times n},\qquad M_R^T=M_R.
\end{aligned}
}
\tag{6.2}\label{eq:6.2}
\end{equation*}
Repeated flavor indices are summed. The Yukawa matrix $Y_\nu$ is dimensionless and $M_R$ has mass dimension one. Under the electroweak transformation of Section 5.3,
\begin{equation*}
\begin{aligned}
\ell_L'&=e^{-i\theta_Y/2}U(\theta)\ell_L,
& (H^c)'&=e^{-i\theta_Y/2}U(\theta)H^c,
& N_R'&=N_R,\\
\overline{\ell_L'}\,Y_\nu(H^c)'N_R'
&=\overline{\ell_L}\,Y_\nu H^cN_R.
\end{aligned}
\tag{6.3}\label{eq:6.3}
\end{equation*}
The Majorana term is gauge invariant. Its flavor matrix can be taken symmetric because $\overline{N_{IR}^c}N_{JR}=\overline{N_{JR}^c}N_{IR}$.

\topic{Mass Terms after Symmetry Breaking}

In the unitary gauge,
\begin{equation*}
H=\frac1{\sqrt2}\begin{pmatrix}0\\v_{\mathrm{SM}}+h\end{pmatrix},\qquad
H^c=\frac1{\sqrt2}\begin{pmatrix}v_{\mathrm{SM}}+h\\0\end{pmatrix},\qquad
M_l=\frac{v_{\mathrm{SM}}}{\sqrt2}Y_l,\qquad M_D:=\frac{v_{\mathrm{SM}}}{\sqrt2}Y_\nu.
\tag{6.4}\label{eq:6.4}
\end{equation*}
Here $v_{\mathrm{SM}}\simeq246\,\mathrm{GeV}$, as in Section 5.3. With $l_L:=(l_{eL},l_{\mu L},l_{\tau L})^T$,
\begin{equation*}
\begin{aligned}
-\mathcal L_{\mathrm{mass}+h}
={}&\overline{l_L}M_lE_R+\overline{\nu_L}M_DN_R
+\frac12\overline{N_R^c}M_RN_R\\
&+\frac h{v_{\mathrm{SM}}}\left(\overline{l_L}M_lE_R+\overline{\nu_L}M_DN_R\right)
+\mathrm{h.c.}
\end{aligned}
\tag{6.5}\label{eq:6.5}
\end{equation*}
Using the charge-conjugation identity
\begin{equation*}
\overline{\nu_L}M_DN_R=\overline{N_R^c}M_D^T\nu_L^c,
\tag{6.6}\label{eq:6.6}
\end{equation*}
define
\begin{equation*}
\Psi_L:=\begin{pmatrix}\nu_L\\N_R^c\end{pmatrix},\qquad
\Psi_L^c=\begin{pmatrix}\nu_L^c\\N_R\end{pmatrix},\qquad
{\displaystyle \mathcal M:=\begin{pmatrix}0_{3\times3}&M_D\\M_D^T&M_R\end{pmatrix}=\mathcal M^T}.
\tag{6.7}\label{eq:6.7}
\end{equation*}
The mass terms become
\begin{equation*}
-\mathcal L_{\mathrm{mass}}
=\overline{l_L}M_lE_R+\frac12\overline{\Psi_L}\,\mathcal M\Psi_L^c+\mathrm{h.c.}
\tag{6.8}\label{eq:6.8}
\end{equation*}

\topic{Flavor Basis and Takagi Factorization}

For constant unitary matrices $R_\ell,R_E\in U(3)$ and $R_N\in U(n)$, write the old fields in terms of new fields as
\begin{equation*}
\begin{aligned}
\ell_L&=R_\ell\ell_L',\qquad E_R=R_EE_R',\qquad N_R=R_NN_R',\\
M_l'&=R_\ell^\dagger M_lR_E,\qquad
M_D'=R_\ell^\dagger M_DR_N,\qquad
M_R'=R_N^TM_RR_N.
\end{aligned}
\tag{6.9}\label{eq:6.9}
\end{equation*}
These transformations preserve the kinetic and gauge terms. We use the charged-lepton mass basis of Section 5.3,
\begin{equation*}
M_l=D_l:=\operatorname{diag}(m_e,m_\mu,m_\tau).
\tag{6.10}\label{eq:6.10}
\end{equation*}
Since $\mathcal M$ is complex symmetric, there exists $\mathcal U\in U(3+n)$ such that
\begin{equation*}
{\displaystyle 
\begin{aligned}
\mathcal M&=\mathcal U D\mathcal U^T,\qquad
\mathcal U^\dagger\mathcal M\mathcal U^*=D,\\
D&:=\operatorname{diag}(D_\nu,D_N),\qquad
D_\nu:=\operatorname{diag}(m_1,m_2,m_3),\qquad
D_N:=\operatorname{diag}(M_1,\ldots,M_n),
\end{aligned}
}
\tag{6.11}\label{eq:6.11}
\end{equation*}
where all masses are nonnegative. We label the three light states by $i=1,2,3$ and the $n$ heavy states by $I=1,\ldots,n$; the light labels allow either normal or inverted ordering. The field transformation and the diagonal mass terms are
\begin{equation*}
\begin{aligned}
\Psi_L&=\mathcal U\Psi_L^m,\qquad
\Psi_L^m:=\begin{pmatrix}\nu_L^m\\(N_R^m)^c\end{pmatrix},\qquad
\Psi_L^c=\mathcal U^*(\Psi_L^m)^c,\\
-\mathcal L_{\mathrm{mass}}
&=\overline{l_L}D_lE_R
+\frac12\overline{\nu_L^m}D_\nu(\nu_L^m)^c
+\frac12\overline{(N_R^m)^c}D_NN_R^m+\mathrm{h.c.}
\end{aligned}
\tag{6.12}\label{eq:6.12}
\end{equation*}
Equivalently, one may first block diagonalize and then Takagi diagonalize each block:
\begin{equation*}
\begin{aligned}
\Psi_L&=W\begin{pmatrix}\nu_L^{\mathrm b}\\(N_R^{\mathrm b})^c\end{pmatrix},\qquad
W^\dagger\mathcal MW^*=\begin{pmatrix}M_\nu&0\\0&M_N\end{pmatrix},\\
M_\nu&=UD_\nu U^T,\qquad M_N=U_ND_NU_N^T,\\
\mathcal U&=W\begin{pmatrix}U&0\\0&U_N\end{pmatrix},\qquad
W\in U(3+n),\quad U\in U(3),\quad U_N\in U(n).
\end{aligned}
\tag{6.13}\label{eq:6.13}
\end{equation*}
The matrices $M_\nu$ and $M_N$ refer to the block-diagonal basis; $M_N$ is generally different from the original block $M_R$.

\topic{Rank and Parameter Counting}

The zero active-active block implies
\begin{equation*}
\mathcal M\begin{pmatrix}a\\0\end{pmatrix}=0
\quad\text{for }a\in\ker M_D^T,\qquad
\operatorname{rank}\mathcal M\le\min(2n,3+n).
\tag{6.14}\label{eq:6.14}
\end{equation*}
If $M_R$ is invertible, block elimination gives the exact rank identity
\begin{equation*}
\begin{aligned}
&\begin{pmatrix}I_3&-M_DM_R^{-1}\\0&I_n\end{pmatrix}
\mathcal M
\begin{pmatrix}I_3&0\\-M_R^{-1}M_D^T&I_n\end{pmatrix}
=\begin{pmatrix}-M_DM_R^{-1}M_D^T&0\\0&M_R\end{pmatrix},\\
&\operatorname{rank}\mathcal M
=n+\operatorname{rank}(-M_DM_R^{-1}M_D^T).
\end{aligned}
\tag{6.15}\label{eq:6.15}
\end{equation*}
This congruence is not a unitary mass diagonalization. For generic matrices, the light rank is $r=\min(n,3)$ and there are $3-r$ massless light states. In particular, $n=2$ gives the minimal seesaw with one massless light neutrino at tree level:
\begin{equation*}
m_1=0\quad\text{(normal ordering)},\qquad
m_3=0\quad\text{(inverted ordering)}.
\tag{6.16}\label{eq:6.16}
\end{equation*}
For the following counts, take $n\ge1$, generic rank and nondegenerate nonzero masses. Before quotienting by field redefinitions, the numbers of independent real parameters are
\begin{equation*}
\begin{aligned}
N(M_D)&=6n,\qquad N(M_R)=n(n+1),\qquad N(\mathcal M)=n^2+7n,\\
N(D)&=n+r=\min(2n,3+n),\\
N(\mathcal U)_{\mathrm{constrained}}
&=n^2+7n-(n+r).
\end{aligned}
\tag{6.17}\label{eq:6.17}
\end{equation*}
The last line counts independent non-mass parameters of $\mathcal M$ at fixed $D$, with the zero-block constraint imposed and Takagi stabilizers removed. It is not the dimension of an unconstrained $U(3+n)$ matrix.

For nondegenerate charged-lepton masses, the remaining basis freedom is
\begin{equation*}
\begin{aligned}
\ell_L&=P\ell_L',\qquad E_R=PE_R',\qquad N_R=R_NN_R',\\
P&:=\operatorname{diag}(e^{i\phi_e},e^{i\phi_\mu},e^{i\phi_\tau}),
\qquad R_N\in U(n),\\
M_D'&=P^\dagger M_DR_N,\qquad M_R'=R_N^TM_RR_N.
\end{aligned}
\tag{6.18}\label{eq:6.18}
\end{equation*}
There are $3+n^2$ redundant real parameters. Excluding the three charged-lepton masses,
\begin{equation*}
{\displaystyle 
\begin{aligned}
N_{\nu,\mathrm{phys}}&=(n^2+7n)-(3+n^2)=7n-3,\\
N_{\mathrm{masses}}&=n+r,\qquad
N_{\mathrm{mixing}}=6n-r-3.
\end{aligned}
}
\tag{6.19}\label{eq:6.19}
\end{equation*}
Thus $n=2$ has four nonzero masses and seven mixing parameters, while $n=3$ has six masses and twelve mixing parameters.

\topic{Charged Currents and the PMNS Matrix}

Write the upper blocks of $\mathcal U$ as
\begin{equation*}
K:=W_{11}U,\qquad R:=W_{12}U_N,\qquad
\nu_{\alpha L}=K_{\alpha i}\nu_{iL}^m+R_{\alpha I}(N_{IR}^m)^c.
\tag{6.20}\label{eq:6.20}
\end{equation*}
In the charged-lepton mass basis, the leptonic charged current of Section 5.3 becomes
\begin{equation*}
{\displaystyle 
\mathcal L_{\mathrm{CC}}^\ell
=\frac{g_2}{\sqrt2}W_\mu^-
\overline{l_{\alpha L}}\gamma^\mu
\left[K_{\alpha i}\nu_{iL}^m+R_{\alpha I}(N_{IR}^m)^c\right]
+\mathrm{h.c.}
}
\tag{6.21}\label{eq:6.21}
\end{equation*}
Following the coefficient definition, set $U_{\mathrm{PMNS}}:=K$. Unitarity of $\mathcal U$ and its zero-block relation give
\begin{equation*}
{\displaystyle 
\begin{aligned}
U_{\mathrm{PMNS}}U_{\mathrm{PMNS}}^\dagger
&=I_3-RR^\dagger,\\
KD_\nu K^T+RD_NR^T&=0_{3\times3}.
\end{aligned}
}
\tag{6.22}\label{eq:6.22}
\end{equation*}
The light charged-current matrix $U_{\mathrm{PMNS}}$ is generally nonunitary; the Takagi matrix $U$ is unitary.

\topic{Exact Parametrization of Light-Heavy Mixing}

Choose the singlet flavor basis so that the heavy block after block diagonalization is $M_N=D_N$, hence $U_N=I_n$. In the coordinate patch where $W_{11}$ is invertible, use the block-unitary freedom to write
\begin{equation*}
\begin{aligned}
X&\in\mathbb C^{3\times n},\qquad
A_X:=(I_3+XX^\dagger)^{-1/2},\qquad
B_X:=(I_n+X^\dagger X)^{-1/2},\\
A_XX&=XB_X,\qquad
{\displaystyle W(X):=\begin{pmatrix}A_X&A_XX\\-X^\dagger A_X&B_X\end{pmatrix}},
\qquad W^\dagger W=WW^\dagger=I_{3+n}.
\end{aligned}
\tag{6.23}\label{eq:6.23}
\end{equation*}
The square roots are the positive Hermitian roots. The upper-left block of $\mathcal M=W\operatorname{diag}(M_\nu,D_N)W^T$ is
\begin{equation*}
0=A_X(M_\nu+XD_NX^T)A_X^T
\quad\Longleftrightarrow\quad
{\displaystyle M_\nu=-XD_NX^T=UD_\nu U^T}.
\tag{6.24}\label{eq:6.24}
\end{equation*}
The other blocks give
\begin{equation*}
{\displaystyle 
\begin{aligned}
M_D&=A_XXD_N(B_X^T)^{-1},\\
M_R&=B_X\left[D_N-(X^\dagger X)D_N(X^\dagger X)^T\right]B_X^T.
\end{aligned}
}
\tag{6.25}\label{eq:6.25}
\end{equation*}
For positive heavy masses, elimination of $X$ from the light block also yields
\begin{equation*}
\begin{aligned}
M_\nu
&=-A_X^{-1}M_DB_X^TD_N^{-1}B_XM_D^T(A_X^T)^{-1}\\
&=-A_X^{-1}M_D
\left(M_R+X^\dagger M_D+M_D^TX^*\right)^{-1}
M_D^T(A_X^T)^{-1}.
\end{aligned}
\tag{6.26}\label{eq:6.26}
\end{equation*}
Here $D_N$ is the physical heavy spectrum. Neither $M_R=D_N$ nor $X=M_DM_R^{-1}$ is an exact identity in this basis.

\topic{Construction from Masses and Mixing Parameters}

Take generic light rank $r=\min(n,3)$ and positive $M_I$. Let $P_\nu$ project onto the nonzero entries of $D_\nu$, and choose a complex $3\times n$ matrix $O$ with
\begin{equation*}
P_\nu:=\operatorname{diag}(p_1,p_2,p_3),\qquad
p_i:=\begin{cases}1,&m_i>0,\\0,&m_i=0,\end{cases}
\qquad P_\nu O=O,\qquad OO^T=P_\nu.
\tag{6.27}\label{eq:6.27}
\end{equation*}
Then
\begin{equation*}
{\displaystyle X=iU\sqrt{D_\nu}\,OD_N^{-1/2}},\qquad
-XD_NX^T=U\sqrt{D_\nu}\,OO^T\sqrt{D_\nu}\,U^T=UD_\nu U^T.
\tag{6.28}\label{eq:6.28}
\end{equation*}
The orthogonality condition uses a transpose, not a Hermitian conjugate. For $n>3$ it does not require $O^TO=I_n$. The exact reconstruction is
\begin{equation*}
(D_\nu,D_N,U,O)\ \longrightarrow\ X\ \longrightarrow\ W(X)
\ \longrightarrow\
\mathcal M=W\begin{pmatrix}UD_\nu U^T&0\\0&D_N\end{pmatrix}W^T.
\tag{6.29}\label{eq:6.29}
\end{equation*}
After removing the three charged-lepton phases and the $U(3-r)$ freedom in the massless subspace, the parameter counts are
\begin{equation*}
\begin{aligned}
N(U)&=9-(3-r)^2-3,\qquad
N(O)=2rn-r(r+1),\\
N(U)+N(O)&=6n-r-3,\qquad
(n+r)+N(U)+N(O)=7n-3.
\end{aligned}
\tag{6.30}\label{eq:6.30}
\end{equation*}
For $n=2$ and normal ordering,
\begin{equation*}
\begin{aligned}
D_\nu&=\operatorname{diag}(0,m_2,m_3),\qquad
D_N=\operatorname{diag}(M_1,M_2),\\
O(z)&=\begin{pmatrix}0&0\\\cos z&-\sin z\\\sin z&\cos z\end{pmatrix},\qquad
z\in\mathbb C,\qquad O^TO=I_2,\qquad OO^T=\operatorname{diag}(0,1,1).
\end{aligned}
\tag{6.31}\label{eq:6.31}
\end{equation*}
For inverted ordering one may use
\begin{equation*}
D_\nu=\operatorname{diag}(m_1,m_2,0),\qquad
O(z)=\begin{pmatrix}\cos z&-\sin z\\\sin z&\cos z\\0&0\end{pmatrix},\qquad
OO^T=\operatorname{diag}(1,1,0).
\tag{6.32}\label{eq:6.32}
\end{equation*}
Discrete column signs can be absorbed into the heavy Majorana fields. For three massive light states, a phase convention for the unitary matrix $U$ is
\begin{equation*}
{\displaystyle 
U=\begin{pmatrix}
c_{12}c_{13}&s_{12}c_{13}&s_{13}e^{-i\delta}\\
-s_{12}c_{23}-c_{12}s_{23}s_{13}e^{i\delta}
&c_{12}c_{23}-s_{12}s_{23}s_{13}e^{i\delta}&s_{23}c_{13}\\
s_{12}s_{23}-c_{12}c_{23}s_{13}e^{i\delta}
&-c_{12}s_{23}-s_{12}c_{23}s_{13}e^{i\delta}&c_{23}c_{13}
\end{pmatrix}
\operatorname{diag}(e^{i\rho},e^{i\sigma},1)
}.
\tag{6.33}\label{eq:6.33}
\end{equation*}
Here $c_{ij}:=\cos\theta_{ij}$ and $s_{ij}:=\sin\theta_{ij}$.

For $n=2$, one Majorana phase is removable by rephasing the massless state. One can take $\rho=0$ in either ordering, after a common charged-lepton rephasing when needed. Thus $U$ contains three angles, one Dirac phase and one Majorana phase; together with the complex $z$ there are seven mixing parameters. The exact charged-current matrix is
\begin{equation*}
{\displaystyle 
\begin{aligned}
U_{\mathrm{PMNS}}&=A_XU\\
&=U\left(I_3+\sqrt{D_\nu}\,OD_N^{-1}O^\dagger\sqrt{D_\nu}\right)^{-1/2}.
\end{aligned}
}
\tag{6.34}\label{eq:6.34}
\end{equation*}

\topic{Seesaw Limit and Casas–Ibarra Form}

The conventional seesaw regime has invertible $M_R$ and
\begin{equation*}
\frac{\sigma_{\max}(M_D)}{\sigma_{\min}(M_R)}\ll1,\qquad
\epsilon_X:=\sigma_{\max}(X)\ll1,
\tag{6.35}\label{eq:6.35}
\end{equation*}
where $\sigma_{\max}$ and $\sigma_{\min}$ denote the largest and smallest singular values. Expanding the exact parametrization,
\begin{equation*}
\begin{aligned}
A_X&=I_3-\frac12XX^\dagger+\mathcal O(\epsilon_X^4),\qquad
B_X=I_n-\frac12X^\dagger X+\mathcal O(\epsilon_X^4),\\
W&=\begin{pmatrix}I_3-\frac12XX^\dagger&X\\-X^\dagger&I_n-\frac12X^\dagger X\end{pmatrix}
+\mathcal O(\epsilon_X^3).
\end{aligned}
\tag{6.36}\label{eq:6.36}
\end{equation*}
To leading order in this basis,
\begin{equation*}
M_D\simeq XD_N,\qquad M_R\simeq D_N,\qquad
X\simeq M_DM_R^{-1}\simeq M_DD_N^{-1}.
\tag{6.37}\label{eq:6.37}
\end{equation*}
The zero-block identity therefore gives
\begin{equation*}
{\displaystyle M_\nu=UD_\nu U^T\simeq-M_DM_R^{-1}M_D^T
\simeq-\frac{v_{\mathrm{SM}}^2}{2}\,Y_\nu D_N^{-1}Y_\nu^T}.
\tag{6.38}\label{eq:6.38}
\end{equation*}
The charged-current matrix satisfies
\begin{equation*}
\begin{aligned}
U_{\mathrm{PMNS}}&=\left(I_3-\frac12XX^\dagger\right)U+\mathcal O(\epsilon_X^4),\\
U_{\mathrm{PMNS}}U_{\mathrm{PMNS}}^\dagger
&=I_3-XX^\dagger+\mathcal O(\epsilon_X^4).
\end{aligned}
\tag{6.39}\label{eq:6.39}
\end{equation*}
Thus $U_{\mathrm{PMNS}}\simeq U$ only after neglecting second-order light-heavy mixing. Substituting the exact expression for $X$ into $M_D\simeq XD_N$ gives the Casas–Ibarra form,
\begin{equation*}
{\displaystyle 
\begin{aligned}
M_D&\simeq iU\sqrt{D_\nu}\,O\sqrt{D_N},\qquad M_R\simeq D_N,\\
Y_\nu&\simeq\frac{i\sqrt2}{v_{\mathrm{SM}}}\,U\sqrt{D_\nu}\,O\sqrt{D_N},\qquad
P_\nu O=O,\quad OO^T=P_\nu.
\end{aligned}
}
\tag{6.40}\label{eq:6.40}
\end{equation*}
The construction of $X$ is exact; these expressions for $M_D$ and $Y_\nu$ are leading-order approximations. In particular,
\begin{equation*}
\epsilon_X=\sigma_{\max}\!\left(\sqrt{D_\nu}\,OD_N^{-1/2}\right).
\tag{6.41}\label{eq:6.41}
\end{equation*}
Since $O$ is complex orthogonal rather than unitary, $m_i\ll M_I$ alone does not ensure small light-heavy mixing.

\subsection{Type-II Seesaw}\label{book-subsection-section-6.1-2}\label{type-ii-seesaw}

\subsection{Type-III Seesaw}\label{book-subsection-section-6.1-3}\label{type-iii-seesaw}

\subsection{Leptogenesis}\label{book-subsection-section-6.1-4}\label{leptogenesis}

\setcounter{section}{1}
\section{\texorpdfstring{$U(1)_{L_\mu-L_\tau}$}{U(1)Lμ-Lτ} Gauge Symmetry}\label{book-section-6.2}\label{lmu-ltau-gauge-symmetry}

The lepton doublets $\ell_{\alpha L}$, charged singlets $E_{\alpha R}$, Higgs doublet $H$, gauge fields $W_\mu^a,B_\mu$ with couplings $g_2,g_1$, and Yukawa matrix $Y_l$ are those of Section 5.3.

The terms of the Standard Model Lagrangian density that contain lepton fields are
\begin{equation*}
{\displaystyle 
\begin{aligned}
\mathcal L_{\mathrm{SM}}^\ell
&=\sum_{\alpha=e,\mu,\tau}\Bigl(\overline{\ell_{\alpha L}}\,i\gamma^\mu D_\mu\ell_{\alpha L}
+\overline{E_{\alpha R}}\,i\gamma^\mu D_\mu E_{\alpha R}\Bigr)\\
&\quad-\sum_{\alpha,\beta=e,\mu,\tau}\Bigl(\overline{\ell_{\alpha L}}\,(Y_l)_{\alpha\beta}\,H\,E_{\beta R}+\mathrm{h.c.}\Bigr),\\
D_\mu\ell_{\alpha L}&=\Bigl(\partial_\mu-ig_2W_\mu^a\frac{\sigma^a}2+\frac i2g_1B_\mu\Bigr)\ell_{\alpha L},\\
D_\mu E_{\alpha R}&=\bigl(\partial_\mu+ig_1B_\mu\bigr)E_{\alpha R}.
\end{aligned}
}
\tag{6.42}\label{eq:6.42}
\end{equation*}
In the flavor basis of Section 5.3,
\begin{equation*}
Y_l=\frac{\sqrt2}{v_{\mathrm{SM}}}\operatorname{diag}(m_e,m_\mu,m_\tau).
\tag{6.43}\label{eq:6.43}
\end{equation*}

\subsection{\texorpdfstring{$U(1)_{L_\mu-L_\tau}$}{U(1)Lμ-Lτ} Gauge Symmetry Transformations}\label{book-subsection-section-6.2-1}

Let $\theta:\mathbb R^4\to\mathbb R$ and introduce a real gauge field $Z'_\mu$ with coupling $g_{Z'}>0$. The lepton fields and $Z'_\mu$ transform as
\begin{equation*}
\begin{gathered}
\begin{aligned}
\ell_{\alpha L}'(x)&=e^{iQ_\alpha\theta(x)}\ell_{\alpha L}(x),\\
E_{\alpha R}'(x)&=e^{iQ_\alpha\theta(x)}E_{\alpha R}(x),\\
(Z'_\mu)'(x)&=Z'_\mu(x)+\frac1{g_{Z'}}\partial_\mu\theta(x),
\end{aligned}\\[6pt]
\text{where the }U(1)_{L_\mu-L_\tau}\text{ charges are}\quad
Q_\alpha=\begin{cases}
0, & \alpha=e,\\
1, & \alpha=\mu,\\
-1, & \alpha=\tau.
\end{cases}
\end{gathered}
\tag{6.44}\label{eq:6.44}
\end{equation*}
The remaining Standard Model fields are unchanged.

The covariant derivatives in \eqref{eq:6.42} become
\begin{equation*}
\begin{aligned}
D_\mu\ell_{\alpha L}&=\Bigl(\partial_\mu-ig_2W_\mu^a\frac{\sigma^a}2+\frac i2g_1B_\mu-ig_{Z'}Q_\alpha Z'_\mu\Bigr)\ell_{\alpha L},\\
D_\mu E_{\alpha R}&=\bigl(\partial_\mu+ig_1B_\mu-ig_{Z'}Q_\alpha Z'_\mu\bigr)E_{\alpha R}.
\end{aligned}
\tag{6.45}\label{eq:6.45}
\end{equation*}

\subsection{Lagrangian}\label{book-subsection-section-6.2-2}

\setcounter{chapter}{6}
\chapter{Renormalization}\label{book-chapter-7}\label{qft-ch7}

\setcounter{section}{0}
\section{Forest Formula}\label{book-section-7.1}\label{forest-formula}

\subsection{Basic Definitions of Graphs}\label{book-subsection-section-7.1-1}\label{graph-definitions}
\begin{itemize}


\item \textbf{\boldmath Graph}: a Feynman graph $\mathcal G$ consists of a finite set $V(\mathcal G)$ of vertices, a finite set $I(\mathcal G)$ of internal lines, each joining two vertices (possibly the same vertex), and a finite set $E(\mathcal G)$ of external lines, each attached to one vertex.
\begin{itemize}


\item Several internal lines may join the same pair of vertices.

\item $\mathcal G$ is connected if any two vertices are joined by a chain of internal lines. A loop is a closed chain of internal lines.

\item For connected $\mathcal G$, the number of loops is $L(\mathcal G)=|I(\mathcal G)|-|V(\mathcal G)|+1$, where $|\cdot|$ is the number of elements; for a disconnected graph, $L$ is the sum over its connected components.
\end{itemize}


\item \textbf{\boldmath 1PI}: $\mathcal G$ is one-particle irreducible (1PI) if it is connected and stays connected when any one internal line is cut.
\begin{itemize}


\item Equivalently, $\mathcal G$ is connected and every internal line lies on a loop.
\end{itemize}


\item \textbf{\boldmath Subgraph}: a nonempty set $I(\gamma)\subseteq I(\mathcal G)$ of internal lines, together with the set $V(\gamma)$ of their endpoints, defines a subgraph $\gamma$ of $\mathcal G$, written $\gamma\subseteq\mathcal G$.
\begin{itemize}


\item The external lines of $\gamma$ are the external lines of $\mathcal G$ attached to $V(\gamma)$ and the internal lines of $\mathcal G$ outside $I(\gamma)$ with an endpoint in $V(\gamma)$; such an internal line with both endpoints in $V(\gamma)$ counts as two external lines of $\gamma$.
\end{itemize}


\item \textbf{\boldmath Spanning tree}: for connected $\mathcal G$, a maximal set of internal lines containing no loop.
\begin{itemize}


\item Its lines connect all vertices of $\mathcal G$, and there are $|V(\mathcal G)|-1$ of them.

\item Each of the remaining $L(\mathcal G)$ internal lines closes exactly one loop with lines of the tree. With independent loop momenta assigned to these lines, momentum conservation at the vertices fixes the momenta of the tree lines.
\end{itemize}


\item \textbf{\boldmath Reduced graph}: for $\gamma\subseteq\mathcal G$, the reduced graph $\mathcal G/\gamma$ is obtained by contracting each connected component of $\gamma$ to a point: the lines $I(\gamma)$ are removed, the vertices of each connected component of $\gamma$ are merged into a single vertex, and all other internal and external lines of $\mathcal G$ are kept.
\begin{itemize}


\item $I(\mathcal G/\gamma)=I(\mathcal G)\setminus I(\gamma)$ and $L(\mathcal G/\gamma)=L(\mathcal G)-L(\gamma)$.

\item If $\mathcal G$ is 1PI, so is $\mathcal G/\gamma$.
\end{itemize}


\item \textbf{\boldmath Renormalization part}: a 1PI subgraph $\gamma\subseteq\mathcal G$ whose superficial degree of divergence $\omega(\gamma)$, namely $d\,L(\gamma)$ plus the degree of the integrand of $\gamma$ in its loop momenta, satisfies $\omega(\gamma)\ge0$; $\gamma=\mathcal G$ is allowed.

\item \textbf{\boldmath Nested, disjoint, overlapping}: two subgraphs $\gamma_1,\gamma_2\subseteq\mathcal G$ are nested if $\gamma_1\subseteq\gamma_2$ or $\gamma_2\subseteq\gamma_1$, disjoint if $V(\gamma_1)\cap V(\gamma_2)=\emptyset$, and overlapping otherwise.

\item \textbf{\boldmath Forest}: a set of renormalization parts of $\mathcal G$, any two of which are nested or disjoint; the empty set is a forest.
\end{itemize}


\subsection{\texorpdfstring{$\mathcal{L}_{CT}$}{LCT} and R-Operation Recursion Formula}\label{book-subsection-section-7.1-2}\label{r-operation-recursion}

Original Lagrangian : $\mathcal L$

\topic{1-Loop Level}
\begin{itemize}


\item \textbf{\boldmath Divergent part}: for a 1PI graph $\mathcal G$, $T\circ\mathcal G$ is the divergent part of $\mathcal G$, a polynomial in the external momenta of $\mathcal G$.

\item \textbf{\boldmath Counterterm Feynman rule}: the counterterm vertex associated with $\mathcal G$ has the Feynman rule $C(\mathcal G):=-T\circ\mathcal G$.

\item \textbf{\boldmath Counterterm Lagrangian}: $\mathcal L_{CT}^{(\mathcal G)}$ is the term of the counterterm Lagrangian whose Feynman rule is $C(\mathcal G)$.
\end{itemize}


\setcounter{section}{1}
\section{QED Renormalization}\label{book-section-7.2}\label{qed-renormalization}

\topic{Gauge-Fixed Lagrangian}

With the covariant gauge fixing of \eqref{eq:4.15}, the QED Lagrangian \eqref{eq:5.67} for Dirac fields $\psi_f$ of charge $Q_fe$ and mass $m_f$ is
\begin{equation*}
\mathcal L
=\sum_f\bar\psi_f\bigl(i\gamma^\mu\partial_\mu-m_f\bigr)\psi_f
-\frac14F_{\mu\nu}F^{\mu\nu}
-\frac1{2\xi}(\partial_\mu A^\mu)^2
+eA_\mu\sum_fQ_f\,\bar\psi_f\gamma^\mu\psi_f .
\tag{7.1}\label{eq:7.1}
\end{equation*}
The Faddeev--Popov ghosts decouple, as noted after \eqref{eq:3.38}.

\topic{Feynman Rules}

The momentum-space factors contributing to $i\mathcal M$ are
\begin{equation*}
\begin{array}{l|l}
\text{element}&\text{factor}\\ \hline
\text{fermion propagator, momentum }p\text{ along the arrow}&\dfrac{i(\slashed{p}+m_f)}{p^2-m_f^2+i\varepsilon}\\[12pt]
\text{photon propagator}&\dfrac{-i}{p^2+i\varepsilon}\left[\eta_{\mu\nu}-(1-\xi)\dfrac{p_\mu p_\nu}{p^2+i\varepsilon}\right]\\[12pt]
\text{fermion-photon vertex}&ieQ_f\gamma^\mu\\[4pt]
\text{incoming / outgoing fermion}&u^s(p)\ /\ \bar u^s(p)\\
\text{incoming / outgoing antifermion}&\bar v^s(p)\ /\ v^s(p)\\
\text{incoming / outgoing photon}&\epsilon^r_\mu(p)\ /\ \epsilon^{r*}_\mu(p)
\end{array}
\tag{7.2}\label{eq:7.2}
\end{equation*}
For the electron, $Q_f=-1$ and the vertex is $-ie\gamma^\mu$. Along each fermion line, write the factors in the order opposite to the arrow, so that the spinor indices form a matrix product. Impose four-momentum conservation at every vertex and integrate each loop momentum with $\int d^4\ell/(2\pi)^4$, continued to $\int d^d\ell/(2\pi)^d$ in $d=4-\epsilon$ dimensions. Each closed fermion loop gives a factor $-1$ and a Dirac trace; exchanging identical external fermions changes the sign, as in \hyperref[feynman-diagrams]{Section 4.3}.

\topic{Tree Level: $e^+e^-\to q\bar q$}

% ee-to-qqbar
\qftdiagram{282.24}{133.01}{%
q 0.1 0 0 0.1 0 0 cm 1 g 0 0 2822.4 1330.15 re f q 72 72 2678.4 1186.15 re W n 0 g 697.941 868.359 m
647.746 952.617 l 604.719 898.836 l f 632.207 409.203 m 725.43 439.68 l 682.402 493.461 l f 2190.19
920.945 m 2096.97 890.469 l 2140 836.688 l f 2124.46 461.789 m 2174.66 377.535 l 2217.68 431.313 l f 14 w
1 J 1 j 0 G 378.102 1124.23 m 952.047 665.074 l S 952.047 665.074 m 378.102 205.922 l S 1870.36 665.074 m
2444.3 1124.23 l S 2444.3 205.922 m 1870.36 665.074 l S 2 J 952.047 665.074 m 975.063 687.613 l 986.57
696.082 l 995.773 700.691 l 1004.98 703.051 l 1014.19 703.016 l 1023.39 700.582 l 1032.6 695.906 l
1044.11 687.371 l 1060.22 671.965 l 1090.14 642.293 l 1101.64 633.891 l 1110.85 629.348 l 1120.06 627.059
l 1129.26 627.176 l 1138.47 629.68 l 1147.68 634.422 l 1159.18 643.023 l 1175.29 658.48 l 1205.21 688.098
l 1216.72 696.434 l 1225.93 700.906 l 1235.13 703.121 l 1244.34 702.934 l 1253.55 700.352 l 1262.75
695.543 l 1274.26 686.875 l 1290.37 671.371 l 1320.29 641.813 l 1331.8 633.547 l 1341 629.137 l 1350.21
626.996 l 1359.41 627.262 l 1368.62 629.914 l 1377.83 634.789 l 1389.34 643.52 l 1405.45 659.074 l
1435.37 688.574 l 1446.88 696.773 l 1456.08 701.113 l 1465.29 703.18 l 1474.49 702.84 l 1483.7 700.113 l
1492.9 695.176 l 1504.41 686.379 l 1520.52 670.777 l 1548.14 643.273 l 1559.65 634.605 l 1568.86 629.797
l 1578.06 627.215 l 1587.27 627.027 l 1596.47 629.242 l 1605.68 633.719 l 1617.19 642.051 l 1633.3
657.297 l 1665.52 689.047 l 1677.03 697.105 l 1686.23 701.313 l 1695.44 703.23 l 1704.64 702.738 l
1713.85 699.867 l 1723.06 694.801 l 1734.56 685.875 l 1750.68 670.18 l 1778.29 642.781 l 1789.8 634.242 l
1799.01 629.566 l 1808.21 627.133 l 1817.42 627.098 l 1826.63 629.457 l 1835.83 634.066 l 1847.34 642.535
l 1863.45 657.887 l 1870.36 665.074 l 1870.36 665.074 l S 10 w 0 J 952.047 642.574 m 958.012 642.574
963.734 644.945 967.957 649.164 c 972.176 653.383 974.547 659.105 974.547 665.074 c 974.547 671.043
972.176 676.766 967.957 680.984 c 963.734 685.203 958.012 687.574 952.047 687.574 c 946.078 687.574
940.355 685.203 936.137 680.984 c 931.918 676.766 929.547 671.043 929.547 665.074 c 929.547 659.105
931.918 653.383 936.137 649.164 c 940.355 644.945 946.078 642.574 952.047 642.574 c h B 1870.36 642.574 m
1876.32 642.574 1882.04 644.945 1886.27 649.164 c 1890.48 653.383 1892.86 659.105 1892.86 665.074 c
1892.86 671.043 1890.48 676.766 1886.27 680.984 c 1882.04 685.203 1876.32 687.574 1870.36 687.574 c
1864.39 687.574 1858.66 685.203 1854.45 680.984 c 1850.23 676.766 1847.86 671.043 1847.86 665.074 c
1847.86 659.105 1850.23 653.383 1854.45 649.164 c 1858.66 644.945 1864.39 642.574 1870.36 642.574 c h B Q
9 w 1 J 1 j 0.449951 G 625.652 1048.69 m 691.777 995.789 755.281 944.984 816.168 896.277 c S 0.449951 g
799.301 932.82 m 816.168 896.277 l 776.813 904.711 l 799.301 932.82 l h B q 727.758 1007.91 65 79 re W n
729.121 1008.84 m 728.398 1008.84 728.039 1009.44 728.039 1010.4 c 728.281 1012.2 728.879 1013.04 729.602
1013.04 c 732 1013.04 733.68 1013.28 734.641 1013.76 c 735.602 1014.24 736.441 1015.32 737.039 1017.12 c
750.602 1071 l 751.078 1072.44 751.441 1074.36 751.441 1076.76 c 751.441 1080.12 750.238 1081.92 747.961
1081.92 c 745.441 1081.92 743.52 1080.36 742.32 1077.36 c 741.121 1074.36 739.922 1070.64 738.84 1066.08
c 738.84 1065.6 738.48 1065.36 737.879 1065.36 c 736.441 1065.36 l 736.199 1065.36 736.078 1065.48 735.84
1065.72 c 735.602 1065.96 735.48 1066.32 735.48 1066.56 c 736.32 1070.04 737.16 1073.04 738 1075.44 c
738.719 1077.84 740.039 1080 741.719 1082.04 c 743.398 1083.96 745.441 1085.04 748.078 1085.04 c 750.719
1085.04 753.121 1084.2 755.281 1082.64 c 757.32 1081.08 758.641 1079.04 759.238 1076.4 c 761.398 1078.92
763.801 1080.96 766.441 1082.64 c 769.078 1084.2 771.84 1085.04 774.719 1085.04 c 778.078 1085.04 780.961
1084.08 783.48 1082.28 c 785.879 1080.48 787.801 1078.08 789 1075.2 c 790.199 1072.32 790.922 1069.2
790.922 1065.84 c 790.922 1060.68 789.602 1055.4 786.961 1049.88 c 784.32 1044.36 780.719 1039.8 776.281
1036.2 c 771.719 1032.6 766.922 1030.8 761.762 1030.8 c 759.359 1030.8 757.32 1031.52 755.398 1032.72 c
753.48 1033.92 751.922 1035.72 750.961 1037.88 c 745.559 1016.64 l 745.32 1015.44 745.199 1014.84 745.199
1014.6 c 745.199 1013.64 747.84 1013.04 753.121 1013.04 c 753.84 1013.04 754.32 1012.56 754.32 1011.48 c
754.078 1010.52 753.84 1009.8 753.602 1009.44 c 753.359 1009.08 753 1008.84 752.281 1008.84 c 729.121
1008.84 l h 761.879 1033.8 m 764.641 1033.8 767.16 1034.88 769.559 1037.16 c 771.961 1039.44 773.879
1042.08 775.441 1044.96 c 776.398 1047.12 777.48 1049.64 778.441 1052.76 c 779.398 1055.88 780.238
1059.12 780.961 1062.72 c 781.68 1066.2 782.039 1068.96 782.039 1071.12 c 782.039 1072.92 781.801 1074.6
781.32 1076.28 c 780.84 1077.84 780 1079.16 778.922 1080.24 c 777.719 1081.32 776.281 1081.92 774.48
1081.92 c 771.48 1081.92 768.719 1080.84 766.078 1078.8 c 763.441 1076.64 761.16 1074.12 759.121 1071 c
759.121 1070.64 l 752.52 1044.12 l 753 1041.24 754.078 1038.72 755.641 1036.8 c 757.199 1034.76 759.238
1033.8 761.879 1033.8 c f Q q 799.109 1011.5 31 58 re W n 799.813 1012 m 799.813 1014.94 l 806.785
1014.94 810.313 1015.78 810.313 1017.63 c 810.313 1061.73 l 807.371 1060.3 803.676 1059.63 799.309
1059.63 c 799.309 1062.57 l 806.113 1062.57 811.32 1064.33 814.848 1067.95 c 816.023 1067.95 l 816.191
1067.95 816.445 1067.86 816.613 1067.69 c 816.781 1067.52 816.863 1067.36 816.863 1067.19 c 816.863
1017.63 l 816.863 1015.78 820.309 1014.94 827.363 1014.94 c 827.363 1012 l 799.813 1012 l f Q 625.652
281.457 m 691.777 334.359 755.281 385.164 816.168 433.871 c S 776.813 425.438 m 816.168 433.871 l 799.301
397.328 l 776.813 425.438 l h B q 727.758 232.047 65 79 re W n 729.121 232.84 m 728.398 232.84 728.039
233.441 728.039 234.398 c 728.281 236.199 728.879 237.039 729.602 237.039 c 732 237.039 733.68 237.281
734.641 237.762 c 735.602 238.238 736.441 239.32 737.039 241.121 c 750.602 295 l 751.078 296.441 751.441
298.359 751.441 300.762 c 751.441 304.121 750.238 305.922 747.961 305.922 c 745.441 305.922 743.52
304.359 742.32 301.359 c 741.121 298.359 739.922 294.641 738.84 290.078 c 738.84 289.602 738.48 289.359
737.879 289.359 c 736.441 289.359 l 736.199 289.359 736.078 289.48 735.84 289.719 c 735.602 289.961
735.48 290.32 735.48 290.559 c 736.32 294.039 737.16 297.039 738 299.441 c 738.719 301.84 740.039 304
741.719 306.039 c 743.398 307.961 745.441 309.039 748.078 309.039 c 750.719 309.039 753.121 308.199
755.281 306.641 c 757.32 305.078 758.641 303.039 759.238 300.398 c 761.398 302.922 763.801 304.961
766.441 306.641 c 769.078 308.199 771.84 309.039 774.719 309.039 c 778.078 309.039 780.961 308.078 783.48
306.281 c 785.879 304.48 787.801 302.078 789 299.199 c 790.199 296.32 790.922 293.199 790.922 289.84 c
790.922 284.68 789.602 279.398 786.961 273.879 c 784.32 268.359 780.719 263.801 776.281 260.199 c 771.719
256.602 766.922 254.801 761.762 254.801 c 759.359 254.801 757.32 255.52 755.398 256.719 c 753.48 257.922
751.922 259.719 750.961 261.879 c 745.559 240.641 l 745.32 239.441 745.199 238.84 745.199 238.602 c
745.199 237.641 747.84 237.039 753.121 237.039 c 753.84 237.039 754.32 236.559 754.32 235.48 c 754.078
234.52 753.84 233.801 753.602 233.441 c 753.359 233.078 753 232.84 752.281 232.84 c 729.121 232.84 l h
761.879 257.801 m 764.641 257.801 767.16 258.879 769.559 261.16 c 771.961 263.441 773.879 266.078 775.441
268.961 c 776.398 271.121 777.48 273.641 778.441 276.762 c 779.398 279.879 780.238 283.121 780.961
286.719 c 781.68 290.199 782.039 292.961 782.039 295.121 c 782.039 296.922 781.801 298.602 781.32 300.281
c 780.84 301.84 780 303.16 778.922 304.238 c 777.719 305.32 776.281 305.922 774.48 305.922 c 771.48
305.922 768.719 304.84 766.078 302.801 c 763.441 300.641 761.16 298.121 759.121 295 c 759.121 294.641 l
752.52 268.121 l 753 265.238 754.078 262.719 755.641 260.801 c 757.199 258.762 759.238 257.801 761.879
257.801 c f Q q 796.109 235.641 36 58 re W n 796.199 236 m 796.199 238.27 l 796.199 238.352 796.199
238.52 796.367 238.688 c 809.387 253.137 l 811.32 255.234 812.918 257 814.176 258.512 c 815.352 259.941
816.613 261.621 817.789 263.469 c 818.965 265.316 819.973 267.246 820.645 269.266 c 821.316 271.195
821.734 273.297 821.734 275.48 c 821.734 277.746 821.316 279.934 820.477 282.031 c 819.637 284.133
818.375 285.813 816.695 287.07 c 815.016 288.332 812.918 289.004 810.563 289.004 c 808.129 289.004
805.945 288.246 804.012 286.82 c 801.996 285.309 800.652 283.461 799.895 281.109 c 800.063 281.109
800.398 281.191 800.82 281.191 c 801.996 281.191 803.09 280.773 804.012 279.934 c 804.852 279.094 805.355
278 805.355 276.656 c 805.355 275.313 804.852 274.219 804.012 273.379 c 803.09 272.457 801.996 272.035
800.82 272.035 c 799.477 272.035 798.383 272.457 797.543 273.379 c 796.621 274.305 796.199 275.395
796.199 276.656 c 796.199 278.672 796.535 280.688 797.375 282.535 c 798.133 284.383 799.223 285.98
800.734 287.492 c 802.246 288.922 803.93 290.012 805.777 290.77 c 807.625 291.523 809.641 291.945 811.824
291.945 c 815.016 291.945 817.957 291.273 820.727 289.93 c 823.5 288.586 825.684 286.652 827.281 284.133
c 828.875 281.613 829.715 278.758 829.715 275.48 c 829.715 273.043 829.129 270.777 828.121 268.594 c
827.027 266.406 825.684 264.391 824.004 262.629 c 822.324 260.781 820.223 258.766 817.621 256.496 c
815.016 254.227 813.336 252.715 812.496 251.961 c 803.004 242.805 l 811.066 242.805 l 815.016 242.805
818.293 242.805 820.98 242.887 c 823.668 242.887 825.098 243.055 825.266 243.223 c 825.938 243.895
826.609 246.5 827.281 250.953 c 829.715 250.953 l 827.363 236 l 796.199 236 l f Q 1998.39 890.004 m
2064.52 942.902 2128.02 993.707 2188.91 1042.42 c S 2149.55 1033.98 m 2188.91 1042.42 l 2172.04 1005.87 l
2149.55 1033.98 l h B q 1976.64 1027.73 57 87 re W n 1977.36 1032.16 m 1977.36 1032.64 1977.36 1033
1977.48 1033.24 c 1994.88 1102.6 l 1995.12 1103.92 1995.36 1105 1995.48 1105.84 c 1995.48 1107.04 1992.84
1107.76 1987.68 1107.76 c 1986.84 1107.76 1986.48 1108.24 1986.48 1109.32 c 1986.48 1109.56 1986.48
1109.92 1986.72 1110.4 c 1986.84 1110.88 1986.96 1111.24 1987.2 1111.48 c 1987.44 1111.72 1987.68 1111.96
1988.16 1111.96 c 2004.36 1113.28 l 2004.72 1113.28 l 2004.72 1113.16 2004.84 1113.04 2005.08 1112.92 c
2005.32 1112.68 2005.56 1112.68 y 2005.68 1112.2 2005.8 1111.96 2005.8 1111.72 c 1993.2 1061.56 l 1996.08
1062.76 1999.68 1065.64 2004.12 1070.2 c 2008.44 1074.76 2011.92 1078 2014.44 1080.04 c 2016.84 1081.96
2019.96 1083.04 2023.8 1083.04 c 2025.96 1083.04 2027.88 1082.2 2029.44 1080.76 c 2031 1079.2 2031.84
1077.28 2031.84 1075.12 c 2031.84 1073.68 2031.48 1072.36 2030.88 1071.16 c 2030.28 1069.96 2029.44 1069
2028.48 1068.28 c 2027.4 1067.56 2026.2 1067.2 2024.76 1067.2 c 2023.44 1067.2 2022.36 1067.56 2021.52
1068.4 c 2020.56 1069.12 2020.2 1070.2 2020.2 1071.52 c 2020.2 1073.44 2020.8 1075 2022.24 1076.32 c
2023.56 1077.64 2025.24 1078.36 2027.16 1078.36 c 2026.32 1079.32 2025.12 1079.92 2023.56 1079.92 c
2021.16 1079.92 2018.76 1079.2 2016.6 1077.88 c 2014.44 1076.44 2012.04 1074.52 2009.52 1072 c 2006.88
1069.36 2004.6 1067.08 2002.44 1064.92 c 2000.28 1062.76 1998.36 1061.2 1996.68 1060.12 c 2001.36 1059.52
2005.44 1058.32 2008.92 1056.4 c 2012.4 1054.48 2014.2 1051.48 2014.2 1047.52 c 2014.2 1046.68 2013.96
1045.6 2013.72 1044.4 c 2013 1041.4 2012.64 1039 2012.64 1037.2 c 2012.64 1033.6 2013.84 1031.8 2016.36
1031.8 c 2019.24 1031.8 2021.4 1033.36 2022.96 1036.48 c 2024.4 1039.6 2025.72 1043.32 2026.68 1047.64 c
2026.8 1048.12 2027.16 1048.36 2027.64 1048.36 c 2029.08 1048.36 l 2029.32 1048.36 2029.56 1048.24 2029.8
1048 c 2030.04 1047.76 2030.16 1047.52 2030.16 1047.28 c 2030.16 1047.16 2030.04 1047.04 2030.04 1046.92
c 2027.04 1034.8 2022.36 1028.8 2016.12 1028.8 c 2013.84 1028.8 2011.8 1029.28 2010.12 1030.24 c 2008.32
1031.2 2007 1032.64 2006.04 1034.56 c 2005.08 1036.36 2004.72 1038.4 2004.72 1040.68 c 2004.72 1041.88
2004.84 1043.2 2005.2 1044.64 c 2005.44 1045.6 2005.56 1046.44 2005.56 1047.28 c 2005.56 1050.4 2004.12
1052.68 2001.48 1054.36 c 1998.72 1055.92 1995.6 1056.88 1992.12 1057.24 c 1986 1032.76 l 1985.64 1031.56
1985.04 1030.6 1984.2 1029.88 c 1983.24 1029.16 1982.28 1028.8 1981.08 1028.8 c 1980 1028.8 1979.16
1029.16 1978.44 1029.76 c 1977.72 1030.36 1977.36 1031.2 1977.36 1032.16 c f Q q 2040.04 1009.32 31 58 re
W n 2040.81 1009 m 2040.81 1011.94 l 2047.79 1011.94 2051.31 1012.78 2051.31 1014.63 c 2051.31 1058.73 l
2048.37 1057.3 2044.68 1056.63 2040.31 1056.63 c 2040.31 1059.57 l 2047.11 1059.57 2052.32 1061.33
2055.85 1064.95 c 2057.02 1064.95 l 2057.19 1064.95 2057.45 1064.86 2057.61 1064.69 c 2057.78 1064.52
2057.86 1064.36 2057.86 1064.19 c 2057.86 1014.63 l 2057.86 1012.78 2061.31 1011.94 2068.36 1011.94 c
2068.36 1009 l 2040.81 1009 l f Q 1998.39 440.145 m 2064.52 387.246 2128.02 336.441 2188.91 287.73 c S
2172.04 324.277 m 2188.91 287.73 l 2149.55 296.164 l 2172.04 324.277 l h B q 1976.64 251.859 57 87 re W n
1977.36 256.16 m 1977.36 256.641 1977.36 257 1977.48 257.238 c 1994.88 326.602 l 1995.12 327.922 1995.36
329 1995.48 329.84 c 1995.48 331.039 1992.84 331.762 1987.68 331.762 c 1986.84 331.762 1986.48 332.238
1986.48 333.32 c 1986.48 333.559 1986.48 333.922 1986.72 334.398 c 1986.84 334.879 1986.96 335.238 1987.2
335.48 c 1987.44 335.719 1987.68 335.961 1988.16 335.961 c 2004.36 337.281 l 2004.72 337.281 l 2004.72
337.16 2004.84 337.039 2005.08 336.922 c 2005.32 336.68 2005.56 336.68 y 2005.68 336.199 2005.8 335.961
2005.8 335.719 c 1993.2 285.559 l 1996.08 286.762 1999.68 289.641 2004.12 294.199 c 2008.44 298.762
2011.92 302 2014.44 304.039 c 2016.84 305.961 2019.96 307.039 2023.8 307.039 c 2025.96 307.039 2027.88
306.199 2029.44 304.762 c 2031 303.199 2031.84 301.281 2031.84 299.121 c 2031.84 297.68 2031.48 296.359
2030.88 295.16 c 2030.28 293.961 2029.44 293 2028.48 292.281 c 2027.4 291.559 2026.2 291.199 2024.76
291.199 c 2023.44 291.199 2022.36 291.559 2021.52 292.398 c 2020.56 293.121 2020.2 294.199 2020.2 295.52
c 2020.2 297.441 2020.8 299 2022.24 300.32 c 2023.56 301.641 2025.24 302.359 2027.16 302.359 c 2026.32
303.32 2025.12 303.922 2023.56 303.922 c 2021.16 303.922 2018.76 303.199 2016.6 301.879 c 2014.44 300.441
2012.04 298.52 2009.52 296 c 2006.88 293.359 2004.6 291.078 2002.44 288.922 c 2000.28 286.762 1998.36
285.199 1996.68 284.121 c 2001.36 283.52 2005.44 282.32 2008.92 280.398 c 2012.4 278.48 2014.2 275.48
2014.2 271.52 c 2014.2 270.68 2013.96 269.602 2013.72 268.398 c 2013 265.398 2012.64 263 2012.64 261.199
c 2012.64 257.602 2013.84 255.801 2016.36 255.801 c 2019.24 255.801 2021.4 257.359 2022.96 260.48 c
2024.4 263.602 2025.72 267.32 2026.68 271.641 c 2026.8 272.121 2027.16 272.359 2027.64 272.359 c 2029.08
272.359 l 2029.32 272.359 2029.56 272.238 2029.8 272 c 2030.04 271.762 2030.16 271.52 2030.16 271.281 c
2030.16 271.16 2030.04 271.039 2030.04 270.922 c 2027.04 258.801 2022.36 252.801 2016.12 252.801 c
2013.84 252.801 2011.8 253.281 2010.12 254.238 c 2008.32 255.199 2007 256.641 2006.04 258.559 c 2005.08
260.359 2004.72 262.398 2004.72 264.68 c 2004.72 265.879 2004.84 267.199 2005.2 268.641 c 2005.44 269.602
2005.56 270.441 2005.56 271.281 c 2005.56 274.398 2004.12 276.68 2001.48 278.359 c 1998.72 279.922 1995.6
280.879 1992.12 281.238 c 1986 256.762 l 1985.64 255.559 1985.04 254.602 1984.2 253.879 c 1983.24 253.16
1982.28 252.801 1981.08 252.801 c 1980 252.801 1979.16 253.16 1978.44 253.762 c 1977.72 254.359 1977.36
255.199 1977.36 256.16 c f Q q 2037.04 233.453 36 58 re W n 2037.2 233 m 2037.2 235.27 l 2037.2 235.352
2037.2 235.52 2037.37 235.688 c 2050.39 250.137 l 2052.32 252.234 2053.92 254 2055.18 255.512 c 2056.35
256.941 2057.61 258.621 2058.79 260.469 c 2059.96 262.316 2060.97 264.246 2061.64 266.266 c 2062.32
268.195 2062.73 270.297 2062.73 272.48 c 2062.73 274.746 2062.32 276.934 2061.48 279.031 c 2060.64
281.133 2059.38 282.813 2057.7 284.07 c 2056.02 285.332 2053.92 286.004 2051.56 286.004 c 2049.13 286.004
2046.95 285.246 2045.01 283.82 c 2043 282.309 2041.65 280.461 2040.89 278.109 c 2041.06 278.109 2041.4
278.191 2041.82 278.191 c 2043 278.191 2044.09 277.773 2045.01 276.934 c 2045.85 276.094 2046.36 275
2046.36 273.656 c 2046.36 272.313 2045.85 271.219 2045.01 270.379 c 2044.09 269.457 2043 269.035 2041.82
269.035 c 2040.48 269.035 2039.38 269.457 2038.54 270.379 c 2037.62 271.305 2037.2 272.395 2037.2 273.656
c 2037.2 275.672 2037.54 277.688 2038.38 279.535 c 2039.13 281.383 2040.22 282.98 2041.73 284.492 c
2043.25 285.922 2044.93 287.012 2046.78 287.77 c 2048.63 288.523 2050.64 288.945 2052.82 288.945 c
2056.02 288.945 2058.96 288.273 2061.73 286.93 c 2064.5 285.586 2066.68 283.652 2068.28 281.133 c 2069.88
278.613 2070.71 275.758 2070.71 272.48 c 2070.71 270.043 2070.13 267.777 2069.12 265.594 c 2068.03
263.406 2066.68 261.391 2065 259.629 c 2063.32 257.781 2061.22 255.766 2058.62 253.496 c 2056.02 251.227
2054.34 249.715 2053.5 248.961 c 2044 239.805 l 2052.07 239.805 l 2056.02 239.805 2059.29 239.805 2061.98
239.887 c 2064.67 239.887 2066.1 240.055 2066.27 240.223 c 2066.94 240.895 2067.61 243.5 2068.28 247.953
c 2070.71 247.953 l 2068.36 233 l 2037.2 233 l f Q q 1373.2 752.211 76 94 re W n 0 g 1412.3 755.699 m
1412.3 757.941 1412.72 761.02 1413.7 764.941 c 1414.54 769 1415.66 773.199 1417.06 777.68 c 1418.46
782.16 1419.58 785.801 1420.56 788.738 c 1420.84 793.359 1421.12 796.859 1421.12 799.238 c 1421.12
805.262 1420.42 810.859 1419.3 816.039 c 1418.18 821.219 1415.94 825.559 1412.58 829.059 c 1409.22
832.422 1404.74 834.238 1398.86 834.238 c 1395.78 834.238 1392.7 833.539 1389.62 832.281 c 1386.54
830.879 1383.88 829.059 1381.64 826.68 c 1379.4 824.301 1378 821.641 1377.16 818.699 c 1376.88 818.141
1376.6 817.859 1376.04 817.859 c 1374.36 817.859 l 1373.52 817.859 1373.24 818.281 1373.24 819.121 c
1373.24 819.539 l 1374.22 823.738 1376.04 827.801 1378.7 831.719 c 1381.36 835.641 1384.58 838.719 1388.5
841.238 c 1392.42 843.621 1396.48 844.879 1400.82 844.879 c 1405.16 844.879 1408.94 843.34 1412.16
840.398 c 1415.24 837.32 1417.76 833.539 1419.72 829.059 c 1421.68 824.578 1423.08 819.961 1423.92 815.34
c 1424.76 810.578 1425.18 806.238 1425.18 802.18 c 1429.94 814.922 1435.12 826.961 1441 838.441 c 1443.1
842.922 l 1443.38 843.199 1443.66 843.34 1444.08 843.34 c 1445.76 843.34 l 1446.46 843.34 1446.88 842.922
1446.88 842.078 c 1446.88 841.801 1446.74 841.52 1446.6 841.238 c 1444.5 836.898 l 1440.3 828.781 1436.66
820.52 1433.3 812.262 c 1429.94 804 1427 796.02 1424.62 788.32 c 1424.2 785.102 1423.5 780.898 1422.52
775.441 c 1421.4 770.121 1420.14 765.078 1418.74 760.32 c 1417.34 755.559 1415.8 753.18 1414.4 752.898 c
1413 752.898 1412.3 753.879 1412.3 755.699 c f Q 1278.14 542.633 m 1366.84 542.633 1452.18 542.633
1534.18 542.633 c S 1498.18 560.633 m 1534.18 542.633 l 1498.18 524.633 l 1498.18 560.633 l h B q 1227.2
383.371 65 79 re W n 1228.12 383.84 m 1227.4 383.84 1227.04 384.441 1227.04 385.398 c 1227.28 387.199
1227.88 388.039 1228.6 388.039 c 1231 388.039 1232.68 388.281 1233.64 388.762 c 1234.6 389.238 1235.44
390.32 1236.04 392.121 c 1249.6 446 l 1250.08 447.441 1250.44 449.359 1250.44 451.762 c 1250.44 455.121
1249.24 456.922 1246.96 456.922 c 1244.44 456.922 1242.52 455.359 1241.32 452.359 c 1240.12 449.359
1238.92 445.641 1237.84 441.078 c 1237.84 440.602 1237.48 440.359 1236.88 440.359 c 1235.44 440.359 l
1235.2 440.359 1235.08 440.48 1234.84 440.719 c 1234.6 440.961 1234.48 441.32 1234.48 441.559 c 1235.32
445.039 1236.16 448.039 1237 450.441 c 1237.72 452.84 1239.04 455 1240.72 457.039 c 1242.4 458.961
1244.44 460.039 1247.08 460.039 c 1249.72 460.039 1252.12 459.199 1254.28 457.641 c 1256.32 456.078
1257.64 454.039 1258.24 451.398 c 1260.4 453.922 1262.8 455.961 1265.44 457.641 c 1268.08 459.199 1270.84
460.039 1273.72 460.039 c 1277.08 460.039 1279.96 459.078 1282.48 457.281 c 1284.88 455.48 1286.8 453.078
1288 450.199 c 1289.2 447.32 1289.92 444.199 1289.92 440.84 c 1289.92 435.68 1288.6 430.398 1285.96
424.879 c 1283.32 419.359 1279.72 414.801 1275.28 411.199 c 1270.72 407.602 1265.92 405.801 1260.76
405.801 c 1258.36 405.801 1256.32 406.52 1254.4 407.719 c 1252.48 408.922 1250.92 410.719 1249.96 412.879
c 1244.56 391.641 l 1244.32 390.441 1244.2 389.84 1244.2 389.602 c 1244.2 388.641 1246.84 388.039 1252.12
388.039 c 1252.84 388.039 1253.32 387.559 1253.32 386.48 c 1253.08 385.52 1252.84 384.801 1252.6 384.441
c 1252.36 384.078 1252 383.84 1251.28 383.84 c 1228.12 383.84 l h 1260.88 408.801 m 1263.64 408.801
1266.16 409.879 1268.56 412.16 c 1270.96 414.441 1272.88 417.078 1274.44 419.961 c 1275.4 422.121 1276.48
424.641 1277.44 427.762 c 1278.4 430.879 1279.24 434.121 1279.96 437.719 c 1280.68 441.199 1281.04
443.961 1281.04 446.121 c 1281.04 447.922 1280.8 449.602 1280.32 451.281 c 1279.84 452.84 1279 454.16
1277.92 455.238 c 1276.72 456.32 1275.28 456.922 1273.48 456.922 c 1270.48 456.922 1267.72 455.84 1265.08
453.801 c 1262.44 451.641 1260.16 449.121 1258.12 446 c 1258.12 445.641 l 1251.52 419.121 l 1252 416.238
1253.08 413.719 1254.64 411.801 c 1256.2 409.762 1258.24 408.801 1260.88 408.801 c f Q q 1298.55 386.965
31 58 re W n 1299.81 387 m 1299.81 389.941 l 1306.79 389.941 1310.31 390.781 1310.31 392.629 c 1310.31
436.727 l 1307.37 435.301 1303.68 434.629 1299.31 434.629 c 1299.31 437.566 l 1306.11 437.566 1311.32
439.332 1314.85 442.945 c 1316.02 442.945 l 1316.19 442.945 1316.45 442.859 1316.61 442.691 c 1316.78
442.523 1316.86 442.355 1316.86 442.188 c 1316.86 392.629 l 1316.86 390.781 1320.31 389.941 1327.36
389.941 c 1327.36 387 l 1299.81 387 l f Q q 1368.27 397.371 82 82 re W n 1371 434.602 m 1370.28 434.602
1369.8 434.84 1369.32 435.32 c 1368.84 435.801 1368.72 436.398 1368.72 437 c 1368.72 437.48 1368.84
438.078 1369.32 438.559 c 1369.8 439.039 1370.28 439.398 1371 439.398 c 1406.28 439.398 l 1406.28 474.801
l 1406.28 475.398 1406.52 475.879 1407 476.359 c 1407.48 476.719 1407.96 476.961 1408.68 476.961 c
1409.16 476.961 1409.76 476.719 1410.24 476.359 c 1410.72 475.879 1411.08 475.398 1411.08 474.801 c
1411.08 439.398 l 1446.36 439.398 l 1446.84 439.398 1447.32 439.039 1447.8 438.559 c 1448.28 438.078
1448.52 437.48 1448.52 437 c 1448.52 436.398 1448.28 435.801 1447.8 435.32 c 1447.32 434.84 1446.84
434.602 1446.36 434.602 c 1411.08 434.602 l 1411.08 399.32 l 1411.08 398.719 1410.72 398.238 1410.24
397.762 c 1409.76 397.398 1409.16 397.16 1408.68 397.16 c 1407.96 397.16 1407.48 397.398 1407 397.762 c
1406.52 398.238 1406.28 398.719 1406.28 399.32 c 1406.28 434.602 l 1371 434.602 l f Q q 1472.57 383.371
65 79 re W n 1474.12 383.84 m 1473.4 383.84 1473.04 384.441 1473.04 385.398 c 1473.28 387.199 1473.88
388.039 1474.6 388.039 c 1477 388.039 1478.68 388.281 1479.64 388.762 c 1480.6 389.238 1481.44 390.32
1482.04 392.121 c 1495.6 446 l 1496.08 447.441 1496.44 449.359 1496.44 451.762 c 1496.44 455.121 1495.24
456.922 1492.96 456.922 c 1490.44 456.922 1488.52 455.359 1487.32 452.359 c 1486.12 449.359 1484.92
445.641 1483.84 441.078 c 1483.84 440.602 1483.48 440.359 1482.88 440.359 c 1481.44 440.359 l 1481.2
440.359 1481.08 440.48 1480.84 440.719 c 1480.6 440.961 1480.48 441.32 1480.48 441.559 c 1481.32 445.039
1482.16 448.039 1483 450.441 c 1483.72 452.84 1485.04 455 1486.72 457.039 c 1488.4 458.961 1490.44
460.039 1493.08 460.039 c 1495.72 460.039 1498.12 459.199 1500.28 457.641 c 1502.32 456.078 1503.64
454.039 1504.24 451.398 c 1506.4 453.922 1508.8 455.961 1511.44 457.641 c 1514.08 459.199 1516.84 460.039
1519.72 460.039 c 1523.08 460.039 1525.96 459.078 1528.48 457.281 c 1530.88 455.48 1532.8 453.078 1534
450.199 c 1535.2 447.32 1535.92 444.199 1535.92 440.84 c 1535.92 435.68 1534.6 430.398 1531.96 424.879 c
1529.32 419.359 1525.72 414.801 1521.28 411.199 c 1516.72 407.602 1511.92 405.801 1506.76 405.801 c
1504.36 405.801 1502.32 406.52 1500.4 407.719 c 1498.48 408.922 1496.92 410.719 1495.96 412.879 c 1490.56
391.641 l 1490.32 390.441 1490.2 389.84 1490.2 389.602 c 1490.2 388.641 1492.84 388.039 1498.12 388.039 c
1498.84 388.039 1499.32 387.559 1499.32 386.48 c 1499.08 385.52 1498.84 384.801 1498.6 384.441 c 1498.36
384.078 1498 383.84 1497.28 383.84 c 1474.12 383.84 l h 1506.88 408.801 m 1509.64 408.801 1512.16 409.879
1514.56 412.16 c 1516.96 414.441 1518.88 417.078 1520.44 419.961 c 1521.4 422.121 1522.48 424.641 1523.44
427.762 c 1524.4 430.879 1525.24 434.121 1525.96 437.719 c 1526.68 441.199 1527.04 443.961 1527.04
446.121 c 1527.04 447.922 1526.8 449.602 1526.32 451.281 c 1525.84 452.84 1525 454.16 1523.92 455.238 c
1522.72 456.32 1521.28 456.922 1519.48 456.922 c 1516.48 456.922 1513.72 455.84 1511.08 453.801 c 1508.44
451.641 1506.16 449.121 1504.12 446 c 1504.12 445.641 l 1497.52 419.121 l 1498 416.238 1499.08 413.719
1500.64 411.801 c 1502.2 409.762 1504.24 408.801 1506.88 408.801 c f Q q 1540.92 386.965 36 58 re W n
1541.2 387 m 1541.2 389.27 l 1541.2 389.352 1541.2 389.52 1541.37 389.688 c 1554.39 404.137 l 1556.32
406.234 1557.92 408 1559.18 409.512 c 1560.35 410.941 1561.61 412.621 1562.79 414.469 c 1563.96 416.316
1564.97 418.246 1565.64 420.266 c 1566.32 422.195 1566.73 424.297 1566.73 426.48 c 1566.73 428.746
1566.32 430.934 1565.48 433.031 c 1564.64 435.133 1563.38 436.813 1561.7 438.07 c 1560.02 439.332 1557.92
440.004 1555.56 440.004 c 1553.13 440.004 1550.95 439.246 1549.01 437.82 c 1547 436.309 1545.65 434.461
1544.89 432.109 c 1545.06 432.109 1545.4 432.191 1545.82 432.191 c 1547 432.191 1548.09 431.773 1549.01
430.934 c 1549.85 430.094 1550.36 429 1550.36 427.656 c 1550.36 426.313 1549.85 425.219 1549.01 424.379 c
1548.09 423.457 1547 423.035 1545.82 423.035 c 1544.48 423.035 1543.38 423.457 1542.54 424.379 c 1541.62
425.305 1541.2 426.395 1541.2 427.656 c 1541.2 429.672 1541.54 431.688 1542.38 433.535 c 1543.13 435.383
1544.22 436.98 1545.73 438.492 c 1547.25 439.922 1548.93 441.012 1550.78 441.77 c 1552.63 442.523 1554.64
442.945 1556.82 442.945 c 1560.02 442.945 1562.96 442.273 1565.73 440.93 c 1568.5 439.586 1570.68 437.652
1572.28 435.133 c 1573.88 432.613 1574.71 429.758 1574.71 426.48 c 1574.71 424.043 1574.13 421.777
1573.12 419.594 c 1572.03 417.406 1570.68 415.391 1569 413.629 c 1567.32 411.781 1565.22 409.766 1562.62
407.496 c 1560.02 405.227 1558.34 403.715 1557.5 402.961 c 1548 393.805 l 1556.07 393.805 l 1560.02
393.805 1563.29 393.805 1565.98 393.887 c 1568.67 393.887 1570.1 394.055 1570.27 394.223 c 1570.94
394.895 1571.61 397.5 1572.28 401.953 c 1574.71 401.953 l 1572.36 387 l 1541.2 387 l f Q q 991.57 721.473
74 88 re W n 0 g 992.641 724.691 m 992.641 725.211 992.641 725.602 992.77 725.73 c 1011.88 802.301 l
1012.14 803.469 1012.79 804.512 1013.83 805.289 c 1014.87 806.07 1016.04 806.461 1017.34 806.461 c
1018.38 806.461 1019.29 806.07 1020.07 805.422 c 1020.85 804.77 1021.24 803.859 1021.24 802.82 c 1021.24
802.559 1021.24 802.301 1021.24 802.172 c 1021.24 802.039 1021.11 801.91 1021.11 801.781 c 1013.44
771.359 l 1012.53 767.719 1012.14 764.602 1012.14 762 c 1012.14 758.75 1012.79 756.148 1014.35 754.07 c
1015.91 751.988 1018.25 750.949 1021.37 750.949 c 1027.61 750.949 1033.2 754.852 1038.01 762.781 c
1038.01 762.91 l 1038.14 763.039 1038.14 763.172 v 1047.5 800.871 l 1047.76 801.91 1048.54 802.949
1049.58 803.73 c 1050.62 804.512 1051.66 805.031 1052.96 805.031 c 1054 805.031 1054.91 804.641 1055.69
803.988 c 1056.47 803.34 1056.86 802.43 1056.86 801.391 c 1056.86 800.871 1056.73 800.48 1056.73 800.352
c 1047.37 762.91 l 1046.72 760.441 1046.46 758.359 1046.46 756.539 c 1046.46 752.77 1047.63 750.949
1050.23 750.949 c 1052.83 750.949 1054.91 752.512 1056.34 755.891 c 1057.64 759.27 1058.94 763.301
1059.98 768.109 c 1060.11 768.629 1060.5 768.891 1061.02 768.891 c 1062.58 768.891 l 1062.84 768.891
1063.1 768.629 1063.36 768.371 c 1063.62 768.109 1063.75 767.852 1063.75 767.59 c 1062.19 761.48 1060.63
756.672 1058.81 753.031 c 1056.99 749.391 1054 747.699 1049.97 747.699 c 1047.11 747.699 1044.51 748.48
1042.3 750.039 c 1040.09 751.602 1038.66 753.809 1038.01 756.672 c 1035.8 753.941 1033.2 751.73 1030.34
750.172 c 1027.48 748.609 1024.36 747.699 1021.24 747.699 c 1015.78 747.699 1011.62 749 1008.5 751.34 c
1001.87 725.211 l 1001.61 724.039 1000.96 723 999.922 722.219 c 998.879 721.441 997.84 721.051 996.539
721.051 c 995.5 721.051 994.59 721.441 993.809 722.09 c 993.031 722.738 992.641 723.648 992.641 724.691 c
f Q q 1764.83 746.504 64 60 re W n 0 g 1767.19 747 m 1766.15 747 1765.63 747.52 1765.63 748.691 c 1776.68
792.891 l 1776.94 794.32 1777.2 795.488 1777.33 796.398 c 1777.33 797.699 1774.47 798.48 1768.88 798.48 c
1767.97 798.48 1767.58 799 1767.58 800.172 c 1767.58 800.43 1767.58 800.82 1767.84 801.34 c 1767.97
801.859 1768.1 802.25 1768.36 802.512 c 1768.62 802.77 1768.88 803.031 1769.4 803.031 c 1786.95 804.461 l
1787.99 804.461 1788.51 803.809 1788.51 802.77 c 1775.77 751.809 l 1782.53 754.02 1788.9 757.531 1794.62
762.469 c 1800.34 767.41 1805.28 773.129 1809.31 779.762 c 1813.34 786.262 1816.2 793.148 1818.02 800.301
c 1818.28 801.469 1818.93 802.512 1819.97 803.289 c 1821.01 804.07 1822.18 804.461 1823.48 804.461 c
1824.52 804.461 1825.43 804.07 1826.21 803.422 c 1826.99 802.641 1827.38 801.73 1827.38 800.691 c 1827.38
799.262 1826.73 797.18 1825.69 794.32 c 1824.52 791.461 1823.35 788.859 1822.31 786.648 c 1818.8 779.891
1814.25 773.648 1808.66 767.93 c 1805.02 764.289 1801.25 761.172 1797.22 758.441 c 1793.06 755.711
1788.77 753.371 1784.35 751.422 c 1779.93 749.469 1775.25 748.039 1770.31 747 c 1767.19 747 l f Q q
147.578 1139.43 57 66 re W n 0 g 170.141 1139.6 m 165.66 1139.6 161.738 1140.72 158.379 1142.96 c 155.02
1145.2 152.5 1148.42 150.82 1152.34 c 149 1156.26 148.16 1160.32 148.16 1164.8 c 148.16 1171.38 149.84
1177.68 153.199 1183.56 c 156.559 1189.44 161.039 1194.06 166.781 1197.56 c 172.52 1201.06 178.68 1202.88
185.262 1202.88 c 189.18 1202.88 192.68 1201.76 195.762 1199.66 c 198.699 1197.56 200.238 1194.62 200.238
1190.84 c 200.238 1185.38 198 1181.32 193.801 1178.8 c 189.461 1176.14 184.559 1174.46 179.238 1173.9 c
173.922 1173.34 167.898 1173.06 161.461 1173.06 c 161.039 1173.06 l 159.5 1167.46 158.801 1162.7 158.801
1158.78 c 158.801 1154.58 159.641 1150.94 161.602 1147.86 c 163.559 1144.64 166.5 1143.1 170.422 1143.1 c
176.02 1143.1 181.34 1144.36 186.52 1146.88 c 191.559 1149.4 195.762 1153.04 198.98 1157.52 c 199.262
1157.8 199.539 1157.94 200.102 1157.94 c 200.52 1157.94 200.941 1157.66 201.5 1157.1 c 201.922 1156.54
202.199 1156.12 202.199 1155.7 c 202.199 1155.28 202.059 1155 201.922 1154.86 c 198.559 1150.1 193.941
1146.32 188.199 1143.66 c 182.32 1141 176.301 1139.6 170.141 1139.6 c h 161.879 1176.7 m 167.199 1176.7
172.238 1176.84 176.859 1177.4 c 181.48 1177.96 185.68 1179.22 189.461 1181.18 c 193.102 1183.14 195.059
1186.36 195.059 1190.7 c 195.059 1192.38 194.5 1193.92 193.52 1195.18 c 192.539 1196.44 191.281 1197.42
189.879 1198.12 c 188.34 1198.82 186.801 1199.24 185.121 1199.24 c 181.199 1199.24 177.559 1198.12 174.48
1196.02 c 171.262 1193.78 168.602 1190.98 166.641 1187.62 c 164.539 1184.12 163 1180.48 161.879 1176.7 c
f Q q 223.484 1217 62 6 re W n 0 g 224.996 1217.54 m 224.406 1217.54 224.016 1217.73 223.625 1218.13 c
223.23 1218.52 223.133 1219.01 223.133 1219.5 c 223.133 1219.89 223.23 1220.38 223.625 1220.77 c 224.016
1221.17 224.406 1221.46 224.996 1221.46 c 281.246 1221.46 l 281.738 1221.46 282.129 1221.17 282.523
1220.77 c 282.816 1220.38 283.012 1219.89 283.012 1219.5 c 283.012 1219.01 282.816 1218.52 282.523
1218.13 c 282.129 1217.73 281.738 1217.54 281.246 1217.54 c 224.996 1217.54 l f Q q 147.578 128.668 57 66
re W n 0 g 170.141 129.602 m 165.66 129.602 161.738 130.719 158.379 132.961 c 155.02 135.199 152.5
138.422 150.82 142.34 c 149 146.262 148.16 150.32 148.16 154.801 c 148.16 161.379 149.84 167.68 153.199
173.559 c 156.559 179.441 161.039 184.059 166.781 187.559 c 172.52 191.059 178.68 192.879 185.262 192.879
c 189.18 192.879 192.68 191.762 195.762 189.66 c 198.699 187.559 200.238 184.621 200.238 180.84 c 200.238
175.379 198 171.32 193.801 168.801 c 189.461 166.141 184.559 164.461 179.238 163.898 c 173.922 163.34
167.898 163.059 161.461 163.059 c 161.039 163.059 l 159.5 157.461 158.801 152.699 158.801 148.781 c
158.801 144.578 159.641 140.941 161.602 137.859 c 163.559 134.641 166.5 133.102 170.422 133.102 c 176.02
133.102 181.34 134.359 186.52 136.879 c 191.559 139.398 195.762 143.039 198.98 147.52 c 199.262 147.801
199.539 147.941 200.102 147.941 c 200.52 147.941 200.941 147.66 201.5 147.102 c 201.922 146.539 202.199
146.121 202.199 145.699 c 202.199 145.281 202.059 145 201.922 144.859 c 198.559 140.102 193.941 136.32
188.199 133.66 c 182.32 131 176.301 129.602 170.141 129.602 c h 161.879 166.699 m 167.199 166.699 172.238
166.84 176.859 167.398 c 181.48 167.961 185.68 169.219 189.461 171.18 c 193.102 173.141 195.059 176.359
195.059 180.699 c 195.059 182.379 194.5 183.922 193.52 185.18 c 192.539 186.441 191.281 187.422 189.879
188.121 c 188.34 188.82 186.801 189.238 185.121 189.238 c 181.199 189.238 177.559 188.121 174.48 186.02 c
171.262 183.781 168.602 180.98 166.641 177.621 c 164.539 174.121 163 170.48 161.879 166.699 c f Q q
220.484 175.238 68 68 re W n 0 g 222.352 206.539 m 221.762 206.539 221.371 206.734 220.977 207.129 c
220.586 207.52 220.488 208.012 220.488 208.5 c 220.488 208.891 220.586 209.383 220.977 209.773 c 221.371
210.164 221.762 210.461 222.352 210.461 c 251.16 210.461 l 251.16 239.371 l 251.16 239.859 251.359
240.254 251.75 240.645 c 252.141 240.938 252.535 241.133 253.121 241.133 c 253.516 241.133 254.004
240.938 254.395 240.645 c 254.789 240.254 255.082 239.859 255.082 239.371 c 255.082 210.461 l 283.895
210.461 l 284.285 210.461 284.68 210.164 285.07 209.773 c 285.461 209.383 285.656 208.891 285.656 208.5 c
285.656 208.012 285.461 207.52 285.07 207.129 c 284.68 206.734 284.285 206.539 283.895 206.539 c 255.082
206.539 l 255.082 177.727 l 255.082 177.238 254.789 176.848 254.395 176.453 c 254.004 176.16 253.516
175.965 253.121 175.965 c 252.535 175.965 252.141 176.16 251.75 176.453 c 251.359 176.848 251.16 177.238
251.16 177.727 c 251.16 206.539 l 222.352 206.539 l f Q q 2525.82 1112.54 61 91 re W n 0 g 2541.58 1115.8
m 2541.86 1117.9 2542.42 1118.88 2543.4 1118.88 c 2546.76 1118.88 2549.42 1119.16 2551.38 1119.72 c
2553.34 1120.28 2554.6 1121.68 2555.16 1123.64 c 2561.18 1147.3 l 2555.86 1142.12 2550.54 1139.6 2545.22
1139.6 c 2539.34 1139.6 2534.72 1141.7 2531.36 1146.18 c 2528 1150.52 2526.32 1155.7 2526.32 1161.86 c
2526.32 1167.74 2527.86 1174.04 2530.94 1180.48 c 2534.02 1186.92 2538.08 1192.24 2543.4 1196.44 c
2548.72 1200.64 2554.32 1202.88 2560.48 1202.88 c 2563.28 1202.88 2565.94 1202.04 2568.18 1200.36 c
2570.42 1198.68 2572.1 1196.44 2573.36 1193.64 c 2573.92 1194.76 2575.32 1196.58 2577.84 1199.1 c 2580.36
1201.62 2582.04 1202.88 2583.16 1202.88 c 2584.14 1202.88 2584.7 1202.18 2584.7 1201.06 c 2565.1 1123.08
l 2564.82 1121.68 2564.68 1120.98 2564.68 1120.7 c 2564.68 1119.58 2567.76 1118.88 2573.92 1118.88 c
2574.34 1118.88 2574.62 1118.74 2574.9 1118.32 c 2575.18 1118.04 2575.32 1117.62 2575.32 1117.06 c
2575.04 1115.94 2574.76 1115.1 2574.48 1114.68 c 2574.2 1114.26 2573.78 1113.98 2572.94 1113.98 c 2542.98
1113.98 l 2542 1113.98 2541.58 1114.68 2541.58 1115.8 c h 2545.5 1143.1 m 2548.72 1143.1 2551.94 1144.22
2555.02 1146.74 c 2558.1 1149.26 2560.9 1152.2 2563.28 1155.7 c 2571.12 1187.2 l 2570.56 1190.42 2569.3
1193.22 2567.48 1195.6 c 2565.66 1197.98 2563.28 1199.24 2560.2 1199.24 c 2556.98 1199.24 2554.04 1197.84
2551.38 1195.32 c 2548.58 1192.8 2546.34 1189.72 2544.52 1186.22 c 2542.56 1182.44 2540.88 1177.26 2539.2
1170.82 c 2537.52 1164.38 2536.82 1159.34 2536.82 1155.7 c 2536.82 1152.34 2537.52 1149.4 2538.92 1146.88
c 2540.32 1144.36 2542.42 1143.1 2545.5 1143.1 c f Q q 2534.71 206.613 53 7 re W n 0 g 2535.66 207.699 m
2535.66 212.602 l 2586.2 212.602 l 2586.2 207.699 l 2535.66 207.699 l f Q q 2525.82 99.8945 61 91 re W n
0 g 2541.58 102.801 m 2541.86 104.898 2542.42 105.879 2543.4 105.879 c 2546.76 105.879 2549.42 106.16
2551.38 106.719 c 2553.34 107.281 2554.6 108.68 2555.16 110.641 c 2561.18 134.301 l 2555.86 129.121
2550.54 126.602 2545.22 126.602 c 2539.34 126.602 2534.72 128.699 2531.36 133.18 c 2528 137.52 2526.32
142.699 2526.32 148.859 c 2526.32 154.738 2527.86 161.039 2530.94 167.48 c 2534.02 173.922 2538.08
179.238 2543.4 183.441 c 2548.72 187.641 2554.32 189.879 2560.48 189.879 c 2563.28 189.879 2565.94
189.039 2568.18 187.359 c 2570.42 185.68 2572.1 183.441 2573.36 180.641 c 2573.92 181.762 2575.32 183.578
2577.84 186.102 c 2580.36 188.621 2582.04 189.879 2583.16 189.879 c 2584.14 189.879 2584.7 189.18 2584.7
188.059 c 2565.1 110.078 l 2564.82 108.68 2564.68 107.98 2564.68 107.699 c 2564.68 106.578 2567.76
105.879 2573.92 105.879 c 2574.34 105.879 2574.62 105.738 2574.9 105.32 c 2575.18 105.039 2575.32 104.621
2575.32 104.059 c 2575.04 102.941 2574.76 102.102 2574.48 101.68 c 2574.2 101.262 2573.78 100.98 2572.94
100.98 c 2542.98 100.98 l 2542 100.98 2541.58 101.68 2541.58 102.801 c h 2545.5 130.102 m 2548.72 130.102
2551.94 131.219 2555.02 133.738 c 2558.1 136.262 2560.9 139.199 2563.28 142.699 c 2571.12 174.199 l
2570.56 177.422 2569.3 180.219 2567.48 182.602 c 2565.66 184.98 2563.28 186.238 2560.2 186.238 c 2556.98
186.238 2554.04 184.84 2551.38 182.32 c 2548.58 179.801 2546.34 176.719 2544.52 173.219 c 2542.56 169.441
2540.88 164.262 2539.2 157.82 c 2537.52 151.379 2536.82 146.34 2536.82 142.699 c 2536.82 139.34 2537.52
136.398 2538.92 133.879 c 2540.32 131.359 2542.42 130.102 2545.5 130.102 c f Q Q
}

Time runs from left to right. Arrows on the fermion lines show the fermion-number flow; grey arrows show the four-momenta.

Label the spins of $e^-,e^+,q,\bar q$ by $s,s',r,r'$ and the colors of $q,\bar q$ by $c,c'\in\{1,2,3\}$; the photon vertex is diagonal in color, and $p_1+p_2=k_1+k_2$. Reading each fermion line against its arrow with \eqref{eq:7.2},
\begin{equation*}
\begin{aligned}
i\mathcal M
={}&\bigl[\bar v^{s'}(p_2)\,(-ie\gamma^\mu)\,u^s(p_1)\bigr]\,
\frac{-i}{(p_1+p_2)^2+i\varepsilon}
\left[\eta_{\mu\nu}-(1-\xi)\frac{(p_1+p_2)_\mu(p_1+p_2)_\nu}{(p_1+p_2)^2+i\varepsilon}\right]\\[6pt]
&\times\bigl[\bar u^r(k_1)\,(ieQ_q\gamma^\nu)\,v^{r'}(k_2)\bigr]\,\delta_{cc'} .
\end{aligned}
\tag{7.3}\label{eq:7.3}
\end{equation*}

\topic{Superficial Degree of Divergence}

For a connected QED graph, let $V$ be the number of vertices, $I_\gamma,I_F$ the numbers of internal photon and fermion lines, and $E_\gamma,E_F$ the numbers of external photon and fermion lines. Each vertex of \eqref{eq:7.1} joins two fermion lines and one photon line and carries no derivative. Counting line ends, and counting powers of loop momenta at $d=4$ with photon propagators $\sim\ell^{-2}$ and fermion propagators $\sim\ell^{-1}$,
\begin{equation*}
\begin{aligned}
2V&=2I_F+E_F,\qquad V=2I_\gamma+E_\gamma,\qquad L=I_\gamma+I_F-V+1,\\[6pt]
\omega&=4L-2I_\gamma-I_F=2I_\gamma+3I_F-4V+4 .
\end{aligned}
\tag{7.4}\label{eq:7.4}
\end{equation*}
Eliminating $I_\gamma=\tfrac12(V-E_\gamma)$ and $I_F=V-\tfrac12E_F$,
\begin{equation*}
{\displaystyle \omega=4-E_\gamma-\frac32E_F .}
\tag{7.5}\label{eq:7.5}
\end{equation*}
The dependence on $V$ cancels because $[e]=0$ in four dimensions.

\topic{Two External Legs}

Every vertex conserves fermion number, so $E_F$ is even and two external legs means $(E_F,E_\gamma)=(2,0)$ or $(0,2)$. Let $-i\Sigma(p)$ and $i\Pi^{\mu\nu}(p)$ be the sums of amputated 1PI graphs with two external fermion lines and with two external photon lines, where $p$ is the external momentum.

\topic{(1) Fermion Propagator}
\begin{equation*}
E_F=2,\qquad E_\gamma=0\qquad\Longrightarrow\qquad\omega=1 .
\tag{7.6}\label{eq:7.6}
\end{equation*}

% fermion-self-energy
\qftdiagram{206.62}{105.22}{%
q 0.1 0 0 0.1 0 0 cm 1 g 0 0 2066.2 1052.23 re f 0 g 405.586 306.102 m 313.754 340.539 l 313.754 271.668
l f 1075.19 306.102 m 983.355 340.539 l 983.355 271.668 l f 1744.79 306.102 m 1652.96 340.539 l 1652.96
271.668 l f 14 w 1 J 1 j 0 G 114.789 306.102 m 612.207 306.102 l S 612.207 306.102 m 1453.99 306.102 l S
1453.99 306.102 m 1951.41 306.102 l S 2 J 612.207 306.102 m 595.133 323.355 l 586.848 334.25 l 580.922
345.363 l 577.895 356.484 l 578.074 367.375 l 581.527 377.805 l 588.074 387.586 l 597.293 396.586 l
612.648 407.305 l 634.148 418.914 l 661.637 433.625 l 670.496 440.238 l 676.914 447.516 l 680.707 455.738
l 681.922 465.113 l 680.84 475.742 l 676.676 491.793 l 662.898 537.699 l 660.855 551.395 l 660.895
564.184 l 663.398 575.598 l 668.555 585.258 l 676.34 592.914 l 686.531 598.461 l 698.73 601.965 l 712.391
603.672 l 731.805 603.859 l 780.867 602.723 l 790.969 604.863 l 799.258 608.992 l 805.797 615.355 l
810.813 624.004 l 815.746 638.785 l 832.438 702.613 l 838.629 714.02 l 846.262 723.09 l 855.328 729.375 l
865.699 732.625 l 877.148 732.824 l 889.371 730.195 l 902.02 725.184 l 918.957 715.902 l 964.652 688.398
l 974.367 685.617 l 983.527 685.508 l 992.328 688.219 l 1000.98 693.633 l 1012.7 704.422 l 1052.14
745.582 l 1063.09 753.52 l 1074.22 758.996 l 1085.32 761.504 l 1096.14 760.777 l 1106.48 756.785 l
1116.14 749.762 l 1125 740.152 l 1135.54 724.465 l 1149.12 698.641 l 1161.69 675.934 l 1168.39 667.457 l
1175.81 661.469 l 1184.21 658.117 l 1193.79 657.309 l 1204.63 658.734 l 1220.93 663.195 l 1262.36 675.797
l 1280.6 678.492 l 1293.19 678.066 l 1304.33 675.129 l 1313.67 669.527 l 1320.98 661.32 l 1326.18 650.762
l 1329.36 638.281 l 1330.96 619.621 l 1330.27 595.199 l 1329.27 564.035 l 1330.56 553.059 l 1333.66
543.859 l 1338.88 536.465 l 1346.39 530.723 l 1356.14 526.348 l 1372.12 521.926 l 1418.77 510.906 l
1431.65 505.824 l 1442.7 499.398 l 1451.34 491.527 l 1457.13 482.23 l 1459.86 471.66 l 1459.57 460.059 l
1456.51 447.742 l 1451.15 435.055 l 1441.61 418.152 l 1416.09 376.23 l 1412.9 366.41 l 1412.33 357.168 l
1414.57 348.324 l 1419.55 339.656 l 1429.88 327.996 l 1453.99 306.102 l 1453.99 306.102 l S 10 w 0 J
612.207 283.602 m 618.172 283.602 623.895 285.973 628.117 290.191 c 632.336 294.414 634.707 300.137
634.707 306.102 c 634.707 312.07 632.336 317.793 628.117 322.012 c 623.895 326.23 618.172 328.602 612.207
328.602 c 606.238 328.602 600.516 326.23 596.297 322.012 c 592.078 317.793 589.707 312.07 589.707 306.102
c 589.707 300.137 592.078 294.414 596.297 290.191 c 600.516 285.973 606.238 283.602 612.207 283.602 c h B
1453.99 283.602 m 1459.96 283.602 1465.68 285.973 1469.9 290.191 c 1474.12 294.414 1476.49 300.137
1476.49 306.102 c 1476.49 312.07 1474.12 317.793 1469.9 322.012 c 1465.68 326.23 1459.96 328.602 1453.99
328.602 c 1448.02 328.602 1442.3 326.23 1438.08 322.012 c 1433.86 317.793 1431.49 312.07 1431.49 306.102
c 1431.49 300.137 1433.86 294.414 1438.08 290.191 c 1442.3 285.973 1448.02 283.602 1453.99 283.602 c h B
q 468.07 362.5 74 88 re W n 468.641 366.691 m 468.641 367.211 468.641 367.602 468.77 367.73 c 487.879
444.301 l 488.141 445.469 488.789 446.512 489.828 447.289 c 490.871 448.07 492.039 448.461 493.34 448.461
c 494.379 448.461 495.289 448.07 496.07 447.422 c 496.852 446.77 497.238 445.859 497.238 444.82 c 497.238
444.559 497.238 444.301 497.238 444.172 c 497.238 444.039 497.109 443.91 497.109 443.781 c 489.441
413.359 l 488.531 409.719 488.141 406.602 488.141 404 c 488.141 400.75 488.789 398.148 490.352 396.07 c
491.91 393.988 494.25 392.949 497.371 392.949 c 503.609 392.949 509.199 396.852 514.012 404.781 c 514.012
404.91 l 514.141 405.039 514.141 405.172 v 523.5 442.871 l 523.762 443.91 524.539 444.949 525.578 445.73
c 526.621 446.512 527.66 447.031 528.961 447.031 c 530 447.031 530.91 446.641 531.691 445.988 c 532.469
445.34 532.859 444.43 532.859 443.391 c 532.859 442.871 532.73 442.48 532.73 442.352 c 523.371 404.91 l
522.719 402.441 522.461 400.359 522.461 398.539 c 522.461 394.77 523.629 392.949 526.23 392.949 c 528.828
392.949 530.91 394.512 532.34 397.891 c 533.641 401.27 534.941 405.301 535.98 410.109 c 536.109 410.629
536.5 410.891 537.02 410.891 c 538.578 410.891 l 538.84 410.891 539.102 410.629 539.359 410.371 c 539.621
410.109 539.75 409.852 539.75 409.59 c 538.191 403.48 536.629 398.672 534.809 395.031 c 532.988 391.391
530 389.699 525.969 389.699 c 523.109 389.699 520.512 390.48 518.301 392.039 c 516.09 393.602 514.66
395.809 514.012 398.672 c 511.801 395.941 509.199 393.73 506.34 392.172 c 503.48 390.609 500.359 389.699
497.238 389.699 c 491.781 389.699 487.621 391 484.5 393.34 c 477.871 367.211 l 477.609 366.039 476.961
365 475.922 364.219 c 474.879 363.441 473.84 363.051 472.539 363.051 c 471.5 363.051 470.59 363.441
469.809 364.09 c 469.031 364.738 468.641 365.648 468.641 366.691 c f Q q 1532.13 387.531 64 60 re W n
1534.19 388 m 1533.15 388 1532.63 388.52 1532.63 389.691 c 1543.68 433.891 l 1543.94 435.32 1544.2
436.488 1544.33 437.398 c 1544.33 438.699 1541.47 439.48 1535.88 439.48 c 1534.97 439.48 1534.58 440
1534.58 441.172 c 1534.58 441.43 1534.58 441.82 1534.84 442.34 c 1534.97 442.859 1535.1 443.25 1535.36
443.512 c 1535.62 443.77 1535.88 444.031 1536.4 444.031 c 1553.95 445.461 l 1554.99 445.461 1555.51
444.809 1555.51 443.77 c 1542.77 392.809 l 1549.53 395.02 1555.9 398.531 1561.62 403.469 c 1567.34 408.41
1572.28 414.129 1576.31 420.762 c 1580.34 427.262 1583.2 434.148 1585.02 441.301 c 1585.28 442.469
1585.93 443.512 1586.97 444.289 c 1588.01 445.07 1589.18 445.461 1590.48 445.461 c 1591.52 445.461
1592.43 445.07 1593.21 444.422 c 1593.99 443.641 1594.38 442.73 1594.38 441.691 c 1594.38 440.262 1593.73
438.18 1592.69 435.32 c 1591.52 432.461 1590.35 429.859 1589.31 427.648 c 1585.8 420.891 1581.25 414.648
1575.66 408.93 c 1572.02 405.289 1568.25 402.172 1564.22 399.441 c 1560.06 396.711 1555.77 394.371
1551.35 392.422 c 1546.93 390.469 1542.25 389.039 1537.31 388 c 1534.19 388 l f Q 9 w 1 J 0.449951 G
230.449 191.313 m 312.777 191.313 391.754 191.313 467.379 191.313 c S 0.449951 g 431.379 209.313 m
467.379 191.313 l 431.379 173.313 l 431.379 209.313 l h B q 313.02 39.0859 65 79 re W n 314.121 39.8398 m
313.398 39.8398 313.039 40.4414 313.039 41.3984 c 313.281 43.1992 313.879 44.0391 314.602 44.0391 c 317
44.0391 318.68 44.2813 319.641 44.7617 c 320.602 45.2383 321.441 46.3203 322.039 48.1211 c 335.602 102 l
336.078 103.441 336.441 105.359 336.441 107.762 c 336.441 111.121 335.238 112.922 332.961 112.922 c
330.441 112.922 328.52 111.359 327.32 108.359 c 326.121 105.359 324.922 101.641 323.84 97.0781 c 323.84
96.6016 323.48 96.3594 322.879 96.3594 c 321.441 96.3594 l 321.199 96.3594 321.078 96.4805 320.84 96.7188
c 320.602 96.9609 320.48 97.3203 320.48 97.5586 c 321.32 101.039 322.16 104.039 323 106.441 c 323.719
108.84 325.039 111 326.719 113.039 c 328.398 114.961 330.441 116.039 333.078 116.039 c 335.719 116.039
338.121 115.199 340.281 113.641 c 342.32 112.078 343.641 110.039 344.238 107.398 c 346.398 109.922
348.801 111.961 351.441 113.641 c 354.078 115.199 356.84 116.039 359.719 116.039 c 363.078 116.039
365.961 115.078 368.48 113.281 c 370.879 111.48 372.801 109.078 374 106.199 c 375.199 103.32 375.922
100.199 375.922 96.8398 c 375.922 91.6797 374.602 86.3984 371.961 80.8789 c 369.32 75.3594 365.719
70.8008 361.281 67.1992 c 356.719 63.6016 351.922 61.8008 346.762 61.8008 c 344.359 61.8008 342.32
62.5195 340.398 63.7188 c 338.48 64.9219 336.922 66.7188 335.961 68.8789 c 330.559 47.6406 l 330.32
46.4414 330.199 45.8398 330.199 45.6016 c 330.199 44.6406 332.84 44.0391 338.121 44.0391 c 338.84 44.0391
339.32 43.5586 339.32 42.4805 c 339.078 41.5195 338.84 40.8008 338.602 40.4414 c 338.359 40.0781 338
39.8398 337.281 39.8398 c 314.121 39.8398 l h 346.879 64.8008 m 349.641 64.8008 352.16 65.8789 354.559
68.1602 c 356.961 70.4414 358.879 73.0781 360.441 75.9609 c 361.398 78.1211 362.48 80.6406 363.441
83.7617 c 364.398 86.8789 365.238 90.1211 365.961 93.7188 c 366.68 97.1992 367.039 99.9609 367.039
102.121 c 367.039 103.922 366.801 105.602 366.32 107.281 c 365.84 108.84 365 110.16 363.922 111.238 c
362.719 112.32 361.281 112.922 359.48 112.922 c 356.48 112.922 353.719 111.84 351.078 109.801 c 348.441
107.641 346.16 105.121 344.121 102 c 344.121 101.641 l 337.52 75.1211 l 338 72.2383 339.078 69.7188
340.641 67.8008 c 342.199 65.7617 344.238 64.8008 346.879 64.8008 c f Q 1588.78 191.313 m 1671.11 191.313
1750.09 191.313 1825.71 191.313 c S 1789.71 209.313 m 1825.71 191.313 l 1789.71 173.313 l 1789.71 209.313
l h B q 1675.18 39.0859 65 79 re W n 1676.12 39.8398 m 1675.4 39.8398 1675.04 40.4414 1675.04 41.3984 c
1675.28 43.1992 1675.88 44.0391 1676.6 44.0391 c 1679 44.0391 1680.68 44.2813 1681.64 44.7617 c 1682.6
45.2383 1683.44 46.3203 1684.04 48.1211 c 1697.6 102 l 1698.08 103.441 1698.44 105.359 1698.44 107.762 c
1698.44 111.121 1697.24 112.922 1694.96 112.922 c 1692.44 112.922 1690.52 111.359 1689.32 108.359 c
1688.12 105.359 1686.92 101.641 1685.84 97.0781 c 1685.84 96.6016 1685.48 96.3594 1684.88 96.3594 c
1683.44 96.3594 l 1683.2 96.3594 1683.08 96.4805 1682.84 96.7188 c 1682.6 96.9609 1682.48 97.3203 1682.48
97.5586 c 1683.32 101.039 1684.16 104.039 1685 106.441 c 1685.72 108.84 1687.04 111 1688.72 113.039 c
1690.4 114.961 1692.44 116.039 1695.08 116.039 c 1697.72 116.039 1700.12 115.199 1702.28 113.641 c
1704.32 112.078 1705.64 110.039 1706.24 107.398 c 1708.4 109.922 1710.8 111.961 1713.44 113.641 c 1716.08
115.199 1718.84 116.039 1721.72 116.039 c 1725.08 116.039 1727.96 115.078 1730.48 113.281 c 1732.88
111.48 1734.8 109.078 1736 106.199 c 1737.2 103.32 1737.92 100.199 1737.92 96.8398 c 1737.92 91.6797
1736.6 86.3984 1733.96 80.8789 c 1731.32 75.3594 1727.72 70.8008 1723.28 67.1992 c 1718.72 63.6016
1713.92 61.8008 1708.76 61.8008 c 1706.36 61.8008 1704.32 62.5195 1702.4 63.7188 c 1700.48 64.9219
1698.92 66.7188 1697.96 68.8789 c 1692.56 47.6406 l 1692.32 46.4414 1692.2 45.8398 1692.2 45.6016 c
1692.2 44.6406 1694.84 44.0391 1700.12 44.0391 c 1700.84 44.0391 1701.32 43.5586 1701.32 42.4805 c
1701.08 41.5195 1700.84 40.8008 1700.6 40.4414 c 1700.36 40.0781 1700 39.8398 1699.28 39.8398 c 1676.12
39.8398 l h 1708.88 64.8008 m 1711.64 64.8008 1714.16 65.8789 1716.56 68.1602 c 1718.96 70.4414 1720.88
73.0781 1722.44 75.9609 c 1723.4 78.1211 1724.48 80.6406 1725.44 83.7617 c 1726.4 86.8789 1727.24 90.1211
1727.96 93.7188 c 1728.68 97.1992 1729.04 99.9609 1729.04 102.121 c 1729.04 103.922 1728.8 105.602
1728.32 107.281 c 1727.84 108.84 1727 110.16 1725.92 111.238 c 1724.72 112.32 1723.28 112.922 1721.48
112.922 c 1718.48 112.922 1715.72 111.84 1713.08 109.801 c 1710.44 107.641 1708.16 105.121 1706.12 102 c
1706.12 101.641 l 1699.52 75.1211 l 1700 72.2383 1701.08 69.7188 1702.64 67.8008 c 1704.2 65.7617 1706.24
64.8008 1708.88 64.8008 c f Q q 1009.1 138.543 49 89 re W n 1010.88 149.879 m 1010.52 149.879 1010.16 150
1009.8 150.48 c 1009.32 150.961 1009.2 151.32 1009.2 151.801 c 1009.2 152.398 1012.44 155.762 1019.16
161.879 c 1019.16 162.121 1019.04 162.48 1019.04 162.719 c 1019.04 162.961 1019.04 163.199 1019.04
163.441 c 1019.04 166.801 1019.28 170.281 1019.76 173.879 c 1020.24 177.48 1021.08 180.961 1022.04
184.559 c 1023 188.16 1024.2 191.879 1025.64 195.719 c 1027.08 199.559 1028.64 203.039 1030.08 205.922 c
1032.6 211.199 1035.36 215.762 1038.48 219.719 c 1041.48 223.559 1044.84 225.602 1048.68 225.602 c 1051.2
225.602 1053 224.641 1054.08 222.961 c 1055.16 221.281 1055.76 219.121 1055.76 216.48 c 1055.76 201.84
1046.04 184.559 1026.6 164.762 c 1026.24 161.281 1026.12 158.879 1026.12 157.441 c 1026.12 147.359 1028.4
142.441 1032.96 142.441 c 1034.88 142.441 1036.8 142.922 1038.72 144 c 1040.64 144.961 1042.32 146.16
1043.76 147.359 c 1045.08 148.559 1047 150.359 1049.52 152.762 c 1049.64 152.879 1049.88 153 1050.24 153
c 1050.6 153 1050.96 152.762 1051.44 152.281 c 1051.8 151.801 1052.04 151.32 1052.04 150.961 c 1052.04
150.719 1051.8 150.359 1051.56 150.121 c 1047.96 146.641 1044.72 144 1041.84 142.078 c 1038.96 140.281
1035.96 139.32 1032.72 139.32 c 1028.4 139.32 1025.28 141.121 1023.24 144.719 c 1021.2 148.32 1019.88
152.52 1019.4 157.559 c 1017.24 155.398 1014.6 152.879 1011.6 150.121 c 1011.48 149.879 1011.24 149.879
1010.88 149.879 c h 1027.56 170.52 m 1031.76 175.078 1035.72 179.762 1039.44 184.801 c 1043.16 189.84
1046.28 195.121 1048.8 200.641 c 1051.2 206.16 1052.52 211.559 1052.52 216.84 c 1052.52 220.559 1051.2
222.48 1048.56 222.48 c 1045.68 222.48 1042.68 219 1039.8 212.281 c 1036.8 205.441 1034.28 198.121
1032.12 190.199 c 1029.96 182.281 1028.52 175.68 1027.56 170.52 c f Q 919.152 841.781 m 995.109 841.781
1067.71 841.781 1136.96 841.781 c S 1100.96 859.781 m 1136.96 841.781 l 1100.96 823.781 l 1100.96 859.781
l h B q 904.098 900.23 65 79 re W n 905.121 900.84 m 904.398 900.84 904.039 901.441 904.039 902.398 c
904.281 904.199 904.879 905.039 905.602 905.039 c 908 905.039 909.68 905.281 910.641 905.762 c 911.602
906.238 912.441 907.32 913.039 909.121 c 926.602 963 l 927.078 964.441 927.441 966.359 927.441 968.762 c
927.441 972.121 926.238 973.922 923.961 973.922 c 921.441 973.922 919.52 972.359 918.32 969.359 c 917.121
966.359 915.922 962.641 914.84 958.078 c 914.84 957.602 914.48 957.359 913.879 957.359 c 912.441 957.359
l 912.199 957.359 912.078 957.48 911.84 957.719 c 911.602 957.961 911.48 958.32 911.48 958.559 c 912.32
962.039 913.16 965.039 914 967.441 c 914.719 969.84 916.039 972 917.719 974.039 c 919.398 975.961 921.441
977.039 924.078 977.039 c 926.719 977.039 929.121 976.199 931.281 974.641 c 933.32 973.078 934.641
971.039 935.238 968.398 c 937.398 970.922 939.801 972.961 942.441 974.641 c 945.078 976.199 947.84
977.039 950.719 977.039 c 954.078 977.039 956.961 976.078 959.48 974.281 c 961.879 972.48 963.801 970.078
965 967.199 c 966.199 964.32 966.922 961.199 966.922 957.84 c 966.922 952.68 965.602 947.398 962.961
941.879 c 960.32 936.359 956.719 931.801 952.281 928.199 c 947.719 924.602 942.922 922.801 937.762
922.801 c 935.359 922.801 933.32 923.52 931.398 924.719 c 929.48 925.922 927.922 927.719 926.961 929.879
c 921.559 908.641 l 921.32 907.441 921.199 906.84 921.199 906.602 c 921.199 905.641 923.84 905.039
929.121 905.039 c 929.84 905.039 930.32 904.559 930.32 903.48 c 930.078 902.52 929.84 901.801 929.602
901.441 c 929.359 901.078 929 900.84 928.281 900.84 c 905.121 900.84 l h 937.879 925.801 m 940.641
925.801 943.16 926.879 945.559 929.16 c 947.961 931.441 949.879 934.078 951.441 936.961 c 952.398 939.121
953.48 941.641 954.441 944.762 c 955.398 947.879 956.238 951.121 956.961 954.719 c 957.68 958.199 958.039
960.961 958.039 963.121 c 958.039 964.922 957.801 966.602 957.32 968.281 c 956.84 969.84 956 971.16
954.922 972.238 c 953.719 973.32 952.281 973.922 950.48 973.922 c 947.48 973.922 944.719 972.84 942.078
970.801 c 939.441 968.641 937.16 966.121 935.121 963 c 935.121 962.641 l 928.52 936.121 l 929 933.238
930.078 930.719 931.641 928.801 c 933.199 926.762 935.238 925.801 937.879 925.801 c f Q q 998.52 951.23
76 7 re W n 1002.24 951.602 m 1001.52 951.602 1001.04 951.84 1000.56 952.32 c 1000.08 952.801 999.961
953.398 999.961 954 c 999.961 954.48 1000.08 955.078 1000.56 955.559 c 1001.04 956.039 1001.52 956.398
1002.24 956.398 c 1071.12 956.398 l 1071.72 956.398 1072.2 956.039 1072.68 955.559 c 1073.04 955.078
1073.28 954.48 1073.28 954 c 1073.28 953.398 1073.04 952.801 1072.68 952.32 c 1072.2 951.84 1071.72
951.602 1071.12 951.602 c 1002.24 951.602 l f Q q 1104.81 922.23 49 89 re W n 1106.88 932.879 m 1106.52
932.879 1106.16 933 1105.8 933.48 c 1105.32 933.961 1105.2 934.32 1105.2 934.801 c 1105.2 935.398 1108.44
938.762 1115.16 944.879 c 1115.16 945.121 1115.04 945.48 1115.04 945.719 c 1115.04 945.961 1115.04
946.199 1115.04 946.441 c 1115.04 949.801 1115.28 953.281 1115.76 956.879 c 1116.24 960.48 1117.08
963.961 1118.04 967.559 c 1119 971.16 1120.2 974.879 1121.64 978.719 c 1123.08 982.559 1124.64 986.039
1126.08 988.922 c 1128.6 994.199 1131.36 998.762 1134.48 1002.72 c 1137.48 1006.56 1140.84 1008.6 1144.68
1008.6 c 1147.2 1008.6 1149 1007.64 1150.08 1005.96 c 1151.16 1004.28 1151.76 1002.12 1151.76 999.48 c
1151.76 984.84 1142.04 967.559 1122.6 947.762 c 1122.24 944.281 1122.12 941.879 1122.12 940.441 c 1122.12
930.359 1124.4 925.441 1128.96 925.441 c 1130.88 925.441 1132.8 925.922 1134.72 927 c 1136.64 927.961
1138.32 929.16 1139.76 930.359 c 1141.08 931.559 1143 933.359 1145.52 935.762 c 1145.64 935.879 1145.88
936 1146.24 936 c 1146.6 936 1146.96 935.762 1147.44 935.281 c 1147.8 934.801 1148.04 934.32 1148.04
933.961 c 1148.04 933.719 1147.8 933.359 1147.56 933.121 c 1143.96 929.641 1140.72 927 1137.84 925.078 c
1134.96 923.281 1131.96 922.32 1128.72 922.32 c 1124.4 922.32 1121.28 924.121 1119.24 927.719 c 1117.2
931.32 1115.88 935.52 1115.4 940.559 c 1113.24 938.398 1110.6 935.879 1107.6 933.121 c 1107.48 932.879
1107.24 932.879 1106.88 932.879 c h 1123.56 953.52 m 1127.76 958.078 1131.72 962.762 1135.44 967.801 c
1139.16 972.84 1142.28 978.121 1144.8 983.641 c 1147.2 989.16 1148.52 994.559 1148.52 999.84 c 1148.52
1003.56 1147.2 1005.48 1144.56 1005.48 c 1141.68 1005.48 1138.68 1002 1135.8 995.281 c 1132.8 988.441
1130.28 981.121 1128.12 973.199 c 1125.96 965.281 1124.52 958.68 1123.56 953.52 c f Q Q
}

The one-loop graph of $-i\Sigma(p)$ for a fermion $f$: the internal fermion carries $\ell$ along its arrow and the photon carries $p-\ell$.

Let $-i\Sigma_2(p)$ be the $O(e^2)$ part of $-i\Sigma(p)$. Starting at the vertex $\nu$ and reading the line against the arrow with \eqref{eq:7.2},
\begin{equation*}
\begin{aligned}
-i\Sigma_2(p)
=\int\frac{d^4\ell}{(2\pi)^4}\,
&(ieQ_f\gamma^\nu)\,\frac{i(\not\ell+m_f)}{\ell^2-m_f^2+i\varepsilon}\,(ieQ_f\gamma^\mu)\\[6pt]
&\times\frac{-i}{(p-\ell)^2+i\varepsilon}\left[\eta_{\mu\nu}-(1-\xi)\frac{(p-\ell)_\mu(p-\ell)_\nu}{(p-\ell)^2+i\varepsilon}\right].
\end{aligned}
\tag{7.7}\label{eq:7.7}
\end{equation*}
For a quark, \eqref{eq:7.7} multiplies the identity in color.

In $d=4-\epsilon$ dimensions, replace $d^4\ell/(2\pi)^4$ by $d^d\ell/(2\pi)^d$, let the coupling carry the factor of $\mu$, $e\to\mu^{\epsilon/2}e$, and continue the $\gamma$ matrices to $d$ dimensions with $\{\gamma^\mu,\gamma^\nu\}=2\eta^{\mu\nu}I_4$ and $\eta^\mu{}_\mu=d$. Collecting $(ieQ_f)^2\,i\,(-i)=-e^2Q_f^2$ and using $\gamma^\nu\gamma_\nu=d\,I_4$, $\gamma^\nu\slashed{a}\gamma_\nu=-(d-2)\slashed{a}$ and, for $\ell=p-k$, $\slashed{k}(\not\ell+m_f)\slashed{k}=2(k\cdot p)\slashed{k}+k^2(m_f-\slashed{p}-\slashed{k})$, where $k^2$ cancels one factor $\mathcal P_\gamma$,
\begin{equation*}
\begin{aligned}
-i\Sigma_2(p)
={}&-\mu^{\epsilon}e^2Q_f^2\int\frac{d^d\ell}{(2\pi)^d}\left[
\frac{\gamma^\nu(\not\ell+m_f)\gamma_\nu}{\mathcal P_1\mathcal P_\gamma}
-(1-\xi)\,\frac{\slashed{k}(\not\ell+m_f)\slashed{k}}{\mathcal P_1\mathcal P_\gamma^2}\right]\\[6pt]
={}&-\mu^{\epsilon}e^2Q_f^2\int\frac{d^d\ell}{(2\pi)^d}\biggl\{
\frac{d\,m_f-(d-2)\not\ell}{\mathcal P_1\mathcal P_\gamma}
-(1-\xi)\biggl[\frac{m_f-\slashed{p}-\slashed{k}}{\mathcal P_1\mathcal P_\gamma}
+\frac{2(k\cdot p)\slashed{k}}{\mathcal P_1\mathcal P_\gamma^2}\biggr]\biggr\},\\[6pt]
\mathcal P_1:={}&\ell^2-m_f^2+i\varepsilon,\qquad
\mathcal P_\gamma:=k^2+i\varepsilon,\qquad
k:=p-\ell .
\end{aligned}
\tag{7.8}\label{eq:7.8}
\end{equation*}
Combine the denominators with a Feynman parameter $x$:
\begin{equation*}
\begin{aligned}
\frac{1}{\mathcal P_1\mathcal P_\gamma}
&=\int_0^1\frac{dx}{\bigl[x\mathcal P_\gamma+(1-x)\mathcal P_1\bigr]^2},\qquad
\frac{1}{\mathcal P_1\mathcal P_\gamma^2}
=\int_0^1\frac{2x\,dx}{\bigl[x\mathcal P_\gamma+(1-x)\mathcal P_1\bigr]^3},\\[6pt]
x\mathcal P_\gamma+(1-x)\mathcal P_1
&=(\ell-xp)^2-\Delta_{\Sigma,f}+i\varepsilon,\qquad
\Delta_{\Sigma,f}:=(1-x)\bigl(m_f^2-xp^2\bigr).
\end{aligned}
\tag{7.9}\label{eq:7.9}
\end{equation*}
Shift $\ell\to\ell+xp$, so that $k\to(1-x)p-\ell$, drop the terms odd in $\ell$, and replace $\ell^\mu\ell^\nu$ by $\eta^{\mu\nu}\ell^2/d$. The numerators of \eqref{eq:7.8} become
\begin{equation*}
\begin{aligned}
d\,m_f-(d-2)\not\ell\ &\to\ d\,m_f-(d-2)x\slashed{p},\\[4pt]
m_f-\slashed{p}-\slashed{k}\ &\to\ m_f-(2-x)\slashed{p},\\[4pt]
2(k\cdot p)\slashed{k}\ &\to\ 2(1-x)^2p^2\slashed{p}+2(p\cdot\ell)\not\ell
\ \to\ 2(1-x)^2p^2\slashed{p}+\frac2d\,\ell^2\slashed{p},
\end{aligned}
\tag{7.10}\label{eq:7.10}
\end{equation*}
so that
\begin{equation*}
\begin{aligned}
-i\Sigma_2(p)
={}&-\mu^{\epsilon}e^2Q_f^2\int_0^1dx\int\frac{d^d\ell}{(2\pi)^d}\,\biggl\{
\frac{d\,m_f-(d-2)x\slashed{p}-(1-\xi)\bigl[m_f-(2-x)\slashed{p}\bigr]}{\bigl(\ell^2-\Delta_{\Sigma,f}+i\varepsilon\bigr)^2}\\[6pt]
&-(1-\xi)\,\frac{2x\bigl[2(1-x)^2p^2+\frac2d\,\ell^2\bigr]\slashed{p}}{\bigl(\ell^2-\Delta_{\Sigma,f}+i\varepsilon\bigr)^3}\biggr\}.
\end{aligned}
\tag{7.11}\label{eq:7.11}
\end{equation*}
After Wick rotation, for $\Delta>0$ and continued to $\Delta<0$ through $\Delta-i\varepsilon$, with $n=2,3$,
\begin{equation*}
\begin{aligned}
\int\frac{d^d\ell}{(2\pi)^d}\,\frac{1}{(\ell^2-\Delta+i\varepsilon)^n}
&=\frac{(-1)^n\,i}{(4\pi)^{d/2}}\,\frac{\Gamma\bigl(n-\frac d2\bigr)}{\Gamma(n)}\,(\Delta-i\varepsilon)^{\frac d2-n},\\[6pt]
\int\frac{d^d\ell}{(2\pi)^d}\,\frac{\ell^2}{(\ell^2-\Delta+i\varepsilon)^n}
&=\frac{(-1)^{n-1}\,i}{(4\pi)^{d/2}}\,\frac d2\,\frac{\Gamma\bigl(n-1-\frac d2\bigr)}{\Gamma(n)}\,(\Delta-i\varepsilon)^{\frac d2-n+1}.
\end{aligned}
\tag{7.12}\label{eq:7.12}
\end{equation*}
Applying \eqref{eq:7.12} to \eqref{eq:7.11},
\begin{equation*}
\begin{aligned}
\Sigma_2(p)
={}&\frac{\mu^{\epsilon}e^2Q_f^2}{(4\pi)^{d/2}}\int_0^1dx\,\biggl\{
\Bigl[d\,m_f-(d-2)x\slashed{p}-(1-\xi)\bigl(m_f-2(1-x)\slashed{p}\bigr)\Bigr]\\[6pt]
&\times\Gamma\Bigl(2-\frac d2\Bigr)(\Delta_{\Sigma,f}-i\varepsilon)^{\frac d2-2}\\[6pt]
&+2(1-\xi)\,x(1-x)^2p^2\slashed{p}\;
\Gamma\Bigl(3-\frac d2\Bigr)(\Delta_{\Sigma,f}-i\varepsilon)^{\frac d2-3}\biggr\}.
\end{aligned}
\tag{7.13}\label{eq:7.13}
\end{equation*}
With $d=4-\epsilon$, $\Gamma(\epsilon/2)=2/\epsilon-\gamma_E+O(\epsilon)$ and $2/\bar\epsilon:=2/\epsilon-\gamma_E+\ln4\pi$,
\begin{equation*}
{\displaystyle 
\begin{aligned}
\Sigma_2(p)={}&\frac{e^2Q_f^2}{16\pi^2}\int_0^1dx\,\biggl\{
\Bigl[4m_f-2x\slashed{p}-(1-\xi)\bigl(m_f-2(1-x)\slashed{p}\bigr)\Bigr]
\left[\frac2{\bar\epsilon}-\ln\frac{\Delta_{\Sigma,f}-i\varepsilon}{\mu^2}\right]\\[6pt]
&-2m_f+2x\slashed{p}+\frac{2(1-\xi)\,x(1-x)\,p^2\slashed{p}}{m_f^2-xp^2-i\varepsilon}\biggr\}+O(\epsilon)\\[6pt]
={}&\frac{e^2Q_f^2}{16\pi^2}\,\frac2\epsilon\,\bigl[(3+\xi)m_f-\xi\slashed{p}\bigr]+(\text{finite as }\epsilon\to0).
\end{aligned}
}
\tag{7.14}\label{eq:7.14}
\end{equation*}

\topic{(2) Photon Propagator}
\begin{equation*}
E_\gamma=2,\qquad E_F=0\qquad\Longrightarrow\qquad\omega=2 .
\tag{7.15}\label{eq:7.15}
\end{equation*}

% photon-self-energy
\qftdiagram{191.31}{118.61}{%
q 0.1 0 0 0.1 0 0 cm 1 g 0 0 1913.14 1186.15 re f 0 g 998.66 937.441 m 906.828 971.875 l 906.828 903.004
l f 914.48 248.707 m 1006.31 214.273 l 1006.31 283.145 l f 14 w 2 J 1 j 0 G 114.789 593.074 m 138.078
614.313 l 151.387 623.621 l 163.031 628.992 l 173.016 631.148 l 182.996 630.895 l 192.977 628.25 l
202.961 623.379 l 214.605 615.305 l 231.242 600.664 l 266.176 568.926 l 279.484 560.488 l 289.469 556.516
l 299.449 554.855 l 309.43 555.613 l 319.41 558.738 l 331.059 565.109 l 344.367 575.293 l 364.328 593.879
l 387.621 614.977 l 400.93 624.102 l 412.574 629.262 l 422.555 631.219 l 432.535 630.762 l 442.52 627.926
l 452.5 622.883 l 464.145 614.645 l 480.781 599.871 l 514.055 569.555 l 527.363 560.914 l 537.344 556.758
l 547.324 554.902 l 557.309 555.457 l 567.289 558.391 l 577.27 563.52 l 590.578 573.195 l 608.879 589.863
l 612.207 593.074 l 612.207 593.074 l S 1300.94 593.074 m 1324.23 614.313 l 1337.54 623.621 l 1349.18
628.992 l 1359.16 631.148 l 1369.14 630.895 l 1379.13 628.25 l 1389.11 623.379 l 1400.75 615.305 l
1417.39 600.664 l 1452.32 568.926 l 1465.63 560.488 l 1475.62 556.516 l 1485.6 554.855 l 1495.58 555.613
l 1505.56 558.738 l 1517.21 565.109 l 1530.52 575.293 l 1550.48 593.879 l 1573.77 614.977 l 1587.08
624.102 l 1598.72 629.262 l 1608.7 631.219 l 1618.68 630.762 l 1628.67 627.926 l 1638.65 622.883 l
1650.29 614.645 l 1666.93 599.871 l 1700.2 569.555 l 1713.51 560.914 l 1723.49 556.758 l 1733.47 554.902
l 1743.46 555.457 l 1753.44 558.391 l 1763.42 563.52 l 1776.73 573.195 l 1795.03 589.863 l 1798.36
593.074 l 1798.36 593.074 l S 1300.94 593.074 m 1299.87 620.16 l 1296.68 647.078 l 1291.38 673.664 l 1284
699.746 l 1274.59 725.172 l 1263.22 749.773 l 1249.94 773.41 l 1234.85 795.926 l 1218.03 817.188 l
1199.59 837.055 l 1179.65 855.414 l 1158.32 872.148 l 1135.75 887.152 l 1112.06 900.336 l 1087.41 911.613
l 1061.95 920.918 l 1035.84 928.191 l 1009.24 933.391 l 982.305 936.477 l 955.215 937.438 l 928.133
936.266 l 901.227 932.965 l 874.664 927.559 l 848.609 920.078 l 823.223 910.574 l 798.664 899.102 l
775.082 885.734 l 752.625 870.551 l 731.434 853.652 l 711.633 835.137 l 693.355 815.121 l 676.703 793.73
l 661.789 771.094 l 648.699 747.355 l 637.52 722.664 l 628.313 697.164 l 621.145 671.023 l 616.051
644.398 l 613.07 617.457 l 612.215 590.363 l 613.496 563.285 l 616.902 536.395 l 622.414 509.852 l
629.996 483.828 l 639.598 458.477 l 651.168 433.965 l 664.629 410.434 l 679.898 388.039 l 696.883 366.91
l 715.477 347.188 l 735.563 328.984 l 757.02 312.418 l 779.715 297.594 l 803.504 284.598 l 828.242
273.516 l 853.773 264.41 l 879.941 257.344 l 906.586 252.355 l 933.543 249.48 l 960.637 248.734 l 987.711
250.121 l 1014.59 253.633 l 1041.11 259.246 l 1067.11 266.93 l 1092.41 276.633 l 1116.88 288.301 l
1140.36 301.852 l 1162.7 317.211 l 1183.75 334.277 l 1203.41 352.949 l 1221.53 373.109 l 1238.01 394.629
l 1252.75 417.383 l 1265.65 441.223 l 1276.63 466.004 l 1285.64 491.57 l 1292.6 517.77 l 1297.48 544.434
l 1300.25 571.398 l 1300.94 593.074 l 1300.94 593.074 l S 10 w 0 J 612.207 570.574 m 618.172 570.574
623.895 572.945 628.117 577.164 c 632.336 581.383 634.707 587.105 634.707 593.074 c 634.707 599.043
632.336 604.766 628.117 608.984 c 623.895 613.203 618.172 615.574 612.207 615.574 c 606.238 615.574
600.516 613.203 596.297 608.984 c 592.078 604.766 589.707 599.043 589.707 593.074 c 589.707 587.105
592.078 581.383 596.297 577.164 c 600.516 572.945 606.238 570.574 612.207 570.574 c h B 1300.94 570.574 m
1306.9 570.574 1312.63 572.945 1316.85 577.164 c 1321.07 581.383 1323.44 587.105 1323.44 593.074 c
1323.44 599.043 1321.07 604.766 1316.85 608.984 c 1312.63 613.203 1306.9 615.574 1300.94 615.574 c
1294.97 615.574 1289.25 613.203 1285.03 608.984 c 1280.81 604.766 1278.44 599.043 1278.44 593.074 c
1278.44 587.105 1280.81 581.383 1285.03 577.164 c 1289.25 572.945 1294.97 570.574 1300.94 570.574 c h B q
479.547 649.473 74 88 re W n 480.641 652.691 m 480.641 653.211 480.641 653.602 480.77 653.73 c 499.879
730.301 l 500.141 731.469 500.789 732.512 501.828 733.289 c 502.871 734.07 504.039 734.461 505.34 734.461
c 506.379 734.461 507.289 734.07 508.07 733.422 c 508.852 732.77 509.238 731.859 509.238 730.82 c 509.238
730.559 509.238 730.301 509.238 730.172 c 509.238 730.039 509.109 729.91 509.109 729.781 c 501.441
699.359 l 500.531 695.719 500.141 692.602 500.141 690 c 500.141 686.75 500.789 684.148 502.352 682.07 c
503.91 679.988 506.25 678.949 509.371 678.949 c 515.609 678.949 521.199 682.852 526.012 690.781 c 526.012
690.91 l 526.141 691.039 526.141 691.172 v 535.5 728.871 l 535.762 729.91 536.539 730.949 537.578 731.73
c 538.621 732.512 539.66 733.031 540.961 733.031 c 542 733.031 542.91 732.641 543.691 731.988 c 544.469
731.34 544.859 730.43 544.859 729.391 c 544.859 728.871 544.73 728.48 544.73 728.352 c 535.371 690.91 l
534.719 688.441 534.461 686.359 534.461 684.539 c 534.461 680.77 535.629 678.949 538.23 678.949 c 540.828
678.949 542.91 680.512 544.34 683.891 c 545.641 687.27 546.941 691.301 547.98 696.109 c 548.109 696.629
548.5 696.891 549.02 696.891 c 550.578 696.891 l 550.84 696.891 551.102 696.629 551.359 696.371 c 551.621
696.109 551.75 695.852 551.75 695.59 c 550.191 689.48 548.629 684.672 546.809 681.031 c 544.988 677.391
542 675.699 537.969 675.699 c 535.109 675.699 532.512 676.48 530.301 678.039 c 528.09 679.602 526.66
681.809 526.012 684.672 c 523.801 681.941 521.199 679.73 518.34 678.172 c 515.48 676.609 512.359 675.699
509.238 675.699 c 503.781 675.699 499.621 677 496.5 679.34 c 489.871 653.211 l 489.609 652.039 488.961
651 487.922 650.219 c 486.879 649.441 485.84 649.051 484.539 649.051 c 483.5 649.051 482.59 649.441
481.809 650.09 c 481.031 650.738 480.641 651.648 480.641 652.691 c f Q q 1367.59 674.504 64 60 re W n
1370.19 675 m 1369.15 675 1368.63 675.52 1368.63 676.691 c 1379.68 720.891 l 1379.94 722.32 1380.2
723.488 1380.33 724.398 c 1380.33 725.699 1377.47 726.48 1371.88 726.48 c 1370.97 726.48 1370.58 727
1370.58 728.172 c 1370.58 728.43 1370.58 728.82 1370.84 729.34 c 1370.97 729.859 1371.1 730.25 1371.36
730.512 c 1371.62 730.77 1371.88 731.031 1372.4 731.031 c 1389.95 732.461 l 1390.99 732.461 1391.51
731.809 1391.51 730.77 c 1378.77 679.809 l 1385.53 682.02 1391.9 685.531 1397.62 690.469 c 1403.34 695.41
1408.28 701.129 1412.31 707.762 c 1416.34 714.262 1419.2 721.148 1421.02 728.301 c 1421.28 729.469
1421.93 730.512 1422.97 731.289 c 1424.01 732.07 1425.18 732.461 1426.48 732.461 c 1427.52 732.461
1428.43 732.07 1429.21 731.422 c 1429.99 730.641 1430.38 729.73 1430.38 728.691 c 1430.38 727.262 1429.73
725.18 1428.69 722.32 c 1427.52 719.461 1426.35 716.859 1425.31 714.648 c 1421.8 707.891 1417.25 701.648
1411.66 695.93 c 1408.02 692.289 1404.25 689.172 1400.22 686.441 c 1396.06 683.711 1391.77 681.371
1387.35 679.422 c 1382.93 677.469 1378.25 676.039 1373.31 675 c 1370.19 675 l f Q 9 w 1 J 0.449951 G
249.555 470.633 m 325.508 470.633 398.109 470.633 467.355 470.633 c S 0.449951 g 431.355 488.633 m
467.355 470.633 l 431.355 452.633 l 431.355 488.633 l h B q 324.496 318.402 65 79 re W n 325.121 318.84 m
324.398 318.84 324.039 319.441 324.039 320.398 c 324.281 322.199 324.879 323.039 325.602 323.039 c 328
323.039 329.68 323.281 330.641 323.762 c 331.602 324.238 332.441 325.32 333.039 327.121 c 346.602 381 l
347.078 382.441 347.441 384.359 347.441 386.762 c 347.441 390.121 346.238 391.922 343.961 391.922 c
341.441 391.922 339.52 390.359 338.32 387.359 c 337.121 384.359 335.922 380.641 334.84 376.078 c 334.84
375.602 334.48 375.359 333.879 375.359 c 332.441 375.359 l 332.199 375.359 332.078 375.48 331.84 375.719
c 331.602 375.961 331.48 376.32 331.48 376.559 c 332.32 380.039 333.16 383.039 334 385.441 c 334.719
387.84 336.039 390 337.719 392.039 c 339.398 393.961 341.441 395.039 344.078 395.039 c 346.719 395.039
349.121 394.199 351.281 392.641 c 353.32 391.078 354.641 389.039 355.238 386.398 c 357.398 388.922
359.801 390.961 362.441 392.641 c 365.078 394.199 367.84 395.039 370.719 395.039 c 374.078 395.039
376.961 394.078 379.48 392.281 c 381.879 390.48 383.801 388.078 385 385.199 c 386.199 382.32 386.922
379.199 386.922 375.84 c 386.922 370.68 385.602 365.398 382.961 359.879 c 380.32 354.359 376.719 349.801
372.281 346.199 c 367.719 342.602 362.922 340.801 357.762 340.801 c 355.359 340.801 353.32 341.52 351.398
342.719 c 349.48 343.922 347.922 345.719 346.961 347.879 c 341.559 326.641 l 341.32 325.441 341.199
324.84 341.199 324.602 c 341.199 323.641 343.84 323.039 349.121 323.039 c 349.84 323.039 350.32 322.559
350.32 321.48 c 350.078 320.52 349.84 319.801 349.602 319.441 c 349.359 319.078 349 318.84 348.281 318.84
c 325.121 318.84 l h 357.879 343.801 m 360.641 343.801 363.16 344.879 365.559 347.16 c 367.961 349.441
369.879 352.078 371.441 354.961 c 372.398 357.121 373.48 359.641 374.441 362.762 c 375.398 365.879
376.238 369.121 376.961 372.719 c 377.68 376.199 378.039 378.961 378.039 381.121 c 378.039 382.922
377.801 384.602 377.32 386.281 c 376.84 387.84 376 389.16 374.922 390.238 c 373.719 391.32 372.281
391.922 370.48 391.922 c 367.48 391.922 364.719 390.84 362.078 388.801 c 359.441 386.641 357.16 384.121
355.121 381 c 355.121 380.641 l 348.52 354.121 l 349 351.238 350.078 348.719 351.641 346.801 c 353.199
344.762 355.238 343.801 357.879 343.801 c f Q 1435.7 470.633 m 1511.66 470.633 1584.26 470.633 1653.5
470.633 c S 1617.5 488.633 m 1653.5 470.633 l 1617.5 452.633 l 1617.5 488.633 l h B q 1510.64 318.402 65
79 re W n 1512.12 318.84 m 1511.4 318.84 1511.04 319.441 1511.04 320.398 c 1511.28 322.199 1511.88
323.039 1512.6 323.039 c 1515 323.039 1516.68 323.281 1517.64 323.762 c 1518.6 324.238 1519.44 325.32
1520.04 327.121 c 1533.6 381 l 1534.08 382.441 1534.44 384.359 1534.44 386.762 c 1534.44 390.121 1533.24
391.922 1530.96 391.922 c 1528.44 391.922 1526.52 390.359 1525.32 387.359 c 1524.12 384.359 1522.92
380.641 1521.84 376.078 c 1521.84 375.602 1521.48 375.359 1520.88 375.359 c 1519.44 375.359 l 1519.2
375.359 1519.08 375.48 1518.84 375.719 c 1518.6 375.961 1518.48 376.32 1518.48 376.559 c 1519.32 380.039
1520.16 383.039 1521 385.441 c 1521.72 387.84 1523.04 390 1524.72 392.039 c 1526.4 393.961 1528.44
395.039 1531.08 395.039 c 1533.72 395.039 1536.12 394.199 1538.28 392.641 c 1540.32 391.078 1541.64
389.039 1542.24 386.398 c 1544.4 388.922 1546.8 390.961 1549.44 392.641 c 1552.08 394.199 1554.84 395.039
1557.72 395.039 c 1561.08 395.039 1563.96 394.078 1566.48 392.281 c 1568.88 390.48 1570.8 388.078 1572
385.199 c 1573.2 382.32 1573.92 379.199 1573.92 375.84 c 1573.92 370.68 1572.6 365.398 1569.96 359.879 c
1567.32 354.359 1563.72 349.801 1559.28 346.199 c 1554.72 342.602 1549.92 340.801 1544.76 340.801 c
1542.36 340.801 1540.32 341.52 1538.4 342.719 c 1536.48 343.922 1534.92 345.719 1533.96 347.879 c 1528.56
326.641 l 1528.32 325.441 1528.2 324.84 1528.2 324.602 c 1528.2 323.641 1530.84 323.039 1536.12 323.039 c
1536.84 323.039 1537.32 322.559 1537.32 321.48 c 1537.08 320.52 1536.84 319.801 1536.6 319.441 c 1536.36
319.078 1536 318.84 1535.28 318.84 c 1512.12 318.84 l h 1544.88 343.801 m 1547.64 343.801 1550.16 344.879
1552.56 347.16 c 1554.96 349.441 1556.88 352.078 1558.44 354.961 c 1559.4 357.121 1560.48 359.641 1561.44
362.762 c 1562.4 365.879 1563.24 369.121 1563.96 372.719 c 1564.68 376.199 1565.04 378.961 1565.04
381.121 c 1565.04 382.922 1564.8 384.602 1564.32 386.281 c 1563.84 387.84 1563 389.16 1561.92 390.238 c
1560.72 391.32 1559.28 391.922 1557.48 391.922 c 1554.48 391.922 1551.72 390.84 1549.08 388.801 c 1546.44
386.641 1544.16 384.121 1542.12 381 c 1542.12 380.641 l 1535.52 354.121 l 1536 351.238 1537.08 348.719
1538.64 346.801 c 1540.2 344.762 1542.24 343.801 1544.88 343.801 c f Q q 832.57 1033.2 49 89 re W n
834.879 1043.88 m 834.52 1043.88 834.16 1044 833.801 1044.48 c 833.32 1044.96 833.199 1045.32 833.199
1045.8 c 833.199 1046.4 836.441 1049.76 843.16 1055.88 c 843.16 1056.12 843.039 1056.48 843.039 1056.72 c
843.039 1056.96 843.039 1057.2 843.039 1057.44 c 843.039 1060.8 843.281 1064.28 843.762 1067.88 c 844.238
1071.48 845.078 1074.96 846.039 1078.56 c 847 1082.16 848.199 1085.88 849.641 1089.72 c 851.078 1093.56
852.641 1097.04 854.078 1099.92 c 856.602 1105.2 859.359 1109.76 862.48 1113.72 c 865.48 1117.56 868.84
1119.6 872.68 1119.6 c 875.199 1119.6 877 1118.64 878.078 1116.96 c 879.16 1115.28 879.762 1113.12
879.762 1110.48 c 879.762 1095.84 870.039 1078.56 850.602 1058.76 c 850.238 1055.28 850.121 1052.88
850.121 1051.44 c 850.121 1041.36 852.398 1036.44 856.961 1036.44 c 858.879 1036.44 860.801 1036.92
862.719 1038 c 864.641 1038.96 866.32 1040.16 867.762 1041.36 c 869.078 1042.56 871 1044.36 873.52
1046.76 c 873.641 1046.88 873.879 1047 874.238 1047 c 874.602 1047 874.961 1046.76 875.441 1046.28 c
875.801 1045.8 876.039 1045.32 876.039 1044.96 c 876.039 1044.72 875.801 1044.36 875.559 1044.12 c
871.961 1040.64 868.719 1038 865.84 1036.08 c 862.961 1034.28 859.961 1033.32 856.719 1033.32 c 852.398
1033.32 849.281 1035.12 847.238 1038.72 c 845.199 1042.32 843.879 1046.52 843.398 1051.56 c 841.238
1049.4 838.602 1046.88 835.602 1044.12 c 835.48 1043.88 835.238 1043.88 834.879 1043.88 c h 851.559
1064.52 m 855.762 1069.08 859.719 1073.76 863.441 1078.8 c 867.16 1083.84 870.281 1089.12 872.801 1094.64
c 875.199 1100.16 876.52 1105.56 876.52 1110.84 c 876.52 1114.56 875.199 1116.48 872.559 1116.48 c 869.68
1116.48 866.68 1113 863.801 1106.28 c 860.801 1099.44 858.281 1092.12 856.121 1084.2 c 853.961 1076.28
852.52 1069.68 851.559 1064.52 c f Q q 908.563 1025.2 82 82 re W n 912 1062.6 m 911.281 1062.6 910.801
1062.84 910.32 1063.32 c 909.84 1063.8 909.719 1064.4 909.719 1065 c 909.719 1065.48 909.84 1066.08
910.32 1066.56 c 910.801 1067.04 911.281 1067.4 912 1067.4 c 947.281 1067.4 l 947.281 1102.8 l 947.281
1103.4 947.52 1103.88 948 1104.36 c 948.48 1104.72 948.961 1104.96 949.68 1104.96 c 950.16 1104.96
950.762 1104.72 951.238 1104.36 c 951.719 1103.88 952.078 1103.4 952.078 1102.8 c 952.078 1067.4 l
987.359 1067.4 l 987.84 1067.4 988.32 1067.04 988.801 1066.56 c 989.281 1066.08 989.52 1065.48 989.52
1065 c 989.52 1064.4 989.281 1063.8 988.801 1063.32 c 988.32 1062.84 987.84 1062.6 987.359 1062.6 c
952.078 1062.6 l 952.078 1027.32 l 952.078 1026.72 951.719 1026.24 951.238 1025.76 c 950.762 1025.4
950.16 1025.16 949.68 1025.16 c 948.961 1025.16 948.48 1025.4 948 1025.76 c 947.52 1026.24 947.281
1026.72 947.281 1027.32 c 947.281 1062.6 l 912 1062.6 l f Q q 1012.86 1011.2 65 79 re W n 1014.12 1011.84
m 1013.4 1011.84 1013.04 1012.44 1013.04 1013.4 c 1013.28 1015.2 1013.88 1016.04 1014.6 1016.04 c 1017
1016.04 1018.68 1016.28 1019.64 1016.76 c 1020.6 1017.24 1021.44 1018.32 1022.04 1020.12 c 1035.6 1074 l
1036.08 1075.44 1036.44 1077.36 1036.44 1079.76 c 1036.44 1083.12 1035.24 1084.92 1032.96 1084.92 c
1030.44 1084.92 1028.52 1083.36 1027.32 1080.36 c 1026.12 1077.36 1024.92 1073.64 1023.84 1069.08 c
1023.84 1068.6 1023.48 1068.36 1022.88 1068.36 c 1021.44 1068.36 l 1021.2 1068.36 1021.08 1068.48 1020.84
1068.72 c 1020.6 1068.96 1020.48 1069.32 1020.48 1069.56 c 1021.32 1073.04 1022.16 1076.04 1023 1078.44 c
1023.72 1080.84 1025.04 1083 1026.72 1085.04 c 1028.4 1086.96 1030.44 1088.04 1033.08 1088.04 c 1035.72
1088.04 1038.12 1087.2 1040.28 1085.64 c 1042.32 1084.08 1043.64 1082.04 1044.24 1079.4 c 1046.4 1081.92
1048.8 1083.96 1051.44 1085.64 c 1054.08 1087.2 1056.84 1088.04 1059.72 1088.04 c 1063.08 1088.04 1065.96
1087.08 1068.48 1085.28 c 1070.88 1083.48 1072.8 1081.08 1074 1078.2 c 1075.2 1075.32 1075.92 1072.2
1075.92 1068.84 c 1075.92 1063.68 1074.6 1058.4 1071.96 1052.88 c 1069.32 1047.36 1065.72 1042.8 1061.28
1039.2 c 1056.72 1035.6 1051.92 1033.8 1046.76 1033.8 c 1044.36 1033.8 1042.32 1034.52 1040.4 1035.72 c
1038.48 1036.92 1036.92 1038.72 1035.96 1040.88 c 1030.56 1019.64 l 1030.32 1018.44 1030.2 1017.84 1030.2
1017.6 c 1030.2 1016.64 1032.84 1016.04 1038.12 1016.04 c 1038.84 1016.04 1039.32 1015.56 1039.32 1014.48
c 1039.08 1013.52 1038.84 1012.8 1038.6 1012.44 c 1038.36 1012.08 1038 1011.84 1037.28 1011.84 c 1014.12
1011.84 l h 1046.88 1036.8 m 1049.64 1036.8 1052.16 1037.88 1054.56 1040.16 c 1056.96 1042.44 1058.88
1045.08 1060.44 1047.96 c 1061.4 1050.12 1062.48 1052.64 1063.44 1055.76 c 1064.4 1058.88 1065.24 1062.12
1065.96 1065.72 c 1066.68 1069.2 1067.04 1071.96 1067.04 1074.12 c 1067.04 1075.92 1066.8 1077.6 1066.32
1079.28 c 1065.84 1080.84 1065 1082.16 1063.92 1083.24 c 1062.72 1084.32 1061.28 1084.92 1059.48 1084.92
c 1056.48 1084.92 1053.72 1083.84 1051.08 1081.8 c 1048.44 1079.64 1046.16 1077.12 1044.12 1074 c 1044.12
1073.64 l 1037.52 1047.12 l 1038 1044.24 1039.08 1041.72 1040.64 1039.8 c 1042.2 1037.76 1044.24 1036.8
1046.88 1036.8 c f Q q 932.57 96.4531 49 89 re W n 934.879 106.879 m 934.52 106.879 934.16 107 933.801
107.48 c 933.32 107.961 933.199 108.32 933.199 108.801 c 933.199 109.398 936.441 112.762 943.16 118.879 c
943.16 119.121 943.039 119.48 943.039 119.719 c 943.039 119.961 943.039 120.199 943.039 120.441 c 943.039
123.801 943.281 127.281 943.762 130.879 c 944.238 134.48 945.078 137.961 946.039 141.559 c 947 145.16
948.199 148.879 949.641 152.719 c 951.078 156.559 952.641 160.039 954.078 162.922 c 956.602 168.199
959.359 172.762 962.48 176.719 c 965.48 180.559 968.84 182.602 972.68 182.602 c 975.199 182.602 977
181.641 978.078 179.961 c 979.16 178.281 979.762 176.121 979.762 173.48 c 979.762 158.84 970.039 141.559
950.602 121.762 c 950.238 118.281 950.121 115.879 950.121 114.441 c 950.121 104.359 952.398 99.4414
956.961 99.4414 c 958.879 99.4414 960.801 99.9219 962.719 101 c 964.641 101.961 966.32 103.16 967.762
104.359 c 969.078 105.559 971 107.359 973.52 109.762 c 973.641 109.879 973.879 110 974.238 110 c 974.602
110 974.961 109.762 975.441 109.281 c 975.801 108.801 976.039 108.32 976.039 107.961 c 976.039 107.719
975.801 107.359 975.559 107.121 c 971.961 103.641 968.719 101 965.84 99.0781 c 962.961 97.2813 959.961
96.3203 956.719 96.3203 c 952.398 96.3203 949.281 98.1211 947.238 101.719 c 945.199 105.32 943.879 109.52
943.398 114.559 c 941.238 112.398 938.602 109.879 935.602 107.121 c 935.48 106.879 935.238 106.879
934.879 106.879 c h 951.559 127.52 m 955.762 132.078 959.719 136.762 963.441 141.801 c 967.16 146.84
970.281 152.121 972.801 157.641 c 975.199 163.16 976.52 168.559 976.52 173.84 c 976.52 177.559 975.199
179.48 972.559 179.48 c 969.68 179.48 966.68 176 963.801 169.281 c 960.801 162.441 958.281 155.121
956.121 147.199 c 953.961 139.281 952.52 132.68 951.559 127.52 c f Q Q
}

The one-loop graph of $i\Pi^{\mu\nu}(p)$, with a closed loop of a fermion $f$: the loop momenta $\ell+p$ and $\ell$ run along the arrows.

Let $i\Pi_2^{\mu\nu}(p)$ be the $O(e^2)$ part of $i\Pi^{\mu\nu}(p)$. Starting at the vertex $\mu$ and reading the loop against the arrows with \eqref{eq:7.2},
\begin{equation*}
\begin{aligned}
i\Pi_2^{\mu\nu}(p)
=\sum_fN_{\mathrm c,f}\,(-1)\int\frac{d^4\ell}{(2\pi)^4}\,
\operatorname{tr}\biggl[
&(ieQ_f\gamma^\mu)\,\frac{i(\not\ell+m_f)}{\ell^2-m_f^2+i\varepsilon}\\[6pt]
&\times(ieQ_f\gamma^\nu)\,\frac{i(\not\ell+\slashed{p}+m_f)}{(\ell+p)^2-m_f^2+i\varepsilon}\biggr],
\end{aligned}
\tag{7.16}\label{eq:7.16}
\end{equation*}
where the factor $-1$ and the Dirac trace come from the closed fermion loop, and $N_{\mathrm c,f}$ is the number of colors of $f$ ($1$ for a lepton, $3$ for a quark) from the color trace.

In $d=4-\epsilon$ dimensions, replace $d^4\ell/(2\pi)^4$ by $d^d\ell/(2\pi)^d$ and let the coupling carry the factor of $\mu$, $e\to\mu^{\epsilon/2}e$. Collecting $(-1)(ieQ_f)^2i^2=-e^2Q_f^2$ and using $\operatorname{tr}(\gamma^\mu\gamma^\nu)=4\eta^{\mu\nu}$, $\operatorname{tr}(\gamma^\mu\gamma^\rho\gamma^\nu\gamma^\sigma)=4(\eta^{\mu\rho}\eta^{\nu\sigma}-\eta^{\mu\nu}\eta^{\rho\sigma}+\eta^{\mu\sigma}\eta^{\rho\nu})$ and the vanishing of traces of an odd number of $\gamma$ matrices,
\begin{equation*}
\begin{aligned}
i\Pi_2^{\mu\nu}(p)
={}&-\mu^{\epsilon}e^2\sum_fN_{\mathrm c,f}Q_f^2\int\frac{d^d\ell}{(2\pi)^d}\,
\frac{\operatorname{tr}\bigl[\gamma^\mu(\not\ell+m_f)\gamma^\nu(\not\ell+\slashed{p}+m_f)\bigr]}{\mathcal P_1\mathcal P_2}\\[6pt]
={}&-4\mu^{\epsilon}e^2\sum_fN_{\mathrm c,f}Q_f^2\int\frac{d^d\ell}{(2\pi)^d}\,\frac{1}{\mathcal P_1\mathcal P_2}\\
&\times\Bigl\{\ell^\mu(\ell+p)^\nu+\ell^\nu(\ell+p)^\mu-\eta^{\mu\nu}\bigl[\ell\cdot(\ell+p)-m_f^2\bigr]\Bigr\},\\[6pt]
\mathcal P_1:={}&\ell^2-m_f^2+i\varepsilon,\qquad
\mathcal P_2:=(\ell+p)^2-m_f^2+i\varepsilon .
\end{aligned}
\tag{7.17}\label{eq:7.17}
\end{equation*}
Combine the denominators with a Feynman parameter $x$:
\begin{equation*}
\begin{aligned}
\frac{1}{\mathcal P_1\mathcal P_2}
&=\int_0^1\frac{dx}{\bigl[x\mathcal P_2+(1-x)\mathcal P_1\bigr]^2}
=\int_0^1\frac{dx}{\bigl[(\ell+xp)^2-\Delta_f+i\varepsilon\bigr]^2},\\[6pt]
\Delta_f&:=m_f^2-x(1-x)p^2 .
\end{aligned}
\tag{7.18}\label{eq:7.18}
\end{equation*}
Shift $\ell\to\ell-xp$ and drop the terms odd in $\ell$, which integrate to zero; then replace $\ell^\mu\ell^\nu$ by $\eta^{\mu\nu}\ell^2/d$ under the integral. The numerator of \eqref{eq:7.17} becomes
\begin{equation*}
\begin{aligned}
&\ell^\mu(\ell+p)^\nu+\ell^\nu(\ell+p)^\mu-\eta^{\mu\nu}\bigl[\ell\cdot(\ell+p)-m_f^2\bigr]\\[4pt]
&\quad\to\ 2\ell^\mu\ell^\nu-\eta^{\mu\nu}\ell^2-2x(1-x)p^\mu p^\nu+\eta^{\mu\nu}\bigl[x(1-x)p^2+m_f^2\bigr]\\[4pt]
&\quad\to\ 2x(1-x)\bigl(p^2\eta^{\mu\nu}-p^\mu p^\nu\bigr)+\eta^{\mu\nu}\Bigl[\frac{2-d}{d}\,\ell^2+\Delta_f\Bigr],
\end{aligned}
\tag{7.19}\label{eq:7.19}
\end{equation*}
so that
\begin{equation*}
\begin{aligned}
i\Pi_2^{\mu\nu}(p)
={}&-4\mu^{\epsilon}e^2\sum_fN_{\mathrm c,f}Q_f^2\int_0^1dx\int\frac{d^d\ell}{(2\pi)^d}\,
\frac{1}{\bigl(\ell^2-\Delta_f+i\varepsilon\bigr)^2}\\[6pt]
&\times\Bigl\{2x(1-x)\bigl(p^2\eta^{\mu\nu}-p^\mu p^\nu\bigr)+\eta^{\mu\nu}\Bigl[\frac{2-d}{d}\,\ell^2+\Delta_f\Bigr]\Bigr\}.
\end{aligned}
\tag{7.20}\label{eq:7.20}
\end{equation*}
After Wick rotation, for $\Delta>0$ and continued to $\Delta<0$ through $\Delta-i\varepsilon$,
\begin{equation*}
\begin{aligned}
\int\frac{d^d\ell}{(2\pi)^d}\,\frac{1}{(\ell^2-\Delta+i\varepsilon)^2}
&=\frac{i}{(4\pi)^{d/2}}\,\Gamma\Bigl(2-\frac d2\Bigr)(\Delta-i\varepsilon)^{\frac d2-2},\\[6pt]
\int\frac{d^d\ell}{(2\pi)^d}\,\frac{\ell^2}{(\ell^2-\Delta+i\varepsilon)^2}
&=-\frac{i}{(4\pi)^{d/2}}\,\frac d2\,\Gamma\Bigl(1-\frac d2\Bigr)(\Delta-i\varepsilon)^{\frac d2-1},\\[6pt]
\int\frac{d^d\ell}{(2\pi)^d}\,\frac{\frac{2-d}{d}\,\ell^2+\Delta}{(\ell^2-\Delta+i\varepsilon)^2}
&=\frac{i\,(\Delta-i\varepsilon)^{\frac d2-1}}{(4\pi)^{d/2}}\Bigl[\Bigl(\frac d2-1\Bigr)\Gamma\Bigl(1-\frac d2\Bigr)+\Gamma\Bigl(2-\frac d2\Bigr)\Bigr]=0,
\end{aligned}
\tag{7.21}\label{eq:7.21}
\end{equation*}
where the last line uses $\Gamma(2-\tfrac d2)=(1-\tfrac d2)\Gamma(1-\tfrac d2)$. Only the transverse term of \eqref{eq:7.20} survives:
\begin{equation*}
\begin{aligned}
i\Pi_2^{\mu\nu}(p)
={}&-4\mu^{\epsilon}e^2\sum_fN_{\mathrm c,f}Q_f^2\int_0^1dx\,2x(1-x)\bigl(p^2\eta^{\mu\nu}-p^\mu p^\nu\bigr)\\[6pt]
&\times\frac{i}{(4\pi)^{d/2}}\,\Gamma\Bigl(2-\frac d2\Bigr)(\Delta_f-i\varepsilon)^{\frac d2-2} .
\end{aligned}
\tag{7.22}\label{eq:7.22}
\end{equation*}
With $d=4-\epsilon$, $\Gamma(\epsilon/2)=2/\epsilon-\gamma_E+O(\epsilon)$ and $2/\bar\epsilon$ as in \eqref{eq:7.14},
\begin{equation*}
{\displaystyle 
\begin{aligned}
\Pi_2^{\mu\nu}(p)&=\bigl(p^2\eta^{\mu\nu}-p^\mu p^\nu\bigr)\,\Pi_2(p^2),\\[6pt]
\Pi_2(p^2)&=-\frac{e^2}{2\pi^2}\sum_fN_{\mathrm c,f}Q_f^2\int_0^1dx\,x(1-x)\left[\frac2{\bar\epsilon}-\ln\frac{\Delta_f-i\varepsilon}{\mu^2}\right]+O(\epsilon).
\end{aligned}
}
\tag{7.23}\label{eq:7.23}
\end{equation*}

\topic{Three External Legs}

Three external legs means $(E_F,E_\gamma)=(2,1)$ or $(0,3)$. Let $ieQ_f\Gamma^\mu(p',p)$ be the sum of amputated 1PI graphs with an incoming fermion of momentum $p$, an outgoing fermion of momentum $p':=p+q$, and an incoming photon of momentum $q$ and index $\mu$. Write $\Gamma^\mu=\gamma^\mu+\Lambda^\mu$.

\topic{(1) Vertex}
\begin{equation*}
E_F=2,\qquad E_\gamma=1\qquad\Longrightarrow\qquad\omega=0 .
\tag{7.24}\label{eq:7.24}
\end{equation*}

% vertex
\qftdiagram{237.23}{156.88}{%
q 0.1 0 0 0.1 0 0 cm 1 g 0 0 2372.3 1568.78 re f 0 g 405.586 669.602 m 313.754 704.035 l 313.754 635.164
l f 941.266 669.602 m 849.434 704.035 l 849.434 635.164 l f 1515.21 669.602 m 1423.38 704.035 l 1423.38
635.164 l f 2050.89 669.602 m 1959.06 704.035 l 1959.06 635.164 l f 14 w 1 J 1 j 0 G 114.789 669.602 m
612.207 669.602 l S 612.207 669.602 m 1186.15 669.602 l S 1186.15 669.602 m 1760.09 669.602 l S 1760.09
669.602 m 2257.51 669.602 l S 2 J 612.207 669.602 m 594.023 688.254 l 583.586 702.844 l 579.066 712.691 l
576.887 722.5 l 577.23 732.145 l 580.125 741.508 l 585.43 750.492 l 592.883 759.02 l 607.164 770.906 l
629.191 785.238 l 654.898 801.906 l 666.203 812.215 l 671.27 819.582 l 674.164 827.527 l 674.891 836.164
l 673.602 845.559 l 668.535 861.078 l 644.289 920.113 l 642.09 931.605 l 641.813 942.379 l 643.691
952.211 l 647.844 960.914 l 654.234 968.379 l 662.699 974.566 l 672.945 979.523 l 690.789 984.949 l
716.648 989.699 l 746.637 995.313 l 760.879 1000.6 l 768.27 1005.55 l 773.871 1011.87 l 777.742 1019.66 l
780.055 1028.93 l 781.219 1045.32 l 781.512 1109.15 l 783.965 1120.54 l 787.93 1130.5 l 793.535 1138.75 l
800.805 1145.08 l 809.66 1149.4 l 819.93 1151.76 l 831.355 1152.31 l 849.977 1150.38 l 875.676 1144.81 l
905.418 1138.63 l 920.496 1138.2 l 929.145 1140.07 l 936.676 1143.89 l 943.176 1149.73 l 948.809 1157.52
l 956.109 1172.36 l 980.91 1231.19 l 987.586 1240.66 l 995.117 1248.22 l 1003.5 1253.55 l 1012.69 1256.46
l 1022.57 1256.91 l 1032.98 1255.02 l 1043.77 1251.04 l 1060.23 1242 l 1087.04 1223.18 l 1106.87 1210.02
l 1120.6 1204.03 l 1129.28 1202.58 l 1137.68 1203.37 l 1145.92 1206.42 l 1154.11 1211.59 l 1166.55
1222.65 l 1207.21 1263.72 l 1216.94 1270.7 l 1226.78 1275.76 l 1236.61 1278.55 l 1246.3 1278.84 l 1255.75
1276.58 l 1264.84 1271.86 l 1273.48 1264.91 l 1285.55 1251.16 l 1300.07 1229.36 l 1316.79 1203.11 l 1327
1191.14 l 1334.24 1185.55 l 1342.03 1182.1 l 1350.48 1180.84 l 1359.68 1181.65 l 1374.93 1186.16 l
1403.54 1198.83 l 1427.73 1208.77 l 1445.34 1213.28 l 1456.32 1214.06 l 1466.4 1212.74 l 1475.4 1209.16 l
1483.18 1203.31 l 1489.69 1195.35 l 1494.94 1185.51 l 1500.74 1168.08 l 1505.74 1142.34 l 1511.18 1111.82
l 1516.11 1096.99 l 1520.74 1089.16 l 1526.7 1083.11 l 1534.13 1078.82 l 1543.03 1076.14 l 1558.95
1074.57 l 1590.34 1075.38 l 1616.46 1075.22 l 1634.36 1072.53 l 1644.71 1068.96 l 1653.41 1063.77 l
1660.23 1056.91 l 1665.06 1048.44 l 1667.9 1038.51 l 1668.88 1027.34 l 1667.43 1008.96 l 1662.13 983.254
l 1655.55 953.074 l 1654.57 937.594 l 1655.97 928.656 l 1659.29 920.855 l 1664.63 914.109 l 1671.95
908.281 l 1686.21 900.789 l 1715.63 889.527 l 1739.64 879.32 l 1755 869.91 l 1763.07 862.594 l 1768.98
854.418 l 1772.51 845.426 l 1773.57 835.707 l 1772.25 825.41 l 1768.76 814.699 l 1760.25 798.258 l
1745.45 776.566 l 1727.95 751.262 l 1721.28 737.375 l 1719.28 728.605 l 1719.5 720.145 l 1722 711.879 l
1726.66 703.691 l 1737.13 691.301 l 1760.09 669.602 l 1760.09 669.602 l S 1186.15 114.789 m 1164.91
140.766 l 1155.6 155.609 l 1150.23 168.602 l 1148.08 179.734 l 1148.33 190.867 l 1150.97 202 l 1155.84
213.133 l 1163.92 226.121 l 1178.56 244.676 l 1209.03 281.789 l 1217.86 296.633 l 1222.2 307.766 l
1224.26 318.898 l 1223.91 330.031 l 1221.16 341.168 l 1216.21 352.301 l 1208.05 365.289 l 1193.34 383.844
l 1162.95 420.957 l 1154.21 435.801 l 1149.96 446.934 l 1148 458.066 l 1148.46 469.199 l 1151.3 480.332 l
1156.34 491.465 l 1164.58 504.457 l 1179.35 523.012 l 1209.67 560.121 l 1218.31 574.965 l 1222.46 586.102
l 1224.32 597.234 l 1223.77 608.367 l 1220.83 619.5 l 1215.7 630.633 l 1207.39 643.621 l 1192.55 662.18 l
1186.15 669.602 l 1186.15 669.602 l S 10 w 0 J 612.207 647.102 m 618.172 647.102 623.895 649.473 628.117
653.691 c 632.336 657.91 634.707 663.633 634.707 669.602 c 634.707 675.566 632.336 681.289 628.117
685.512 c 623.895 689.73 618.172 692.102 612.207 692.102 c 606.238 692.102 600.516 689.73 596.297 685.512
c 592.078 681.289 589.707 675.566 589.707 669.602 c 589.707 663.633 592.078 657.91 596.297 653.691 c
600.516 649.473 606.238 647.102 612.207 647.102 c h B 1186.15 647.102 m 1192.12 647.102 1197.84 649.473
1202.06 653.691 c 1206.28 657.91 1208.65 663.633 1208.65 669.602 c 1208.65 675.566 1206.28 681.289
1202.06 685.512 c 1197.84 689.73 1192.12 692.102 1186.15 692.102 c 1180.18 692.102 1174.46 689.73 1170.24
685.512 c 1166.02 681.289 1163.65 675.566 1163.65 669.602 c 1163.65 663.633 1166.02 657.91 1170.24
653.691 c 1174.46 649.473 1180.18 647.102 1186.15 647.102 c h B 1760.09 647.102 m 1766.06 647.102 1771.78
649.473 1776 653.691 c 1780.22 657.91 1782.59 663.633 1782.59 669.602 c 1782.59 675.566 1780.22 681.289
1776 685.512 c 1771.78 689.73 1766.06 692.102 1760.09 692.102 c 1754.13 692.102 1748.4 689.73 1744.18
685.512 c 1739.96 681.289 1737.59 675.566 1737.59 669.602 c 1737.59 663.633 1739.96 657.91 1744.18
653.691 c 1748.4 649.473 1754.13 647.102 1760.09 647.102 c h B q 473.07 725.996 64 88 re W n 473.641
729.691 m 473.641 730.211 473.641 730.602 473.77 730.73 c 486.641 782.469 l 487.809 787.148 489.891
791.699 492.879 796.121 c 495.871 800.539 499.512 804.309 503.801 807.172 c 507.961 810.031 512.512
811.461 517.191 811.461 c 520.828 811.461 524.078 810.422 526.941 808.469 c 529.801 806.391 531.879
803.789 533.309 800.672 c 534.738 797.551 535.52 794.039 535.52 790.27 c 535.52 785.98 534.609 781.559
532.922 777.141 c 531.102 772.719 528.762 768.559 525.77 764.922 c 522.781 761.148 519.27 758.16 515.371
755.949 c 511.469 753.871 507.309 752.699 503.02 752.699 c 500.422 752.699 498.078 753.48 495.871 754.781
c 493.531 756.211 491.84 758.16 490.539 760.762 c 482.871 730.211 l 482.609 729.039 481.961 728 480.922
727.219 c 479.879 726.441 478.84 726.051 477.539 726.051 c 476.5 726.051 475.59 726.441 474.809 727.09 c
474.031 727.738 473.641 728.648 473.641 729.691 c h 503.281 755.949 m 507.828 755.949 511.859 758.289
515.238 762.969 c 518.621 767.648 521.219 773.238 523.039 779.609 c 524.859 785.98 525.77 791.441 525.77
795.859 c 525.77 799.109 524.988 801.969 523.559 804.441 c 522.129 806.781 519.922 808.078 516.93 808.078
c 513.031 808.078 509.781 806.648 506.922 803.789 c 504.059 800.93 501.719 797.551 500.031 793.648 c
498.211 789.75 496.781 785.852 495.738 781.691 c 492.23 767.129 l 492.75 764.012 493.922 761.281 495.871
759.199 c 497.82 756.988 500.289 755.949 503.281 755.949 c f Q q 1149.15 733.652 74 88 re W n 1149.64
737.691 m 1149.64 738.211 1149.64 738.602 1149.77 738.73 c 1168.88 815.301 l 1169.14 816.469 1169.79
817.512 1170.83 818.289 c 1171.87 819.07 1173.04 819.461 1174.34 819.461 c 1175.38 819.461 1176.29 819.07
1177.07 818.422 c 1177.85 817.77 1178.24 816.859 1178.24 815.82 c 1178.24 815.559 1178.24 815.301 1178.24
815.172 c 1178.24 815.039 1178.11 814.91 1178.11 814.781 c 1170.44 784.359 l 1169.53 780.719 1169.14
777.602 1169.14 775 c 1169.14 771.75 1169.79 769.148 1171.35 767.07 c 1172.91 764.988 1175.25 763.949
1178.37 763.949 c 1184.61 763.949 1190.2 767.852 1195.01 775.781 c 1195.01 775.91 l 1195.14 776.039
1195.14 776.172 v 1204.5 813.871 l 1204.76 814.91 1205.54 815.949 1206.58 816.73 c 1207.62 817.512
1208.66 818.031 1209.96 818.031 c 1211 818.031 1211.91 817.641 1212.69 816.988 c 1213.47 816.34 1213.86
815.43 1213.86 814.391 c 1213.86 813.871 1213.73 813.48 1213.73 813.352 c 1204.37 775.91 l 1203.72
773.441 1203.46 771.359 1203.46 769.539 c 1203.46 765.77 1204.63 763.949 1207.23 763.949 c 1209.83
763.949 1211.91 765.512 1213.34 768.891 c 1214.64 772.27 1215.94 776.301 1216.98 781.109 c 1217.11
781.629 1217.5 781.891 1218.02 781.891 c 1219.58 781.891 l 1219.84 781.891 1220.1 781.629 1220.36 781.371
c 1220.62 781.109 1220.75 780.852 1220.75 780.59 c 1219.19 774.48 1217.63 769.672 1215.81 766.031 c
1213.99 762.391 1211 760.699 1206.97 760.699 c 1204.11 760.699 1201.51 761.48 1199.3 763.039 c 1197.09
764.602 1195.66 766.809 1195.01 769.672 c 1192.8 766.941 1190.2 764.73 1187.34 763.172 c 1184.48 761.609
1181.36 760.699 1178.24 760.699 c 1172.78 760.699 1168.62 762 1165.5 764.34 c 1158.87 738.211 l 1158.61
737.039 1157.96 736 1156.92 735.219 c 1155.88 734.441 1154.84 734.051 1153.54 734.051 c 1152.5 734.051
1151.59 734.441 1150.81 735.09 c 1150.03 735.738 1149.64 736.648 1149.64 737.691 c f Q q 1838.23 751.027
64 60 re W n 1840.19 751 m 1839.15 751 1838.63 751.52 1838.63 752.691 c 1849.68 796.891 l 1849.94 798.32
1850.2 799.488 1850.33 800.398 c 1850.33 801.699 1847.47 802.48 1841.88 802.48 c 1840.97 802.48 1840.58
803 1840.58 804.172 c 1840.58 804.43 1840.58 804.82 1840.84 805.34 c 1840.97 805.859 1841.1 806.25
1841.36 806.512 c 1841.62 806.77 1841.88 807.031 1842.4 807.031 c 1859.95 808.461 l 1860.99 808.461
1861.51 807.809 1861.51 806.77 c 1848.77 755.809 l 1855.53 758.02 1861.9 761.531 1867.62 766.469 c
1873.34 771.41 1878.28 777.129 1882.31 783.762 c 1886.34 790.262 1889.2 797.148 1891.02 804.301 c 1891.28
805.469 1891.93 806.512 1892.97 807.289 c 1894.01 808.07 1895.18 808.461 1896.48 808.461 c 1897.52
808.461 1898.43 808.07 1899.21 807.422 c 1899.99 806.641 1900.38 805.73 1900.38 804.691 c 1900.38 803.262
1899.73 801.18 1898.69 798.32 c 1897.52 795.461 1896.35 792.859 1895.31 790.648 c 1891.8 783.891 1887.25
777.648 1881.66 771.93 c 1878.02 768.289 1874.25 765.172 1870.22 762.441 c 1866.06 759.711 1861.77
757.371 1857.35 755.422 c 1852.93 753.469 1848.25 752.039 1843.31 751 c 1840.19 751 l f Q 9 w 1 J
0.449951 G 230.449 554.813 m 312.777 554.813 391.754 554.813 467.379 554.813 c S 0.449951 g 431.379
572.813 m 467.379 554.813 l 431.379 536.813 l 431.379 572.813 l h B q 313.02 402.582 65 79 re W n 314.121
403.84 m 313.398 403.84 313.039 404.441 313.039 405.398 c 313.281 407.199 313.879 408.039 314.602 408.039
c 317 408.039 318.68 408.281 319.641 408.762 c 320.602 409.238 321.441 410.32 322.039 412.121 c 335.602
466 l 336.078 467.441 336.441 469.359 336.441 471.762 c 336.441 475.121 335.238 476.922 332.961 476.922 c
330.441 476.922 328.52 475.359 327.32 472.359 c 326.121 469.359 324.922 465.641 323.84 461.078 c 323.84
460.602 323.48 460.359 322.879 460.359 c 321.441 460.359 l 321.199 460.359 321.078 460.48 320.84 460.719
c 320.602 460.961 320.48 461.32 320.48 461.559 c 321.32 465.039 322.16 468.039 323 470.441 c 323.719
472.84 325.039 475 326.719 477.039 c 328.398 478.961 330.441 480.039 333.078 480.039 c 335.719 480.039
338.121 479.199 340.281 477.641 c 342.32 476.078 343.641 474.039 344.238 471.398 c 346.398 473.922
348.801 475.961 351.441 477.641 c 354.078 479.199 356.84 480.039 359.719 480.039 c 363.078 480.039
365.961 479.078 368.48 477.281 c 370.879 475.48 372.801 473.078 374 470.199 c 375.199 467.32 375.922
464.199 375.922 460.84 c 375.922 455.68 374.602 450.398 371.961 444.879 c 369.32 439.359 365.719 434.801
361.281 431.199 c 356.719 427.602 351.922 425.801 346.762 425.801 c 344.359 425.801 342.32 426.52 340.398
427.719 c 338.48 428.922 336.922 430.719 335.961 432.879 c 330.559 411.641 l 330.32 410.441 330.199
409.84 330.199 409.602 c 330.199 408.641 332.84 408.039 338.121 408.039 c 338.84 408.039 339.32 407.559
339.32 406.48 c 339.078 405.52 338.84 404.801 338.602 404.441 c 338.359 404.078 338 403.84 337.281 403.84
c 314.121 403.84 l h 346.879 428.801 m 349.641 428.801 352.16 429.879 354.559 432.16 c 356.961 434.441
358.879 437.078 360.441 439.961 c 361.398 442.121 362.48 444.641 363.441 447.762 c 364.398 450.879
365.238 454.121 365.961 457.719 c 366.68 461.199 367.039 463.961 367.039 466.121 c 367.039 467.922
366.801 469.602 366.32 471.281 c 365.84 472.84 365 474.16 363.922 475.238 c 362.719 476.32 361.281
476.922 359.48 476.922 c 356.48 476.922 353.719 475.84 351.078 473.801 c 348.441 471.641 346.16 469.121
344.121 466 c 344.121 465.641 l 337.52 439.121 l 338 436.238 339.078 433.719 340.641 431.801 c 342.199
429.762 344.238 428.801 346.879 428.801 c f Q 1894.88 554.813 m 1977.21 554.813 2056.19 554.813 2131.81
554.813 c S 2095.81 572.813 m 2131.81 554.813 l 2095.81 536.813 l 2095.81 572.813 l h B q 1966.28 395.836
65 79 re W n 1967.12 396.84 m 1966.4 396.84 1966.04 397.441 1966.04 398.398 c 1966.28 400.199 1966.88
401.039 1967.6 401.039 c 1970 401.039 1971.68 401.281 1972.64 401.762 c 1973.6 402.238 1974.44 403.32
1975.04 405.121 c 1988.6 459 l 1989.08 460.441 1989.44 462.359 1989.44 464.762 c 1989.44 468.121 1988.24
469.922 1985.96 469.922 c 1983.44 469.922 1981.52 468.359 1980.32 465.359 c 1979.12 462.359 1977.92
458.641 1976.84 454.078 c 1976.84 453.602 1976.48 453.359 1975.88 453.359 c 1974.44 453.359 l 1974.2
453.359 1974.08 453.48 1973.84 453.719 c 1973.6 453.961 1973.48 454.32 1973.48 454.559 c 1974.32 458.039
1975.16 461.039 1976 463.441 c 1976.72 465.84 1978.04 468 1979.72 470.039 c 1981.4 471.961 1983.44
473.039 1986.08 473.039 c 1988.72 473.039 1991.12 472.199 1993.28 470.641 c 1995.32 469.078 1996.64
467.039 1997.24 464.398 c 1999.4 466.922 2001.8 468.961 2004.44 470.641 c 2007.08 472.199 2009.84 473.039
2012.72 473.039 c 2016.08 473.039 2018.96 472.078 2021.48 470.281 c 2023.88 468.48 2025.8 466.078 2027
463.199 c 2028.2 460.32 2028.92 457.199 2028.92 453.84 c 2028.92 448.68 2027.6 443.398 2024.96 437.879 c
2022.32 432.359 2018.72 427.801 2014.28 424.199 c 2009.72 420.602 2004.92 418.801 1999.76 418.801 c
1997.36 418.801 1995.32 419.52 1993.4 420.719 c 1991.48 421.922 1989.92 423.719 1988.96 425.879 c 1983.56
404.641 l 1983.32 403.441 1983.2 402.84 1983.2 402.602 c 1983.2 401.641 1985.84 401.039 1991.12 401.039 c
1991.84 401.039 1992.32 400.559 1992.32 399.48 c 1992.08 398.52 1991.84 397.801 1991.6 397.441 c 1991.36
397.078 1991 396.84 1990.28 396.84 c 1967.12 396.84 l h 1999.88 421.801 m 2002.64 421.801 2005.16 422.879
2007.56 425.16 c 2009.96 427.441 2011.88 430.078 2013.44 432.961 c 2014.4 435.121 2015.48 437.641 2016.44
440.762 c 2017.4 443.879 2018.24 447.121 2018.96 450.719 c 2019.68 454.199 2020.04 456.961 2020.04
459.121 c 2020.04 460.922 2019.8 462.602 2019.32 464.281 c 2018.84 465.84 2018 467.16 2016.92 468.238 c
2015.72 469.32 2014.28 469.922 2012.48 469.922 c 2009.48 469.922 2006.72 468.84 2004.08 466.801 c 2001.44
464.641 1999.16 462.121 1997.12 459 c 1997.12 458.641 l 1990.52 432.121 l 1991 429.238 1992.08 426.719
1993.64 424.801 c 1995.2 422.762 1997.24 421.801 1999.88 421.801 c f Q q 2040.18 468.754 22 46 re W n
2040.94 470.535 m 2040.61 470.621 2040.44 470.871 2040.44 471.375 c 2050.68 509.512 l 2051.02 510.52
2051.52 511.359 2052.36 512.031 c 2053.12 512.621 2054.04 512.957 2055.05 512.957 c 2056.39 512.957
2057.49 512.453 2058.5 511.613 c 2059.5 510.773 2060.01 509.68 2060.01 508.422 c 2060.01 507.832 2059.84
507.16 2059.59 506.488 c 2044.38 470.199 l 2044.21 469.781 2043.96 469.613 2043.63 469.613 c 2043.29
469.613 2042.87 469.695 2042.2 470.031 c 2041.53 470.285 2041.11 470.453 2040.94 470.535 c f Q q 875.176
502.043 49 89 re W n 876.879 512.879 m 876.52 512.879 876.16 513 875.801 513.48 c 875.32 513.961 875.199
514.32 875.199 514.801 c 875.199 515.398 878.441 518.762 885.16 524.879 c 885.16 525.121 885.039 525.48
885.039 525.719 c 885.039 525.961 885.039 526.199 885.039 526.441 c 885.039 529.801 885.281 533.281
885.762 536.879 c 886.238 540.48 887.078 543.961 888.039 547.559 c 889 551.16 890.199 554.879 891.641
558.719 c 893.078 562.559 894.641 566.039 896.078 568.922 c 898.602 574.199 901.359 578.762 904.48
582.719 c 907.48 586.559 910.84 588.602 914.68 588.602 c 917.199 588.602 919 587.641 920.078 585.961 c
921.16 584.281 921.762 582.121 921.762 579.48 c 921.762 564.84 912.039 547.559 892.602 527.762 c 892.238
524.281 892.121 521.879 892.121 520.441 c 892.121 510.359 894.398 505.441 898.961 505.441 c 900.879
505.441 902.801 505.922 904.719 507 c 906.641 507.961 908.32 509.16 909.762 510.359 c 911.078 511.559 913
513.359 915.52 515.762 c 915.641 515.879 915.879 516 916.238 516 c 916.602 516 916.961 515.762 917.441
515.281 c 917.801 514.801 918.039 514.32 918.039 513.961 c 918.039 513.719 917.801 513.359 917.559
513.121 c 913.961 509.641 910.719 507 907.84 505.078 c 904.961 503.281 901.961 502.32 898.719 502.32 c
894.398 502.32 891.281 504.121 889.238 507.719 c 887.199 511.32 885.879 515.52 885.398 520.559 c 883.238
518.398 880.602 515.879 877.602 513.121 c 877.48 512.879 877.238 512.879 876.879 512.879 c h 893.559
533.52 m 897.762 538.078 901.719 542.762 905.441 547.801 c 909.16 552.84 912.281 558.121 914.801 563.641
c 917.199 569.16 918.52 574.559 918.52 579.84 c 918.52 583.559 917.199 585.48 914.559 585.48 c 911.68
585.48 908.68 582 905.801 575.281 c 902.801 568.441 900.281 561.121 898.121 553.199 c 895.961 545.281
894.52 538.68 893.559 533.52 c f Q q 1354.12 505.168 49 89 re W n 1355.88 515.879 m 1355.52 515.879
1355.16 516 1354.8 516.48 c 1354.32 516.961 1354.2 517.32 1354.2 517.801 c 1354.2 518.398 1357.44 521.762
1364.16 527.879 c 1364.16 528.121 1364.04 528.48 1364.04 528.719 c 1364.04 528.961 1364.04 529.199
1364.04 529.441 c 1364.04 532.801 1364.28 536.281 1364.76 539.879 c 1365.24 543.48 1366.08 546.961
1367.04 550.559 c 1368 554.16 1369.2 557.879 1370.64 561.719 c 1372.08 565.559 1373.64 569.039 1375.08
571.922 c 1377.6 577.199 1380.36 581.762 1383.48 585.719 c 1386.48 589.559 1389.84 591.602 1393.68
591.602 c 1396.2 591.602 1398 590.641 1399.08 588.961 c 1400.16 587.281 1400.76 585.121 1400.76 582.48 c
1400.76 567.84 1391.04 550.559 1371.6 530.762 c 1371.24 527.281 1371.12 524.879 1371.12 523.441 c 1371.12
513.359 1373.4 508.441 1377.96 508.441 c 1379.88 508.441 1381.8 508.922 1383.72 510 c 1385.64 510.961
1387.32 512.16 1388.76 513.359 c 1390.08 514.559 1392 516.359 1394.52 518.762 c 1394.64 518.879 1394.88
519 1395.24 519 c 1395.6 519 1395.96 518.762 1396.44 518.281 c 1396.8 517.801 1397.04 517.32 1397.04
516.961 c 1397.04 516.719 1396.8 516.359 1396.56 516.121 c 1392.96 512.641 1389.72 510 1386.84 508.078 c
1383.96 506.281 1380.96 505.32 1377.72 505.32 c 1373.4 505.32 1370.28 507.121 1368.24 510.719 c 1366.2
514.32 1364.88 518.52 1364.4 523.559 c 1362.24 521.398 1359.6 518.879 1356.6 516.121 c 1356.48 515.879
1356.24 515.879 1355.88 515.879 c h 1372.56 536.52 m 1376.76 541.078 1380.72 545.762 1384.44 550.801 c
1388.16 555.84 1391.28 561.121 1393.8 566.641 c 1396.2 572.16 1397.52 577.559 1397.52 582.84 c 1397.52
586.559 1396.2 588.48 1393.56 588.48 c 1390.68 588.48 1387.68 585 1384.8 578.281 c 1381.8 571.441 1379.28
564.121 1377.12 556.199 c 1374.96 548.281 1373.52 541.68 1372.56 536.52 c f Q q 1430.11 497.168 82 82 re
W n 1433 534.602 m 1432.28 534.602 1431.8 534.84 1431.32 535.32 c 1430.84 535.801 1430.72 536.398 1430.72
537 c 1430.72 537.48 1430.84 538.078 1431.32 538.559 c 1431.8 539.039 1432.28 539.398 1433 539.398 c
1468.28 539.398 l 1468.28 574.801 l 1468.28 575.398 1468.52 575.879 1469 576.359 c 1469.48 576.719
1469.96 576.961 1470.68 576.961 c 1471.16 576.961 1471.76 576.719 1472.24 576.359 c 1472.72 575.879
1473.08 575.398 1473.08 574.801 c 1473.08 539.398 l 1508.36 539.398 l 1508.84 539.398 1509.32 539.039
1509.8 538.559 c 1510.28 538.078 1510.52 537.48 1510.52 537 c 1510.52 536.398 1510.28 535.801 1509.8
535.32 c 1509.32 534.84 1508.84 534.602 1508.36 534.602 c 1473.08 534.602 l 1473.08 499.32 l 1473.08
498.719 1472.72 498.238 1472.24 497.762 c 1471.76 497.398 1471.16 497.16 1470.68 497.16 c 1469.96 497.16
1469.48 497.398 1469 497.762 c 1468.52 498.238 1468.28 498.719 1468.28 499.32 c 1468.28 534.602 l 1433
534.602 l f Q q 1542.41 483.168 53 79 re W n 1555.64 485.398 m 1555.88 487.199 1556.36 488.039 1557.2
488.039 c 1560.08 488.039 1562.36 488.281 1564.04 488.762 c 1565.72 489.238 1566.8 490.441 1567.28
492.121 c 1572.44 512.398 l 1567.88 507.961 1563.32 505.801 1558.76 505.801 c 1553.72 505.801 1549.76
507.602 1546.88 511.441 c 1544 515.16 1542.56 519.602 1542.56 524.879 c 1542.56 529.922 1543.88 535.32
1546.52 540.84 c 1549.16 546.359 1552.64 550.922 1557.2 554.52 c 1561.76 558.121 1566.56 560.039 1571.84
560.039 c 1574.24 560.039 1576.52 559.32 1578.44 557.879 c 1580.36 556.441 1581.8 554.52 1582.88 552.121
c 1583.36 553.078 1584.56 554.641 1586.72 556.801 c 1588.88 558.961 1590.32 560.039 1591.28 560.039 c
1592.12 560.039 1592.6 559.441 1592.6 558.48 c 1575.8 491.641 l 1575.56 490.441 1575.44 489.84 1575.44
489.602 c 1575.44 488.641 1578.08 488.039 1583.36 488.039 c 1583.72 488.039 1583.96 487.922 1584.2
487.559 c 1584.44 487.32 1584.56 486.961 1584.56 486.48 c 1584.32 485.52 1584.08 484.801 1583.84 484.441
c 1583.6 484.078 1583.24 483.84 1582.52 483.84 c 1556.84 483.84 l 1556 483.84 1555.64 484.441 1555.64
485.398 c h 1559 508.801 m 1561.76 508.801 1564.52 509.762 1567.16 511.922 c 1569.8 514.078 1572.2
516.602 1574.24 519.602 c 1580.96 546.602 l 1580.48 549.359 1579.4 551.762 1577.84 553.801 c 1576.28
555.84 1574.24 556.922 1571.6 556.922 c 1568.84 556.922 1566.32 555.719 1564.04 553.559 c 1561.64 551.398
1559.72 548.762 1558.16 545.762 c 1556.48 542.52 1555.04 538.078 1553.6 532.559 c 1552.16 527.039 1551.56
522.719 1551.56 519.602 c 1551.56 516.719 1552.16 514.199 1553.36 512.039 c 1554.56 509.879 1556.36
508.801 1559 508.801 c f Q 1331.55 192.184 m 1331.55 274.516 1331.55 353.492 1331.55 429.113 c S 1313.55
393.113 m 1331.55 429.113 l 1349.55 393.113 l 1313.55 393.113 l h B q 1412.68 272.488 53 79 re W n
1426.64 274.398 m 1426.88 276.199 1427.36 277.039 1428.2 277.039 c 1431.08 277.039 1433.36 277.281
1435.04 277.762 c 1436.72 278.238 1437.8 279.441 1438.28 281.121 c 1443.44 301.398 l 1438.88 296.961
1434.32 294.801 1429.76 294.801 c 1424.72 294.801 1420.76 296.602 1417.88 300.441 c 1415 304.16 1413.56
308.602 1413.56 313.879 c 1413.56 318.922 1414.88 324.32 1417.52 329.84 c 1420.16 335.359 1423.64 339.922
1428.2 343.52 c 1432.76 347.121 1437.56 349.039 1442.84 349.039 c 1445.24 349.039 1447.52 348.32 1449.44
346.879 c 1451.36 345.441 1452.8 343.52 1453.88 341.121 c 1454.36 342.078 1455.56 343.641 1457.72 345.801
c 1459.88 347.961 1461.32 349.039 1462.28 349.039 c 1463.12 349.039 1463.6 348.441 1463.6 347.48 c 1446.8
280.641 l 1446.56 279.441 1446.44 278.84 1446.44 278.602 c 1446.44 277.641 1449.08 277.039 1454.36
277.039 c 1454.72 277.039 1454.96 276.922 1455.2 276.559 c 1455.44 276.32 1455.56 275.961 1455.56 275.48
c 1455.32 274.52 1455.08 273.801 1454.84 273.441 c 1454.6 273.078 1454.24 272.84 1453.52 272.84 c 1427.84
272.84 l 1427 272.84 1426.64 273.441 1426.64 274.398 c h 1430 297.801 m 1432.76 297.801 1435.52 298.762
1438.16 300.922 c 1440.8 303.078 1443.2 305.602 1445.24 308.602 c 1451.96 335.602 l 1451.48 338.359
1450.4 340.762 1448.84 342.801 c 1447.28 344.84 1445.24 345.922 1442.6 345.922 c 1439.84 345.922 1437.32
344.719 1435.04 342.559 c 1432.64 340.398 1430.72 337.762 1429.16 334.762 c 1427.48 331.52 1426.04
327.078 1424.6 321.559 c 1423.16 316.039 1422.56 311.719 1422.56 308.602 c 1422.56 305.719 1423.16
303.199 1424.36 301.039 c 1425.56 298.879 1427.36 297.801 1430 297.801 c f Q 1072.21 1358.33 m 1148.16
1358.33 1220.76 1358.33 1290.01 1358.33 c S 1254.01 1376.33 m 1290.01 1358.33 l 1254.01 1340.33 l 1254.01
1376.33 l h B q 1057.15 1416.78 65 79 re W n 1058.12 1417.84 m 1057.4 1417.84 1057.04 1418.44 1057.04
1419.4 c 1057.28 1421.2 1057.88 1422.04 1058.6 1422.04 c 1061 1422.04 1062.68 1422.28 1063.64 1422.76 c
1064.6 1423.24 1065.44 1424.32 1066.04 1426.12 c 1079.6 1480 l 1080.08 1481.44 1080.44 1483.36 1080.44
1485.76 c 1080.44 1489.12 1079.24 1490.92 1076.96 1490.92 c 1074.44 1490.92 1072.52 1489.36 1071.32
1486.36 c 1070.12 1483.36 1068.92 1479.64 1067.84 1475.08 c 1067.84 1474.6 1067.48 1474.36 1066.88
1474.36 c 1065.44 1474.36 l 1065.2 1474.36 1065.08 1474.48 1064.84 1474.72 c 1064.6 1474.96 1064.48
1475.32 1064.48 1475.56 c 1065.32 1479.04 1066.16 1482.04 1067 1484.44 c 1067.72 1486.84 1069.04 1489
1070.72 1491.04 c 1072.4 1492.96 1074.44 1494.04 1077.08 1494.04 c 1079.72 1494.04 1082.12 1493.2 1084.28
1491.64 c 1086.32 1490.08 1087.64 1488.04 1088.24 1485.4 c 1090.4 1487.92 1092.8 1489.96 1095.44 1491.64
c 1098.08 1493.2 1100.84 1494.04 1103.72 1494.04 c 1107.08 1494.04 1109.96 1493.08 1112.48 1491.28 c
1114.88 1489.48 1116.8 1487.08 1118 1484.2 c 1119.2 1481.32 1119.92 1478.2 1119.92 1474.84 c 1119.92
1469.68 1118.6 1464.4 1115.96 1458.88 c 1113.32 1453.36 1109.72 1448.8 1105.28 1445.2 c 1100.72 1441.6
1095.92 1439.8 1090.76 1439.8 c 1088.36 1439.8 1086.32 1440.52 1084.4 1441.72 c 1082.48 1442.92 1080.92
1444.72 1079.96 1446.88 c 1074.56 1425.64 l 1074.32 1424.44 1074.2 1423.84 1074.2 1423.6 c 1074.2 1422.64
1076.84 1422.04 1082.12 1422.04 c 1082.84 1422.04 1083.32 1421.56 1083.32 1420.48 c 1083.08 1419.52
1082.84 1418.8 1082.6 1418.44 c 1082.36 1418.08 1082 1417.84 1081.28 1417.84 c 1058.12 1417.84 l h
1090.88 1442.8 m 1093.64 1442.8 1096.16 1443.88 1098.56 1446.16 c 1100.96 1448.44 1102.88 1451.08 1104.44
1453.96 c 1105.4 1456.12 1106.48 1458.64 1107.44 1461.76 c 1108.4 1464.88 1109.24 1468.12 1109.96 1471.72
c 1110.68 1475.2 1111.04 1477.96 1111.04 1480.12 c 1111.04 1481.92 1110.8 1483.6 1110.32 1485.28 c
1109.84 1486.84 1109 1488.16 1107.92 1489.24 c 1106.72 1490.32 1105.28 1490.92 1103.48 1490.92 c 1100.48
1490.92 1097.72 1489.84 1095.08 1487.8 c 1092.44 1485.64 1090.16 1483.12 1088.12 1480 c 1088.12 1479.64 l
1081.52 1453.12 l 1082 1450.24 1083.08 1447.72 1084.64 1445.8 c 1086.2 1443.76 1088.24 1442.8 1090.88
1442.8 c f Q q 1151.57 1467.78 76 7 re W n 1155.24 1468.6 m 1154.52 1468.6 1154.04 1468.84 1153.56
1469.32 c 1153.08 1469.8 1152.96 1470.4 1152.96 1471 c 1152.96 1471.48 1153.08 1472.08 1153.56 1472.56 c
1154.04 1473.04 1154.52 1473.4 1155.24 1473.4 c 1224.12 1473.4 l 1224.72 1473.4 1225.2 1473.04 1225.68
1472.56 c 1226.04 1472.08 1226.28 1471.48 1226.28 1471 c 1226.28 1470.4 1226.04 1469.8 1225.68 1469.32 c
1225.2 1468.84 1224.72 1468.6 1224.12 1468.6 c 1155.24 1468.6 l f Q q 1257.86 1438.78 49 89 re W n
1259.88 1449.88 m 1259.52 1449.88 1259.16 1450 1258.8 1450.48 c 1258.32 1450.96 1258.2 1451.32 1258.2
1451.8 c 1258.2 1452.4 1261.44 1455.76 1268.16 1461.88 c 1268.16 1462.12 1268.04 1462.48 1268.04 1462.72
c 1268.04 1462.96 1268.04 1463.2 1268.04 1463.44 c 1268.04 1466.8 1268.28 1470.28 1268.76 1473.88 c
1269.24 1477.48 1270.08 1480.96 1271.04 1484.56 c 1272 1488.16 1273.2 1491.88 1274.64 1495.72 c 1276.08
1499.56 1277.64 1503.04 1279.08 1505.92 c 1281.6 1511.2 1284.36 1515.76 1287.48 1519.72 c 1290.48 1523.56
1293.84 1525.6 1297.68 1525.6 c 1300.2 1525.6 1302 1524.64 1303.08 1522.96 c 1304.16 1521.28 1304.76
1519.12 1304.76 1516.48 c 1304.76 1501.84 1295.04 1484.56 1275.6 1464.76 c 1275.24 1461.28 1275.12
1458.88 1275.12 1457.44 c 1275.12 1447.36 1277.4 1442.44 1281.96 1442.44 c 1283.88 1442.44 1285.8 1442.92
1287.72 1444 c 1289.64 1444.96 1291.32 1446.16 1292.76 1447.36 c 1294.08 1448.56 1296 1450.36 1298.52
1452.76 c 1298.64 1452.88 1298.88 1453 1299.24 1453 c 1299.6 1453 1299.96 1452.76 1300.44 1452.28 c
1300.8 1451.8 1301.04 1451.32 1301.04 1450.96 c 1301.04 1450.72 1300.8 1450.36 1300.56 1450.12 c 1296.96
1446.64 1293.72 1444 1290.84 1442.08 c 1287.96 1440.28 1284.96 1439.32 1281.72 1439.32 c 1277.4 1439.32
1274.28 1441.12 1272.24 1444.72 c 1270.2 1448.32 1268.88 1452.52 1268.4 1457.56 c 1266.24 1455.4 1263.6
1452.88 1260.6 1450.12 c 1260.48 1449.88 1260.24 1449.88 1259.88 1449.88 c h 1276.56 1470.52 m 1280.76
1475.08 1284.72 1479.76 1288.44 1484.8 c 1292.16 1489.84 1295.28 1495.12 1297.8 1500.64 c 1300.2 1506.16
1301.52 1511.56 1301.52 1516.84 c 1301.52 1520.56 1300.2 1522.48 1297.56 1522.48 c 1294.68 1522.48
1291.68 1519 1288.8 1512.28 c 1285.8 1505.44 1283.28 1498.12 1281.12 1490.2 c 1278.96 1482.28 1277.52
1475.68 1276.56 1470.52 c f Q Q
}

The one-loop graph of $ieQ_f\Lambda^\mu(p',p)$ for a fermion $f$: the internal fermions carry $\ell$ and $\ell+q$ along their arrows and the internal photon carries $p-\ell$.

Let $\Lambda_2^\mu(p',p)$ be the $O(e^2)$ part of $\Lambda^\mu(p',p)$. Starting at the vertex $\nu$ and reading the line against the arrow with \eqref{eq:7.2},
\begin{equation*}
\begin{aligned}
ieQ_f\Lambda_2^\mu(p',p)
=\int\frac{d^4\ell}{(2\pi)^4}\,
&(ieQ_f\gamma^\nu)\,\frac{i(\not\ell+\slashed{q}+m_f)}{(\ell+q)^2-m_f^2+i\varepsilon}\,(ieQ_f\gamma^\mu)\\[6pt]
&\times\frac{i(\not\ell+m_f)}{\ell^2-m_f^2+i\varepsilon}\,(ieQ_f\gamma^\rho)\\[6pt]
&\times\frac{-i}{(p-\ell)^2+i\varepsilon}\left[\eta_{\nu\rho}-(1-\xi)\frac{(p-\ell)_\nu(p-\ell)_\rho}{(p-\ell)^2+i\varepsilon}\right].
\end{aligned}
\tag{7.25}\label{eq:7.25}
\end{equation*}
For a quark, \eqref{eq:7.25} multiplies the identity in color.

In $d=4-\epsilon$ dimensions, continue as for \eqref{eq:7.8}, with the stripped coupling $ieQ_f$ also carrying $\mu^{\epsilon/2}$. Collecting $(ieQ_f)^3\,i\cdot i\cdot(-i)=(ieQ_f)(-ie^2Q_f^2)$,
\begin{equation*}
\begin{aligned}
\Lambda_2^\mu(p',p)
={}&-i\mu^{\epsilon}e^2Q_f^2\int\frac{d^d\ell}{(2\pi)^d}\biggl[
\frac{\gamma^\nu(\not\ell'+m_f)\gamma^\mu(\not\ell+m_f)\gamma_\nu}{\mathcal P_1\mathcal P_1'\mathcal P_\gamma}\\[6pt]
&-(1-\xi)\,\frac{\slashed{k}(\not\ell'+m_f)\gamma^\mu(\not\ell+m_f)\slashed{k}}{\mathcal P_1\mathcal P_1'\mathcal P_\gamma^2}\biggr],\\[6pt]
\ell':={}&\ell+q,\qquad k:=p-\ell,\\[6pt]
\mathcal P_1:={}&\ell^2-m_f^2+i\varepsilon,\qquad
\mathcal P_1':=\ell'^2-m_f^2+i\varepsilon,\qquad
\mathcal P_\gamma:=k^2+i\varepsilon .
\end{aligned}
\tag{7.26}\label{eq:7.26}
\end{equation*}
Combine the denominators with Feynman parameters $x,y,z$:
\begin{equation*}
\begin{aligned}
\frac{1}{\mathcal P_1\mathcal P_1'\mathcal P_\gamma}
&=2\int dF\,\frac{1}{\bigl[x\mathcal P_1+y\mathcal P_1'+z\mathcal P_\gamma\bigr]^3},\\[6pt]
\frac{1}{\mathcal P_1\mathcal P_1'\mathcal P_\gamma^2}
&=6\int dF\,\frac{z}{\bigl[x\mathcal P_1+y\mathcal P_1'+z\mathcal P_\gamma\bigr]^4},\\[6pt]
\int dF&:=\int_0^1dx\int_0^1dy\int_0^1dz\;\delta(x+y+z-1),\\[6pt]
x\mathcal P_1+y\mathcal P_1'+z\mathcal P_\gamma
&=(\ell+yq-zp)^2-\Delta_{V,f}+i\varepsilon,\\[6pt]
\Delta_{V,f}&:=(1-z)m_f^2-xz\,p^2-yz\,p'^2-xy\,q^2 .
\end{aligned}
\tag{7.27}\label{eq:7.27}
\end{equation*}
Shift $\ell\to\ell+A$, so that $\ell'\to\ell+A'$ and $k\to K-\ell$, with $A:=(1-x)p-yp'$, $A':=(1-y)p'-xp$ and $K:=xp+yp'$. Drop the terms odd in $\ell$ and replace $\ell^\alpha\ell^\beta$ by $\eta^{\alpha\beta}\ell^2/d$; with $\not\ell\not\ell=\ell^2$, the numerators of \eqref{eq:7.26} become
\begin{equation*}
\begin{aligned}
\gamma^\nu(\not\ell'+m_f)\gamma^\mu(\not\ell+m_f)\gamma_\nu\ &\to\ \frac{(d-2)^2}{d}\,\ell^2\gamma^\mu+N^\mu,\\[4pt]
\slashed{k}(\not\ell'+m_f)\gamma^\mu(\not\ell+m_f)\slashed{k}\ &\to\ (\ell^2)^2\gamma^\mu+\ell^2N_2^\mu+N_0^\mu,\\[10pt]
N^\mu&:=\gamma^\nu(\slashed{A}'+m_f)\gamma^\mu(\slashed{A}+m_f)\gamma_\nu,\\[6pt]
N_0^\mu&:=\slashed{K}(\slashed{A}'+m_f)\gamma^\mu(\slashed{A}+m_f)\slashed{K},\\[6pt]
N_2^\mu&:=-\gamma^\mu(\slashed{A}+m_f)\slashed{K}-\slashed{K}(\slashed{A}'+m_f)\gamma^\mu\\
&\quad\ +\frac1d\Bigl[N^\mu-(d-2)\slashed{K}\gamma^\mu\slashed{K}\\
&\qquad\quad-\gamma^\nu(\slashed{A}'+m_f)\gamma^\mu\gamma_\nu\slashed{K}-\slashed{K}\gamma^\nu\gamma^\mu(\slashed{A}+m_f)\gamma_\nu\Bigr],
\end{aligned}
\tag{7.28}\label{eq:7.28}
\end{equation*}
so that
\begin{equation*}
\begin{aligned}
\Lambda_2^\mu(p',p)
={}&-i\mu^{\epsilon}e^2Q_f^2\int dF\int\frac{d^d\ell}{(2\pi)^d}\,\biggl\{
\frac{2\bigl[\frac{(d-2)^2}{d}\,\ell^2\gamma^\mu+N^\mu\bigr]}{\bigl(\ell^2-\Delta_{V,f}+i\varepsilon\bigr)^3}\\[6pt]
&-6(1-\xi)\,z\,\frac{(\ell^2)^2\gamma^\mu+\ell^2N_2^\mu+N_0^\mu}{\bigl(\ell^2-\Delta_{V,f}+i\varepsilon\bigr)^4}\biggr\}.
\end{aligned}
\tag{7.29}\label{eq:7.29}
\end{equation*}
After Wick rotation, for $\Delta>0$ and continued to $\Delta<0$ through $\Delta-i\varepsilon$, the integrals \eqref{eq:7.12} generalize to
\begin{equation*}
\int\frac{d^d\ell}{(2\pi)^d}\,\frac{(\ell^2)^a}{(\ell^2-\Delta+i\varepsilon)^n}
=\frac{(-1)^{n+a}\,i}{(4\pi)^{d/2}}\,
\frac{\Gamma\bigl(a+\frac d2\bigr)\,\Gamma\bigl(n-a-\frac d2\bigr)}{\Gamma\bigl(\frac d2\bigr)\,\Gamma(n)}\,
(\Delta-i\varepsilon)^{\frac d2+a-n}.
\tag{7.30}\label{eq:7.30}
\end{equation*}
Applying \eqref{eq:7.30} to \eqref{eq:7.29},
\begin{equation*}
\begin{aligned}
\Lambda_2^\mu(p',p)
={}&\frac{\mu^{\epsilon}e^2Q_f^2}{(4\pi)^{d/2}}\int dF\,\biggl\{
\Bigl[\frac{(d-2)^2}{2}-(1-\xi)\,z\,\frac{d(d+2)}{4}\Bigr]
\Gamma\Bigl(2-\frac d2\Bigr)(\Delta_{V,f}-i\varepsilon)^{\frac d2-2}\,\gamma^\mu\\[6pt]
&-\Gamma\Bigl(3-\frac d2\Bigr)(\Delta_{V,f}-i\varepsilon)^{\frac d2-3}
\Bigl[N^\mu-(1-\xi)\,z\,\frac d2\,N_2^\mu\Bigr]\\[6pt]
&-(1-\xi)\,z\,\Gamma\Bigl(4-\frac d2\Bigr)(\Delta_{V,f}-i\varepsilon)^{\frac d2-4}\,N_0^\mu\biggr\}.
\end{aligned}
\tag{7.31}\label{eq:7.31}
\end{equation*}
With $d=4-\epsilon$, $\Gamma(\epsilon/2)=2/\epsilon-\gamma_E+O(\epsilon)$, $2/\bar\epsilon$ as in \eqref{eq:7.14}, and $N^\mu$, $N_2^\mu$ in the finite terms taken at $d=4$,
\begin{equation*}
{\displaystyle 
\begin{aligned}
\Lambda_2^\mu(p',p)={}&\frac{e^2Q_f^2}{16\pi^2}\int dF\,\biggl\{
\bigl[2-6(1-\xi)z\bigr]\gamma^\mu\left[\frac2{\bar\epsilon}-\ln\frac{\Delta_{V,f}-i\varepsilon}{\mu^2}\right]
-\bigl[4-5(1-\xi)z\bigr]\gamma^\mu\\[6pt]
&-\frac{N^\mu-2(1-\xi)z\,N_2^\mu}{\Delta_{V,f}-i\varepsilon}
-\frac{(1-\xi)z\,N_0^\mu}{(\Delta_{V,f}-i\varepsilon)^2}\biggr\}+O(\epsilon)\\[6pt]
={}&\frac{e^2Q_f^2}{16\pi^2}\,\frac2\epsilon\,\xi\gamma^\mu+(\text{finite as }\epsilon\to0),\\[10pt]
N^\mu\big|_{d=4}={}&-2\slashed{A}\gamma^\mu\slashed{A}'+4m_f(A+A')^\mu-2m_f^2\gamma^\mu,\\[6pt]
N_2^\mu\big|_{d=4}={}&\frac14N^\mu\big|_{d=4}-\gamma^\mu\slashed{A}\slashed{K}-\slashed{K}\slashed{A}'\gamma^\mu-\frac12\slashed{K}\gamma^\mu\slashed{K}-(A+A')^\mu\slashed{K}-m_fK^\mu .
\end{aligned}
}
\tag{7.32}\label{eq:7.32}
\end{equation*}
The parameter integrals converge unless $p^2=p'^2=m_f^2$; in that on-shell case they diverge logarithmically at $z\to1$. With \eqref{eq:7.14}, the results satisfy $q_\mu\Lambda_2^\mu(p+q,p)=\Sigma_2(p)-\Sigma_2(p+q)$, as required by the Ward identity.

\topic{(2) Three Photons}
\begin{equation*}
E_\gamma=3,\qquad E_F=0\qquad\Longrightarrow\qquad\omega=1 .
\tag{7.33}\label{eq:7.33}
\end{equation*}
The graphs with the two orientations of the fermion loop cancel (Furry's theorem), so the sum of 1PI graphs with three external photons vanishes.

\setcounter{chapter}{7}
\chapter{Cosmology}\label{book-chapter-8}\label{qft-ch8}

\setcounter{section}{0}
\section{Geometry of the Universe}\label{book-section-8.1}\label{geometry-of-the-universe}

\subsection{Experimental Facts}\label{book-subsection-section-8.1-1}\label{cosmology-experimental-facts}
\begin{enumerate}


\item Hubble's law: the universe is expanding.

\item The cosmic microwave background (CMB) is isotropic.

\item The large-scale distribution of galaxies is homogeneous.
\end{enumerate}


\subsection{FLRW Metric}\label{book-subsection-section-8.1-2}\label{flrw-metric}
\begin{equation*}
ds^2=dt^2-a(t)^2\left[\frac{dr^2}{1-kr^2}+r^2\bigl(d\theta^2+\sin^2\theta\,d\phi^2\bigr)\right]
\tag{8.1}\label{eq:8.1}
\end{equation*}

\setcounter{section}{1}
\section{Evolution of the Universe}\label{book-section-8.2}\label{evolution-of-the-universe}

\subsection{Dynamics}\label{book-subsection-section-8.2-1}\label{cosmology-dynamics}
\begin{equation*}
\left(\frac{\dot a}{a}\right)^2=\frac{8\pi G}{3}\rho-\frac{k}{a^2}
\tag{8.2}\label{eq:8.2}
\end{equation*}

\subsection{Dark Matter}\label{book-subsection-section-8.2-2}\label{dark-matter}

\topic{Experimental Results}

\subsection{Dark Energy}\label{book-subsection-section-8.2-3}\label{dark-energy}

\subsection{Cosmic Evolution}\label{book-subsection-section-8.2-4}\label{cosmic-evolution}
\begin{enumerate}


\item Radiation-dominated era (the early universe)

\item Matter-dominated era

\item Dark-energy-dominated era
\end{enumerate}


\setcounter{section}{2}
\section{Boltzmann Equation}\label{book-section-8.3}\label{boltzmann-equation}

\appendix
\renewcommand{\chaptermark}[1]{\markboth{\appendixname\ \thechapter\enspace #1}{}}
\setcounter{chapter}{0}
\chapter{Summary of Feynman Rules}\label{book-chapter-A}

\setcounter{section}{0}
\section{Feynman Diagrams of the Standard Model}\label{book-section-A.1}\label{sm-feynman-rules}

\topic{Scope and Gauge}
\begin{itemize}


\item \textbf{\boldmath Theory}: the minimal Standard Model of \hyperref[electroweak-theory]{Section 5.3} and \hyperref[qcd]{Section 5.4}, with gauge group $SU(3)_c\times SU(2)_L\times U(1)_Y$, three generations of quarks and leptons in the mass bases of Section 5.3, and the Higgs doublet $H$. There are no right-handed neutrinos, so the neutrinos are massless; the seesaw fields of \hyperref[qft-ch6]{Chapter 6} are not included.

\item \textbf{\boldmath Electroweak gauge}: the unitary gauge $H=(0,v_{\mathrm{SM}}+h)^T/\sqrt2$ of Section 5.3, with $v_{\mathrm{SM}}\simeq246\,\mathrm{GeV}$. The electroweak bosons are $W^\pm_\mu$, $Z_\mu$, $A^{\mathrm{em}}_\mu$ and $h$; there are no Goldstone fields.

\item \textbf{\boldmath Photon and gluons}: the covariant gauge fixing of \hyperref[feynman-diagrams]{Section 4.3}, with one parameter $\xi$ for $A^{\mathrm{em}}_\mu$ and for the gluons $A^a_\mu$ of Section 5.4 (separate parameters may be used),
\begin{equation*}
\mathcal L_{\mathrm{gf}}=-\frac1{2\xi}\bigl(\partial^\mu A^{\mathrm{em}}_\mu\bigr)^2-\frac1{2\xi}\bigl(\partial^\mu A^a_\mu\bigr)^2 .
\tag{A.1}\label{eq:A.1}
\end{equation*}

\item \textbf{\boldmath Order}: tree level, with the tree-level masses of \hyperref[sm-fermions]{Section 5.2} and Section 5.3.
\begin{itemize}


\item Tree graphs contain no closed loops, so the Faddeev--Popov ghosts of \hyperref[path-integrals-in-field-theory]{Section 3.2} and the counterterms of \hyperref[qft-ch7]{Chapter 7} do not enter.

\item These rules are not a loop-ready set: loop calculations in a renormalizable gauge need additional Goldstone, ghost and counterterm rules, which are not given here.

\item Widths are not included. The propagators below are the fixed-order tree-level ones; predictions near the $W^\pm$, $Z$, $h$ or $t$ resonance require an appropriate treatment of the width, such as a resummation of the self-energy.
\end{itemize}

\end{itemize}


These choices hold only in this subsection.

\topic{Fields and Free Lagrangian}
\begin{itemize}


\item \textbf{\boldmath Fermions}: $\psi_f$ as in the neutral currents of Section 5.3: $\psi_f=q_f$ for $f=u,d,s,c,b,t$ (color triplets), $\psi_f=l_f$ for $f=e,\mu,\tau$, and $\psi_{\nu_\alpha}=\nu_{\alpha L}$. Here $m_f$ is the mass of $\psi_f$, with $m_{\nu_\alpha}=0$; $Q_f$ is the charge of \hyperref[qed]{Section 5.3}, and $T_f^3$ is that of Section 5.3.

\item \textbf{\boldmath Vector fields}: $A^{\mathrm{em}}_\mu$, $Z_\mu$, $W^\pm_\mu$ of Section 5.3 and the gluons $A^a_\mu$, $a=1,\ldots,8$. For each of them write $X_{\mu\nu}:=\partial_\mu X_\nu-\partial_\nu X_\mu$.

\item \textbf{\boldmath Scalar}: the Higgs field $h$.
\end{itemize}


Inverting the definitions of $W^\pm_\mu$, $Z_\mu$ and $A^{\mathrm{em}}_\mu$ in Section 5.3,
\begin{equation*}
\begin{aligned}
W^1_\mu&=\frac{W^+_\mu+W^-_\mu}{\sqrt2},&\qquad 
W^2_\mu&=\frac{i\bigl(W^+_\mu-W^-_\mu\bigr)}{\sqrt2},\\[3pt]
W^3_\mu&=c_WZ_\mu+s_WA^{\mathrm{em}}_\mu,&\qquad 
B_\mu&=-s_WZ_\mu+c_WA^{\mathrm{em}}_\mu .
\end{aligned}
\tag{A.2}\label{eq:A.2}
\end{equation*}
Since $W^1_\mu W^{1\nu}+W^2_\mu W^{2\nu}=W^+_\mu W^{-\nu}+W^-_\mu W^{+\nu}$ and $(W^3_\mu,B_\mu)\mapsto(Z_\mu,A^{\mathrm{em}}_\mu)$ is a rotation, the quadratic parts of the gauge kinetic terms, the mass terms of the Gauge-Boson Masses part of Section 5.3, the Higgs mass of Section 5.2 and $\mathcal L_{\mathrm{gf}}$ give
\begin{equation*}
\begin{aligned}
\mathcal L_0={}&\sum_f\bar\psi_f\bigl(i\gamma^\mu\partial_\mu-m_f\bigr)\psi_f
-\frac14A^{\mathrm{em}}_{\mu\nu}A^{\mathrm{em}\,\mu\nu}-\frac1{2\xi}\bigl(\partial^\mu A^{\mathrm{em}}_\mu\bigr)^2
-\frac14A^a_{\mu\nu}A^{a\mu\nu}-\frac1{2\xi}\bigl(\partial^\mu A^a_\mu\bigr)^2\\
&-\frac12W^+_{\mu\nu}W^{-\mu\nu}+m_W^2W^+_\mu W^{-\mu}
-\frac14Z_{\mu\nu}Z^{\mu\nu}+\frac12m_Z^2Z_\mu Z^\mu
+\frac12\partial_\mu h\,\partial^\mu h-\frac12m_h^2h^2 .
\end{aligned}
\tag{A.3}\label{eq:A.3}
\end{equation*}

\topic{Interaction Lagrangian}

Write $\mathcal L_{\mathrm{SM}}+\mathcal L_{\mathrm{gf}}=\mathcal L_0+\mathcal L_{\mathrm{int}}$. From Sections 5.2--5.4, with $\lambda_H=m_h^2/(2v_{\mathrm{SM}}^2)$ and the constant $\lambda_Hv_{\mathrm{SM}}^4/4$ dropped,
\begin{equation*}
\begin{aligned}
\mathcal L_{\mathrm{int}}={}&\mathcal L_{\mathrm{NC}}+\mathcal L_{\mathrm{CC}}
+g_sA^a_\mu\sum_{f}\bar q_f\gamma^\mu T^aq_f
-\sum_f\frac{m_f}{v_{\mathrm{SM}}}\,h\,\bar\psi_f\psi_f\\
&+\Bigl(\frac{2h}{v_{\mathrm{SM}}}+\frac{h^2}{v_{\mathrm{SM}}^2}\Bigr)\Bigl(m_W^2W^+_\mu W^{-\mu}+\frac12m_Z^2Z_\mu Z^\mu\Bigr)
-\frac{m_h^2}{2v_{\mathrm{SM}}}h^3-\frac{m_h^2}{8v_{\mathrm{SM}}^2}h^4\\
&+\mathcal L_{3g}+\mathcal L_{4g}+\mathcal L_{3W}+\mathcal L_{4W},
\end{aligned}
\tag{A.4}\label{eq:A.4}
\end{equation*}
where $\mathcal L_{\mathrm{NC}}$ and $\mathcal L_{\mathrm{CC}}$ are the currents of Section 5.3, and the cubic and quartic parts of $-\frac14F^a_{\mu\nu}F^{a\mu\nu}$ in Section 5.4 are
\begin{equation*}
\mathcal L_{3g}=-g_sf^{abc}\bigl(\partial_\mu A^a_\nu\bigr)A^{b\mu}A^{c\nu},\qquad
\mathcal L_{4g}=-\frac{g_s^2}4f^{abc}f^{ade}A^b_\mu A^c_\nu A^{d\mu}A^{e\nu},
\tag{A.5}\label{eq:A.5}
\end{equation*}
using the antisymmetry of $f^{abc}A^{b\mu}A^{c\nu}$ in $\mu,\nu$. Since $W^a_{\mu\nu}$ of Section 5.3 has the form of $F^a_{\mu\nu}$ with $g_s\to g_2$ and $f^{abc}\to\epsilon^{abc}$, $\mathcal L_{3W}$ and $\mathcal L_{4W}$ are $\mathcal L_{3g}$ and $\mathcal L_{4g}$ with these replacements, $A^a_\mu\to W^a_\mu$ and $a=1,2,3$; $B_\mu$ has no self-interactions.

\topic{From the Lagrangian to the Rules}

By \hyperref[interaction-picture]{Section 4.1} and \hyperref[s-matrix-and-wick-theorem]{Section 4.2}, for interactions without derivatives $\hat{\mathcal H}_{\mathrm{int}}=-\hat{\mathcal L}_{\mathrm{int}}$, and the order-$n$ term of $\hat S$ is
\begin{equation*}
\frac{i^n}{n!}\int d^4x_1\cdots d^4x_n\,\mathcal T\bigl[\hat{\mathcal L}_{\mathrm{int}}(x_1)\cdots\hat{\mathcal L}_{\mathrm{int}}(x_n)\bigr].
\tag{A.6}\label{eq:A.6}
\end{equation*}
Wick's theorem pairs the fields: a contraction between two factors gives a propagator of $\mathcal L_0$, a contraction with an external state gives an external factor, and each factor $i\int d^4x\,\hat{\mathcal L}_{\mathrm{int}}$ gives a vertex, its $x$-integral imposing four-momentum conservation. The $1/n!$ cancels against the permutations of the $n$ vertices.

$\mathcal L_{3g}$ and $\mathcal L_{3W}$ contain derivatives, and the massive vector fields are constrained (Section 1.8), so $\hat{\mathcal H}_{\mathrm{int}}\neq-\hat{\mathcal L}_{\mathrm{int}}$ and the time-ordered products of $\partial_\mu\hat A_\nu$ and of massive vector fields contain noncovariant contact terms. At tree level these two effects cancel (not derived here). The rules below are the covariant ones, which also follow from the functional integral of \hyperref[qft-ch3]{Chapter 3} with weight $e^{iS}$.

\topic{Vertex Rule}
\begin{itemize}


\item \textbf{\boldmath Legs}: leg $j=1,\ldots,n$ of a vertex sits on a boson field of $\mathcal L_{\mathrm{int}}$ ($h$, $A^{\mathrm{em}}_\mu$, $Z_\mu$, $W^+_\mu$, $W^-_\mu$ or $A^a_\mu$), with four-momentum $k_j$ flowing into the vertex and wave function $\varepsilon_j$: a vector $\varepsilon_j^\mu$, and $1$ for $h$. A gluon leg also carries an adjoint index $a_j$.

\item \textbf{\boldmath Vertex factor}: replace each boson field by the plane waves of its legs,
\begin{equation*}
\phi(x)\to\sum_{j\text{ on }\phi}c_j\,\varepsilon_j\,e^{-ik_j\cdot x},\qquad
A^a_\mu(x)\to\sum_{j\text{ on }A}c_j\,\delta^{aa_j}\varepsilon_{j\mu}\,e^{-ik_j\cdot x},
\tag{A.7}\label{eq:A.7}
\end{equation*}
so that $\partial_\mu\to-ik_{j\mu}$ on leg $j$, and keep the term proportional to $c_1\cdots c_n$:
\begin{equation*}
{\displaystyle i\int d^4x\,\mathcal L_{\mathrm{int}}\Big|_{c_1c_2\cdots c_n}=(2\pi)^4\delta^{(4)}\Bigl(\sum_{j=1}^nk_j\Bigr)\,i\mathcal V(\varepsilon_1,\ldots,\varepsilon_n).}
\tag{A.8}\label{eq:A.8}
\end{equation*}
The vertex factor is the tensor $i\mathcal V$, obtained by removing the $\varepsilon_j$; purely bosonic vertices below are written with lower Lorentz indices. The term linear in every $c_j$ sums over all ways of attaching the legs to identical fields, which gives the combinatorial factors of identical fields.

\item \textbf{\boldmath Fermion bilinears}: a term $\bar\psi_{f'}\Gamma\psi_f$, with $\Gamma$ a matrix in spinor and color space multiplying boson fields, gives the matrix $i\Gamma$, with its boson legs as above; the fermion line enters the vertex as $\psi_f$ and leaves it as $\psi_{f'}$.
\end{itemize}


\topic{Charged and Colored Lines}
\begin{itemize}


\item \textbf{\boldmath Charge of $W^\pm$}: for $\theta^3=\theta_Y=\theta$ constant and $\theta^1=\theta^2=0$, the transformation of Section 5.3 is $e^{i\theta Q}$. The infinitesimal transformation of the gauge field in Section 5.4, with $f^{abc}\to\epsilon^{abc}$, gives $\delta W^1_\mu=\theta W^2_\mu$ and $\delta W^2_\mu=-\theta W^1_\mu$, so
\begin{equation*}
\delta W^\pm_\mu=\pm i\theta\,W^\pm_\mu .
\tag{A.9}\label{eq:A.9}
\end{equation*}
Thus $W^+_\mu$ has charge $+e$ in the sense of Section 5.3, and, like a field $\psi_f$ of charge $Q_fe$, it annihilates the particle of that charge: a leg on $W^+_\mu$ is an incoming $W^+$ or an outgoing $W^-$, and a leg on $W^-_\mu$ is an incoming $W^-$ or an outgoing $W^+$.

\item \textbf{\boldmath Internal $W$ lines}: an internal $W$ line joins a $W^+_\mu$ leg of one vertex to a $W^-_\nu$ leg of another. If $k$ flows into the first vertex through its $W^+_\mu$ leg, then $-k$ flows into the second through its $W^-_\nu$ leg.

\item \textbf{\boldmath Fermion lines}: arrows show the flow of fermion number, and the momentum of a fermion propagator points along the arrow. Along each line, write the factors in the order opposite to the arrow, as in \hyperref[qed-renormalization]{Section 7.2}. For an external antifermion the arrow points opposite to its motion.

\item \textbf{\boldmath Color}: quark colors are $c,c'\in\{1,2,3\}$, as in Section 7.2. In $(T^a)_{c'c}$, $c$ is the color of the quark line entering the vertex and $c'$ that of the line leaving it; all other vertices on quark lines carry $\delta_{c'c}$.
\end{itemize}


\topic{Propagators}

The contractions of the fields of $\mathcal L_0$, Fourier transformed as in \hyperref[fourier-transform]{Section B.1}, are those of Section 4.3, with $\delta^{ab}$ for gluons. In $\mathcal L_0$, $W^1_\mu$ and $W^2_\mu$ are real Proca fields of mass $m_W$, so
\begin{equation*}
\langle0|\mathcal T[\hat W^+_\mu\hat W^-_\nu]|0\rangle=\frac12\Bigl(\langle0|\mathcal T[\hat W^1_\mu\hat W^1_\nu]|0\rangle+\langle0|\mathcal T[\hat W^2_\mu\hat W^2_\nu]|0\rangle\Bigr),\qquad
\langle0|\mathcal T[\hat W^\pm_\mu\hat W^\pm_\nu]|0\rangle=0 .
\tag{A.10}\label{eq:A.10}
\end{equation*}
Since $\nu_{\alpha L}=P_L\nu_{\alpha L}$ and $\overline{\nu_{\alpha L}}=\overline{\nu_{\alpha L}}P_R$, the neutrino contraction is $P_L\widetilde S_F(p)P_R$ at $m=0$, and $P_L\slashed{p}P_R=P_L\slashed{p}$. With $p$ the momentum of the line,
\begin{equation*}
\begin{array}{l|l}
\text{line}&\text{factor}\\ \hline
\psi_f,\ f\neq\nu_\alpha,\ p\text{ along the arrow}&\dfrac{i(\slashed{p}+m_f)}{p^2-m_f^2+i\varepsilon}\quad(\times\,\delta_{c'c}\text{ for quarks})\\[12pt]
\nu_{\alpha L},\ p\text{ along the arrow}&\dfrac{iP_L\slashed{p}}{p^2+i\varepsilon}\\[12pt]
A^{\mathrm{em}}_\mu\text{ to }A^{\mathrm{em}}_\nu&\dfrac{-i}{p^2+i\varepsilon}\left[\eta_{\mu\nu}-(1-\xi)\dfrac{p_\mu p_\nu}{p^2+i\varepsilon}\right]\\[12pt]
A^a_\mu\text{ to }A^b_\nu&\delta^{ab}\,\dfrac{-i}{p^2+i\varepsilon}\left[\eta_{\mu\nu}-(1-\xi)\dfrac{p_\mu p_\nu}{p^2+i\varepsilon}\right]\\[12pt]
W^+_\mu\text{ to }W^-_\nu&\dfrac{-i}{p^2-m_W^2+i\varepsilon}\left(\eta_{\mu\nu}-\dfrac{p_\mu p_\nu}{m_W^2}\right)\\[12pt]
Z_\mu\text{ to }Z_\nu&\dfrac{-i}{p^2-m_Z^2+i\varepsilon}\left(\eta_{\mu\nu}-\dfrac{p_\mu p_\nu}{m_Z^2}\right)\\[12pt]
h&\dfrac{i}{p^2-m_h^2+i\varepsilon}
\end{array}
\tag{A.11}\label{eq:A.11}
\end{equation*}

\topic{External Lines}

The fields have unit pole residue at tree level, so by Section 4.3 and the charge assignment above,
\begin{equation*}
\begin{array}{l|ll}
\text{external state}&\text{incoming}&\text{outgoing}\\ \hline
\text{fermion}&u^s(p)&\bar u^s(p)\\
\text{antifermion}&\bar v^s(p)&v^s(p)\\
\gamma,\ g,\ Z&\epsilon^r_\mu(p)&\epsilon^{r*}_\mu(p)\\
W^+&\epsilon^r_\mu(p)\text{ on a }W^+_\mu\text{ leg}&\epsilon^{r*}_\mu(p)\text{ on a }W^-_\mu\text{ leg}\\
W^-&\epsilon^r_\mu(p)\text{ on a }W^-_\mu\text{ leg}&\epsilon^{r*}_\mu(p)\text{ on a }W^+_\mu\text{ leg}\\
h&1&1
\end{array}
\tag{A.12}\label{eq:A.12}
\end{equation*}
External quarks and gluons also carry their color and adjoint indices. $W^\pm$ and $Z$ have the three polarizations of \hyperref[proca-field]{Section 1.8}, and $\gamma$ and $g$ the two physical ones. Neutrinos use the $m=0$ spinors of \hyperref[dirac-field]{Section 1.7}; through the $P_L$ in their vertices, only $P_Lu^s$ and $P_Lv^s$ contribute.

\topic{Vertices: Fermions}

The rows follow from $\mathcal L_{\mathrm{NC}}$, $\mathcal L_{\mathrm{CC}}$, the quark--gluon term and the $h\bar\psi_f\psi_f$ term. Since $P_R\gamma^\mu=\gamma^\mu P_L$, $\overline{q_{fL}}\gamma^\mu q_{f'L}=\bar q_f\gamma^\mu P_Lq_{f'}$, and with $(W^+_\mu)^*=W^-_\mu$ the Hermitian conjugate in $\mathcal L_{\mathrm{CC}}$ is
\begin{equation*}
\frac{g_2}{\sqrt2}\,W^-_\mu\Bigl(\sum_{\alpha=e,\mu,\tau}\bar l_\alpha\gamma^\mu P_L\nu_{\alpha L}
+\sum_{f=u,c,t}\ \sum_{f'=d,s,b}V^*_{ff'}\,\bar q_{f'}\gamma^\mu P_Lq_f\Bigr).
\tag{A.13}\label{eq:A.13}
\end{equation*}
With "in" the fermion line entering the vertex along its arrow, "out" the line leaving it, and the boson leg carrying the index $\mu$,
\begin{equation*}
\begin{array}{lll|l}
\text{in}&\text{out}&\text{boson}&i\mathcal V\\ \hline
\psi_f&\psi_f&A^{\mathrm{em}}_\mu&ieQ_f\gamma^\mu\\[4pt]
\psi_f&\psi_f&Z_\mu&i\dfrac{g_2}{c_W}\gamma^\mu\bigl(T_f^3P_L-Q_fs_W^2\bigr)\\[10pt]
q_f,\ \text{color }c&q_f,\ \text{color }c'&A^a_\mu&ig_s\gamma^\mu(T^a)_{c'c}\\[4pt]
q_{f'},\ f'=d,s,b&q_f,\ f=u,c,t&W^+_\mu&i\dfrac{g_2}{\sqrt2}V_{ff'}\gamma^\mu P_L\\[10pt]
q_f,\ f=u,c,t&q_{f'},\ f'=d,s,b&W^-_\mu&i\dfrac{g_2}{\sqrt2}V^*_{ff'}\gamma^\mu P_L\\[10pt]
l_\alpha&\nu_{\alpha L}&W^+_\mu&i\dfrac{g_2}{\sqrt2}\gamma^\mu P_L\\[10pt]
\nu_{\alpha L}&l_\alpha&W^-_\mu&i\dfrac{g_2}{\sqrt2}\gamma^\mu P_L\\[10pt]
\psi_f&\psi_f&h&-i\dfrac{m_f}{v_{\mathrm{SM}}}
\end{array}
\tag{A.14}\label{eq:A.14}
\end{equation*}
The $Z$ vertex equals $i\frac{g_2}{c_W}\gamma^\mu\bigl[(T_f^3-Q_fs_W^2)P_L-Q_fs_W^2P_R\bigr]$; for $\nu_{\alpha L}$ it is $i\frac{g_2}{2c_W}\gamma^\mu P_L$. The photon vertex of the electron is $-ie\gamma^\mu$, as in Section 7.2. The neutral and $h$ vertices are flavor diagonal, and the neutrinos have no $\gamma$ or $h$ vertex.

\topic{Vertices: Higgs and Gauge Bosons}

The $h$ terms of $\mathcal L_{\mathrm{int}}$ contain no derivatives. Attaching the legs to identical fields gives a factor $2$ for $ZZ$ and for $hh$, $2\cdot2$ for $hhZZ$, $3!$ for $h^3$ and $4!$ for $h^4$. With $m_W=g_2v_{\mathrm{SM}}/2$ and $m_Z=m_W/c_W$,
\begin{equation*}
\begin{array}{l|l}
\text{legs}&i\mathcal V\\ \hline
h\,W^+_\mu W^-_\nu&i\dfrac{2m_W^2}{v_{\mathrm{SM}}}\eta_{\mu\nu}=ig_2m_W\eta_{\mu\nu}\\[10pt]
h\,Z_\mu Z_\nu&i\dfrac{2m_Z^2}{v_{\mathrm{SM}}}\eta_{\mu\nu}=i\dfrac{g_2m_Z}{c_W}\eta_{\mu\nu}\\[10pt]
h\,h\,W^+_\mu W^-_\nu&i\dfrac{2m_W^2}{v_{\mathrm{SM}}^2}\eta_{\mu\nu}=i\dfrac{g_2^2}{2}\eta_{\mu\nu}\\[10pt]
h\,h\,Z_\mu Z_\nu&i\dfrac{2m_Z^2}{v_{\mathrm{SM}}^2}\eta_{\mu\nu}=i\dfrac{g_2^2}{2c_W^2}\eta_{\mu\nu}\\[10pt]
h\,h\,h&-i\dfrac{3m_h^2}{v_{\mathrm{SM}}}\\[10pt]
h\,h\,h\,h&-i\dfrac{3m_h^2}{v_{\mathrm{SM}}^2}
\end{array}
\tag{A.15}\label{eq:A.15}
\end{equation*}

\topic{Vertices: Gauge-Boson Self-Interactions}

Define
\begin{equation*}
\begin{aligned}
\mathcal Y^{(3)}_{\mu\nu\rho}(k_1,k_2,k_3)&:=\eta_{\mu\nu}(k_1-k_2)_\rho+\eta_{\nu\rho}(k_2-k_3)_\mu+\eta_{\rho\mu}(k_3-k_1)_\nu,\\
\mathcal Y^{(4)}_{\mu\nu,\rho\sigma}&:=2\eta_{\mu\nu}\eta_{\rho\sigma}-\eta_{\mu\rho}\eta_{\nu\sigma}-\eta_{\mu\sigma}\eta_{\nu\rho}.
\end{aligned}
\tag{A.16}\label{eq:A.16}
\end{equation*}

\topic{(1) Gluons}

Label the legs $j=1,2,3,4$ by $(a_1,\mu),(a_2,\nu),(a_3,\rho),(a_4,\sigma)$ with momenta $k_j$ into the vertex. In $\mathcal L_{3g}$ the rule sums over the assignments $\pi\in S_3$ of the legs to the three fields, the first being the differentiated one. By $f^{a_{\pi(1)}a_{\pi(2)}a_{\pi(3)}}=\operatorname{sgn}(\pi)f^{a_1a_2a_3}$,
\begin{equation*}
\begin{aligned}
i\mathcal V(\varepsilon_1,\varepsilon_2,\varepsilon_3)
&=i\,(-g_s)\,(-i)\,f^{a_1a_2a_3}\sum_{\pi\in S_3}\operatorname{sgn}(\pi)\,\bigl(k_{\pi(1)}\cdot\varepsilon_{\pi(2)}\bigr)\bigl(\varepsilon_{\pi(1)}\cdot\varepsilon_{\pi(3)}\bigr)\\
&=g_sf^{a_1a_2a_3}\,\mathcal Y^{(3)}_{\mu\nu\rho}(k_1,k_2,k_3)\,\varepsilon_1^\mu\varepsilon_2^\nu\varepsilon_3^\rho .
\end{aligned}
\tag{A.17}\label{eq:A.17}
\end{equation*}
In $\mathcal L_{4g}$ the $4!$ assignments fall into six classes of four equal terms, since $f^{abc}f^{ade}A^b_\mu A^c_\nu A^{d\mu}A^{e\nu}$ is invariant under $(b,c)\leftrightarrow(d,e)$ and under the simultaneous exchanges $b\leftrightarrow c$, $d\leftrightarrow e$. Writing the summed adjoint index of the result as $b$,
\begin{equation*}
{\displaystyle 
\begin{aligned}
A^{a_1}_\mu A^{a_2}_\nu A^{a_3}_\rho:\quad i\mathcal V&=g_sf^{a_1a_2a_3}\,\mathcal Y^{(3)}_{\mu\nu\rho}(k_1,k_2,k_3),\\
A^{a_1}_\mu A^{a_2}_\nu A^{a_3}_\rho A^{a_4}_\sigma:\quad i\mathcal V&=-ig_s^2\Bigl[f^{a_1a_2b}f^{a_3a_4b}\bigl(\eta_{\mu\rho}\eta_{\nu\sigma}-\eta_{\mu\sigma}\eta_{\nu\rho}\bigr)
+f^{a_1a_3b}f^{a_2a_4b}\bigl(\eta_{\mu\nu}\eta_{\rho\sigma}-\eta_{\mu\sigma}\eta_{\nu\rho}\bigr)\\
&\qquad+f^{a_1a_4b}f^{a_2a_3b}\bigl(\eta_{\mu\nu}\eta_{\rho\sigma}-\eta_{\mu\rho}\eta_{\nu\sigma}\bigr)\Bigr].
\end{aligned}}
\tag{A.18}\label{eq:A.18}
\end{equation*}

\topic{(2) Electroweak Gauge Bosons}

$\mathcal L_{3W}$ and $\mathcal L_{4W}$ give the vertices of (1) with $g_s\to g_2$, $f^{abc}\to\epsilon^{abc}$, for legs on $W^a_\mu$, $a=1,2,3$. By the inverted definitions above, a leg on $W^+_\mu$, $W^-_\mu$, $Z_\mu$ or $A^{\mathrm{em}}_\mu$ is a leg on $W^a_\mu$ with the isospin vector $n_+^a$, $n_-^a$, $c_Wn_3^a$ or $s_Wn_3^a$ in place of $\delta^{aa_j}$, where
\begin{equation*}
n_\pm:=\frac1{\sqrt2}(1,\pm i,0),\qquad n_3:=(0,0,1),\qquad (x\times y)^a:=\epsilon^{abc}x^by^c .
\tag{A.19}\label{eq:A.19}
\end{equation*}
These satisfy, with no complex conjugation in $x^ay^a$,
\begin{equation*}
n_+\times n_-=-i\,n_3,\qquad n_\pm\times n_3=\pm i\,n_\pm,\qquad
n_+^an_-^a=n_3^an_3^a=1,\qquad n_\pm^an_\pm^a=n_\pm^an_3^a=0 .
\tag{A.20}\label{eq:A.20}
\end{equation*}
Hence $\epsilon^{abc}n_+^an_-^bn_3^c=n_+^a(n_-\times n_3)^a=-i$, so the three-boson vertex is $-ig_2$ times $c_W$ or $s_W$, with $g_2s_W=e$. Let $x_j$ be the isospin vector of leg $j$. In the four-boson vertex the three terms carry $(x_1\times x_2)^b(x_3\times x_4)^b$, $(x_1\times x_3)^b(x_2\times x_4)^b$ and $(x_1\times x_4)^b(x_2\times x_3)^b$. For legs $W^+W^+W^-W^-$, in this order, the first vanishes and the other two equal $(n_+\times n_-)^b(n_+\times n_-)^b=-1$. For legs $W^+W^-$ followed by two neutral legs, the first vanishes and the other two equal $(n_+\times n_3)^b(n_-\times n_3)^b=1$, times $c_W^2$, $s_W^2$ or $c_Ws_W$. With all momenta into the vertex,
\begin{equation*}
\begin{array}{l|l}
\text{legs}&i\mathcal V\\ \hline
W^+_\mu(k_1)\,W^-_\nu(k_2)\,A^{\mathrm{em}}_\rho(k_3)&-ie\,\mathcal Y^{(3)}_{\mu\nu\rho}(k_1,k_2,k_3)\\[4pt]
W^+_\mu(k_1)\,W^-_\nu(k_2)\,Z_\rho(k_3)&-ig_2c_W\,\mathcal Y^{(3)}_{\mu\nu\rho}(k_1,k_2,k_3)\\[4pt]
W^+_\mu\,W^+_\nu\,W^-_\rho\,W^-_\sigma&ig_2^2\,\mathcal Y^{(4)}_{\mu\nu,\rho\sigma}\\[4pt]
W^+_\mu\,W^-_\nu\,Z_\rho\,Z_\sigma&-ig_2^2c_W^2\,\mathcal Y^{(4)}_{\mu\nu,\rho\sigma}\\[4pt]
W^+_\mu\,W^-_\nu\,A^{\mathrm{em}}_\rho\,A^{\mathrm{em}}_\sigma&-ie^2\,\mathcal Y^{(4)}_{\mu\nu,\rho\sigma}\\[4pt]
W^+_\mu\,W^-_\nu\,Z_\rho\,A^{\mathrm{em}}_\sigma&-ieg_2c_W\,\mathcal Y^{(4)}_{\mu\nu,\rho\sigma}
\end{array}
\tag{A.21}\label{eq:A.21}
\end{equation*}
For example, for an incoming $W^+$ of momentum $p$ and an outgoing $W^+$ of momentum $p'$ at a $W^+W^-A^{\mathrm{em}}$ vertex, the $W^+_\mu$ leg has $k_1=p$ and $\varepsilon_1=\epsilon(p)$, and the $W^-_\nu$ leg has $k_2=-p'$ and $\varepsilon_2=\epsilon^*(p')$.

These tables exhaust $\mathcal L_{\mathrm{int}}$. In particular, at tree level no vertex contains only neutral gauge bosons, and $h$ couples neither to $\gamma$ nor to gluons.

\topic{Assembling Diagrams}
\begin{enumerate}


\item Draw every connected tree graph with the given external lines and the vertices above; each internal $W$ line joins a $W^+$ leg to a $W^-$ leg.

\item Fix the internal momenta by four-momentum conservation at each vertex.

\item Multiply the vertex factors, propagators and external factors. Lorentz indices are raised and lowered with $\eta$, so that each repeated Lorentz index occurs once up and once down; e.g. an external vector on the leg $\mu$ of a vertex gives $i\mathcal V_{\mu\cdots}\,\epsilon^{r\mu}(p)$ with $\epsilon^{r\mu}=\eta^{\mu\nu}\epsilon^r_\nu$, and a vector propagator $D_{\alpha\beta}$ of the table joins the vertex factors $i\mathcal V_{\cdots\mu}$ and $i\mathcal V'_{\nu\cdots}$ of its two ends through $\eta^{\mu\alpha}D_{\alpha\beta}\eta^{\beta\nu}$. Contract the spinor indices along each fermion line against the arrow, the colors along each quark line, and the adjoint indices of internal gluons with $\delta^{ab}$.

\item Exchanging identical external fermions changes the sign, as in Section 4.3. Tree graphs have symmetry factor $1$, because the factors from identical fields are already in the vertices.

\item The sum is $i\mathcal M_{fi}$ of Section 4.2. Cross sections and decay rates follow from \hyperref[cross-sections-and-decay-rates]{Section 4.6}, with $C_{\mathrm{id}}$ for identical final particles.
\end{enumerate}


\topic{Example: $t\to bW^+$}

Let the $t$ have momentum $p$, spin $s$, color $c$, the $b$ momentum $p'$, spin $s'$, color $c'$, and the $W^+$ momentum $k$, polarization $r$, with $p=p'+k$. The outgoing $W^+$ sits on a $W^-_\mu$ leg, and the line enters as $q_t$ and leaves as $q_b$, so
\begin{equation*}
i\mathcal M=\bar u^{s'}(p')\Bigl(i\frac{g_2}{\sqrt2}V^*_{tb}\gamma^\mu P_L\Bigr)u^s(p)\,\epsilon^{r*}_\mu(k)\,\delta_{c'c} .
\tag{A.22}\label{eq:A.22}
\end{equation*}

\setcounter{chapter}{1}
\chapter{Mathematical Foundations}\label{book-chapter-B}\label{qft-ch0}

\setcounter{section}{0}
\section{Fourier Transform}\label{book-section-B.1}\label{fourier-transform}

We use the following convention for the four-dimensional Fourier transform, with the factor $(2\pi)^{-4}$ included in the inverse transform:
\begin{equation*}
\begin{aligned}
\widetilde f(p)&=\int d^4x\,e^{+ip\cdot x}f(x),\\
f(x)&=\int\frac{d^4p}{(2\pi)^4}\,e^{-ip\cdot x}\widetilde f(p).
\end{aligned}
\tag{B.1}\label{eq:B.1}
\end{equation*}
Here $x^\mu=(t,\mathbf x)$ and $p^\mu=(p^0,\mathbf p)$, and the phase is given by
\begin{equation*}
p\cdot x=p_\mu x^\mu=p^0t-\mathbf p\cdot\mathbf x.
\tag{B.2}\label{eq:B.2}
\end{equation*}
The Fourier representations of the Dirac $\delta$ function are
\begin{equation*}
\begin{aligned}
\int\frac{d^4p}{(2\pi)^4}\,e^{-ip\cdot(x-y)}&=\delta^{(4)}(x-y),\\[6pt]
\int d^4x\,e^{i(p-q)\cdot x}&=(2\pi)^4\delta^{(4)}(p-q).
\end{aligned}
\tag{B.3}\label{eq:B.3}
\end{equation*}

\setcounter{section}{1}
\section{\texorpdfstring{$\gamma$}{γ} Matrices: Weyl Representation}\label{book-section-B.2}\label{gamma-matrices}

The Pauli matrices are
\begin{equation*}
\sigma^1=\begin{pmatrix}0&1\\1&0\end{pmatrix},\qquad
\sigma^2=\begin{pmatrix}0&-i\\i&0\end{pmatrix},\qquad
\sigma^3=\begin{pmatrix}1&0\\0&-1\end{pmatrix}.
\tag{B.4}\label{eq:B.4}
\end{equation*}

\topic{Properties of the Pauli Matrices}

For each $i=1,2,3$,
\begin{equation*}
(\sigma^i)^\dagger=\sigma^i,\qquad
(\sigma^i)^2=I_2,\qquad
(\sigma^i)^{-1}=\sigma^i,\qquad
\operatorname{tr}\sigma^i=0,\qquad
\det\sigma^i=-1.
\tag{B.5}\label{eq:B.5}
\end{equation*}
Thus, each Pauli matrix is both Hermitian and unitary, with eigenvalues $+1$ and $-1$.

The product identity is
\begin{equation*}
\sigma^i\sigma^j=\delta^{ij}I_2+i\epsilon^{ijk}\sigma^k.
\tag{B.6}\label{eq:B.6}
\end{equation*}
This gives the commutation and anticommutation relations, as well as the trace orthogonality relation:
\begin{equation*}
[\sigma^i,\sigma^j]=2i\epsilon^{ijk}\sigma^k,\qquad
\{\sigma^i,\sigma^j\}=2\delta^{ij}I_2,\qquad
\operatorname{tr}(\sigma^i\sigma^j)=2\delta^{ij}.
\tag{B.7}\label{eq:B.7}
\end{equation*}
For any $\mathbf a,\mathbf b\in\mathbb R^3$, the product identity can also be written as
\begin{equation*}
(\mathbf a\cdot\boldsymbol\sigma)(\mathbf b\cdot\boldsymbol\sigma)
=(\mathbf a\cdot\mathbf b)I_2
+i(\mathbf a\times\mathbf b)\cdot\boldsymbol\sigma,
\qquad \boldsymbol\sigma=(\sigma^1,\sigma^2,\sigma^3).
\tag{B.8}\label{eq:B.8}
\end{equation*}

\topic{$\gamma$ matrices}

Pauli four-vectors
\begin{equation*}
\sigma^\mu:=(I_2,\boldsymbol\sigma),\qquad
\bar\sigma^\mu:=(I_2,-\boldsymbol\sigma),\qquad \mu=0,1,2,3.
\tag{B.9}\label{eq:B.9}
\end{equation*}
Using them, the $\gamma$ matrices in the Weyl representation are defined as:
\begin{equation*}
\gamma^\mu:=
\begin{pmatrix}
0&\sigma^\mu\\
\bar\sigma^\mu&0
\end{pmatrix}_{4\times4}.
\tag{B.10}\label{eq:B.10}
\end{equation*}
These matrices satisfy
\begin{equation*}
\{\gamma^\mu,\gamma^\nu\}=2\eta^{\mu\nu}I_4.
\tag{B.11}\label{eq:B.11}
\end{equation*}
\begin{equation*}
\begin{gathered}
(\gamma^0)^2=I_4,\qquad
(\gamma^i)^2=-I_4,\\
(\gamma^0)^\dagger=\gamma^0,\qquad
(\gamma^i)^\dagger=-\gamma^i.
\end{gathered}
\tag{B.12}\label{eq:B.12}
\end{equation*}
Hence each $\gamma^\mu$ is unitary:
\begin{equation*}
(\gamma^\mu)^\dagger\gamma^\mu=I_4.
\tag{B.13}\label{eq:B.13}
\end{equation*}
Define
\begin{equation*}
\gamma^5:=i\gamma^0\gamma^1\gamma^2\gamma^3
=\begin{pmatrix}-I_2&0\\0&I_2\end{pmatrix}
.
\tag{B.14}\label{eq:B.14}
\end{equation*}
It satisfies
\begin{equation*}
(\gamma^5)^2=I_4,\qquad
\{\gamma^5,\gamma^\mu\}=0,\qquad
(\gamma^5)^\dagger = (\gamma^5).
\tag{B.15}\label{eq:B.15}
\end{equation*}

\setcounter{section}{2}
\section{IO(1,3) Group}\label{book-section-B.3}\label{classical-poincare}

\topic{The Group O(1,3)}

To distinguish a matrix $\Lambda$ from the corresponding tensor $\Lambda$ of coordinate transformation, we denote the matrix by $[\Lambda]$ and the tensor simply by $\Lambda$. We use the same notation for $\eta$.

The matrix $[\Lambda]$
\begin{equation*}
O(1,3):=\{[\Lambda]\in GL(4,\mathbb R)\mid [\Lambda]^T[\eta][\Lambda]=[\eta]\}.
\tag{B.16}\label{eq:B.16}
\end{equation*}
where $[\eta]=\operatorname{diag}(1,-1,-1,-1)$.

The tensor $\Lambda$

Given a specific diffeomorphism of Minkowskispacetime $\phi:x\mapsto x'$, we obtain a tensor
\begin{equation*}
\Lambda^\mu{}_\nu:=\frac{\partial x'^\mu}{\partial x^\nu}.
\tag{B.17}\label{eq:B.17}
\end{equation*}
By the way, we can prove that the isometries of Minkowski spacetime must be linear up to a translation.

The entry of $[\Lambda]$ in row $\mu$ and column $\nu$ is the corresponding component of the tensor $\Lambda$:
\begin{equation*}
[\Lambda]_{\mu\nu}=\Lambda^\mu{}_\nu.
\tag{B.18}\label{eq:B.18}
\end{equation*}
Similarly,
\begin{equation*}
[\eta]_{\mu\nu}=\eta_{\mu\nu}.
\tag{B.19}\label{eq:B.19}
\end{equation*}
Therefore, equivalently, we obtain the following relationship:
\begin{equation*}
\eta_{\mu\nu}\Lambda^\mu{}_\rho\Lambda^\nu{}_\sigma=\eta_{\rho\sigma}.
\tag{B.20}\label{eq:B.20}
\end{equation*}
where we have used the relation $([\Lambda]^T)_{\mu\nu}=\Lambda^\nu{}_\mu$. This is exactly the condition for $\phi$ to be an isometry of Minkowski spacetime.

The inverse and transpose of $[\Lambda]$ satisfy
\begin{equation*}
[\Lambda]^{-1}=[\eta][\Lambda]^T[\eta],\qquad
[\Lambda][\eta][\Lambda]^T=[\eta],
\tag{B.21}\label{eq:B.21}
\end{equation*}
so $[\Lambda]^{-1}$ and $[\Lambda]^T$ also belong to $O(1,3)$.

In components, the first relation reads
\begin{equation*}
(\Lambda^{-1})^\mu{}_\nu=\Lambda_\nu{}^\mu.
\tag{B.22}\label{eq:B.22}
\end{equation*}

\topic{Connectedness}

Taking the determinant of $[\Lambda]^T[\eta][\Lambda]=[\eta]$ gives
\begin{equation*}
(\det[\Lambda])^2=1\quad\Longrightarrow\quad\det[\Lambda]=\pm1.
\tag{B.23}\label{eq:B.23}
\end{equation*}
Setting $\rho=\sigma=0$ in $\eta_{\mu\nu}\Lambda^\mu{}_\rho\Lambda^\nu{}_\sigma=\eta_{\rho\sigma}$ gives
\begin{equation*}
(\Lambda^0{}_0)^2=1+\sum_i(\Lambda^i{}_0)^2\ge1\quad\Longrightarrow\quad\Lambda^0{}_0\ge1\ \text{or}\ \Lambda^0{}_0\le-1.
\tag{B.24}\label{eq:B.24}
\end{equation*}
We can prove that $O(1,3)$ has exactly these four connected components.

Define the time reversal $T$ and the space inversion $P$ by
\begin{equation*}
T:(t,\mathbf x)\mapsto(-t,\mathbf x),\qquad
P:(t,\mathbf x)\mapsto(t,-\mathbf x),
\tag{B.25}\label{eq:B.25}
\end{equation*}
whose matrices are
\begin{equation*}
[T]=\operatorname{diag}(-1,1,1,1),\qquad
[P]=\operatorname{diag}(1,-1,-1,-1).
\tag{B.26}\label{eq:B.26}
\end{equation*}
The corresponding tensors are
\begin{equation*}
T^\mu{}_\nu=\delta^\mu{}_\nu-2\delta^\mu{}_0\delta^0{}_\nu,\qquad
P^\mu{}_\nu=-\delta^\mu{}_\nu+2\delta^\mu{}_0\delta^0{}_\nu.
\tag{B.27}\label{eq:B.27}
\end{equation*}

\begin{center}
\renewcommand{\arraystretch}{1.2}
\begin{tabular}{cccc}
\toprule
Connected component & $\det[\Lambda]$ & $\Lambda^0{}_0$ & Representative \\
\midrule
$\mathcal L^\uparrow_+$ & $+1$ & $\ge1$ & $I$ \\
$\mathcal L^\downarrow_+$ & $+1$ & $\le-1$ & $[PT]$ \\
$\mathcal L^\uparrow_-$ & $-1$ & $\ge1$ & $[P]$ \\
$\mathcal L^\downarrow_-$ & $-1$ & $\le-1$ & $[T]$ \\
\bottomrule
\end{tabular}
\end{center}

$SO(1,3)$ consists of the two classes with $\det[\Lambda]=+1$:
\begin{equation*}
SO(1,3)=\mathcal L^\uparrow_+\cup\mathcal L^\downarrow_+.
\tag{B.28}\label{eq:B.28}
\end{equation*}
Furthermore, $SO(1,3)^+=\mathcal L^\uparrow_+$ is a subgroup of $O(1,3)$.

Physically, $T$ and $P$ correspond to time-reversal symmetry and parity symmetry, respectively. What is usually called Lorentz symmetry is invariance under $SO(1,3)^+$ only. The distinction matters: no violation of Lorentz symmetry has ever been observed, so it is regarded as a fundamental symmetry. Parity, in contrast, is violated in the weak interaction, as proposed by Lee and Yang and confirmed experimentally by Wu et al.; time-reversal symmetry is violated as well.

\topic{Lie Algebra of O(1,3)}

Consider a one-parameter subgroup $e^{tX}$, $t\in\mathbb R$, of $O(1,3)$. It satisfies
\begin{equation*}
(e^{tX})^T[\eta]\,e^{tX}=[\eta].
\tag{B.29}\label{eq:B.29}
\end{equation*}
Expanding the left-hand side to first order in $t$ gives
\begin{equation*}
(I+tX^T)[\eta](I+tX)+O(t^2)=[\eta]+t\left(X^T[\eta]+[\eta]X\right)+O(t^2),
\tag{B.30}\label{eq:B.30}
\end{equation*}
so the first-order term must vanish:
\begin{equation*}
X^T[\eta]+[\eta]X=0.
\tag{B.31}\label{eq:B.31}
\end{equation*}
Therefore,
\begin{equation*}
\mathfrak o(1,3):=\{X\in M_4(\mathbb R)\mid X^T[\eta]+[\eta]X=0\}.
\tag{B.32}\label{eq:B.32}
\end{equation*}
Furthermore, this relation can be written more explicitly as
\begin{equation*}
[\eta]X=-X^T[\eta]=-([\eta]X)^T,
\tag{B.33}\label{eq:B.33}
\end{equation*}
which means that the matrix $[\eta]X$ is antisymmetric.

An antisymmetric $4\times4$ matrix has $4\times3/2=6$ independent real parameters. Hence the general form of $X$ is
\begin{equation*}
X=\begin{pmatrix}
0&b^1&b^2&b^3\\
b^1&0&-\theta^3&\theta^2\\
b^2&\theta^3&0&-\theta^1\\
b^3&-\theta^2&\theta^1&0
\end{pmatrix}
=\sum_i\theta^iJ^i+\sum_ib^iK^i
\tag{B.34}\label{eq:B.34}
\end{equation*}
In physics, the following notation is commonly used:
\begin{equation*}
e^{-\frac i2\omega^{\mu\nu}J^{\mu\nu}}
\tag{B.35}\label{eq:B.35}
\end{equation*}
where $\omega_{\mu\nu}$ is antisymmetric in its indices, and
\begin{equation*}
-iJ^{ij}=\epsilon_{ijk}J^k,\qquad -iJ^{0i}=K^i.
\tag{B.36}\label{eq:B.36}
\end{equation*}
Each $J^{\mu\nu}$ is simply a $4\times4$ matrix, with $\alpha,\beta$ labelling its rows and columns. Its entries are
\begin{equation*}
(J^{\mu\nu})_{\alpha\beta}=i\left(\eta^{\mu\alpha}\delta^{\nu\beta}-\eta^{\nu\alpha}\delta^{\mu\beta}\right).
\tag{B.37}\label{eq:B.37}
\end{equation*}
Explicitly (row index $\alpha$, column index $\beta$; the remaining components follow from $J^{\nu\mu}=-J^{\mu\nu}$), the rotation generators $J^{ij}=i\epsilon_{ijk}J^k$ are Hermitian:
\begin{equation*}
J^{23}=\begin{pmatrix}0&0&0&0\\0&0&0&0\\0&0&0&-i\\0&0&i&0\end{pmatrix},\quad
J^{31}=\begin{pmatrix}0&0&0&0\\0&0&0&i\\0&0&0&0\\0&-i&0&0\end{pmatrix},\quad
J^{12}=\begin{pmatrix}0&0&0&0\\0&0&-i&0\\0&i&0&0\\0&0&0&0\end{pmatrix},
\tag{B.38}\label{eq:B.38}
\end{equation*}
and the boost generators $J^{0i}=iK^i$ are anti-Hermitian:
\begin{equation*}
J^{01}=\begin{pmatrix}0&i&0&0\\i&0&0&0\\0&0&0&0\\0&0&0&0\end{pmatrix},\quad
J^{02}=\begin{pmatrix}0&0&i&0\\0&0&0&0\\i&0&0&0\\0&0&0&0\end{pmatrix},\quad
J^{03}=\begin{pmatrix}0&0&0&i\\0&0&0&0\\0&0&0&0\\i&0&0&0\end{pmatrix}.
\tag{B.39}\label{eq:B.39}
\end{equation*}

\setcounter{section}{3}
\section{Green's Functions}\label{book-section-B.4}\label{greens-functions}

\topic{Definition and Integral Solution}

In the general construction below, $x,y$ denote $d$ independent variables. Let $\mathcal K_x$ be a linear differential operator on $\Omega\subseteq\mathbb R^d$, and let $\mathcal B$ specify linear boundary or initial data:
\begin{equation*}
\mathcal K_xu(x)=f(x),\qquad \mathcal Bu=b.
\tag{B.40}\label{eq:B.40}
\end{equation*}
Assume the data are sufficiently regular and the integrals below exist. Equations involving $\delta$ functions are understood in the sense of distributions.
\begin{itemize}


\item \textbf{\boldmath Green's function}: A kernel satisfying the differential equation and homogeneous boundary or initial conditions,
\begin{equation*}
{\displaystyle \mathcal K_xG(x,y)=\delta^{(d)}(x-y),\qquad \mathcal B_xG(x,y)=0.}
\tag{B.41}\label{eq:B.41}
\end{equation*}
The subscript $x$ specifies the variable on which the operator acts. The choice of boundary or initial conditions is part of the definition of $G$.

\item \textbf{\boldmath Integral solution}: If $u_{\mathrm h}$ solves the homogeneous equation with the prescribed data, then
\begin{equation*}
\begin{aligned}
{\displaystyle u(x)=u_{\mathrm h}(x)+\int_\Omega d^dy\,G(x,y)f(y),}\\
\mathcal K_xu_{\mathrm h}(x)=0,\qquad \mathcal Bu_{\mathrm h}=b.
\end{aligned}
\tag{B.42}\label{eq:B.42}
\end{equation*}
\end{itemize}

Indeed,
\begin{equation*}
\begin{aligned}
\mathcal K_xu(x)
&=\int_\Omega d^dy\,[\mathcal K_xG(x,y)]f(y)
=\int_\Omega d^dy\,\delta^{(d)}(x-y)f(y)=f(x),\\
\mathcal Bu&=\mathcal Bu_{\mathrm h}
+\int_\Omega d^dy\,[\mathcal B_xG(x,y)]f(y)=b.
\end{aligned}
\tag{B.43}\label{eq:B.43}
\end{equation*}
The solution is unique if $\mathcal Ku=0$, $\mathcal Bu=0$ has only the zero solution in the chosen function space. With nontrivial zero modes, uniqueness fails and solvability may impose compatibility conditions; an inverse on arbitrary sources cannot be assumed. Without boundary or initial conditions, two kernels satisfying the first equation in \eqref{eq:B.41} can differ by a solution of the homogeneous equation in $x$.

\topic{Construction by Fourier Transform}

For a constant-coefficient operator on Minkowski spacetime, return to $x^\mu=(t,\mathbf x)$ and the Fourier convention \eqref{eq:B.1}. If the chosen conditions are also translation invariant, $G(x,y)=G(x-y)$, and
\begin{equation*}
\begin{aligned}
\mathcal K(-ip)\widetilde G(p)&=1,\\
G(x-y)&=\int\frac{d^4p}{(2\pi)^4}\,
e^{-ip\cdot(x-y)}\widetilde G(p),\\
\widetilde G(p)&=\left[\frac{1}{\mathcal K(-ip)}\right]_{\mathrm{prescription}}.
\end{aligned}
\tag{B.44}\label{eq:B.44}
\end{equation*}
Here $\mathcal K(-ip)$ is obtained by replacing each $\partial_\mu$ with $-ip_\mu$. At zeros of $\mathcal K(-ip)$, the reciprocal requires a distributional prescription fixed by the chosen conditions; the algebraic equation alone does not select it.

\topic{Poisson Equation and Boundary Data}

First consider a smooth source of compact support in $\mathbb R^3$:
\begin{equation*}
-\nabla^2u(\mathbf x)=f(\mathbf x),\qquad
u(\mathbf x)\to0\quad\text{as }|\mathbf x|\to\infty.
\tag{B.45}\label{eq:B.45}
\end{equation*}
For $r:=|\mathbf x-\mathbf y|>0$, a rotationally invariant kernel that vanishes at infinity obeys
\begin{equation*}
\nabla^2G_0(r)=\frac{1}{r^2}\frac{d}{dr}\left(r^2\frac{dG_0}{dr}\right)=0
\quad\Longrightarrow\quad G_0(r)=\frac{A}{r}.
\tag{B.46}\label{eq:B.46}
\end{equation*}
Integrating $-\nabla^2G_0=\delta^{(3)}$ over a ball $B_\epsilon$ centred at the source fixes $A$:
\begin{equation*}
\begin{aligned}
1&=-\int_{B_\epsilon}d^3x\,\nabla^2G_0
=-4\pi\epsilon^2\left.\frac{dG_0}{dr}\right|_{r=\epsilon}
=4\pi A,\\
G_0(\mathbf x,\mathbf y)&=\frac{1}{4\pi|\mathbf x-\mathbf y|}.
\end{aligned}
\tag{B.47}\label{eq:B.47}
\end{equation*}
Hence
\begin{equation*}
{\displaystyle u(\mathbf x)=\int_{\mathbb R^3}d^3y\,
\frac{f(\mathbf y)}{4\pi|\mathbf x-\mathbf y|}.}
\tag{B.48}\label{eq:B.48}
\end{equation*}
Now let $\Omega\subset\mathbb R^3$ be bounded with smooth boundary, and prescribe $u=b$ on $\partial\Omega$. The Dirichlet Green's function for $-\nabla^2$ is symmetric in its arguments and satisfies
\begin{equation*}
\begin{aligned}
-\nabla_{\mathbf y}^2G_D(\mathbf x,\mathbf y)
&=\delta^{(3)}(\mathbf y-\mathbf x),\qquad \mathbf y\in\Omega,\\
G_D(\mathbf x,\mathbf y)&=0,\qquad \mathbf y\in\partial\Omega.
\end{aligned}
\tag{B.49}\label{eq:B.49}
\end{equation*}
Let $\mathbf n_y$ be the outward unit normal and $\partial_{n_y}:=\mathbf n_y\cdot\nabla_{\mathbf y}$. Green's second identity is
\begin{equation*}
\int_\Omega d^3y\,(v\nabla^2w-w\nabla^2v)
=\int_{\partial\Omega}dS_y\,(v\partial_{n_y}w-w\partial_{n_y}v).
\tag{B.50}\label{eq:B.50}
\end{equation*}
Apply \eqref{eq:B.50} to $v=u$, $w=G_D(\mathbf x,\cdot)$ on $\Omega\setminus B_\epsilon(\mathbf x)$ and take $\epsilon\to0^+$. For $\mathbf x\in\Omega$,
\begin{equation*}
{\displaystyle \begin{aligned}
u(\mathbf x)&=\int_\Omega d^3y\,G_D(\mathbf x,\mathbf y)f(\mathbf y)\\
&\quad-\int_{\partial\Omega}dS_y\,b(\mathbf y)\partial_{n_y}G_D(\mathbf x,\mathbf y).
\end{aligned}}
\tag{B.51}\label{eq:B.51}
\end{equation*}
The second term supplies the nonzero boundary data. The free-space kernel $G_0$ in \eqref{eq:B.47} generally does not satisfy the boundary condition in \eqref{eq:B.49}.

\topic{Klein–Gordon Equation and Initial Data}

With $m\geq0$ and the mostly-minus metric, consider
\begin{equation*}
(\Box+m^2)u(x)=f(x),\qquad
\Box:=\partial_\mu\partial^\mu=\partial_t^2-\nabla^2.
\tag{B.52}\label{eq:B.52}
\end{equation*}
\begin{itemize}


\item \textbf{\boldmath Retarded Green's function}: The translation-invariant kernel with support only at nonnegative time differences,
\begin{equation*}
(\Box+m^2)G_{\mathrm{ret}}(t,\mathbf x)
=\delta(t)\delta^{(3)}(\mathbf x),\qquad
G_{\mathrm{ret}}(t,\mathbf x)=0\quad(t<0).
\tag{B.53}\label{eq:B.53}
\end{equation*}
\end{itemize}

Here the source point has been translated to the origin. Since $\mathcal K(-ip)=m^2-p^2$, the retarded prescription is
\begin{equation*}
\widetilde G_{\mathrm{ret}}(p)
=-\lim_{\epsilon\to0^+}\frac{1}{(p^0+i\epsilon)^2-E_{\mathbf p}^2},\qquad
E_{\mathbf p}:=\sqrt{\mathbf p^2+m^2}.
\tag{B.54}\label{eq:B.54}
\end{equation*}
Both poles lie below the real $p^0$ axis. Closing the contour in the lower half-plane for $t>0$ and the upper half-plane for $t<0$ gives
\begin{equation*}
{\displaystyle G_{\mathrm{ret}}(t,\mathbf x)
=\theta(t)\int\frac{d^3p}{(2\pi)^3}\,
e^{i\mathbf p\cdot\mathbf x}\frac{\sin(E_{\mathbf p}t)}{E_{\mathbf p}}.}
\tag{B.55}\label{eq:B.55}
\end{equation*}
Here $\theta(t)$ is the Heaviside function, and the value at $E_{\mathbf p}=0$ is understood by continuity. The equal-time limits are
\begin{equation*}
G_{\mathrm{ret}}(0^+,\mathbf x)=0,\qquad
\partial_tG_{\mathrm{ret}}(0^+,\mathbf x)=\delta^{(3)}(\mathbf x).
\tag{B.56}\label{eq:B.56}
\end{equation*}
For prescribed initial data $u(t_0,\mathbf x)=u_0(\mathbf x)$ and $\partial_tu(t_0,\mathbf x)=u_1(\mathbf x)$, \eqref{eq:B.56} yields, for $t>t_0$,
\begin{equation*}
{\displaystyle \begin{aligned}
u(t,\mathbf x)
&=\int d^3y\,\bigl[
\partial_tG_{\mathrm{ret}}(t-t_0,\mathbf x-\mathbf y)u_0(\mathbf y)\\
&\hspace{4.5em}+G_{\mathrm{ret}}(t-t_0,\mathbf x-\mathbf y)u_1(\mathbf y)\bigr]\\
&\quad+\int_{t_0}^{t}dt'\int d^3y\,
G_{\mathrm{ret}}(t-t',\mathbf x-\mathbf y)f(t',\mathbf y).
\end{aligned}}
\tag{B.57}\label{eq:B.57}
\end{equation*}
For $m=0$, the spatial integral in \eqref{eq:B.55} gives the wave-equation kernel. Writing $r:=|\mathbf x|$,
\begin{equation*}
\begin{aligned}
G_{\mathrm{ret}}(t,\mathbf x)\big|_{m=0}
&=\frac{\theta(t)}{2\pi^2r}\int_0^\infty dp\,\sin(pr)\sin(pt)\\
&={\displaystyle \frac{\delta(t-r)}{4\pi r}}.
\end{aligned}
\tag{B.58}\label{eq:B.58}
\end{equation*}
With zero initial data, extend $f(t',\mathbf y)$ by zero for $t'<t_0$. Then \eqref{eq:B.57} reduces to
\begin{equation*}
{\displaystyle u(t,\mathbf x)=\int d^3y\,
\frac{f(t-|\mathbf x-\mathbf y|,\mathbf y)}{4\pi|\mathbf x-\mathbf y|}.}
\tag{B.59}\label{eq:B.59}
\end{equation*}
For comparison with the scalar Feynman propagator defined in \eqref{eq:4.12}, the Feynman prescription for the same normalization $\mathcal KG_F=\delta^{(4)}$ gives
\begin{equation*}
\begin{aligned}
G_F(x-y)&=-\lim_{\varepsilon\to0^+}\int\frac{d^4p}{(2\pi)^4}\,
\frac{e^{-ip\cdot(x-y)}}{p^2-m^2+i\varepsilon}
=i\Delta_F(x-y),\\
(\Box_x+m^2)\Delta_F(x-y)&=-i\delta^{(4)}(x-y).
\end{aligned}
\tag{B.60}\label{eq:B.60}
\end{equation*}

\clearpage
\end{CJK*}
\end{document}
